% !TEX program = xelatex
% =====================================================================
%  Lộ trình Data Engineer: xây và vận hành lakehouse mã nguồn mở
%  Phiên bản 2.0 nháp, cập nhật 06/11/2026
%  Biên dịch bằng XeLaTeX (chạy 2 lần để có mục lục).
%  Overleaf: Menu > Settings > Compiler = XeLaTeX.
% =====================================================================
% Biên dịch: xelatex main.tex (chạy 2 lần để có mục lục)
\documentclass[11pt,a4paper,oneside]{report}

% ---------- Font và ngôn ngữ ----------
\usepackage{fontspec}
\usepackage{polyglossia}
\setdefaultlanguage{vietnamese}
\setotherlanguage{english}
% Nạp font theo tên file: Tectonic chỉ thấy font trong cây TeX theo tên file, TeX Live cũng tìm được
\setmainfont{texgyrepagella}[Extension=.otf,UprightFont=*-regular,BoldFont=*-bold,ItalicFont=*-italic,BoldItalicFont=*-bolditalic]
\setsansfont{texgyreheros}[Extension=.otf,UprightFont=*-regular,BoldFont=*-bold,ItalicFont=*-italic,BoldItalicFont=*-bolditalic]
% Font mono phải có đủ dấu tiếng Việt (DejaVu Sans Mono thiếu ổ, ỗ...)
\IfFontExistsTF{NotoSansMono-Regular.ttf}{\setmonofont{NotoSansMono}[Extension=.ttf,UprightFont=*-Regular,BoldFont=*-Bold,Scale=0.86]}{%
  \IfFontExistsTF{Liberation Mono}{\setmonofont{Liberation Mono}[Scale=0.9]}{%
    \setmonofont{texgyrecursor}[Extension=.otf,UprightFont=*-regular,BoldFont=*-bold,ItalicFont=*-italic,BoldItalicFont=*-bolditalic,Scale=1.0]}}

% ---------- Bố cục ----------
\usepackage[a4paper,top=2.2cm,bottom=2.4cm,left=2.3cm,right=2.3cm,headheight=14pt]{geometry}
\usepackage{microtype}
\usepackage{xcolor}
\usepackage{array,booktabs,xltabular,multirow}
\usepackage{enumitem}
\usepackage{fvextra}
\usepackage[most]{tcolorbox}
\usepackage[nobottomtitles*]{titlesec}
\renewcommand{\bottomtitlespace}{0.12\textheight}
\usepackage{fancyhdr}
\usepackage{etoolbox}
\usepackage[colorlinks=true,linkcolor=linkblue,urlcolor=linkblue,citecolor=linkblue,
            pdftitle={Lộ trình Data Engineer: xây lakehouse mã nguồn mở}]{hyperref}
\usepackage{bookmark}

% ---------- Màu ----------
\definecolor{linkblue}{HTML}{1F5FAD}
\definecolor{ink}{HTML}{1E2329}
\definecolor{muted}{HTML}{5B636E}
\definecolor{accent}{HTML}{0F6E56}
\definecolor{accentbg}{HTML}{EAF5F0}
\definecolor{warnline}{HTML}{B5651D}
\definecolor{warnbg}{HTML}{FBF2E8}
\definecolor{infoline}{HTML}{1F5FAD}
\definecolor{infobg}{HTML}{EEF4FB}
\definecolor{rulegray}{HTML}{C9CED6}
\definecolor{codebg}{HTML}{F5F6F8}

\color{ink}
\setlength{\parindent}{0pt}
\setlength{\emergencystretch}{3em}
\setlength{\parskip}{0.55em}
\linespread{1.08}
\setlist{itemsep=0.15em,topsep=0.3em,parsep=0.1em}

% ---------- Tiêu đề chương, mục ----------
\titleformat{\chapter}[display]
  {\sffamily\bfseries\color{ink}}
  {\large\color{accent}\MakeUppercase{\chaptertitlename}\ \thechapter}
  {0.3em}{\LARGE}
\titlespacing*{\chapter}{0pt}{-1.2em}{1.2em}
\titleformat{\section}{\sffamily\bfseries\large\color{ink}}{\thesection}{0.6em}{}
\titleformat{\subsection}{\sffamily\bfseries\normalsize\color{ink}}{\thesubsection}{0.6em}{}
\titlespacing*{\section}{0pt}{1.3em}{0.4em}
\titlespacing*{\subsection}{0pt}{1.0em}{0.3em}
\setcounter{secnumdepth}{0}
\setcounter{tocdepth}{0}
\usepackage{tocloft}
\setlength{\cftbeforechapskip}{0.5em}
\renewcommand{\cfttoctitlefont}{\sffamily\bfseries\LARGE}
\renewcommand{\cftchapfont}{\sffamily\bfseries\small}
\renewcommand{\cftchappagefont}{\sffamily\bfseries\small}

% Chương không thuộc chặng (mở đầu, phụ lục): \plainchapter{tên}
\newcommand{\plainchapter}[1]{%
  \clearpage\phantomsection
  \addcontentsline{toc}{chapter}{#1}%
  \markboth{#1}{}%
  \thispagestyle{plain}%
  {\sffamily\bfseries\LARGE\color{ink}#1\par}\vspace{1.0em}}

% ---------- Header, footer ----------
\pagestyle{fancy}
\fancyhf{}
\fancyhead[L]{\small\sffamily\color{muted}\leftmark}
\fancyhead[R]{\small\sffamily\color{muted}Lộ trình Data Engineer}
\fancyfoot[C]{\small\sffamily\color{muted}\thepage}
\renewcommand{\headrulewidth}{0.3pt}
\renewcommand{\chaptermark}[1]{\markboth{#1}{}}
\fancypagestyle{plain}{\fancyhf{}\fancyfoot[C]{\small\sffamily\color{muted}\thepage}\renewcommand{\headrulewidth}{0pt}}

% ---------- Bảng ----------
% Cột X căn trái, không giãn chữ; P{w}: cột rộng cố định căn trái
\newcolumntype{L}{>{\raggedright\arraybackslash}X}
\newcolumntype{P}[1]{>{\raggedright\arraybackslash}p{#1}}
\newcommand{\thd}[1]{\textbf{\sffamily\small #1}}
\setlength{\tabcolsep}{4.5pt}
\renewcommand{\arraystretch}{1.18}
\AtBeginEnvironment{xltabular}{\small}
\keepXColumns

% ---------- Mã, đường dẫn ----------
% \code{...}: mã ngắn trong dòng. Thoát tay: \_ \% \# \& \$ \{ \} \textbackslash{} \textasciitilde{}
\newcommand{\code}[1]{\texttt{#1}}
% \repo{...}: file hoặc lệnh trong repo lakehouse/ của bạn
\newcommand{\repo}[1]{\texttt{\color{accent}#1}}
\DefineVerbatimEnvironment{codeblock}{Verbatim}{%
  fontsize=\small,breaklines=true,breakanywhere=true,
  frame=leftline,framerule=2pt,rulecolor=\color{rulegray},
  xleftmargin=0.8em,framesep=0.8em}

% ---------- Hộp ----------
\newtcolorbox{luuy}[1][]{enhanced,breakable,colback=warnbg,colframe=warnline,
  boxrule=0pt,leftrule=3pt,arc=0pt,outer arc=0pt,left=8pt,right=8pt,top=6pt,bottom=6pt,#1}
\newtcolorbox{ghichu}[1][]{enhanced,breakable,colback=infobg,colframe=infoline,
  boxrule=0pt,leftrule=3pt,arc=0pt,outer arc=0pt,left=8pt,right=8pt,top=6pt,bottom=6pt,#1}


% ---------- Phần (Part) ----------
% \partpage{I}{Nền tảng}{mô tả ngắn}
\newcommand{\partpage}[3]{%
  \clearpage\phantomsection
  \addcontentsline{toc}{part}{Phần #1\quad #2}%
  \thispagestyle{empty}%
  \vspace*{5cm}
  {\sffamily\bfseries\large\color{accent}PHẦN #1\par}\vspace{0.5em}
  {\sffamily\bfseries\Huge\color{ink}#2\par}\vspace{1em}
  {\color{accent}\rule{0.35\linewidth}{1.2pt}}\par\vspace{1em}
  \begin{minipage}{0.8\linewidth}\large\color{muted}#3\end{minipage}
  \clearpage}

% ---------- Module ----------
% \modulechapter{5}{Open table format: Apache Iceberg}
\newcommand{\modulechapter}[2]{%
  \clearpage\phantomsection
  \addcontentsline{toc}{chapter}{Module #1\quad #2}%
  \markboth{Module #1: #2}{}%
  \thispagestyle{plain}%
  {\sffamily\bfseries\large\color{accent}MODULE #1\par}\vspace{0.35em}%
  {\sffamily\bfseries\LARGE\color{ink}#2\par}\vspace{1.0em}%
  \label{module:#1}}
% \moduleinfo{cấp độ}{thời lượng}{tiên quyết}
\newcommand{\moduleinfo}[3]{%
  \begin{tcolorbox}[enhanced,colback=codebg,colframe=codebg,arc=0pt,boxrule=0pt,
     left=8pt,right=8pt,top=4pt,bottom=4pt,fontupper=\small\sffamily]
  \textbf{Cấp độ:} #1 \quad\textbar\quad \textbf{Thời lượng:} #2 \quad\textbar\quad \textbf{Tiên quyết:} #3
  \end{tcolorbox}}

% Hộp "Thử thách kỹ năng" (giống skill badge challenge lab)
\newtcolorbox{thuthach}[1][]{enhanced,breakable,colback=accentbg,colframe=accent,
  boxrule=0pt,leftrule=3pt,arc=0pt,outer arc=0pt,left=8pt,right=8pt,top=6pt,bottom=6pt,
  fonttitle=\sffamily\bfseries\small,coltitle=accent,colbacktitle=accentbg,
  title={Thử thách kỹ năng},attach title to upper={\par\smallskip},#1}
% Hộp "Hoàn thành khi"
\newcommand{\hoanthanh}[1]{\par\smallskip\textbf{\sffamily\small\color{accent}Hoàn thành khi:} #1\par}

\renewcommand{\cftpartfont}{\sffamily\bfseries\normalsize\color{accent}}
\renewcommand{\cftpartpagefont}{\sffamily\bfseries\normalsize\color{accent}}
\setlength{\cftbeforepartskip}{1.0em}


\begin{document}

\begin{titlepage}
\thispagestyle{empty}
\vspace*{3.5cm}
{\sffamily\bfseries\color{accent}\large LỘ TRÌNH HỌC THEO VAI TRÒ\par}
\vspace{0.6em}
{\sffamily\bfseries\Huge\color{ink} Data Engineer\par}
\vspace{0.3em}
{\sffamily\bfseries\LARGE\color{ink} Xây và vận hành lakehouse mã nguồn mở\par}
\vspace{0.6em}
{\sffamily\large\color{muted} Từ SQL tới production: Iceberg, Spark, Airflow, Trino, Kafka, AWS\par}
\vspace{2.2em}
{\color{accent}\rule{\linewidth}{1.2pt}}\par
\vspace{1.2em}
\begin{minipage}{0.88\linewidth}
\large
\textbf{5 phần \textbar{} 13 module \textbar{} 1 dự án capstone \textbar{} 288 giờ lõi}\par\medskip
Mỗi module là một tutorial thực hành trên dữ liệu Olist và repo \repo{lakehouse}, gồm bài học có code chạy được kèm output thật, lab kết thúc bằng một con số kiểm được, và một bài thử thách kỹ năng. Phiên bản phần mềm được ghim và đối chiếu với tài liệu chính thức của từng dự án.
\end{minipage}
\vfill
{\sffamily\color{muted}Phiên bản 2.0 nháp \textbar{} cập nhật 06/11/2026\par}
\end{titlepage}

\tableofcontents
\clearpage

% ---------------- front ----------------
% ====================== PHẦN MỞ ĐẦU ======================
\plainchapter{Về lộ trình này}

Lộ trình này đưa bạn từ SQL tới việc tự xây và vận hành một \textbf{lakehouse mã nguồn mở}: dữ liệu nằm trên object storage, được tổ chức thành bảng Apache Iceberg, quản lý bằng một REST catalog, xử lý bằng Spark, điều phối bằng Airflow, kiểm tra chất lượng trước khi phục vụ qua Trino và BI, rồi đưa lên cloud.

Cách trình bày theo mô hình của các lộ trình học chính thức. Giống \href{https://www.skills.google/paths/16}{learning path Professional Data Engineer của Google Cloud}, lộ trình đi theo vai trò công việc, gồm nội dung học, lab thực hành và một bài thử thách cuối mỗi phần để chứng minh kỹ năng, rồi ánh xạ sang chứng chỉ. Từ bản 2.0, mỗi module mở bằng một vấn đề thật có số đo, rồi tới các bài học có code chạy được; khung đầy đủ nằm ở mục \textit{Cách tổ chức}.

\section{Dành cho ai}
\begin{itemize}
  \item Sinh viên công nghệ thông tin và kỹ sư phần mềm muốn chuyển sang data engineering. Nếu bạn đã làm phần mềm, đọc mục \textit{Bắt đầu từ nền SWE} ở cuối chương này trước Module 1.
  \item Data analyst muốn tự xây pipeline thay vì chỉ dùng dữ liệu có sẵn.
  \item Kỹ sư dữ liệu đã quen warehouse truyền thống, muốn hiểu lakehouse và open table format từ bên trong.
\end{itemize}
Bạn cần biết lập trình Python ở mức cơ bản, dùng được dòng lệnh và Docker. Không cần biết trước SQL: Module 1 dạy SQL từ câu \code{SELECT}, lấy ví dụ từ ORM cho người đã làm app, và có thêm tới 25 giờ luyện tập cho người mới bắt đầu.

\section{Học xong bạn làm được gì}
\begin{itemize}
  \item Viết SQL phân tích nhiều bảng và thiết kế mô hình dữ liệu dạng star schema đúng grain.
  \item Giải thích một bảng Iceberg được tạo thành từ những file nào, một lần commit diễn ra thế nào, và vì sao catalog là nơi duy nhất giữ trạng thái.
  \item Xây pipeline tăng dần, chạy lại không sinh trùng, có test chặn dữ liệu sai trước khi tới tầng phục vụ.
  \item Điều phối pipeline bằng Airflow 3, phục vụ dữ liệu qua Trino và dashboard, thêm một nhánh streaming với Kafka và Flink.
  \item Đưa lakehouse lên AWS bằng Terraform với giới hạn chi phí, và đối chiếu được với các dịch vụ tương đương trên Google Cloud.
\end{itemize}

\section{Cách tổ chức}
Lộ trình có \textbf{5 phần}: bốn phần học theo module, và Phần V cho capstone, chứng chỉ và nghề nghiệp. Phần lõi gồm \textbf{trang khởi động cho SWE và 13 module}, tổng \textbf{288 giờ}. Dự án capstone được xây dần qua bốn mốc, ở cuối Phần I tới IV. Ngoài lõi có \textbf{3 nhánh tuỳ chọn}, tổng 26 giờ.

Lịch tính 10 giờ mỗi tuần: 9 giờ học module và 1 giờ cho chương nghề nghiệp. Với nhịp đó, phần lõi mất khoảng \textbf{8 tháng}; thêm ôn chứng chỉ và cả ba nhánh thì gần 10 tháng. Mỗi module viết theo kiểu tutorial thực hành và có cùng một khung:

\begin{xltabular}{\textwidth}{@{}P{3.6cm} L@{}}
\toprule
\thd{mục} & \thd{vai trò} \\
\midrule
\endhead
\bottomrule
\endfoot
Thẻ module & cấp độ, tiên quyết, giờ lõi và giờ thêm; file, bảng và service của repo mà module dùng; phiên bản ghim kèm link docs đúng bản; bài nào lướt được nhờ nền SWE, kèm một bài kiểm 15 phút \\
Vấn đề thật & một chuyện có số đo trên Olist hoặc trong production, kết bằng câu hỏi mà module trả lời \\
Từ SWE sang DE & bảng nối thứ bạn đã biết với tên gọi trong DE, ghi giống ở đâu và khác ở đâu; mỗi chỗ khác được xử lý trong một bài \\
Bạn sẽ làm được & bốn tới sáu việc bắt đầu bằng động từ, mỗi việc khớp một điều kiện ``hoàn thành khi'' \\
Bài học & mỗi khái niệm một bài: vấn đề, giải thích bằng lời thường từ thứ bạn đã biết, so với SWE và chỗ so sánh gãy, làm thử bằng code có output thật, bẫy thường gặp \\
Lab & làm từng bước trên repo và dữ liệu Olist; mỗi bước có lệnh, output mong đợi và gợi ý khi lỗi; kết bằng ``hoàn thành khi'' là một lệnh trả về đúng số \\
Thử thách kỹ năng & một bài không có hướng dẫn từng bước, giống challenge lab của skill badge, kèm tiêu chí đạt và phần mang vào capstone \\
Tự kiểm tra & năm tới tám câu, đáp án ngắn ở cuối module \\
Đọc thêm & tối đa 7 mục, mỗi mục ghi vì sao đọc, đọc phần nào, mất bao lâu, phiên bản nào \\
Kế hoạch theo tuần & chia module theo tuần 9 giờ, tách phần lõi và phần thêm nếu còn giờ; tuần cuối làm thử thách \\
\end{xltabular}
Trang khởi động cho SWE không theo khung này: nó chỉ có bốn mục, gồm thứ mang sang thẳng, thói quen phải đổi, lối đi nhanh và danh sách kiểm 1 giờ trên repo.

\section{Chọn điểm bắt đầu}
Không phải ai cũng cần học theo đúng thứ tự từ Module 1. Hãy chọn lối vào hợp với bạn, giống cách các học viện của nhà cung cấp chia track theo người học.

\begin{xltabular}{\textwidth}{@{}P{3.4cm} L L@{}}
\toprule
\thd{bạn là} & \thd{bắt đầu ở} & \thd{ghi chú} \\
\midrule
\endhead
\bottomrule
\endfoot
Mới với dữ liệu & Module 1, đi theo thứ tự & giữ đủ thời lượng của Module 1 và 3, đó là nền cho mọi phần sau \\
Đã vững SQL & làm thử thách kỹ năng của Module 1; đạt thì sang Module 2 và 3 & phần lớn thời gian nên dồn vào Phần II \\
Kỹ sư phần mềm, backend hoặc cloud & mục \textit{Bắt đầu từ nền SWE} ở cuối chương này, rồi Module 1 theo thứ tự & làm danh sách kiểm 1 giờ trên repo trước; ở mỗi module, bài kiểm 15 phút trong thẻ module cho biết bài nào lướt được \\
Đã quen warehouse & Module 4 tới 6 trước, rồi quay lại Module 7 & đọc kỹ phần catalog và commit, đó là chỗ lakehouse khác warehouse nhiều nhất \\
\end{xltabular}

\section{Nguyên tắc về nguồn}
\begin{itemize}
  \item Khẳng định kỹ thuật dẫn tới \textbf{tài liệu chính thức}: docs, spec, release notes của dự án, hướng dẫn của nhà cung cấp cloud, trang của chương trình chứng chỉ.
  \item Bài viết của kỹ sư có tên tuổi và sách cũng được trích, nhất là ở phần so với SWE. Chúng luôn được ghi rõ là ``bài viết'' hay ``sách'', vì đó là ý kiến của tác giả chứ không phải đặc tả.
  \item Phụ lục B liệt kê sách đọc thêm theo module; sách chỉ trích ở một chỗ thì ghi ngay tại chỗ đó.
  \item Mọi con số ghi kèm cách đo: lệnh đã chạy cùng ngày và phiên bản, file và dòng của repo, hoặc link nguồn. Cùng một phép đo có thể ra số khác nhau tuỳ cách dựng, ví dụ đọc từ bảng đã nạp hay đọc thẳng file CSV, nên cách đo quan trọng ngang con số.
  \item Output dưới mỗi khối code là output thật kèm ngày chạy; khối chưa chạy được ghi rõ ``chưa chạy'' và lý do. Nguồn và lệnh kiểm của từng câu nằm trong các comment \code{\% Nguồn} và \code{\% Kiểm} của file \code{.tex}.
  \item Phiên bản phần mềm và thông tin chứng chỉ được kiểm ngày \textbf{02/11/2026}. Phần mềm thay đổi nhanh: hãy ghim phiên bản trong môi trường học và luôn mở bộ docs đúng phiên bản bạn chạy, không mặc định mở bản \code{latest}.
  \item Thời lượng là ước tính cho phần lõi. Mỗi module có phần ``thêm'' để bạn quay lại khi còn thời gian.
\end{itemize}

% Trang khởi động M0 (2 giờ), bản 2.0 nháp: mục cuối của chương "Về lộ trình này".
% Output chạy ngày 06/11/2026 từ thư mục có symlink lakehouse tới repo: uv 0.10.9, Python 3.10.20,
% GNU Make 3.81, openrsync protocol 29, OpenJDK 11.0.30. Lúc kiểm, cache uv, TMPDIR, venv Airflow
% và thư mục Ivy của Spark nằm trong thư mục nháp, nên repo lakehouse không bị ghi gì.
\section{Bắt đầu từ nền SWE}

Nếu bạn đã làm phần mềm vài năm, bạn không bắt đầu từ số 0. Git, review, test, CI, hạ tầng dạng code (IaC) và Docker Compose dùng gần như nguyên vẹn trong lộ trình này. Phần phải học lại là cách nghĩ về dữ liệu, và lỗi ở phần này thường im lặng: ở Module 3, một câu join quen tay cho tổng tiền cao hơn 26,9\% mà database không báo lỗi nào.
% Kiểm: chạy lại khối A của M3 ngày 06/11/2026 bằng lakehouse/.venv/bin/python, DuckDB 1.5.6: 16008872.12 và 20308134.71, cao hơn 26.9%.

Trang này là bản đồ, không phải bài học. Nó tính 2 giờ trong phần lõi: một giờ đọc ba mục đầu, một giờ làm danh sách kiểm trên repo \repo{lakehouse}. Lời giải thích cho từng chỗ khác nằm ở module được trỏ tới.

\subsection{Thứ mang sang thẳng}

Sáu thứ dưới đây bạn dùng như khi làm app. Cột cuối là chỗ dữ liệu làm chúng khác đi.
\begin{xltabular}{\textwidth}{@{}P{2.6cm} L L@{}}
\toprule
\thd{bạn đã biết} & \thd{dùng ở đâu trong lộ trình} & \thd{chỗ khác cần để ý} \\
\midrule
\endhead
\bottomrule
\endfoot
Git và pull request & code của lab, DAG, câu SQL và job Spark; capstone là một repo của bạn & Git giữ code, không giữ dữ liệu: CSV do \repo{make data} tải về, còn bảng nằm trong bucket. Theo bài viết ``Functional Data Engineering'' của Maxime Beauchemin, chạy lại ra cùng số cần cả dữ liệu không đổi lẫn code có version \\
Code review & SQL, DAG và job Spark được review như code của app & bạn review thêm dữ liệu mà code mới sinh ra, bằng cách so số dòng và tổng tiền trước và sau (Module 10) \\
Unit test với pytest & repo có 93 test, trong đó \repo{make test} chạy 88 test mà không cần Docker & test xanh chỉ nói code đúng trên dữ liệu mẫu. Sách \textit{The Site Reliability Workbook} của Google, chương 13, ghi rằng unit test không bắt được mọi lỗi của pipeline; test trên dữ liệu thật nằm ở Module 10 \\
CI và lint & GitHub Actions chạy Ruff và pytest ở Module 10; \repo{make lint} chạy Ruff trên máy & test cần cả stack được tách ra bằng marker \code{e2e} của pytest, nên \repo{make test} chạy được mà không cần Docker, như trên một runner CI chỉ có uv và Java \\
IaC & Terraform dựng một lát cắt AWS ở Module 13 rồi xoá & bucket chứa bảng là dữ liệu, không dựng lại được bằng \code{apply}. Mặc định Terraform không xoá bucket còn object; bật \code{force\_\allowbreak{}destroy} thì lệnh xoá hạ tầng xoá luôn mọi object, không khôi phục được \\
Docker Compose & stack 6 service trong \repo{docker-compose.yml}, từ object storage tới Spark Connect & container dựng lại được, volume thì không: \repo{make down} giữ 3 named volume, còn \code{docker compose down -v} xoá chúng và mọi bảng mất theo \\
\end{xltabular}
% Nguồn: bài viết https://maximebeauchemin.medium.com/functional-data-engineering-a-modern-paradigm-for-batch-data-processing-2327ec32c42a (08/01/2018), đọc qua web.archive.org bản 20260812053740 ngày 06/11/2026: "To put it simply, immutable data along with versioned logic are key to reproducibility."
% Nguồn: sách The Site Reliability Workbook, chương 13 Data Processing Pipelines, https://sre.google/workbook/data-processing/ mục Staging: "Unit tests won't catch all pipeline issues". https://raw.githubusercontent.com/hashicorp/terraform-provider-aws/main/website/docs/r/s3_bucket.html.markdown: force_destroy mặc định false, object bị xoá không khôi phục được.
% Kiểm: collect-only với marker như make test trên bản chép, 06/11/2026: "88/93 tests collected (5 deselected)". lakehouse/pyproject.toml dòng 36 (marker e2e); Makefile dòng 33, 34 (test), 43, 44 (lint); docker-compose.yml dòng 10, 28, 45, 62, 88, 137 (6 service), 153 tới 156 (3 volume); README.md dòng 219, 220 (down giữ volume, down -v mất mọi bảng).

\subsection{Thói quen phải đổi}

Sáu thói quen dưới đây đúng với app nhưng sai với dữ liệu. Mỗi dòng trỏ tới module dạy cách làm mới bằng code chạy được, nên ở đây bạn chỉ cần nhận ra chúng.
\begin{xltabular}{\textwidth}{@{}L L P{2.4cm}@{}}
\toprule
\thd{thói quen SWE} & \thd{trong DE thành} & \thd{học ở} \\
\midrule
\endhead
\bottomrule
\endfoot
vòng lặp gọi một query cho mỗi bản ghi & một câu theo tập cho cả bảng & Module 1 \\
sửa dòng tại chỗ bằng \code{UPDATE} & ghi version mới; bảng hiện tại suy ra từ lịch sử thay đổi & Module 5 và 8 \\
retry an toàn nhờ header \code{Idempotency-Key} & cả job chạy lại không nhân đôi dòng & Module 8 \\
unit test xanh là đủ & test dữ liệu chạy trên dữ liệu thật sau mỗi lần ghi & Module 10 \\
hợp đồng API viết bằng OpenAPI & bảng là API, có hợp đồng dữ liệu (data contract) chạy được như một bài kiểm & Module 10 \\
giá server tính theo tháng & giá tính theo byte quét & Module 13 \\
\end{xltabular}

\subsection{Lối đi nhanh}

Mỗi module có dòng \textit{Lướt được} trong thẻ module. Dòng này ghi các bài bạn có thể đọc nhanh nhờ nền SWE, kèm một bài kiểm 15 phút.
\begin{itemize}
  \item Làm bài kiểm 15 phút trước khi đọc các bài đó.
  \item Đạt thì đọc lướt, chỉ dừng ở đoạn \textit{Chỗ so sánh gãy} và hộp \textit{Bẫy thường gặp}, rồi làm lab như bình thường.
  \item Không đạt thì học đủ. Lab không lướt được, vì điều kiện ``hoàn thành khi'' là con số đúng trên máy của bạn.
\end{itemize}

Bảng dưới là gợi ý cho người đã làm backend; danh sách chính xác, kèm số bài, nằm ở thẻ của từng module. Module 3 và Module 5 không có trong bảng: hãy học đủ hai module này.
\begin{xltabular}{\textwidth}{@{}l L L@{}}
\toprule
\thd{module} & \thd{thường lướt được} & \thd{đừng lướt, dù trông quen} \\
\midrule
\endhead
\bottomrule
\endfoot
1 & khoá, ràng buộc và chuẩn hoá, nếu bạn từng viết migration có khoá ngoại & \code{NULL} và \code{NOT IN}; một câu theo tập thay cho vòng lặp \\
2 & vòng đời DE; warehouse, lake và lakehouse; batch và streaming & viết yêu cầu và định nghĩa metric để hai người tính ra cùng một số \\
4 & S3 là kho key và value: key phẳng, list có phân trang, upload nhiều phần & cách Parquet bỏ qua cả khối dòng nhờ giá trị min và max \\
6 & mã hoá lúc lưu và lúc truyền; secret nằm ngoài code & xoá dữ liệu cá nhân khi bảng còn giữ các bản cũ \\
7 & DataFrame chỉ chạy khi cần kết quả, nếu bạn quen Java Streams & dữ liệu chuyển qua mạng giữa các máy (shuffle), và vài khoá chiếm phần lớn dữ liệu (skew) \\
8 & retry có backoff luỹ thừa; phân trang theo cursor & lưu cursor cùng lần commit dữ liệu; cả job chạy lại không nhân đôi dòng; dữ liệu đến muộn \\
9 & test DAG bằng pytest & lịch chạy theo khoảng dữ liệu, và nạp lại dữ liệu quá khứ \\
10 & runbook và postmortem, nếu bạn từng trực on-call & test dữ liệu; cổng write-audit-publish (ghi, kiểm rồi mới công bố); sao lưu và khôi phục catalog \\
11 & data API có cache & quyền theo dòng và theo cột trên Trino \\
12 & topic, partition, offset và consumer group, nếu bạn đã viết consumer Kafka & giờ của sự kiện (event time) và ngưỡng chờ dữ liệu đến muộn (watermark); ghi đúng một lần vào Iceberg \\
13 & IAM quyền tối thiểu, ngân sách và cảnh báo chi phí, cú pháp Terraform & \code{terraform destroy} có thể xoá cả dữ liệu; giá theo byte quét; giới hạn của S3 Tables \\
\end{xltabular}

\subsection{Danh sách kiểm 1 giờ trên repo}

Mục tiêu là chứng minh repo chạy trên máy bạn mà chưa cần Docker, và biết lệnh nào của \repo{Makefile} cần Docker. Bạn làm trên một bản chép của repo trong \code{lab\_\allowbreak{}m00/}, giống bản checkout sạch mà CI dùng. Nhờ vậy \repo{lakehouse/.venv}, venv mà các module sau dùng, không bị đụng tới.

\begin{ghichu}
\textbf{Cách chạy code trong mục này.}
\begin{itemize}
  \item Lần đầu, clone repo \repo{lakehouse/} từ \href{https://github.com/DPhO05/lakehouse}{github.com/DPhO05/lakehouse} vào một thư mục. Thư mục cha của \repo{lakehouse/} là nơi chạy các khối dưới.
  \item Mọi khối là lệnh bash, chạy từ thư mục cha của \repo{lakehouse/}. Bạn cần \code{uv}, \code{make}, \code{rsync} và mạng ở lần chạy đầu; bước 4 cần mạng ở mọi lần chạy.
  \item Lần đầu, trong \repo{lakehouse/} chạy \code{make data} để tải 9 file CSV của Olist vào \repo{source\_\allowbreak{}sim/data/}. Bước này cần mạng, và M1 cần các file này.
  \item Bước 3 cần JDK 8, 11 hoặc 17 thì mới đạt \code{88 passed}, vì Spark 3.5.5 hỗ trợ ba bản này. Máy kiểm dùng OpenJDK 11.0.30.
  \item Lệnh \code{cd} luôn nằm trong cặp ngoặc tròn, tức chạy trong một shell con (subshell), nên sau mỗi khối terminal của bạn vẫn đứng ở thư mục cha của \repo{lakehouse/}. Nếu \code{cd} lỗi, các lệnh sau nó trong ngoặc không chạy.
  \item Thời gian chạy đổi theo máy và mạng, nên các khối chỉ in những dòng không đổi giữa các lần chạy. Thời gian đo được ghi trong lời giải thích.
  \item Phần ``Kết quả'' là output thật, chạy ngày 06/11/2026 bằng uv 0.10.9 và Python 3.10.20. Lần đầu, máy chạy cả mục mất khoảng 4 tới 5 phút; phần còn lại của giờ là đọc output và tự làm bước 6.
\end{itemize}
\end{ghichu}
% Kiểm: lần đầu với cache rỗng, đo bằng /usr/bin/time ngày 05/11/2026: uv sync 62,48 giây, make test 109,21 giây, make dag-check 103,50 giây, CLI help 0,65 giây; tổng khoảng 276 giây. Lần đầu ngày 06/11/2026, cache rỗng, đo thời gian quanh từng khối: K1 0,05 giây, K2 63,87 giây, K3 111,65 giây, K4 64,20 giây, K5 0,63 giây; tổng khoảng 240 giây.
% Nguồn: https://github.com/apache/spark/blob/v3.5.5/docs/index.md dòng 37, mở ngày 06/11/2026: "Spark runs on Java 8/11/17". Kiểm: check_dags.sh dòng 7, 10 gọi uv pip install với file constraints trên GitHub ở mọi lần chạy.

\textbf{Bước 1: chép repo.} Lệnh dưới xoá bản chép cũ nếu có, rồi chép repo mà bỏ \code{.venv}, cache bytecode và 9 file CSV. Test không cần CSV thật: test của \repo{source\_\allowbreak{}sim} dùng dữ liệu mẫu riêng.

% chạy: K1
\begin{codeblock}
rm -rf lab_m00/lakehouse && mkdir -p lab_m00
rsync -a --exclude .venv --exclude __pycache__ --exclude source_sim/data \
  lakehouse/ lab_m00/lakehouse/
(cd lab_m00/lakehouse &&
  ls Makefile pyproject.toml uv.lock .python-version && cat .python-version)
\end{codeblock}
\textit{Kết quả:}
% kết quả: K1
\begin{codeblock}
.python-version
Makefile
pyproject.toml
uv.lock
3.10
\end{codeblock}

Bốn file này là thứ các bước sau đọc. \repo{.python-version} ghim Python 3.10, còn \repo{uv.lock} ghim đúng phiên bản của từng gói, giống \code{package-lock.json} của npm. Nếu \code{rsync} báo lỗi có chữ \code{No such file or directory}, bạn đang đứng sai thư mục: \repo{lakehouse/} phải nằm ngay trong thư mục hiện tại.
% Kiểm: tests/source_sim/test_app.py dòng 11, 14 đặt SIM_DATA_DIR là tests/fixtures/olist; K3 qua hết trên bản chép không có source_sim/data. Lỗi của openrsync ngày 06/11/2026: "rsync(51602): error: nosuchdir/: (l)stat: No such file or directory".

\textbf{Bước 2: cài môi trường bằng \code{uv sync}.} uv tạo \code{.venv} trong bản chép theo \repo{uv.lock}: 40 gói, gồm cả nhóm dev như pytest, DuckDB và Ruff. Dòng cuối in phiên bản của Python, DuckDB và PySpark, ba thứ mà lab dùng nhiều nhất.

% chạy: K2
\begin{codeblock}
(cd lab_m00/lakehouse && uv sync --quiet &&
  .venv/bin/python -c 'import sys, duckdb, pyspark
print(sys.version.split()[0], duckdb.__version__, pyspark.__version__)')
\end{codeblock}
\textit{Kết quả:}
% kết quả: K2
\begin{codeblock}
3.10.20 1.5.6 3.5.5
\end{codeblock}

Số bản vá của Python có thể khác 3.10.20, vì \repo{.python-version} chỉ ghim 3.10; hai số sau phải đúng là 1.5.6 và 3.5.5.

Lần đầu, uv tải hơn 400 MB, phần lớn là gói nguồn của PySpark 3.5.5 nặng 317 MB, nên lệnh im lặng khoảng 1 phút; bỏ cờ quiet nếu muốn xem tiến độ. Khi gói đã nằm trong cache của uv, lệnh này mất khoảng 4 giây. Bản chép chiếm khoảng 715 MB đĩa, gần hết là \code{.venv}.

Lộ trình ghim DuckDB 2.0.0, nhưng bản này cần Python 3.11 trở lên. Repo dùng Python 3.10, nên venv của nó không cài được bản 2.0.0; \repo{uv.lock} ghim DuckDB 1.5.6, và các module chạy DuckDB ghi rõ output là của bản nào. Nếu uv dừng ở dòng có chữ \code{Failed to download} kèm tên một gói, uv chưa tải được gói đó: kiểm mạng hoặc proxy rồi chạy lại.
% Kiểm: uv sync lần đầu với cache rỗng, 05/11/2026: "Installed 40 packages", real 62,48 giây, 8 wheel lớn nhất cộng 96,5 MiB; pyspark-3.5.5.tar.gz có size = 317214975 ở lakehouse/uv.lock dòng 508. Cache ấm, 06/11/2026: real 4,18 giây. Lần đầu, 06/11/2026, cache rỗng: 63,87 giây. du -sk, 06/11/2026: bản chép 698492 KiB, tức 715,3 MB; .venv 696092 KiB, tức 712,8 MB. Mọi MB trong trang này là 10^6 byte.
% Kiểm: uv.lock dòng 71, 72 (duckdb 1.5.6), 502, 503 (pyspark 3.5.5); pyproject.toml dòng 5 (Python từ 3.10 tới trước 3.12), 14 tới 22 (nhóm dev). Nguồn: https://pypi.org/pypi/duckdb/2.0.0/json, requires_python >=3.11.0, file upload ngày 2026-10-21, mở ngày 06/11/2026. Kiểm, 06/11/2026: stat cho uv.lock sửa lần cuối 19/10/2026, trước ngày đó; uv.lock dòng 4 tới 8 đã tách nhánh python_full_version >= '3.11', nên giới hạn là Python của venv chứ không phải lock. Thiếu mạng, mô phỏng bằng cờ offline và cache rỗng ngày 06/11/2026: "Failed to download `certifi==2026.10.18`" rồi "Network connectivity is disabled, but the requested data wasn't found in the cache".

\textbf{Bước 3: chạy unit test bằng \repo{make test}.} Lệnh này chạy pytest và bỏ hai nhóm test theo marker: test DAG, chạy ở bước 4, và test e2e, cần cả stack. Khối dưới tìm dòng tóm tắt cuối của pytest, dòng bắt đầu bằng dấu \code{=} và có thời gian chạy, rồi chỉ giữ các con số trên dòng đó. Khối không lấy dòng cuối cùng, vì khi có test lỗi, dòng cuối là thông báo lỗi của \code{make}.

% chạy: K3
\begin{codeblock}
(cd lab_m00/lakehouse && make test 2>&1 | grep -E '^=+ .* in [0-9.]+s' \
  | grep -oE '[0-9]+ (passed|failed|skipped|deselected|errors?)')
\end{codeblock}
\textit{Kết quả:}
% kết quả: K3
\begin{codeblock}
88 passed
5 deselected
\end{codeblock}

5 test bị bỏ là 4 test DAG và 1 test e2e. Trên máy kiểm, lệnh mất chưa tới 2 phút, gần hết cho 30 test chạy Spark local; lần đầu, Spark còn tải jar Iceberg runtime khoảng 45 MB.

Nếu kết quả có \code{30 skipped} hoặc \code{30 errors} và chỉ còn \code{58 passed}, máy bạn chưa có JDK chạy được. Khi PATH không có lệnh \code{java}, như trên Linux chưa cài JDK, 30 test Spark local tự bỏ qua với lý do ``cần Java để chạy Spark local'' (\repo{tests/conftest.py}, dòng 17 và 18).

macOS thì luôn có sẵn \code{/usr/bin/java}, kể cả khi chưa cài JDK, nên trên Mac các test này báo lỗi thay vì bỏ qua. Chạy \code{java -version} để kiểm, rồi cài JDK 11 như máy kiểm, hoặc JDK 8 hay 17. Nếu có \code{failed}, hoặc vẫn có \code{errors} khi JDK đã đúng bản, chạy \repo{make test} không qua \code{grep} để xem test nào hỏng.
% Kiểm: 06/11/2026: real 104,96 giây; 05/11/2026 với thư mục Ivy rỗng: real 109,21 giây; 06/11/2026 với thư mục Ivy rỗng: 111,65 giây. Jar iceberg-spark-runtime-3.5_2.12-1.8.1 nặng 45166693 byte theo ls -l trong thư mục Ivy, 06/11/2026, khớp content-length trên Maven Central; khai ở tests/conftest.py dòng 12, nạp ở dòng 27. 4 test DAG ở tests/airflow/test_dags.py; 1 test e2e ở tests/e2e/test_pipeline.py dòng 106.
% Kiểm: 06/11/2026, make test trên bản chép. PATH không có java: "58 passed, 30 skipped, 5 deselected, 1 warning in 5.74s". Một java giả thoát mã 1 đứng đầu PATH: "58 passed, 5 deselected, 1 warning, 30 errors in 6.66s", rồi dòng cuối "make: *** [test] Error 1"; K3 vẫn in 58 passed, 5 deselected, 30 errors.
% Kiểm: /usr/bin/java là file của macOS: ls -lO cho cờ restricted, mount cho / là volume sealed, read-only (Darwin 25.6.0, 06/11/2026); trên máy kiểm, which -a java chỉ ra /usr/bin/java, và nó chạy OpenJDK 11.0.30 của Homebrew.

\textbf{Bước 4: kiểm DAG bằng \repo{make dag-check}.} Script \repo{airflow/check\_\allowbreak{}dags.sh} dựng một venv riêng có Airflow 3.1.6 trên Python 3.11, rồi chạy 4 test parse, serialize và kiểm cấu hình của 3 DAG. Airflow chạy trong venv đó ngay trên máy bạn, không trong container, nên bước này không cần Docker. Venv nằm ngoài repo, mặc định ở \code{/tmp/lakehouse-airflow-check}.

% chạy: K4
\begin{codeblock}
(cd lab_m00/lakehouse && make dag-check 2>&1 | tail -n 1 \
  | grep -oE '[0-9]+ (passed|failed|skipped|errors?)')
\end{codeblock}
\textit{Kết quả:}
% kết quả: K4
\begin{codeblock}
4 passed
\end{codeblock}

Lần đầu mất 1 tới 2 phút vì uv phải giải và tải 130 gói, và venv đó chiếm khoảng 690 MB; các lần sau mất khoảng 3 giây. Script đọc file constraints của Airflow 3.1.6 từ GitHub ở mọi lần chạy, nên luôn cần mạng. Nếu không ghi được vào \code{/tmp}, đặt biến \code{AIRFLOW\_\allowbreak{}CHECK\_\allowbreak{}VENV} trỏ tới thư mục khác.
% Kiểm: airflow/check_dags.sh dòng 2 (3 DAG; parse, serialize, cấu hình), 6 (mặc định /tmp/lakehouse-airflow-check), 7 (constraints-3.1.6 trên raw.githubusercontent.com), 9 (Python 3.11), 10, 11 (apache-airflow==3.1.6). Lần đầu, 05/11/2026: "Resolved 130 packages in 1m 27s", "Installed 130 packages", real 103,50 giây; lần sau, 06/11/2026: real 2,74 giây. Lần đầu, 06/11/2026, cache rỗng: 64,20 giây; du -sk của venv: 674188 KiB, tức 690,4 MB. UV_OFFLINE=1 make dag-check với venv đã có, 06/11/2026: "error: Network connectivity is disabled, but a remote requirements file was requested: https://raw.githubusercontent.com/apache/airflow/constraints-3.1.6/constraints-3.11.txt".

\textbf{Bước 5: gọi CLI.} Lệnh dưới in phần trợ giúp của CLI \code{lakehouse} và chỉ giữ bốn dòng lệnh con. Lệnh này chỉ đọc code, không gửi gì tới Spark, nên chạy được khi Docker tắt.

% chạy: K5
\begin{codeblock}
(cd lab_m00/lakehouse && uv run lakehouse --help | grep -E '^ {4}[a-z]')
\end{codeblock}
\textit{Kết quả:}
% kết quả: K5
\begin{codeblock}
    bootstrap           tạo namespace và table (chạy lại an toàn)
    bronze              REST API -> iceberg.raw.*
    silver              raw -> staging -> semantic -> core
    maintenance         compaction, expire, xoá file mồ côi
\end{codeblock}

Bốn lệnh con là bốn bước của pipeline, và mỗi lệnh có một target cùng tên trong \repo{Makefile}. Nếu uv báo \code{Failed to spawn}, bạn đang đứng ngoài thư mục có \repo{pyproject.toml}: \code{uv run} chỉ thấy lệnh \code{lakehouse} trong project của nó.
% Kiểm: CLI help chạy được ngày 05/11/2026 lúc Docker tắt, real 0,65 giây. Gọi từ lab_m00/ ngày 06/11/2026: "error: Failed to spawn: `lakehouse`", "Caused by: No such file or directory (os error 2)".

\textbf{Bước 6: lệnh nào cần Docker.} Mở \repo{Makefile} của bản chép: 13 target nằm ở dòng 4 tới 44. Tự đánh dấu target nào cần Docker, rồi mới so với bảng dưới.

\begin{xltabular}{\textwidth}{@{}P{3.6cm} P{1.6cm} L@{}}
\toprule
\thd{target} & \thd{cần Docker} & \thd{vì sao} \\
\midrule
\endhead
\bottomrule
\endfoot
\repo{up}, \repo{down}, \repo{ps}, \repo{logs} & có & gọi thẳng \code{docker compose} \\
\repo{bootstrap}, \repo{bronze}, \repo{silver}, \repo{maintenance} & có & CLI chạy trên máy bạn nhưng gửi job tới Spark Connect ở \code{sc://localhost:15002}; Bronze còn gọi API nguồn ở \code{http://localhost:8000}. Cả hai là container của stack \\
\repo{e2e} & có & chạy cả pipeline trên stack đang chạy \\
\repo{data} & không & \code{curl} tải 9 file CSV từ GitHub, chỉ cần mạng \\
\repo{test}, \repo{dag-check}, \repo{lint} & không & pytest, venv Airflow riêng và Ruff đều chạy trên máy \\
\end{xltabular}
% Kiểm: Makefile dòng 4 tới 15 (docker compose), 17, 18 (data), 20 tới 30 (CLI), 33, 34 (test), 37, 38 (e2e), 40, 41 (dag-check), 43, 44 (lint). src/lakehouse/config.py dòng 22, 23 (SPARK_REMOTE mặc định sc://localhost:15002, OLIST_API_BASE_URL mặc định http://localhost:8000); pyproject.toml dòng 36; source_sim/download_olist.sh dòng 5, 27.

\begin{luuy}
\textbf{Bẫy thường gặp.}
\begin{itemize}
  \item \textbf{Nghĩ bản chép tách khỏi stack.} Bản chép chỉ tách code và venv. CLI của nó vẫn nối tới \code{localhost:15002}, nên \repo{make bronze} chạy trong \code{lab\_\allowbreak{}m00/\allowbreak{}lakehouse/} ghi vào đúng các bảng của stack đang chạy. Giống chạy migration từ một nhánh feature vào database dev dùng chung: code nằm ở nhánh riêng, dữ liệu thì không.
  \item \textbf{Chạy \code{uv sync} ngay trong \repo{lakehouse/}.} Mặc định \code{uv sync} xoá mọi gói không khai báo trong project, nên gói bạn tự cài thêm vào \repo{lakehouse/.venv} sẽ mất. Đó là lý do danh sách kiểm này chạy trên bản chép.
\end{itemize}
% Kiểm: uv help sync, uv 0.10.9, 06/11/2026: "By default, an exact sync is performed: uv removes packages that are not declared as dependencies of the project."
\end{luuy}

\hoanthanh{bước 3 in \code{88 passed} và \code{5 deselected}; bước 4 in \code{4 passed}; bước 5 liệt kê bốn lệnh con \code{bootstrap}, \code{bronze}, \code{silver}, \code{maintenance}; và ở bước 6 bạn tự đánh dấu đúng 9 trên 13 target cần Docker trước khi đọc bảng.}


% ====================== BẢN ĐỒ ======================
\plainchapter{Bản đồ lộ trình}

Thứ tự đi từ nền tảng tới vận hành: học cách hỏi dữ liệu và mô hình hoá trước, rồi hiểu dữ liệu được lưu thế nào, sau đó mới xây pipeline, cuối cùng là phục vụ, streaming và cloud. Dự án capstone chạy xuyên suốt, mỗi phần thêm một mốc.

\begin{codeblock}
M0  Trang khởi động cho SWE, mục "Bắt đầu từ nền SWE" ở cuối chương mở đầu
   │
PHẦN I    Nền tảng                                      85 giờ, tính cả M0
            M1  SQL cho người đã quen ORM
            M2  Vòng đời DE, yêu cầu và chọn kho
            M3  Mô hình dữ liệu phân tích: thiết kế Gold cho Olist
            ─► Mốc capstone 1: star schema Gold trên DuckDB, đối soát đúng
   │
PHẦN II   Lưu trữ và lakehouse                                      61 giờ
            M4  File, object storage và Python cho dữ liệu
            M5  Apache Iceberg qua lăng kính Git
            M6  Catalog, bảo mật và dữ liệu cá nhân
            ─► Mốc capstone 2: giải thích mọi file của raw.olist_orders,
               diễn tập xoá dữ liệu cá nhân trên bản sao
   │
PHẦN III  Xử lý và vận hành pipeline                                90 giờ
            M7  Spark cho người quen Java Streams, và dựng Gold
            M8  Ingestion từ REST API và áp thay đổi vào bảng
            M9  Airflow 3: viết và test DAG
            M10 Chất lượng, môi trường và vận hành dữ liệu
            ─► Mốc capstone 3: marts tăng dần chạy bằng Airflow,
               có cổng write-audit-publish và SLO độ tươi
   │
PHẦN IV   Phục vụ, streaming và cloud                               52 giờ
            M11 Phục vụ dữ liệu: Trino, quyền, BI và data API
            M12 Streaming: Kafka API và Flink
            M13 Cloud: mang lakehouse lên AWS
            ─► Mốc capstone 4: dashboard qua Trino, data API, lát cắt AWS
   │
PHẦN V    Capstone, chứng chỉ và nghề nghiệp
            nghề nghiệp: 1 giờ mỗi tuần, song song với module
            ôn chứng chỉ: 30 tới 40 giờ, sau phần lõi

NHÁNH TUỲ CHỌN, sau phần lõi                                        26 giờ
            T1  dbt trên DuckDB                                sau M10
            T2  CDC từ log: Debezium và Kafka Connect          sau M12
            T3  Dữ liệu cho AI trên review Olist          sau capstone
\end{codeblock}

\section{Các module}
\begin{xltabular}{\textwidth}{@{}l L P{2.6cm} r P{2.6cm}@{}}
\toprule
\thd{module} & \thd{tên} & \thd{cấp độ} & \thd{giờ} & \thd{tiên quyết} \\
\midrule
\endhead
\bottomrule
\endfoot
0 & Trang khởi động cho SWE, mục \textit{Bắt đầu từ nền SWE} ở cuối chương ``Về lộ trình này'' & cơ bản & 2 & không \\
1 & SQL cho người đã quen ORM; thêm tới 25 giờ luyện tập & cơ bản & 40 & 0 \\
2 & Vòng đời DE, yêu cầu và chọn kho & cơ bản & 13 & 1; học song song với 3 được \\
3 & Mô hình dữ liệu phân tích: thiết kế Gold cho Olist & cơ bản tới trung cấp & 30 & 1 \\
4 & File, object storage và Python cho dữ liệu & trung cấp & 20 & 1, 2 \\
5 & Apache Iceberg qua lăng kính Git & trung cấp tới nâng cao & 28 & 4 \\
6 & Catalog, bảo mật và dữ liệu cá nhân & trung cấp & 13 & 5 \\
7 & Spark cho người quen Java Streams, và dựng Gold & trung cấp & 28 & 1, 3, 5 \\
8 & Ingestion từ REST API và áp thay đổi vào bảng & trung cấp & 24 & 5, 7 \\
9 & Airflow 3: viết và test DAG & trung cấp & 16 & 8 \\
10 & Chất lượng, môi trường và vận hành dữ liệu & trung cấp & 22 & 7, 9 \\
11 & Phục vụ dữ liệu: Trino, quyền, BI và data API & trung cấp & 14 & 3, 6, 7 \\
12 & Streaming: Kafka API và Flink & nâng cao & 20 & 5, 7, 8 \\
13 & Cloud: mang lakehouse lên AWS & nâng cao & 18 & 5, 6, 9, 10 \\
\midrule
 & \textbf{tổng phần lõi} & & \textbf{288} & \\
\midrule
T1 & Nhánh dbt trên DuckDB, tuỳ chọn & trung cấp & 10 & 3; nên học sau 10 \\
T2 & Nhánh CDC từ log: Debezium và Kafka Connect, tuỳ chọn & nâng cao & 8 & 8, 12 \\
T3 & Nhánh dữ liệu cho AI trên review Olist, tuỳ chọn & trung cấp & 8 & 10; học sau capstone \\
\midrule
 & \textbf{tổng ba nhánh} & & \textbf{26} & \\
\end{xltabular}

\section{Lịch gợi ý theo tháng}
Lịch tính cho 10 giờ mỗi tuần: 9 giờ học module và 1 giờ cho chương nghề nghiệp. Mỗi tháng tính 4 tuần, tức 36 giờ module; tuần dư của tháng dài dùng làm đệm.

\begin{xltabular}{\textwidth}{@{}l P{3.8cm} r L@{}}
\toprule
\thd{tháng} & \thd{học} & \thd{giờ cộng dồn} & \thd{mốc capstone} \\
\midrule
\endhead
\bottomrule
\endfoot
1 & M0 và 34 giờ đầu của M1 & 36 & \\
2 & hết M1; M2; 17 giờ đầu của M3 & 72 & \\
3 & hết M3; M4; 3 giờ đầu của M5 & 108 & mốc 1 khi xong M3, ở giờ 85: bus matrix, grain và star schema Gold trên DuckDB cho năm câu hỏi; phương trình đối soát đúng \\
4 & hết M5; 11 giờ đầu của M6 & 144 & \\
5 & hết M6; M7; 6 giờ đầu của M8 & 180 & mốc 2 khi xong M6, ở giờ 146: giải thích được mọi file của \code{raw.olist\_\allowbreak{}orders} trong bucket; diễn tập xoá dữ liệu cá nhân trên bản sao, có bước khôi phục \\
6 & hết M8; M9; 2 giờ đầu của M10 & 216 & \\
7 & hết M10; M11; 2 giờ đầu của M12 & 252 & mốc 3 khi xong M10, ở giờ 236: \code{marts} tăng dần chạy bằng Airflow, có cổng write-audit-publish, SLO độ tươi có alert, và một lần khôi phục catalog đã thử \\
8 & hết M12; M13 & 288 & mốc 4 khi xong M13, ở giờ 288: dashboard Metabase qua Trino có quyền chỉ đọc, một data API, và một lát cắt AWS dựng rồi xoá bằng Terraform; luồng sự kiện vào Iceberg là phần tuỳ chọn \\
9 & T1, T2; 18 giờ đầu của phần ôn chứng chỉ & 324 & \\
10 & hết ôn chứng chỉ, tổng 30 tới 40 giờ; T3 nếu muốn & 344 tới 354 & \\
\end{xltabular}
\begin{itemize}
  \item Chỉ phần lõi: 288 giờ, 8 tháng. Lõi cộng ôn chứng chỉ: 318 tới 328 giờ, 8,8 tới 9,1 tháng. Lõi, ba nhánh và ôn chứng chỉ: 344 tới 354 giờ, 9,6 tới 9,8 tháng.
  \item Không tháng nào vượt 36 giờ module. Nếu Module 1 dùng hết 25 giờ thêm, mọi mốc sau lùi tối đa khoảng 3 tuần.
  \item Giờ của từng module là ước lượng; chưa ai học thử theo lịch này.
\end{itemize}

\begin{ghichu}
\textbf{Thói quen mỗi tuần.} Một buổi đọc docs, hai buổi làm lab, một buổi viết lại bằng lời của bạn những gì vừa học. Ghi từ mới vào một file thuật ngữ riêng. Cuối mỗi module, làm thử thách kỹ năng mà không mở tài liệu.
\end{ghichu}

% ====================== STACK THAM CHIẾU ======================
\plainchapter{Stack tham chiếu và dữ liệu}

Mọi lab chạy trên một stack mã nguồn mở dựng bằng Docker Compose trên laptop. Bạn không tốn tiền cloud cho tới Module 13.

\section{Thành phần}
Cột ``bản mới nhất'' là bản đã kiểm ngày 02/11/2026. Bạn không cần dùng bản mới nhất; điều quan trọng là ghim một phiên bản và đọc docs của đúng phiên bản đó.

\begin{xltabular}{\textwidth}{@{}P{2.7cm} L P{3.8cm}@{}}
\toprule
\thd{thành phần} & \thd{vai trò trong stack} & \thd{bản mới nhất, docs} \\
\midrule
\endhead
\bottomrule
\endfoot
RustFS & object storage tương thích S3, chứa mọi file dữ liệu và metadata; quickstart của Apache Polaris và bộ test S3 của PyIceberg đều dùng RustFS & \href{https://docs.rustfs.com/}{docs.rustfs.com} \\
Apache Iceberg & open table format: biến các file Parquet thành bảng có giao dịch, time travel, schema evolution & 1.12.0, \href{https://iceberg.apache.org/releases/}{releases}, \href{https://iceberg.apache.org/spec/}{spec} \\
Apache Polaris hoặc Lakekeeper & Iceberg REST catalog: giữ con trỏ tới metadata hiện hành của từng bảng, phân quyền, cấp credential & Polaris 1.8.0, \href{https://polaris.apache.org/downloads/}{downloads}; \href{https://docs.lakekeeper.io/}{Lakekeeper docs} \\
Apache Spark & engine xử lý batch, chạy MERGE và các procedure bảo trì Iceberg & 4.3.0; nhánh 3.5 mới nhất là 3.5.9, \href{https://spark.apache.org/news/index.html}{Spark news} \\
PyIceberg & thư viện Python đọc, ghi bảng Iceberg không cần JVM & 0.12.0, \href{https://py.iceberg.apache.org/}{py.iceberg.apache.org} \\
DuckDB & SQL phân tích trên laptop, đọc Parquet và Iceberg & 2.0.0 LTS, \href{https://duckdb.org/2026/10/21/announcing-duckdb-200}{thông báo 2.0} \\
PostgreSQL & học SQL và hành vi OLTP: ràng buộc, transaction, index & 18, \href{https://www.postgresql.org/docs/current/}{docs current} \\
Apache Airflow & điều phối pipeline & 3.3.2, \href{https://airflow.apache.org/docs/apache-airflow/stable/release_notes.html}{release notes} \\
Trino & engine truy vấn tương tác, phục vụ BI & 484, \href{https://trino.io/docs/current/release.html}{release notes} \\
Metabase & dashboard trên dữ liệu tầng marts & \href{https://www.metabase.com/docs/latest/}{Metabase docs} \\
Apache Kafka & log sự kiện cho nhánh streaming & 4.4.0, \href{https://kafka.apache.org/community/downloads/}{downloads} \\
Apache Flink & xử lý stream có trạng thái, ghi vào Iceberg & 2.3, \href{https://nightlies.apache.org/flink/flink-docs-stable/}{docs stable} \\
Apache Arrow (pyarrow) & đọc, ghi Parquet trong Python, định dạng bộ nhớ dùng chung & 26.0.0, \href{https://arrow.apache.org/release/}{releases} \\
\end{xltabular}

\begin{luuy}
\textbf{Cổng hay bị trùng.} REST catalog thường dùng 8181; object storage dùng 9000 cho API và 9001 cho console; Spark UI dùng 4040; Spark Connect dùng 15002; Metabase, Grafana và giao diện Marquez đều mặc định dùng 3000. Khi gộp nhiều quickstart vào một máy, hãy đổi cổng ngay trong file compose.
\end{luuy}

\section{Dữ liệu: Olist}
Mọi lab và dự án capstone dùng \href{https://www.kaggle.com/datasets/olistbr/brazilian-ecommerce}{Brazilian E-Commerce Public Dataset by Olist} trên Kaggle: khoảng 100 nghìn đơn hàng ẩn danh từ 2016 tới 2018 của một sàn thương mại điện tử Brazil, chia thành nhiều file CSV liên kết với nhau.

\begin{xltabular}{\textwidth}{@{}P{4.6cm} L@{}}
\toprule
\thd{bảng} & \thd{dùng để học} \\
\midrule
\endhead
\bottomrule
\endfoot
\code{orders} & trạng thái đơn và các mốc thời gian; nguồn cho tải tăng dần theo thời điểm mua \\
\code{order\_items} & mỗi dòng một sản phẩm trong đơn; grain của bảng fact chính \\
\code{order\_payments} & nhiều dòng thanh toán cho một đơn; khoá gồm hai cột \\
\code{order\_reviews} & đánh giá; có bản trùng, hợp để luyện bỏ trùng \\
\code{customers} & phân biệt \code{customer\_id} theo đơn và \code{customer\_unique\_id} theo người \\
\code{sellers}, \code{products} & bảng chiều; luyện SCD và chuẩn hoá \\
\code{geolocation} & toạ độ theo mã bưu chính; nhiều dòng trùng, hợp để luyện chất lượng dữ liệu \\
\end{xltabular}
Đọc giấy phép trên trang dataset trước khi công khai lại dữ liệu hoặc dùng cho mục đích thương mại.

\section{Tổ chức tầng dữ liệu}
Lộ trình dùng kiến trúc nhiều tầng theo kiểu medallion, một quy ước thiết kế chứ không phải tính năng của phần mềm nào (\href{https://docs.databricks.com/aws/en/lakehouse/medallion}{Databricks: medallion architecture}).

\begin{codeblock}
nguồn CSV / sự kiện Kafka
      │  ingestion (Module 8, 12)
      ▼
raw       (bronze)  dữ liệu nguyên trạng, chỉ thêm, có cột metadata nạp
      │  ép kiểu, bỏ trùng, MERGE theo khoá (Module 7, 8)
      ▼
staging / core (silver)  sạch, đúng grain, có test (Module 10)
      │  mô hình hoá, tổng hợp (Module 3)
      ▼
marts     (gold)   star schema và bảng tổng hợp cho BI (Module 11)

Xuyên suốt: Parquet trên RustFS · bảng Iceberg · REST catalog · Airflow điều phối
\end{codeblock}


\partpage{I}{Nền tảng}{Hỏi dữ liệu bằng SQL, hiểu vòng đời và kiến trúc dữ liệu, thiết kế mô hình cho phân tích. Kết thúc bằng mốc capstone 1: bus matrix, grain và star schema Gold trên DuckDB cho năm câu hỏi, có phương trình đối soát đúng.}
% ---------------- m01_03 ----------------
% Module 1, 2, 3 của "Lộ trình Data Engineer: xây lakehouse mã nguồn mở"
% =====================================================================
\modulechapter{1}{SQL và mô hình quan hệ}
\moduleinfo{Cơ bản}{35 giờ}{không}

Mọi tầng của một lakehouse, từ staging tới mart, đều được viết và kiểm bằng SQL. Module này dạy bạn đọc, viết và giải thích câu SQL nhiều bảng, thiết kế bảng có khoá và ràng buộc đúng, và hiểu engine chạy câu lệnh ra sao. Bạn học hành vi OLTP trên PostgreSQL 18, rồi SQL phân tích trên DuckDB. Học xong, bạn viết được truy vấn có \code{JOIN}, \code{GROUP BY}, CTE và window function trên dữ liệu Olist, và biết ràng buộc nào engine thực sự kiểm, kể cả trong lakehouse.

\section{Bạn sẽ học được}

\begin{itemize}
  \item Viết câu SQL nhiều bảng có \code{JOIN}, \code{GROUP BY}, \code{HAVING}, CTE và window function, rồi giải thích từng dòng kết quả.
  \item Phân loại câu lệnh thành DDL, DML, DCL, TCL và nhận ra cách phân loại khác nhau giữa các hệ quản trị.
  \item Thiết kế bảng với \code{PRIMARY KEY}, \code{FOREIGN KEY}, \code{UNIQUE}, \code{NOT NULL}, \code{CHECK}; chọn khoá tự nhiên hay khoá thay thế.
  \item Xác định grain và khoá logic của một bảng, rồi kiểm nó bằng truy vấn.
  \item Giải thích NULL và logic ba trị, frame mặc định của window function, và các mức isolation của PostgreSQL.
  \item Đọc kế hoạch truy vấn bằng \code{EXPLAIN} và \code{EXPLAIN ANALYZE}.
  \item Phân biệt OLTP với OLAP, và giải thích vì sao lakehouse giữ khoá bằng pipeline chứ không bằng ràng buộc.
\end{itemize}

\section{Nội dung chính}

\subsection{Truy vấn}

\begin{enumerate}
  \item \textbf{Mô hình quan hệ.} Dữ liệu nằm trong bảng (relation); mỗi dòng (row, tuple) có cùng tập cột (column, attribute); tập bảng gom trong một database và schema. Đọc \href{https://www.postgresql.org/docs/current/tutorial-concepts.html}{PostgreSQL 18: Concepts}.
  \item \textbf{Nhóm câu lệnh.} DDL (Data Definition Language) tạo và đổi cấu trúc: \code{CREATE}, \code{ALTER}, \code{DROP}. DML (Data Manipulation Language) đọc và đổi dữ liệu: \code{SELECT}, \code{INSERT}, \code{UPDATE}, \code{DELETE}, \code{MERGE}. DCL (Data Control Language) cấp quyền: \code{GRANT}, \code{REVOKE}. TCL (Transaction Control Language) điều khiển transaction: \code{BEGIN}, \code{COMMIT}, \code{ROLLBACK}, \code{SAVEPOINT}. Cách chia không thống nhất giữa các hệ quản trị: tài liệu YugabyteDB YSQL (tương thích PostgreSQL) chia DDL, DML, DCL, TCL (\href{https://docs.yugabyte.com/preview/api/ysql/the-sql-language/}{YSQL: The SQL language}), còn Oracle xếp \code{GRANT}, \code{REVOKE} vào DDL, có nhóm Transaction Control riêng, và tự commit transaction trước và sau mỗi câu DDL (\href{https://docs.oracle.com/en/database/oracle/oracle-database/23/sqlrf/Types-of-SQL-Statements.html}{Oracle 23: Types of SQL Statements}). Danh sách đầy đủ câu lệnh của PostgreSQL nằm ở \href{https://www.postgresql.org/docs/current/sql-commands.html}{PostgreSQL 18: SQL Commands}.
  \item \textbf{\code{SELECT} và thứ tự xử lý.} Về logic, PostgreSQL tính \code{WITH}, rồi \code{FROM}, \code{WHERE}, \code{GROUP BY} và \code{HAVING}, rồi danh sách \code{SELECT}, \code{DISTINCT}, phép tập hợp, \code{ORDER BY}, cuối cùng \code{LIMIT}. Vì vậy alias đặt trong \code{SELECT} không dùng được trong \code{WHERE}. Đọc mục Description của \href{https://www.postgresql.org/docs/current/sql-select.html}{PostgreSQL 18: SELECT}.
  \item \textbf{NULL và logic ba trị.} So sánh thường với NULL cho ra NULL, nên \code{WHERE col = NULL} không trả dòng nào. Dùng \code{IS NULL}; khi cần coi hai NULL là bằng nhau, dùng \code{IS NOT DISTINCT FROM} (\href{https://www.postgresql.org/docs/current/functions-comparison.html}{PostgreSQL 18, mục 9.2}). Spark SQL có toán tử tương đương \code{<=>}.
  \item \textbf{\code{JOIN} các loại.} \code{INNER}, \code{LEFT}, \code{RIGHT}, \code{FULL}, \code{CROSS}. Semi join viết bằng \code{EXISTS}, anti join bằng \code{NOT EXISTS}. Cẩn thận \code{NOT IN} với subquery có NULL: kết quả có thể là NULL thay vì true (\href{https://www.postgresql.org/docs/current/functions-subquery.html}{PostgreSQL 18, mục 9.24}).
  \item \textbf{\code{GROUP BY}, \code{HAVING}, \code{FILTER}.} \code{WHERE} lọc dòng trước khi gộp; \code{HAVING} lọc nhóm sau khi gộp. Mệnh đề \code{FILTER (WHERE ...)} cho một aggregate chỉ nhận một phần dòng (\href{https://www.postgresql.org/docs/current/sql-expressions.html}{PostgreSQL 18, mục 4.2.7}).
  \item \textbf{Subquery và CTE.} \code{WITH} đặt tên cho truy vấn con; \code{WITH RECURSIVE} cho truy vấn tự tham chiếu. PostgreSQL gộp một CTE không đệ quy, không có side effect, được tham chiếu một lần vào câu cha; \code{MATERIALIZED} và \code{NOT MATERIALIZED} ghi đè quyết định đó (\href{https://www.postgresql.org/docs/current/queries-with.html}{PostgreSQL 18, mục 7.8}).
  \item \textbf{Window function.} \code{row\_\allowbreak{}number()}, \code{rank()}, \code{dense\_\allowbreak{}rank()}, \code{lag()}, \code{lead()}, aggregate chạy trên cửa sổ. Frame có ba chế độ \code{ROWS}, \code{RANGE}, \code{GROUPS} và mệnh đề \code{EXCLUDE}. Frame mặc định là \code{RANGE BETWEEN UNBOUNDED PRECEDING AND CURRENT ROW}: khi có \code{ORDER BY}, frame chạy từ đầu partition tới dòng ngang hàng (peer) cuối cùng của dòng hiện tại; không có \code{ORDER BY} thì mọi dòng là peer nên frame phủ cả partition (\href{https://www.postgresql.org/docs/current/sql-expressions.html}{PostgreSQL 18, mục 4.2.8}; \href{https://www.postgresql.org/docs/current/tutorial-window.html}{tutorial 3.5}). DuckDB có thêm \code{QUALIFY} để lọc theo kết quả window function, giống \code{HAVING} lọc theo aggregate (\href{https://duckdb.org/docs/2.0/sql/query_syntax/qualify}{DuckDB 2.0: QUALIFY}; \href{https://duckdb.org/docs/2.0/sql/functions/window_functions}{DuckDB 2.0: Window Functions}).
\end{enumerate}

\subsection{Thiết kế bảng}

\begin{enumerate}[resume]
  \item \textbf{Khoá và ràng buộc.} \code{CHECK} nhận dòng khi biểu thức là true hoặc NULL. \code{NOT NULL} tương đương \code{CHECK (col IS NOT NULL)} nhưng hiệu quả hơn. \code{PRIMARY KEY} đòi giá trị vừa duy nhất vừa khác NULL, và tự tạo một unique B-tree index. Với \code{UNIQUE}, mặc định hai NULL không bằng nhau, nên nhiều dòng có NULL vẫn lọt qua; thêm \code{NULLS NOT DISTINCT} để đổi điều đó. \code{FOREIGN KEY} có các hành động \code{NO ACTION} (mặc định), \code{RESTRICT}, \code{CASCADE}, \code{SET NULL}, \code{SET DEFAULT}, và không tự tạo index trên cột tham chiếu (\href{https://www.postgresql.org/docs/current/ddl-constraints.html}{PostgreSQL 18, mục 5.5}).
  \item \textbf{Khoá tự nhiên và khoá thay thế.} Khoá tự nhiên (natural key) có nghĩa nghiệp vụ và đến từ hệ thống nguồn, ví dụ \code{seller\_\allowbreak{}id} của Olist. Khoá thay thế (surrogate key) là định danh bạn thêm vào, theo định nghĩa không có trong dữ liệu nguồn (\href{https://learn.microsoft.com/en-us/power-bi/guidance/star-schema}{Microsoft Learn: star schema}). Nếu bảng chỉ có khoá thay thế mà khoá tự nhiên không có \code{UNIQUE}, không gì ngăn hai dòng mô tả cùng một thực thể; vì vậy thường giữ cả hai.
  \item \textbf{Grain và khoá logic.} Grain là câu trả lời cho “một dòng của bảng này là gì”, ví dụ “một dòng là một lần thanh toán của một đơn”. Trong cách nói thường ngày của data engineering, khoá logic là tập cột làm mỗi dòng duy nhất ở grain đó, dù engine có kiểm hay không. Với bảng payments của Olist, giả thuyết khoá logic là (\code{order\_\allowbreak{}id}, \code{payment\_\allowbreak{}sequential}); bạn kiểm giả thuyết đó bằng truy vấn trong Lab 1.2. Thuật ngữ này không có định nghĩa chuẩn trong docs của PostgreSQL; hãy hiểu nó như một quy ước làm việc.
  \item \textbf{Chuẩn hoá.} Mục tiêu là bỏ dư thừa và phụ thuộc không nhất quán. 1NF: bỏ nhóm lặp, mỗi tập dữ liệu liên quan một bảng có khoá chính. 2NF: giá trị dùng cho nhiều dòng tách thành bảng riêng, nối bằng khoá ngoại. 3NF: bỏ cột không phụ thuộc vào khoá. Trang của Microsoft coi 3NF là mức cao nhất cần cho đa số ứng dụng, và nhắc rằng nhiều bảng nhỏ có thể làm chậm hệ thống (\href{https://learn.microsoft.com/en-us/troubleshoot/office/access/database-normalization-description}{Microsoft Learn: Database normalization description}). Các dạng chuẩn cao hơn (BCNF, 4NF, 5NF) có trong bài giảng CMU 15-445 ở mục Khoá học miễn phí.
  \item \textbf{OLTP và OLAP.} OLTP phục vụ giao dịch thời gian thực; OLAP thường chạy theo lô, cho truy vấn, báo cáo, phân tích trên dữ liệu lịch sử; tách xử lý phân tích khỏi OLTP để không ảnh hưởng tải giao dịch (\href{https://docs.aws.amazon.com/whitepapers/latest/data-warehousing-on-aws/data-processing.html}{AWS: Data processing}). Hệ OLTP như PostgreSQL lưu theo dòng; dữ liệu phân tích thường lưu theo cột, ví dụ Apache Parquet là định dạng file hướng cột (\href{https://parquet.apache.org/docs/overview/}{Parquet: Overview}).
\end{enumerate}

\subsection{Engine chạy câu lệnh thế nào}

\begin{enumerate}[resume]
  \item \textbf{Transaction và isolation.} PostgreSQL nhận cả bốn mức isolation chuẩn nhưng chỉ cài ba mức khác nhau: Read Uncommitted chạy như Read Committed. Read Committed là mặc định và cho phép nonrepeatable read; Repeatable Read chặn nonrepeatable read, và trong PostgreSQL chặn cả phantom read; chỉ Serializable chặn serialization anomaly (\href{https://www.postgresql.org/docs/current/transaction-iso.html}{PostgreSQL 18, mục 13.2}).
  \item \textbf{Index và \code{EXPLAIN}.} Index tăng tốc đọc nhưng thêm chi phí cho cả hệ thống, nên dùng có chọn lọc (\href{https://www.postgresql.org/docs/current/indexes.html}{PostgreSQL 18, chương 11}). \code{EXPLAIN} in kế hoạch với số dòng ước lượng; \code{EXPLAIN ANALYZE} chạy thật câu lệnh rồi in thêm số dòng và thời gian thật. Ở bản 18, \code{ANALYZE} tự bật \code{BUFFERS}. Với câu sửa dữ liệu, bọc trong \code{BEGIN} … \code{ROLLBACK} (\href{https://www.postgresql.org/docs/current/using-explain.html}{PostgreSQL 18, mục 14.1}).
\end{enumerate}

\subsection{Khoá trong lakehouse}

\begin{enumerate}[resume]
  \item \textbf{Ràng buộc chỉ để thông tin.} Trên Databricks, \code{PRIMARY KEY}, \code{FOREIGN KEY} và \code{UNIQUE} là ràng buộc informational, không được kiểm; \code{NOT NULL} và \code{CHECK} được kiểm, và câu ghi vi phạm sẽ thất bại. Unique constraint đang ở Public Preview từ Databricks Runtime 18.2 (\href{https://docs.databricks.com/aws/en/tables/constraints}{Databricks: Constraints}).
  \item \textbf{Identifier field của Iceberg.} Schema Iceberg có thể khai báo \code{identifier-\allowbreak{}field-\allowbreak{}ids}: tập cột nguyên thuỷ xác định một dòng. Engine dùng chúng làm khoá cho upsert và equality delete, nhưng Iceberg không bảo đảm hay đòi hỏi tính duy nhất; engine hoặc bên ghi dữ liệu phải tự giữ (\href{https://iceberg.apache.org/spec/\#identifier-field-ids}{Iceberg spec: Identifier Field IDs}). Trong Spark, bạn khai báo bằng \code{ALTER TABLE ... SET IDENTIFIER FIELDS} (\href{https://iceberg.apache.org/docs/latest/spark-ddl/\#alter-table-set-identifier-fields}{Iceberg: Spark DDL}). Spec hiện đã có phần cho format v4, trong đó không còn cho ghi equality delete mới; trước khi dùng v4, kiểm engine của bạn đã hỗ trợ chưa.
  \item \textbf{Hệ quả.} Trong lakehouse, khoá duy nhất là việc của pipeline: khử trùng theo khoá logic trước khi ghi (thường bằng \code{row\_\allowbreak{}number()} rồi lọc dòng số 1), ghi bằng \code{MERGE} theo khoá logic, và chạy test đếm trùng sau mỗi lần nạp.
\end{enumerate}

\section{Bắt đầu}

\begin{itemize}
  \item \href{https://www.postgresql.org/docs/current/tutorial.html}{PostgreSQL 18: Tutorial}, chương 2 “The SQL Language”: tạo bảng, nạp dữ liệu, truy vấn, join, aggregate.
  \item \href{https://duckdb.org/docs/2.0/sql/introduction}{DuckDB 2.0: SQL Introduction}: chạy SQL đầu tiên trên DuckDB.
  \item \href{https://sqlbolt.com/}{SQLBolt}, bài 1 tới 12: luyện \code{SELECT}, \code{WHERE}, \code{JOIN} ngay trên trình duyệt.
\end{itemize}

\section{Tài liệu chính thức}

\begin{xltabular}{\textwidth}{@{}P{4.2cm} L P{1.8cm}@{}}
\toprule
\thd{tài liệu} & \thd{nội dung cần đọc} & \thd{thời lượng} \\
\midrule
\endhead
\bottomrule
\endfoot
\href{https://www.postgresql.org/docs/current/tutorial.html}{PostgreSQL 18: Tutorial} & chương 2; mục 3.4 Transactions, 3.5 Window Functions & 4 giờ \\
\href{https://www.postgresql.org/docs/current/sql-select.html}{PostgreSQL 18: SELECT} & mục Description: thứ tự xử lý; mệnh đề \code{WINDOW}, \code{DISTINCT ON} & 1 giờ \\
\href{https://www.postgresql.org/docs/current/functions-comparison.html}{PostgreSQL 18: 9.2 Comparison} & \code{IS NULL}, \code{IS [NOT] DISTINCT FROM} & 0,5 giờ \\
\href{https://www.postgresql.org/docs/current/functions-subquery.html}{PostgreSQL 18: 9.24 Subquery Expressions} & \code{EXISTS}, \code{IN}, \code{NOT IN} và NULL & 0,5 giờ \\
\href{https://www.postgresql.org/docs/current/queries-with.html}{PostgreSQL 18: 7.8 WITH Queries} & CTE, \code{WITH RECURSIVE}, \code{MATERIALIZED}, CTE sửa dữ liệu & 1,5 giờ \\
\href{https://www.postgresql.org/docs/current/sql-expressions.html}{PostgreSQL 18: 4.2 Value Expressions} & 4.2.7 Aggregate Expressions (\code{FILTER}); 4.2.8 Window Function Calls (frame) & 1,5 giờ \\
\href{https://www.postgresql.org/docs/current/functions-window.html}{PostgreSQL 18: 9.22 Window Functions} & danh sách hàm, ghi chú về \code{last\_\allowbreak{}value} và frame & 0,5 giờ \\
\href{https://www.postgresql.org/docs/current/ddl-constraints.html}{PostgreSQL 18: 5.5 Constraints} & 5.5.1 tới 5.5.5, gồm \code{NULLS NOT DISTINCT} & 1,5 giờ \\
\href{https://www.postgresql.org/docs/current/transaction-iso.html}{PostgreSQL 18: 13.2 Transaction Isolation} & bảng 13.1; Read Committed và Repeatable Read & 1,5 giờ \\
\href{https://www.postgresql.org/docs/current/indexes.html}{PostgreSQL 18: chương 11 Indexes} & 11.1 tới 11.3, 11.6 Unique Indexes & 1 giờ \\
\href{https://www.postgresql.org/docs/current/using-explain.html}{PostgreSQL 18: 14.1 Using EXPLAIN} & 14.1.1 Basics, 14.1.2 \code{EXPLAIN ANALYZE} & 1,5 giờ \\
\href{https://duckdb.org/docs/2.0/sql/functions/window_functions}{DuckDB 2.0: Window Functions} & framing, \code{GROUPS}, \code{EXCLUDE} & 1 giờ \\
\href{https://duckdb.org/docs/2.0/sql/query_syntax/qualify}{DuckDB 2.0: QUALIFY} & lọc top N mỗi nhóm không cần subquery & 0,3 giờ \\
\href{https://docs.oracle.com/en/database/oracle/oracle-database/23/sqlrf/Types-of-SQL-Statements.html}{Oracle 23: Types of SQL Statements} & cách một hệ quản trị chia nhóm câu lệnh & 0,3 giờ \\
\href{https://learn.microsoft.com/en-us/troubleshoot/office/access/database-normalization-description}{Microsoft Learn: Database normalization} & 1NF, 2NF, 3NF và ví dụ & 0,5 giờ \\
\href{https://docs.aws.amazon.com/whitepapers/latest/data-warehousing-on-aws/data-processing.html}{AWS: Data processing} & OLTP và OLAP & 0,3 giờ \\
\href{https://docs.databricks.com/aws/en/tables/constraints}{Databricks: Constraints} & ràng buộc được kiểm và ràng buộc informational & 0,3 giờ \\
\href{https://iceberg.apache.org/spec/\#identifier-field-ids}{Iceberg spec: Identifier Field IDs} & identifier field và tính duy nhất & 0,3 giờ \\
\end{xltabular}

\section{Khoá học miễn phí}

\begin{itemize}
  \item \href{https://15445.courses.cs.cmu.edu/fall2025/}{CMU 15-445/645 Database Systems, Fall 2025}: trang khoá học có lịch, slide, ghi chú và link playlist YouTube của bài giảng. Nên xem bài \#01 Relational Model \& Algebra, \#02 Modern SQL, \#06 Storage Models \& Compression, \#08 Indexes \& Filters I, \#17 Concurrency Control Theory. Bản mới nhất là \href{https://15445.courses.cs.cmu.edu/fall2026/}{Fall 2026}; lịch hai kỳ đánh số bài khác nhau, đừng trộn. Danh sách mọi kỳ ở \href{https://db.cs.cmu.edu/courses/}{CMU Database Group: Courses}.
  \item \href{https://sqlbolt.com/}{SQLBolt}: bài tập SQL tương tác, miễn phí.
  \item \href{https://pgexercises.com/}{PostgreSQL Exercises}: bài tập trên PostgreSQL theo nhóm Basic, Joins and Subqueries, Aggregates, Recursive; miễn phí.
\end{itemize}

\section{Lab thực hành}

Mọi lab dùng bộ \href{https://www.kaggle.com/datasets/olistbr/brazilian-ecommerce}{Brazilian E-Commerce Public Dataset by Olist}. Tải về, giải nén vào thư mục \code{olist/}. Ghim phiên bản: PostgreSQL 18, DuckDB 2.0.0.

\subsection{Lab 1.1: Bảng có khoá và ràng buộc trên PostgreSQL 18}

\begin{enumerate}
  \item Chạy PostgreSQL và mở \code{psql}:
\begin{codeblock}
docker run -d --name pg18 -e POSTGRES_PASSWORD=postgres \
  -v "$PWD/olist:/data:ro" -p 5432:5432 postgres:18
docker exec -it pg18 psql -U postgres
\end{codeblock}
  \item Tạo bảng \code{customers}, \code{sellers}, \code{products}, \code{orders}, \code{order\_\allowbreak{}items}, \code{order\_\allowbreak{}payments} theo cột của file CSV. Đặt \code{PRIMARY KEY} theo khoá logic bạn đoán, \code{FOREIGN KEY} từ \code{orders} tới \code{customers}, từ \code{order\_\allowbreak{}items} tới \code{orders}, \code{products}, \code{sellers}, và \code{CHECK (price >= 0)}.
\begin{codeblock}
CREATE TABLE order_items (
  order_id      text NOT NULL REFERENCES orders,
  order_item_id int  NOT NULL,
  product_id    text NOT NULL REFERENCES products,
  seller_id     text NOT NULL REFERENCES sellers,
  shipping_limit_date timestamp,
  price         numeric(12,2) CHECK (price >= 0),
  freight_value numeric(12,2),
  PRIMARY KEY (order_id, order_item_id)
);
\copy order_items FROM '/data/olist_order_items_dataset.csv' WITH (FORMAT csv, HEADER true)
\end{codeblock}
  \item Nạp các bảng cha trước, bảng con sau. Nếu một lệnh \code{\textbackslash{}copy} thất bại vì trùng khoá, đọc thông báo lỗi và sửa giả thuyết khoá.
  \item Thử bốn câu sai và đọc lỗi: chèn \code{order\_\allowbreak{}items} có \code{price} âm; chèn đơn với \code{customer\_\allowbreak{}id} không tồn tại; chèn trùng khoá chính; xoá một \code{seller} đang được tham chiếu.
  \item Tạo bảng nhỏ có cột \code{UNIQUE} và cột \code{UNIQUE NULLS NOT DISTINCT}, chèn hai dòng NULL vào mỗi cột, so kết quả.
  \item Tạo index trên \code{order\_\allowbreak{}items(seller\_\allowbreak{}id)} vì khoá ngoại không tự có index.
\end{enumerate}

\hoanthanh{sáu bảng nạp đủ dòng của sáu file CSV với PK và FK bật; bạn có ảnh chụp bốn thông báo lỗi ràng buộc và giải thích được vì sao cột \code{UNIQUE} thường nhận hai NULL còn cột \code{NULLS NOT DISTINCT} thì không.}

\subsection{Lab 1.2: SQL phân tích trên DuckDB}

\begin{enumerate}
  \item Mở DuckDB và tạo bảng từ CSV:
\begin{codeblock}
CREATE TABLE orders   AS SELECT * FROM read_csv('olist/olist_orders_dataset.csv');
CREATE TABLE items    AS SELECT * FROM read_csv('olist/olist_order_items_dataset.csv');
CREATE TABLE payments AS SELECT * FROM read_csv('olist/olist_order_payments_dataset.csv');
CREATE TABLE products AS SELECT * FROM read_csv('olist/olist_products_dataset.csv');
CREATE TABLE customers AS SELECT * FROM read_csv('olist/olist_customers_dataset.csv');
\end{codeblock}
  \item Kiểm khoá logic: với mỗi bảng, so \code{count(*)} với số tổ hợp khác nhau của các cột bạn đoán là khoá.
\begin{codeblock}
SELECT count(*) AS n_rows,
       count(DISTINCT (order_id, payment_sequential)) AS n_keys
FROM payments;
\end{codeblock}
  \item So \code{count(DISTINCT customer\_\allowbreak{}id)} với \code{count(DISTINCT customer\_\allowbreak{}unique\_\allowbreak{}id)} trong \code{customers}. Ghi lại cột nào định danh người mua.
  \item Doanh thu theo tháng và theo bang: join \code{items}, \code{orders}, \code{customers}, \code{GROUP BY}, rồi dùng \code{HAVING} giữ tháng có trên 100 đơn.
  \item Top 3 danh mục mỗi bang theo doanh thu bằng \code{QUALIFY rank() OVER (...) <= 3}; viết lại cùng câu trên PostgreSQL bằng subquery.
  \item Doanh thu cộng dồn 7 ngày bằng \code{ROWS BETWEEN 6 PRECEDING AND CURRENT ROW}; tăng trưởng so với tháng trước bằng \code{lag()}.
  \item Chạy \code{last\_\allowbreak{}value(price) OVER (PARTITION BY order\_\allowbreak{}id ORDER BY order\_\allowbreak{}item\_\allowbreak{}id)} rồi thêm \code{ROWS BETWEEN UNBOUNDED PRECEDING AND UNBOUNDED FOLLOWING}; giải thích khác biệt bằng frame mặc định.
  \item Đoán rồi chạy: \code{NULL = NULL}, \code{NULL IS NOT DISTINCT FROM NULL}, \code{count(*)} so với \code{count(order\_\allowbreak{}delivered\_\allowbreak{}customer\_\allowbreak{}date)}, một \code{NOT IN} với subquery có NULL so với \code{NOT EXISTS}.
\end{enumerate}

\hoanthanh{bạn có một file SQL chạy lại được, trả đúng khoá logic của \code{payments} và \code{items} (số dòng bằng số khoá), câu top 3 cho cùng kết quả trên DuckDB và PostgreSQL, và giải thích bằng frame vì sao \code{last\_\allowbreak{}value} đổi kết quả.}

\subsection{Lab 1.3: Transaction, isolation và EXPLAIN}

\begin{enumerate}
  \item Mở hai phiên \code{psql} tới container của Lab 1.1. Phiên A: \code{BEGIN;} rồi đọc \code{price} của một món hàng. Phiên B: \code{UPDATE} món đó và \code{COMMIT}. Phiên A đọc lại: ở Read Committed, giá trị đổi (nonrepeatable read).
  \item Lặp lại với phiên A mở bằng \code{BEGIN ISOLATION LEVEL REPEATABLE READ;}: lần đọc thứ hai vẫn thấy giá trị cũ.
  \item Thử \code{BEGIN ISOLATION LEVEL READ UNCOMMITTED;} và chứng minh phiên A không đọc được thay đổi chưa commit của phiên B.
  \item Tạo bảng lớn bằng \code{generate\_\allowbreak{}series} (1 triệu dòng), chạy \code{EXPLAIN ANALYZE} một câu lọc theo một cột trước và sau \code{CREATE INDEX}; so loại node quét, số dòng ước lượng và số dòng thật.
  \item Chạy \code{EXPLAIN ANALYZE} một câu \code{UPDATE} bên trong \code{BEGIN} … \code{ROLLBACK} và kiểm dữ liệu không đổi.
\end{enumerate}

\hoanthanh{bạn có bản ghi hai phiên cho thấy nonrepeatable read ở Read Committed và không xảy ra ở Repeatable Read, cùng hai kế hoạch truy vấn trước và sau khi tạo index, mỗi kế hoạch có chú thích loại scan và chênh lệch giữa số dòng ước lượng và số dòng thật.}

\section{Kế hoạch theo tuần}

\begin{xltabular}{\textwidth}{@{}P{1.4cm} L P{1.6cm}@{}}
\toprule
\thd{tuần} & \thd{việc} & \thd{giờ} \\
\midrule
\endhead
\bottomrule
\endfoot
1 & PostgreSQL Tutorial chương 2; SQLBolt; CMU bài \#01, \#02; mục \code{SELECT}, 9.2, 9.24 & 10 \\
2 & mục 7.8, 4.2.7, 4.2.8, 3.5; DuckDB Window Functions, QUALIFY; Lab 1.2; pgexercises nhóm Joins, Aggregates & 10 \\
3 & mục 5.5; chuẩn hoá; khoá trong lakehouse; Lab 1.1 & 10 \\
4 & mục 13.2, 14.1, chương 11; CMU bài \#06, \#08; Lab 1.3; thử thách kỹ năng & 5 \\
thêm & CMU bài \#17; pgexercises nhóm Recursive; homework 1 của CMU Fall 2026 & 6 tới 10 \\
\end{xltabular}

\begin{thuthach}
Không có hướng dẫn từng bước. Trên dữ liệu Olist, viết một file SQL trả lời năm câu hỏi: (1) tỉ lệ đơn giao trễ so với ngày hẹn theo từng tháng; (2) mười người bán có doanh thu cao nhất và tỉ lệ giao trễ của họ; (3) phần trăm người mua (theo \code{customer\_\allowbreak{}unique\_\allowbreak{}id}) có từ hai đơn trở lên; (4) với mỗi đơn có nhiều lần thanh toán, kiểu thanh toán chiếm giá trị lớn nhất; (5) danh mục sản phẩm có doanh thu tăng ba tháng liên tiếp. Chạy file trên cả DuckDB và PostgreSQL. Ghi grain của mỗi kết quả trong một dòng chú thích.
\hoanthanh{năm câu cho cùng kết quả trên hai engine (chấp nhận khác biệt làm tròn), mỗi câu có chú thích grain, và không câu nào cộng \code{payment\_\allowbreak{}value} sau khi join với bảng items.}
\end{thuthach}

\section{Câu hỏi tự kiểm tra}

\begin{enumerate}
  \item DDL, DML, DCL, TCL là gì? Cho hai câu lệnh mỗi nhóm, và nêu một hệ quản trị xếp \code{GRANT} vào nhóm khác.
  \item Vì sao alias đặt trong danh sách \code{SELECT} không dùng được trong \code{WHERE}?
  \item Vì sao \code{WHERE col = NULL} không trả dòng nào, và nên viết gì thay?
  \item Vì sao \code{NOT IN (SELECT ...)} có thể không trả dòng nào khi subquery chứa NULL?
  \item Khi nào điều kiện phải nằm trong \code{HAVING} thay vì \code{WHERE}?
  \item \code{rank()}, \code{dense\_\allowbreak{}rank()} và \code{row\_\allowbreak{}number()} khác nhau thế nào khi có giá trị bằng nhau?
  \item Frame mặc định khi có \code{ORDER BY} là gì, và nó ảnh hưởng \code{last\_\allowbreak{}value()} ra sao?
  \item Primary key, unique constraint và unique index khác nhau thế nào trong PostgreSQL? \code{NULLS NOT DISTINCT} đổi điều gì?
  \item Khoá tự nhiên, khoá thay thế và khoá logic là gì? Grain của bảng payments Olist là gì?
  \item Read Committed của PostgreSQL cho phép hiện tượng nào mà Repeatable Read chặn? Read Uncommitted chạy thế nào?
  \item Trong \code{EXPLAIN ANALYZE}, số dòng ước lượng và số dòng thật cho bạn biết gì? Vì sao phải cẩn thận khi chạy nó với \code{UPDATE}?
  \item Trên Databricks và trên Iceberg, khai báo khoá chính có chặn được dòng trùng không? Pipeline phải làm gì thay?
\end{enumerate}

\section{Điểm dễ nhầm}

\begin{xltabular}{\textwidth}{@{}P{3.8cm} L P{3.4cm}@{}}
\toprule
\thd{điểm dễ nhầm} & \thd{thực tế} & \thd{docs chính thức} \\
\midrule
\endhead
\bottomrule
\endfoot
Nhiều người nghĩ \code{NULL = NULL} là true & So sánh thường với NULL ra NULL; dùng \code{IS NULL} hoặc \code{IS NOT DISTINCT FROM} & \href{https://www.postgresql.org/docs/current/functions-comparison.html}{PostgreSQL 18, mục 9.2} \\
Cột \code{UNIQUE} chỉ nhận một NULL & Mặc định hai NULL không bằng nhau, nên nhiều dòng NULL lọt qua; \code{NULLS NOT DISTINCT} đổi hành vi này & \href{https://www.postgresql.org/docs/current/ddl-constraints.html}{PostgreSQL 18, mục 5.5.3} \\
Không ghi frame thì window luôn phủ cả partition & Khi có \code{ORDER BY}, frame mặc định chỉ tới peer cuối cùng của dòng hiện tại; chỉ khi không có \code{ORDER BY} mới phủ cả partition & \href{https://www.postgresql.org/docs/current/sql-expressions.html}{PostgreSQL 18, mục 4.2.8} \\
\code{EXPLAIN ANALYZE} chỉ in kế hoạch & Nó chạy thật câu lệnh, side effect vẫn xảy ra; bọc trong \code{BEGIN} … \code{ROLLBACK} & \href{https://www.postgresql.org/docs/current/using-explain.html}{PostgreSQL 18, mục 14.1.2} \\
READ UNCOMMITTED của PostgreSQL cho đọc dữ liệu chưa commit & Nó chạy như Read Committed; PostgreSQL chỉ cài ba mức khác nhau & \href{https://www.postgresql.org/docs/current/transaction-iso.html}{PostgreSQL 18, mục 13.2} \\
CTE luôn được tính riêng một lần và chặn tối ưu & CTE không đệ quy, không side effect, tham chiếu một lần được gộp vào câu cha; \code{[NOT] MATERIALIZED} điều khiển việc đó & \href{https://www.postgresql.org/docs/current/queries-with.html}{PostgreSQL 18, mục 7.8.3} \\
Khai báo \code{FOREIGN KEY} là có index trên cột tham chiếu & PostgreSQL không tự tạo index đó; tự tạo khi cần join hoặc xoá bảng cha & \href{https://www.postgresql.org/docs/current/ddl-constraints.html}{PostgreSQL 18, mục 5.5.5} \\
Càng nhiều index càng nhanh & Index thêm chi phí cho cả hệ thống, cần dùng có chọn lọc & \href{https://www.postgresql.org/docs/current/indexes.html}{PostgreSQL 18, chương 11} \\
\code{GRANT} là DCL ở mọi hệ quản trị & Oracle xếp \code{GRANT}, \code{REVOKE} vào DDL và tự commit quanh mỗi câu DDL & \href{https://docs.oracle.com/en/database/oracle/oracle-database/23/sqlrf/Types-of-SQL-Statements.html}{Oracle 23: Types of SQL Statements} \\
Khai báo \code{PRIMARY KEY} trong lakehouse là có khoá duy nhất & Databricks không kiểm PK, FK, UNIQUE; Iceberg không bảo đảm tính duy nhất của identifier field; giữ bằng khử trùng, \code{MERGE} và test & \href{https://docs.databricks.com/aws/en/tables/constraints}{Databricks: Constraints}, \href{https://iceberg.apache.org/spec/\#identifier-field-ids}{Iceberg spec} \\
\end{xltabular}


% =====================================================================
\modulechapter{2}{Vòng đời DE, yêu cầu và chọn kho}
\moduleinfo{Cơ bản}{13 giờ}{Module 1; học song song với Module 3 được}
% =====================================================================
%  Bản viết lại theo kiểu tutorial cho Module 2, theo mẫu Module 3 bản thử.
%  Mỗi khối code có dòng "% chạy: <id>"; output thật nằm sau "% kết quả: <id>".
%  Khối "% chạy: <id> cần-lời-giải" chỉ chạy sau khi người học viết file của mình;
%  lời giải tham khảo nằm ở solutions/m02/.
%  Output chạy ngày 06/11/2026 bằng lakehouse/.venv/bin/python, Python 3.10.20,
%  DuckDB 1.5.6, SQLite 3.50.4, thư mục ~/Desktop/DE; cùng output trên DuckDB 2.0.0, Python 3.11.15.
% =====================================================================

\begin{xltabular}{\textwidth}{@{}P{3.4cm} L@{}}
\toprule
\thd{thẻ module} & \thd{nội dung} \\
\midrule
\endhead
\bottomrule
\endfoot
Giờ & 13 giờ lõi, trong đó 3 giờ cho chọn kho ở Bài 2 và Lab 2.3; thêm khoảng 6 giờ nếu đọc hết phần Đọc thêm \\
Dùng trong repo & sơ đồ ``Dữ liệu đi qua các namespace'' và lịch chạy trong \repo{README.md}; mục 1, 9, 10, 11 của \repo{docs/data-model.md}; mục 7 của \repo{docs/source-api-contract.md}; \repo{source\_\allowbreak{}sim/replay.py}; 9 file CSV trong \repo{lakehouse/source\_\allowbreak{}sim/data/}. Không cần service nào của repo chạy \\
Phiên bản & DuckDB 1.5.6 và SQLite 3.50.4 trong \repo{lakehouse/.venv}, Python 3.10.20. Docs: \href{https://duckdb.org/docs/1.5/}{DuckDB 1.5}, \href{https://docs.python.org/3.10/library/sqlite3.html}{Python 3.10 sqlite3} \\
Lướt được & Bài 6, vì repo \repo{lakehouse/} đã chia sẵn các tầng. Bài kiểm 15 phút, không mở tài liệu: ghi tầng medallion và kiểu ghi của 5 namespace \code{raw}, \code{staging}, \code{semantic}, \code{core}, \code{metadata}; kể 5 giai đoạn vòng đời, mỗi giai đoạn kèm phần repo làm việc đó; kể 6 undercurrent; viết một dòng cho mỗi khái niệm data warehouse, data lake, lakehouse; viết một câu cho Lambda, một câu cho Kappa, mỗi câu kèm lý do repo không cần nó. So với bảng ở Bài 6 và \repo{docs/data-model.md} mục 1; đúng 22 trên 26 ý là đạt \\
\end{xltabular}
% Kiểm: 26 ý = 5 namespace x 2 thuộc tính + 5 giai đoạn + 6 undercurrent + 3 khái niệm kho + 2 kiến trúc. Đáp án phần namespace ở lakehouse/docs/data-model.md dòng 9 tới 15 (cột lớp và ghi kiểu); phần còn lại ở bảng của Bài 6.

Module 1 dạy bạn viết SQL. Module này dạy hai việc phải làm trước câu SQL đầu tiên của một yêu cầu: biến câu hỏi bằng lời thành một định nghĩa chỉ cho ra một con số, và chọn kho hợp với cách dữ liệu sẽ được đọc.

Phần lớn giờ dành cho hai việc đó, ở Bài 1 và Bài 2. Các bài còn lại gắn tên cho những thứ có sẵn trong repo \repo{lakehouse}: hệ thống nguồn, vai trò trong team, ETL và ELT, các tầng dữ liệu. Repo \repo{lakehouse/} đã chia sẵn các tầng, nên Bài 6 chỉ cần lướt.

\begin{ghichu}
\textbf{Cách chạy code trong module này.}
\begin{itemize}
  \item Lời giải mẫu nằm trong \repo{solutions/m02/}. Đọc sau khi tự làm, để đối chiếu.
  \item Mọi khối code là lệnh bash, chạy từ thư mục cha của \repo{lakehouse/} bằng \repo{lakehouse/\allowbreak{}.venv/\allowbreak{}bin/\allowbreak{}python}. Venv này có DuckDB 1.5.6, và Python 3.10.20 của nó có sẵn SQLite 3.50.4 cùng phần tìm theo chữ FTS5.
  \item Lộ trình ghim DuckDB 2.0.0, nhưng bản này cần Python 3.11 trở lên, nên output trong module là của 1.5.6. Code chỉ dùng API có ở cả hai bản: lấy một dòng bằng \code{fetchall()[0]}, vì DuckDB 2.0 bỏ \code{fetchone()} trên relation.
  % Nguồn: https://pypi.org/pypi/duckdb/2.0.0/json (requires_python >=3.11.0); chạy thử của tác giả trên DuckDB 2.0.0
  \item Lab ghi file vào thư mục \code{lab\_\allowbreak{}m02/} và không ghi gì vào \repo{lakehouse/}. Khối của Bài 3 import code của repo, nên nó chạy Python với cờ \code{-B} để không sinh file \code{.pyc} trong repo.
  \item Phần ``Kết quả'' dưới mỗi khối là output thật, chạy ngày 06/11/2026. Thời gian đo ở Bài 2 đổi theo máy, nên khối đó chỉ in bậc của tỉ số và ghi thời gian ra file \code{lab\_\allowbreak{}m02/bai2\_\allowbreak{}thoi\_\allowbreak{}gian.txt}.
\end{itemize}
\end{ghichu}

% =====================================================================
\section{Vấn đề thật: 746 hay 673 đơn, cùng một câu hỏi}

Một BA, tức người phân tích nghiệp vụ, nhắn: ``Trong các đơn mua tháng 2/2017, bao nhiêu đơn giao tới khách sau 01/03/2017?'' Câu hỏi nghe rõ ràng, và bạn viết SQL trong hai phút. Đồng nghiệp ngồi cạnh cũng viết, và ra số khác.

% chạy: VD
\begin{codeblock}
lakehouse/.venv/bin/python - <<'EOF'
import duckdb
D = 'lakehouse/source_sim/data/'
con = duckdb.connect()
con.sql(f"CREATE VIEW orders AS SELECT * FROM '{D}olist_orders_dataset.csv'")
one = lambda q: con.sql(q).fetchall()[0][0]
thang_2 = """FROM orders WHERE order_purchase_timestamp >= TIMESTAMP '2017-02-01'
                         AND order_purchase_timestamp <  TIMESTAMP '2017-03-01'"""
print('đơn mua trong tháng 2/2017:', one(f'SELECT count(*) {thang_2}'))
for moc in ['2017-03-01', '2017-03-02']:
    n = one(f"""SELECT count(*) {thang_2}
                AND order_delivered_customer_date >= TIMESTAMP '{moc}'""")
    print(f'  giao tới khách từ 00:00 ngày {moc}:', n)
EOF
\end{codeblock}
\textit{Kết quả:}
% kết quả: VD
\begin{codeblock}
đơn mua trong tháng 2/2017: 1780
  giao tới khách từ 00:00 ngày 2017-03-01: 746
  giao tới khách từ 00:00 ngày 2017-03-02: 673
\end{codeblock}

Bạn đọc ``sau 01/03'' là từ 00:00 ngày 01/03, nên ra 746. Đồng nghiệp đọc là sau khi hết ngày 01/03, nên ra 673. Cả hai câu SQL đều đúng với cách hiểu của người viết. Không có lỗi hay cảnh báo nào cho biết hai người đang trả lời hai câu hỏi khác nhau.

Lab 2.1 tìm thêm các cách hiểu khác, như tính cả đơn chưa giao, hay đọc mốc theo giờ UTC. Repo \repo{lakehouse} cũng chưa có định nghĩa metric nào cho năm câu hỏi Gold ở \repo{docs/data-model.md} mục 11, nên câu hỏi Gold nào cũng có thể gặp chuyện này.
% Nguồn: lakehouse/docs/data-model.md dòng 372 tới 384: mục 11 chỉ liệt kê câu hỏi và cột nguồn, chưa có định nghĩa metric.

Module này trả lời hai câu hỏi. Viết yêu cầu thế nào để hai người đọc tính ra cùng một con số? Và khi đã biết dữ liệu sẽ được đọc ra sao, chọn kho nào để việc đọc đó nhanh?

% =====================================================================
\section{Từ SWE sang DE}

Bảng dưới nối thứ bạn đã biết với tên gọi trong data engineering, theo thứ tự bài học. Cột cuối là chỗ so sánh gãy, và bài học ghi ở cuối ô xử lý chỗ gãy đó trên dữ liệu thật.

\begin{xltabular}{\textwidth}{@{}P{3.2cm} P{2.8cm} L L@{}}
\toprule
\thd{bạn đã biết (SWE)} & \thd{trong DE gọi là} & \thd{giống ở đâu} & \thd{khác ở đâu} \\
\midrule
\endhead
\bottomrule
\endfoot
user story có tiêu chí chấp nhận & định nghĩa metric & viết trước khi code, kiểm được bằng ví dụ & hai người đọc phải tính ra cùng một con số trên dữ liệu còn đổi; Bài 1 \\
SLO độ trễ p99 của API & độ tươi của bảng & một con số cam kết, đo liên tục, có cảnh báo & đo bằng tuổi của dữ liệu mới nhất, không đo thời gian trả lời; pipeline trễ không báo lỗi nào; Bài 1 \\
đọc từ read replica còn trễ & số liệu chưa chốt & đọc sớm thì thấy giá trị cũ & trễ có thể tới hàng tháng, nên con số cần một ngày chốt đi kèm; Bài 1 \\
chọn database cho microservice theo query pattern & chọn kho theo kiểu truy cập & hỏi dữ liệu được đọc thế nào trước khi chọn kho & cùng dữ liệu nằm ở hai kho, và truy vấn phân tích quét cả cột của cả bảng; Bài 2 \\
database của chính service bạn & hệ thống nguồn & dữ liệu do app sinh ra & bạn không sở hữu nó; schema và nhãn thời gian có thể đổi mà không báo; Bài 3 \\
backend engineer & data engineer, analytics engineer & viết code, test, review, deploy & sản phẩm là bảng đúng và đúng giờ, người dùng là người phân tích; Bài 4 \\
event sourcing & ELT & giữ bản gốc để dựng lại khi luật đổi & bản gốc do team khác thiết kế; ETL chỉ giữ kết quả nên đổi luật là phải đọc lại nguồn; Bài 5 \\
tầng adapter, domain, read model & tầng raw, staging, core, marts & mỗi tầng một trách nhiệm, phụ thuộc một chiều & tầng sau luôn dựng lại được từ tầng trước; Bài 6 \\
ghi đôi vào database và cache & kiến trúc Lambda & một đường nhanh, một đường chắc & logic nằm ở hai nơi, và hai đường phải ra cùng số; Bài 6 \\
team sở hữu một service & sản phẩm dữ liệu (data product) & một team sở hữu, có hợp đồng và người dùng & người dùng là người phân tích, và mọi domain dùng chung nền tảng và luật quản trị; Bài 7 \\
dùng thư viện hay tự viết & tự viết hay mua công cụ dữ liệu & cùng phép tính chi phí sở hữu & phần đắt là giữ connector chạy đúng khi nguồn đổi; Bài 8 \\
\end{xltabular}

\textbf{Thói quen SWE gây hại nhất ở module này: nhận câu hỏi bằng lời rồi viết SQL ngay, như nhận ticket rồi code ngay.} Câu SQL sai định nghĩa vẫn chạy, vẫn trả về một con số trông hợp lý, và không test nào đỏ.

% =====================================================================
\section{Bạn sẽ làm được}

\begin{itemize}
  \item Liệt kê mọi cách hiểu của một câu hỏi có mốc thời gian, đếm từng cách, rồi viết lại câu hỏi để chỉ còn một đáp án.
  \item Viết thẻ yêu cầu gồm người dùng, grain, định nghĩa metric, mốc cắt và độ tươi, đủ rõ để người khác tính ra cùng con số.
  \item Chọn kho cho năm việc của Olist theo kiểu truy cập, mỗi lựa chọn có một số đo làm bằng chứng.
  \item Gắn đúng tên chuẩn cho các phần của repo trong 15 phút: tầng medallion, kiểu ghi, giai đoạn vòng đời, undercurrent, Lambda và Kappa.
  \item Quyết định ETL hay ELT, batch hay streaming, tự viết hay mua cho một nhu cầu, kèm lý do kiểm được.
\end{itemize}

% =====================================================================
\section{Bài 1: Từ câu hỏi nghiệp vụ tới định nghĩa metric}

\textbf{Vấn đề.} BA hỏi ``doanh thu tháng 2/2017''. Có ít nhất sáu cách tính hợp lý, mỗi cách ra một số khác, và số lớn nhất hơn số nhỏ nhất 47\%. Nếu yêu cầu không nói cách nào, bạn chọn thay BA mà không ai biết.

\textbf{Giải thích.} Bên SWE, bạn không code một user story khi chưa có tiêu chí chấp nhận. Trong DE, tiêu chí chấp nhận của một câu hỏi là một \textit{thẻ yêu cầu}: một trang đặc tả ngắn có năm phần. Thiếu phần nào thì câu SQL phải tự đoán phần đó.

\begin{xltabular}{\textwidth}{@{}P{2.6cm} P{4.2cm} L@{}}
\toprule
\thd{phần} & \thd{trả lời câu hỏi} & \thd{ví dụ cho doanh thu tháng} \\
\midrule
\endhead
\bottomrule
\endfoot
người dùng & ai đọc, để quyết định gì, xem lúc nào & BA xem dashboard mỗi sáng thứ Hai để chọn danh mục chạy khuyến mãi \\
grain & một dòng của kết quả là gì & một dòng mỗi tháng đặt hàng và mỗi danh mục \\
định nghĩa metric & công thức, cột nguồn, điều kiện lọc, đơn vị & tổng \code{price} của dòng hàng, không cộng \code{freight\_\allowbreak{}value}, bỏ đơn \code{canceled} và \code{unavailable}, đơn vị BRL \\
mốc cắt & cột thời gian nào, khoảng nào, múi giờ nào, chốt số vào ngày nào & \code{order\_\allowbreak{}purchase\_\allowbreak{}timestamp} từ 00:00 ngày 01/02 tới trước 00:00 ngày 01/03, giờ São Paulo; số của tháng chốt vào ngày 01 của tháng thứ hai sau đó \\
độ tươi & dữ liệu mới nhất được phép cũ bao lâu khi người dùng mở & lúc 08:00 thứ Hai, dữ liệu đã có tới 00:00 cùng ngày \\
\end{xltabular}

Grain là câu trả lời cho ``một dòng là gì''. Module 3 dạy grain kỹ; ở đây chỉ cần viết nó ra trước khi viết SQL.

Mốc cắt gồm hai thứ. Thứ nhất là một \textit{khoảng nửa mở} (half-open interval): lấy từ mốc đầu, tới trước mốc cuối, như \code{range()} của Python. Viết như vậy thì hai tháng liền nhau không chồng lên nhau và không hở. Thứ hai là \textit{ngày chốt} (as-of date): ngày bạn tính con số. Dữ liệu của một tháng đã qua vẫn còn tới thêm, nên cùng một câu hỏi tính vào hai ngày khác nhau ra hai số khác nhau.

Độ tươi (freshness) là tuổi của dữ liệu mới nhất mà người dùng thấy. Yêu cầu độ tươi phải so với lịch chạy của pipeline. Repo chạy Silver 6 lần mỗi ngày, nên dữ liệu Silver có thể trễ tới khoảng 4 giờ; một yêu cầu ``trễ không quá 1 giờ'' sẽ không đạt nếu chưa đổi lịch.
% Nguồn: lakehouse/airflow/README.md dòng 186 (cron hiện tại: 00, 04, ..., 20 giờ UTC; "dữ liệu Silver trễ tới khoảng 4 giờ")

\begin{ghichu}
\textbf{So với SWE.} Thẻ yêu cầu giống user story kèm tiêu chí chấp nhận viết theo kiểu Given-When-Then. Martin Fowler mô tả Given-When-Then là cách đặc tả hành vi của hệ thống bằng ví dụ. Câu kiểm của thẻ, như ``tháng 2/2017, chốt ngày 01/04/2017, ra 244.709,43 BRL'', chính là ví dụ đó.
% Nguồn: https://martinfowler.com/bliki/GivenWhenThen.html
% Kiểm: DuckDB 1.5.6 ngày 06/11/2026: sum(price) của đơn mua tháng 2/2017, bỏ canceled và unavailable, thêm điều kiện coalesce(order_approved_at, order_purchase_timestamp) < TIMESTAMP '2017-04-01', ra 244709.43; bỏ điều kiện đó ra 244959.35. Mô phỏng bằng source_sim.replay tại 2017-04-01 00:00 cũng ra 244709.43.
Độ tươi giống SLO độ trễ. Sách SRE của Google cho ví dụ ``99\% of Get RPC calls will complete in less than 100 ms''. SRE Workbook viết SLO độ tươi của pipeline theo cùng kiểu, ví dụ ``The oldest data is no older than Y''. Số chưa chốt thì giống đọc từ read replica còn trễ: đọc sớm thì thấy giá trị cũ.
% Nguồn: https://sre.google/sre-book/service-level-objectives/ ; https://sre.google/workbook/data-processing/

\textbf{Chỗ so sánh gãy.} Tiêu chí chấp nhận, kể cả khi đã thành test tự động, chạy trên dữ liệu mẫu cố định nên lần nào cũng ra cùng kết quả. Định nghĩa metric phải cho cùng một số mỗi lần chạy, trên dữ liệu vẫn còn đổi, nên thẻ phải có ngày chốt. Replica trễ vì chưa chép xong những gì đã xảy ra, còn số của tháng 2/2017 đổi vì sự việc chưa xảy ra: tới 19/09/2017 vẫn còn đơn tháng 2 được giao. Độ tươi đo bằng tuổi của dữ liệu, không đo thời gian trả lời: khi pipeline chạy trễ, dashboard vẫn mở nhanh và không có lỗi nào, chỉ có con số là cũ.
\end{ghichu}

\subsection{Làm thử}

Khối dưới tính ``doanh thu tháng 2/2017'' theo sáu cách, rồi đếm số đơn tháng 2 đã giao nếu bạn tính vào năm ngày khác nhau.

% chạy: A
\begin{codeblock}
lakehouse/.venv/bin/python - <<'EOF'
import duckdb
D = 'lakehouse/source_sim/data/'
con = duckdb.connect()
for t in ['orders', 'order_items', 'order_payments']:
    con.sql(f"CREATE VIEW {t} AS SELECT * FROM '{D}olist_{t}_dataset.csv'")
one = lambda q: con.sql(q).fetchall()[0][0]
thang_2 = lambda c: f"{c} >= TIMESTAMP '2017-02-01' AND {c} < TIMESTAMP '2017-03-01'"
mua = thang_2('order_purchase_timestamp')
hang = """SELECT o.*, CAST(i.price AS DECIMAL(14,2)) AS price,
                 CAST(i.freight_value AS DECIMAL(14,2)) AS freight
          FROM orders o JOIN order_items i USING (order_id)"""
cach = {
  'giá hàng, mọi trạng thái': f"SELECT sum(price) FROM ({hang}) WHERE {mua}",
  'giá hàng, bỏ canceled và unavailable': f"""SELECT sum(price) FROM ({hang})
      WHERE {mua} AND order_status NOT IN ('canceled', 'unavailable')""",
  'giá hàng, chỉ delivered': f"""SELECT sum(price) FROM ({hang})
      WHERE {mua} AND order_status = 'delivered'""",
  'giá hàng cộng phí vận chuyển': f"""SELECT sum(price + freight) FROM ({hang})
      WHERE {mua}""",
  'tiền khách trả, bảng order_payments': f"""SELECT
      sum(CAST(p.payment_value AS DECIMAL(14,2)))
      FROM orders o JOIN order_payments p USING (order_id) WHERE {mua}""",
  'giá hàng, theo ngày giao tới khách': f"""SELECT sum(price) FROM ({hang})
      WHERE {thang_2('order_delivered_customer_date')}""",
}
print('Doanh thu tháng 2/2017, BRL:')
for ten, q in cach.items():
    print(f'{one(q):>11}  {ten}')
print('Đơn mua tháng 2/2017, hai cách viết mốc cắt:')
print('  nửa mở [01/02, 01/03):', one(f'SELECT count(*) FROM orders WHERE {mua}'))
print("  BETWEEN DATE '2017-02-01' AND DATE '2017-02-28':", one(
    """SELECT count(*) FROM orders WHERE order_purchase_timestamp
       BETWEEN DATE '2017-02-01' AND DATE '2017-02-28'"""))
print('Đơn mua tháng 2/2017 đã giao tới khách, đếm vào lúc:')
for ngay in ['2017-03-01', '2017-03-15', '2017-04-01', '2017-07-01', '2018-12-31']:
    n = one(f"""SELECT count(*) FROM orders WHERE {mua}
                AND order_delivered_customer_date < TIMESTAMP '{ngay}'""")
    print(f'  00:00 ngày {ngay}: {n}')
EOF
\end{codeblock}
\textit{Kết quả:}
% kết quả: A
\begin{codeblock}
Doanh thu tháng 2/2017, BRL:
  247303.02  giá hàng, mọi trạng thái
  244959.35  giá hàng, bỏ canceled và unavailable
  234223.40  giá hàng, chỉ delivered
  286280.62  giá hàng cộng phí vận chuyển
  291908.01  tiền khách trả, bảng order_payments
  198909.29  giá hàng, theo ngày giao tới khách
Đơn mua tháng 2/2017, hai cách viết mốc cắt:
  nửa mở [01/02, 01/03): 1780
  BETWEEN DATE '2017-02-01' AND DATE '2017-02-28': 1725
Đơn mua tháng 2/2017 đã giao tới khách, đếm vào lúc:
  00:00 ngày 2017-03-01: 907
  00:00 ngày 2017-03-15: 1524
  00:00 ngày 2017-04-01: 1633
  00:00 ngày 2017-07-01: 1651
  00:00 ngày 2018-12-31: 1653
\end{codeblock}

Sáu con số đi từ 198.909,29 tới 291.908,01 BRL, và mỗi số đúng với một cách hiểu chữ ``doanh thu''. Tiền khách trả còn lớn hơn giá hàng cộng phí vận chuyển 5.627,39 BRL. Phần lớn khoản chênh, 5.567,14 BRL, là tiền của 47 đơn tháng 2 có thanh toán nhưng không có dòng hàng nào, gần hết là đơn \code{unavailable}; 60,25 BRL còn lại là chênh lệch bên trong 14 đơn. Thanh toán là một quy trình nghiệp vụ riêng, nên Module 3 tách nó thành một fact riêng.
% Kiểm: DuckDB 1.5.6 ngày 06/11/2026, đơn mua tháng 2/2017: 47 đơn, với 50 dòng thanh toán, có thanh toán mà không có dòng hàng nào (45 unavailable, 2 canceled), tổng 5567.14; 14 đơn có dòng hàng mà tổng thanh toán khác tổng price cộng freight_value, lệch tổng 60.25. Trên toàn bộ dữ liệu có 775 đơn không có dòng hàng (lakehouse/docs/data-model.md mục 9 hàng 5).

Phần cuối cho thấy vì sao thẻ cần ngày chốt. Nếu tính vào 00:00 ngày 01/03, chỉ có 907 đơn tháng 2 đã giao. Tới 15/03 là 1.524, và con số chỉ đứng yên ở 1.653 từ 19/09/2017, ngày đơn tháng 2 cuối cùng được giao. Theo luật phát lại của mock \repo{source\_\allowbreak{}sim}, version \code{delivered} của một đơn chỉ hiện ra khi đồng hồ mô phỏng tới ngày giao, nên một dashboard chạy lúc 00:00 ngày 15/03, chưa tính độ trễ của pipeline, thấy đúng 1.524.
% Nguồn: lakehouse/source_sim/replay.py dòng 59 tới 73 (mỗi mốc vòng đời là một version, updated_at là chính mốc đó)
% Kiểm: 291908.01 - 286280.62 = 5627.39; 291908.01 / 198909.29 = 1.4675, tức hơn 47%.
% Kiểm: DuckDB 1.5.6 ngày 06/11/2026, đơn tháng 2/2017 đã giao trước các mốc: 15/03 08:00 có 1525, 01/08 có 1651, 19/09 00:00 có 1652, 20/09 có 1653; lần giao cuối lúc 2017-09-19 14:36:39.

Doanh thu cũng đổi theo ngày chốt, vì dòng hàng chỉ hiện ra khi đơn được duyệt. Một đơn mua ngày 03/02/2017 tới 04/04/2017 mới được duyệt, nên số chốt ngày 01/04/2017 là 244.709,43 BRL, còn khối trên tính trên toàn bộ dữ liệu ra 244.959,35 BRL. Khoản chênh đúng bằng 249,92 BRL của đơn đó.
% Nguồn: lakehouse/source_sim/replay.py dòng 86 tới 88 (dòng hàng có updated_at là order_approved_at, không có thì order_purchase_timestamp); lakehouse/source_sim/app.py dòng 127 tới 136 (API chỉ trả version có updated_at <= giờ mô phỏng)
% Kiểm: đơn 2e5dc86c8c4aa663549caf5e31de840d mua 2017-02-03 00:04:49, duyệt 2017-04-04 10:56:48, trạng thái processing, tổng price 249.92; 244959.35 - 244709.43 = 249.92.

\begin{luuy}
\textbf{Bẫy thường gặp.}
\begin{itemize}
  \item \textbf{Viết mốc cắt bằng \code{BETWEEN} với ngày.} Với cột timestamp, \code{BETWEEN DATE '2017-02-01' AND DATE '2017-02-28'} dừng ở 00:00 ngày 28/02, nên mất cả ngày cuối tháng: 1.725 đơn thay vì 1.780. Luôn dùng khoảng nửa mở.
  \item \textbf{Để định nghĩa nằm trong câu SQL của từng người.} Hai người viết hai câu là hai định nghĩa. Viết định nghĩa một lần trong thẻ, rồi một lần trong bảng Gold mà mọi người cùng đọc, như Module 3 và Module 7 làm.
  \item \textbf{Hứa độ tươi chặt hơn lịch pipeline.} Với Silver chạy mỗi 4 giờ, lời hứa ``trễ dưới 1 giờ'' chỉ đúng vào vài thời điểm trong ngày.
  \item \textbf{Hứa ``đủ đơn'' rồi tưởng là đủ doanh thu.} Dòng hàng chỉ hiện ra khi đơn được duyệt. Theo luật phát lại, lúc 08:00 ngày 25/11/2017 dashboard chỉ thấy doanh thu ngày 24/11 của 1.029 đơn, 130.347,74 BRL; số cuối cùng là 1.166 đơn, 152.653,74 BRL.
  % Kiểm: mô phỏng bằng source_sim.replay ngày 06/11/2026, Python 3.10.20: đơn mua ngày 24/11/2017, bỏ canceled và unavailable, giờ mô phỏng 2017-11-25 08:00 thấy 1029 đơn có dòng hàng, tổng price 130347.74; trên toàn bộ dữ liệu 1166 đơn, 152653.74.
\end{itemize}
\end{luuy}

% =====================================================================
\section{Bài 2: Chọn kho theo kiểu truy cập}

\textbf{Vấn đề.} Olist cần hai việc trên cùng bảng \code{orders}. API của app tra một đơn theo \code{order\_\allowbreak{}id}, nhiều lần mỗi giây. Dashboard đếm đơn và số ngày giao trung bình theo tháng, trên cả 99.441 đơn. Chọn một kho cho cả hai thì một việc sẽ chậm, và bài này đo xem chậm bao nhiêu.

\textbf{Giải thích.} Bên SWE, bạn đã chọn database cho từng microservice theo cách nó được đọc. DE làm đúng việc đó với năm \textit{kiểu truy cập} (access pattern) hay gặp:

\begin{xltabular}{\textwidth}{@{}P{2.6cm} P{3.6cm} L P{3.2cm}@{}}
\toprule
\thd{kiểu truy cập} & \thd{ví dụ Olist} & \thd{thứ làm nó nhanh} & \thd{kho hay dùng} \\
\midrule
\endhead
\bottomrule
\endfoot
tra theo khoá & API tra một đơn theo \code{order\_\allowbreak{}id} & index B-tree hoặc bảng băm dẫn thẳng tới dòng & PostgreSQL, SQLite \\
quét phân tích & đếm đơn theo tháng trên mọi đơn & lưu theo cột, nén, tính trên từng lô giá trị & DuckDB, Trino trên Iceberg, BigQuery \\
tìm theo chữ & review có cụm ``não recebi'' & \textit{index đảo ngược} (inverted index): mỗi từ trỏ tới các dòng chứa nó & SQLite FTS5, OpenSearch \\
tìm gần nhất theo vector & review giống nghĩa một câu cho trước & index HNSW hoặc IVF trên vector & pgvector; nhánh T3 thử kiểu truy cập này \\
chuỗi thời gian & độ tươi của từng bảng, đo mỗi 15 phút & ghi nối đuôi theo thời gian, đọc theo khoảng & Prometheus \\
\end{xltabular}
% Nguồn: https://www.sqlite.org/fts5.html ; https://github.com/pgvector/pgvector (README: HNSW, IVFFlat) ; https://prometheus.io/docs/introduction/overview/ ("stores its metrics as time series data")

Hai kiểu đầu là hai cách xếp dữ liệu. \textit{Kho theo dòng} (row store) để mọi cột của một dòng nằm cạnh nhau. Tra một đơn qua index B-tree chỉ đọc vài trang, và có ngay đủ cả dòng. \textit{Kho theo cột} (column store) để mọi giá trị của một cột nằm cạnh nhau. Câu gộp chỉ đọc đúng những cột nó cần và tính trên cả lô giá trị một lần, nhưng muốn lấy đủ một dòng thì phải ghép từ mọi cột.

\begin{ghichu}
\textbf{So với SWE.} Mẫu ``database per service'' trên microservices.io viết: ``Different services have different data storage requirements. For some services, a relational database is the best choice.'' Chọn kho trong DE dùng đúng phép suy luận đó: hỏi dữ liệu được đọc thế nào trước, rồi mới chọn kho.
% Nguồn: https://microservices.io/patterns/data/database-per-service.html

\textbf{Chỗ so sánh gãy.} Với microservice, mỗi service giữ dữ liệu của mình và tự quản các bản sao như cache. Trong DE, dữ liệu đơn hàng nằm ở hai kho: database của app cho việc tra, và lakehouse cho việc quét, nơi pipeline của team dữ liệu chép sang theo lịch. Engine phân tích còn có chi phí cố định cho mỗi câu lệnh: trên máy viết bài, \code{SELECT 1} mất khoảng 33 µs với DuckDB và dưới 1 µs với SQLite. Vì vậy hàng nghìn câu nhỏ là cách dùng sai kho theo cột.
% Kiểm: ngày 06/11/2026, DuckDB 1.5.6, Python 3.10.20, Apple M5: 1.000 lần execute('SELECT 1').fetchall(), trung vị 7 lần đo: DuckDB 32,8 µs mỗi câu, SQLite 0,4 µs; chạy lại cùng ngày: DuckDB 33,4 µs, SQLite 0,8 µs.
\end{ghichu}

\begin{ghichu}
\textbf{Ba đề thi chấm đúng kỹ năng này,} dù tên dịch vụ khác nhau theo cloud. AWS DEA-C01 có Task 2.1 ``Choose a data store'' trong Domain 2, phần chiếm 26\% điểm; Skill 2.1.2 đòi cấu hình dịch vụ lưu trữ phù hợp ``for specific access patterns and requirements''. Hướng dẫn 4.2 của Google Professional Data Engineer mở mục 3.1 ``Selecting storage systems'' bằng ý ``Analyzing data access patterns''. Microsoft DP-700 mở nhóm kỹ năng ``Ingest and transform batch data'' bằng ``Choose an appropriate data store''.
% Nguồn: https://docs.aws.amazon.com/aws-certification/latest/data-engineer-associate-01/data-engineer-associate-01.html (Domain 2, 26%) và trang ...-domain2.html (Task 2.1, Skill 2.1.2); https://services.google.com/fh/files/misc/professional_data_engineer_exam_guide_english.pdf (tiêu đề PDF "v4.2 Professional Data Engineer Exam Guide", mục 3.1)
% Nguồn: https://learn.microsoft.com/en-us/credentials/certifications/resources/study-guides/dp-700 (updated_at 2026-11-05). Kiểm: mở ba trang ngày 06/11/2026 và chép đúng chữ; Skill 2.1.2 viết "Configure the appropriate storage services for specific access patterns and requirements".
\end{ghichu}

\subsection{Làm thử}

Khối dưới nạp cùng file \code{orders} vào SQLite, có index B-tree trên \code{order\_\allowbreak{}id}, và vào DuckDB, không có index. Nó in plan của mỗi engine, kiểm hai kho cho cùng kết quả, rồi đo hai việc: tra 1.000 \code{order\_\allowbreak{}id}, và gộp cả bảng theo tháng. Mỗi việc chạy nóng một lần, rồi đo xen kẽ hai kho 7 lần và lấy trung vị. DuckDB chỉ dùng 1 luồng, như SQLite, để số luồng không làm lệch phép so. Hai kho vẫn khác nhau ở index, ở kiểu dữ liệu của cột thời gian và ở cách chạy câu lệnh, nên phần sau khối xét từng chỗ khác đó.

% chạy: B
\begin{codeblock}
mkdir -p lab_m02
lakehouse/.venv/bin/python - <<'EOF'
import csv, json, math, sqlite3, statistics, time, duckdb
F = 'lakehouse/source_sim/data/olist_orders_dataset.csv'
rows = list(csv.reader(open(F, newline='')))
cols = rows[0]
# Kho theo dòng: SQLite; PRIMARY KEY tạo một index B-tree trên order_id
lite = sqlite3.connect(':memory:')
lite.execute(f'CREATE TABLE orders ({cols[0]} TEXT PRIMARY KEY, '
             + ', '.join(f'{c} TEXT' for c in cols[1:]) + ')')
lite.executemany(f"INSERT INTO orders VALUES ({', '.join('?' * len(cols))})",
                 [[v or None for v in r] for r in rows[1:]])
# Kho theo cột: DuckDB, cùng dữ liệu, không index; 1 luồng như SQLite
duck = duckdb.connect(config={'threads': 1})
duck.sql(f"CREATE TABLE orders AS SELECT * FROM '{F}'")
kho = {'SQLite': lite, 'DuckDB': duck}
tra = 'SELECT * FROM orders WHERE order_id = ?'
gop = {'SQLite': """SELECT substr(order_purchase_timestamp, 1, 7), count(*),
         round(avg(julianday(order_delivered_customer_date)
                   - julianday(order_purchase_timestamp)), 6)
       FROM orders GROUP BY 1 ORDER BY 1""",
       'DuckDB': """SELECT strftime(order_purchase_timestamp, '%Y-%m'), count(*),
         round(avg(date_diff('second', order_purchase_timestamp,
                             order_delivered_customer_date)) / 86400, 6)
       FROM orders GROUP BY 1 ORDER BY 1"""}
ids = sorted(r[0] for r in rows[1:])[::99][:1000]

def trung_vi(viec, lan=7):  # đo xen kẽ hai kho; vòng 0 chạy nóng, không tính
    giay = {k: [] for k in viec}
    for vong in range(lan + 1):
        for k, f in viec.items():
            t = time.perf_counter(); f(); t = time.perf_counter() - t
            if vong: giay[k].append(t)
    return {k: statistics.median(v) for k, v in giay.items()}

def quet(sql):  # nút quét bảng, ở đáy plan JSON của DuckDB
    nut = json.loads(duck.sql('EXPLAIN (FORMAT json) ' + sql).fetchall()[0][1])[0]
    while nut['children']:
        nut = nut['children'][0]
    return nut['extra_info']

plan = lambda sql, *p: '; '.join(
    r[3] for r in lite.execute('EXPLAIN QUERY PLAN ' + sql, p))
print(f'orders: {len(rows) - 1} dòng, {len(cols)} cột')
print('SQLite, tra:', plan(tra, 'x'))
print('SQLite, gộp:', plan(gop['SQLite']))
q = quet(tra.replace('?', f"'{ids[0]}'"))
print('DuckDB, tra:', q['Type'], '+ lọc',
      q['Filters'].strip('(').split('=')[0].strip())
q = quet(gop['DuckDB'])
print('DuckDB, gộp:', q['Type'], '+ chỉ đọc', len(q['Projections']), 'trên',
      len(cols), 'cột:')
print('  ' + ', '.join(q['Projections']))
kq = {k: [tuple(r) for r in c.execute(gop[k]).fetchall()] for k, c in kho.items()}
print(f"hai kho cho cùng {len(kq['SQLite'])} dòng kết quả:",
      kq['SQLite'] == kq['DuckDB'])

t_tra = trung_vi({k: lambda c=c: [c.execute(tra, [i]).fetchall() for i in ids]
                  for k, c in kho.items()})
t_gop = trung_vi({k: lambda c=c, k=k: c.execute(gop[k]).fetchall()
                  for k, c in kho.items()})
with open('lab_m02/bai2_thoi_gian.txt', 'w') as f:
    for k in kho:
        print(f'{k}: tra {t_tra[k] / len(ids) * 1e6:.1f} µs mỗi lần,',
              f'gộp {t_gop[k] * 1e3:.1f} ms', file=f)
bac = lambda r: 10 ** round(math.log10(r))  # làm tròn tỉ số về lũy thừa 10 gần nhất
print(f'tra {len(ids)} order_id: SQLite nhanh hơn DuckDB khoảng',
      bac(t_tra['DuckDB'] / t_tra['SQLite']), 'lần')
print('gộp theo tháng: DuckDB nhanh hơn SQLite khoảng',
      bac(t_gop['SQLite'] / t_gop['DuckDB']), 'lần')
EOF
\end{codeblock}
\textit{Kết quả:}
% kết quả: B
\begin{codeblock}
orders: 99441 dòng, 8 cột
SQLite, tra: SEARCH orders USING INDEX sqlite_autoindex_orders_1 (order_id=?)
SQLite, gộp: SCAN orders; USE TEMP B-TREE FOR GROUP BY
DuckDB, tra: Sequential Scan + lọc order_id
DuckDB, gộp: Sequential Scan + chỉ đọc 2 trên 8 cột:
  order_purchase_timestamp, order_delivered_customer_date
hai kho cho cùng 25 dòng kết quả: True
tra 1000 order_id: SQLite nhanh hơn DuckDB khoảng 100 lần
gộp theo tháng: DuckDB nhanh hơn SQLite khoảng 10 lần
\end{codeblock}

Mở \code{lab\_\allowbreak{}m02/bai2\_\allowbreak{}thoi\_\allowbreak{}gian.txt} để xem số của máy bạn. Trên máy viết bài, trung vị của 7 lần đo là: tra một \code{order\_\allowbreak{}id} mất 2,2 µs với SQLite và 159,6 µs với DuckDB; gộp 99.441 đơn theo tháng mất 28,0 ms với SQLite và 3,7 ms với DuckDB. Tỉ số thật vì vậy là khoảng 70 lần và khoảng 8 lần; khối chỉ in bậc độ lớn của chúng, vì con số cụ thể đổi theo máy.
% Kiểm: lab_m02/bai2_thoi_gian.txt sau khối B, ngày 06/11/2026, Apple M5 10 nhân, macOS 26, DuckDB 1.5.6, SQLite 3.50.4. Năm lần chạy lại cùng ngày: tra 2,1 tới 2,2 µs và 159,6 tới 161,1 µs; gộp 27,2 tới 28,9 ms và 3,5 tới 3,7 ms. Trên DuckDB 2.0.0 với Python 3.11.15, DuckDB tra khoảng 283 µs và gộp khoảng 3 ms; hai dòng bậc tỉ số in ra giống hệt.

Các dòng plan cho thấy mỗi engine đi đường nào. Khi tra theo khoá, SQLite đi \code{SEARCH ... USING INDEX}: B-tree dẫn thẳng tới một dòng. DuckDB không có index nên đi \code{Sequential Scan} và lọc \code{order\_\allowbreak{}id} trên cả 99.441 giá trị, cộng thêm chi phí cố định của mỗi câu lệnh. Thêm \code{PRIMARY KEY} cho \code{order\_\allowbreak{}id} trong DuckDB thì mỗi lần tra còn khoảng 62 µs, vẫn chậm hơn SQLite khoảng 30 lần: index giúp được, nhưng chi phí cố định của mỗi câu vẫn giữ kho theo cột ở xa phía sau.
% Kiểm: ngày 06/11/2026, DuckDB 1.5.6, 1 luồng, 1.000 lần tra, trung vị 7 lần đo: không index 161,0 µs, PRIMARY KEY 62,3 µs, CREATE INDEX 67,8 µs, SQLite 1,5 µs; 62,3 µs so với 2,2 µs của khối B là 28 lần. EXPLAIN của bảng có PRIMARY KEY vẫn in Sequential Scan.

Khi gộp, hai kho khác nhau ở ba chỗ. SQLite đi \code{SCAN orders}: nó xử lý từng dòng một, và với mỗi dòng gọi \code{julianday()} để đọc hai mốc thời gian lưu dạng chữ. DuckDB tính trên từng lô giá trị kiểu \code{TIMESTAMP}, và theo mục \code{Projections} của plan, nó chỉ đọc 2 trên 8 cột. Hai kho cho đúng cùng 25 dòng kết quả, nên chỗ khác chỉ là tốc độ.

Ở đây, đọc ít cột giúp ít nhất, vì dữ liệu nằm sẵn trong RAM. Khi mốc thời gian đã được đổi sẵn ra số, cùng câu gộp trên SQLite chạy 19,2 ms với bảng 3 cột và 20,1 ms với bảng 11 cột: thêm 8 cột chỉ tốn khoảng 4\%. Đọc mốc từ chữ tốn khoảng 8 ms, vì câu gốc chạy 28 ms. Phần chênh còn lại, từ khoảng 20 ms xuống 3,7 ms, đến từ cách chạy câu lệnh: SQLite xử lý từng dòng và gom nhóm bằng một B-tree tạm, còn DuckDB xử lý theo lô.
% Kiểm: ngày 06/11/2026, SQLite 3.50.4, Python 3.10.20, Apple M5, trung vị 7 lần đo sau một lần chạy nóng, cùng phép gộp avg(d - p) theo tháng trên mốc đã đổi bằng julianday() ra REAL: bảng 3 cột 19,2 và 19,3 ms; bảng 11 cột, gồm 8 cột chữ và 3 cột số đặt cuối, 20,1 ms; câu gốc của khối B trên mốc dạng chữ 28,0 và 28,1 ms; DuckDB 1.5.6 cùng phép gộp 3,4 ms.
Đọc ít cột thắng lớn khi dữ liệu nằm trên đĩa hay trên kho đối tượng, nơi mỗi cột bỏ qua là một phần byte không phải đọc. Ở đó chi phí đi theo byte: với giá on-demand, BigQuery tính tiền theo số byte câu truy vấn đọc.
% Nguồn: https://docs.cloud.google.com/bigquery/docs/best-practices-costs ("When using the on-demand pricing model, queries are billed according to the number of bytes read."; WebFetch ngày 06/11/2026)

Quy luật rút ra: kho nào thắng phụ thuộc vào việc, không có kho nào nhanh hơn cho mọi việc. Với bảng lớn hơn, chênh lệch ở việc tra còn tăng, vì khi bảng lớn gấp đôi, mỗi B-tree đọc thêm nhiều nhất một trang, còn quét phải đọc gấp đôi số giá trị.
% Lý luận, chưa đo: độ sâu B-tree tăng theo log của số dòng.

\begin{luuy}
\textbf{Bẫy thường gặp.}
\begin{itemize}
  \item \textbf{Chạy dashboard thẳng trên database của app.} Câu quét cả bảng giành CPU, bộ nhớ đệm và đĩa với giao dịch của người dùng. Đây là lý do có pipeline chép dữ liệu sang kho phân tích.
  \item \textbf{Dùng kho theo cột làm backend cho API tra theo khoá.} Ở Làm thử, DuckDB chậm hơn SQLite khoảng 70 lần cho việc này, và vẫn chậm hơn khoảng 30 lần khi đã có index.
  \item \textbf{Đo một lần rồi kết luận.} Lần chạy đầu gồm cả thời gian khởi động và nạp bộ nhớ đệm. Khối trên chạy nóng một lần, đo xen kẽ hai kho để chúng chịu cùng tải của máy, rồi lấy trung vị.
\end{itemize}
\end{luuy}

% =====================================================================
\section{Bài 3: Hệ thống nguồn}

\textbf{Vấn đề.} Bạn không sở hữu dữ liệu nguồn. Khi team server đổi API, đổi cách phân trang hay đổi nhãn thời gian, pipeline của bạn chịu hậu quả mà không ai báo trước.

\textbf{Giải thích.} Bên SWE, khi gọi service của team khác, bạn chỉ dựa vào hợp đồng API của họ, còn mọi thứ phía sau hợp đồng là việc của họ. Pipeline dữ liệu đứng đúng ở chỗ đó với mỗi \textit{hệ thống nguồn} (source system), tức nơi dữ liệu sinh ra. Cách đọc thì khác: pipeline đọc mọi bản ghi theo lịch, chứ không chỉ vài bản ghi cho mỗi request. Có bốn loại nguồn hay gặp, mỗi loại một cách đọc:

\begin{itemize}
  \item \textbf{Database giao dịch (OLTP) của app.} Đọc bằng câu truy vấn trên bản sao, hoặc đọc log thay đổi (change data capture, CDC) của database, ở nhánh T2.
  \item \textbf{API.} Kéo từng trang theo cursor, chịu giới hạn tốc độ, và chỉ thấy những gì API cho xem.
  \item \textbf{Log và luồng sự kiện.} Đọc từ một log như Kafka, ở Module 12.
  \item \textbf{File.} CSV hay Parquet được đổ ra định kỳ. Bộ Olist trên Kaggle là loại này.
\end{itemize}

Repo dùng loại thứ hai. Bronze gọi \code{GET /api/v1/<entity>} từng trang, với cursor do server cấp và client chỉ so như chuỗi. Nhóm server viết logic phát lại, còn \repo{source\_\allowbreak{}sim/} là mock làm đúng theo hợp đồng API.
% Nguồn: lakehouse/src/lakehouse/sources/olist_api.py dòng 3 tới 4; lakehouse/README.md dòng 22

Điểm quan trọng nhất: API này không trả sự kiện mà trả \textit{version} của trạng thái. Mỗi khi một đơn tới một mốc mới, API phát lại cả đơn đó, với \code{updated\_\allowbreak{}at} là lúc tới mốc.

\begin{ghichu}
\textbf{So với SWE.} Hệ thống nguồn là database của một service, nhưng không phải service của bạn. Với service của mình, bạn viết migration, biết từng cột nghĩa là gì, và đổi nhãn thời gian thì sửa luôn code đọc.

\textbf{Chỗ so sánh gãy.} Ở đây bạn chỉ có hợp đồng và dữ liệu. Nhãn còn có thể sai: hợp đồng coi mọi mốc là nhãn không có múi giờ, README của mock gọi là UTC, nhưng các chữ số thật ra là giờ São Paulo. Đọc quá nhiều còn làm chậm app của team khác, nên Bronze đọc theo trang và theo lịch.
% Nguồn: lakehouse/docs/source-api-contract.md dòng 279 tới 283 (mục 7 Múi giờ)
\end{ghichu}

\subsection{Làm thử}

Khối dưới dùng đúng hàm phát lại của mock để đếm số version API phát ra cho bảng \code{orders}, rồi in các version của một đơn.

% chạy: C
\begin{codeblock}
lakehouse/.venv/bin/python -B - <<'EOF'
import collections, sys
sys.path.insert(0, 'lakehouse')  # dùng đúng luật phát lại của mock trong repo
from source_sim.replay import order_versions, read_csv
orders = list(read_csv('lakehouse/source_sim/data/olist_orders_dataset.csv'))
so_version = collections.Counter(len(order_versions(o)) for o in orders)
print('số đơn trong CSV:', len(orders))
print('số version API phát ra:', sum(k * n for k, n in so_version.items()))
print('số đơn theo số version:', dict(sorted(so_version.items())))
don = next(o for o in orders if o['order_id'] == 'e481f51cbdc54678b7cc49136f2d6af7')
for updated_at, ban in order_versions(don):
    print(updated_at, ban['order_status'].ljust(9),
          ban['order_delivered_customer_date'])
EOF
\end{codeblock}
\textit{Kết quả:}
% kết quả: C
\begin{codeblock}
số đơn trong CSV: 99441
số version API phát ra: 391551
số đơn theo số version: {1: 177, 2: 1626, 3: 2430, 4: 95208}
2017-10-02 10:56:33 created   None
2017-10-02 11:07:15 approved  None
2017-10-04 19:55:00 shipped   None
2017-10-10 21:25:13 delivered 2017-10-10 21:25:13
\end{codeblock}

99.441 đơn thành 391.551 version, đúng bằng số dòng của \repo{raw.olist\_\allowbreak{}orders}. Staging giữ version mới nhất của mỗi khoá, nên \repo{stg\_\allowbreak{}olist\_\allowbreak{}orders} còn 99.441 dòng. Đơn mẫu đi qua bốn version, và chỉ version cuối có ngày giao.
% Nguồn: lakehouse/docs/data-model.md mục 10 (raw.olist_orders 391.551, staging.stg_olist_orders 99.441) và mục 9 hàng 17

\begin{luuy}
\textbf{Bẫy thường gặp.}
\begin{itemize}
  \item \textbf{Đếm đơn trên bảng raw.} Mỗi version là một dòng, nên đếm trên \repo{raw.olist\_\allowbreak{}orders} ra 391.551 thay vì 99.441.
  \item \textbf{Tin nhãn múi giờ.} Lab 2.1 chỉ cách kiểm múi giờ bằng chính dữ liệu.
  \item \textbf{So hai trục thời gian khác nhau.} \code{\_ingested\_\allowbreak{}at} của raw là giờ UTC thật của Spark, còn \code{updated\_\allowbreak{}at} là giờ mô phỏng. Hợp đồng dặn đừng so hai cột này.
  % Nguồn: lakehouse/docs/source-api-contract.md dòng 285
\end{itemize}
\end{luuy}

% =====================================================================
\section{Bài 4: Ai làm gì trong một team dữ liệu}

\textbf{Vấn đề.} Một con số trên dashboard sai. Ai sửa? Nếu không rõ ranh giới giữa các vai, lỗi định nghĩa như ở Bài 1 rơi vào khoảng trống giữa hai team.

\textbf{Giải thích.} Bên SWE, product owner viết tiêu chí chấp nhận, còn engineer là người tìm ra các trường hợp biên khi viết code. Team dữ liệu chia việc giống vậy, và repo có sẵn một ví dụ. Ngày hẹn giao của Olist luôn là 00:00, nên ``giao trễ'' có hai cách tính: so theo giờ thì có 7.827 đơn trễ, so theo ngày thì có 6.535 đơn. Người viết Silver thấy trường hợp biên này và chọn so theo ngày; quyết định đó nằm trong cột \repo{core.orders.is\_\allowbreak{}delivered\_\allowbreak{}late}, kèm lý do trong \repo{docs/data-model.md}. Ở team lớn hơn, analytics engineer giữ định nghĩa này, còn BA xác nhận nó đúng nghĩa nghiệp vụ.
% Nguồn: lakehouse/docs/data-model.md dòng 233 và 239 (is_delivered_late so theo ngày, vì ngày hẹn luôn là 00:00:00) và dòng 332 (mục 9 hàng 19: 6.535 đơn trễ); lakehouse/src/lakehouse/silver/core.py dòng 190 tới 191
% Kiểm: DuckDB 1.5.6 ngày 06/11/2026: so order_delivered_customer_date > order_estimated_delivery_date theo timestamp ra 7827 đơn trễ, so theo ngày ra 6535.

Bốn vai dưới hay gặp quanh một nền tảng dữ liệu. Cột cuối chỉ chỗ mỗi vai chạm vào repo.

\begin{xltabular}{\textwidth}{@{}P{2.8cm} L P{3.2cm} P{3.4cm}@{}}
\toprule
\thd{vai} & \thd{làm ra gì} & \thd{ai dùng trực tiếp} & \thd{trong repo} \\
\midrule
\endhead
\bottomrule
\endfoot
backend engineer & service, API và database của app & người dùng app, service khác & team server viết API Olist \\
data engineer & pipeline, bảng raw và Silver, lịch chạy, kiểm chất lượng & analytics engineer, hệ thống khác & Bronze, Silver, 3 DAG trong \repo{airflow/dags/} \\
analytics engineer & bảng Gold đã mô hình hoá, định nghĩa metric dùng chung & data analyst & Gold, chưa làm; Module 3 và Module 7 \\
data analyst & dashboard, phân tích, câu trả lời cho nghiệp vụ & BA, quản lý & Metabase, Module 11 \\
\end{xltabular}

Claire Carroll của dbt Labs viết: ``Analytics engineers create clean data sets for end users.'' Theo bà, vai này nối khoảng cách giữa data engineering và phân tích. Repo còn cho thấy một vai thứ năm: team nền tảng (platform team). Ở đây đó là nhóm khác vận hành Airflow, nên repo chỉ giao 3 file DAG và tài liệu.
% Nguồn: https://www.getdbt.com/what-is-analytics-engineering (Claire Carroll, 26/03/2019); lakehouse/README.md dòng 4

Ở công ty nhỏ, một người có thể làm cả ba vai dữ liệu: data engineer, analytics engineer và data analyst. Thứ cần chốt là việc, không phải chức danh: ai giữ định nghĩa metric, và ai trực khi pipeline hỏng.

\begin{ghichu}
\textbf{So với SWE.} Data engineer gần với backend engineer nhất: cùng viết code, viết test, review và deploy.

\textbf{Chỗ so sánh gãy.} Sản phẩm của data engineer là một bảng đúng và đúng giờ. Người dùng là người đọc số, không phải client gọi API theo hợp đồng chặt, nên lỗi thường lộ ra muộn, khi ai đó thấy một con số lạ.
\end{ghichu}

\begin{luuy}
\textbf{Bẫy thường gặp.}
\begin{itemize}
  \item \textbf{Nghĩ định nghĩa metric là việc riêng của BA.} BA biết nghiệp vụ, còn người viết SQL biết dữ liệu có những cách hiểu nào. Định nghĩa chỉ xong khi cả hai cùng ký vào thẻ yêu cầu.
\end{itemize}
\end{luuy}

% =====================================================================
\section{Bài 5: ETL và ELT}

\textbf{Vấn đề.} Tháng sau, BA muốn doanh thu cộng thêm phí vận chuyển. Nếu kho chỉ giữ bảng kết quả, bạn phải kéo lại dữ liệu từ nguồn, và nguồn có thể không còn giữ bản cũ.

\textbf{Giải thích.} Bên SWE, bạn từng chọn giữa hai cách lưu: chỉ giữ một bảng tổng đã tính sẵn, hay giữ cả log sự kiện để tính lại khi cần. Bảng tổng nhẹ hơn, nhưng khi công thức đổi thì không còn gì để tính lại. ETL và ELT là đúng lựa chọn đó cho cả một kho dữ liệu. ETL (extract, transform, load) biến đổi dữ liệu trước, rồi chỉ nạp kết quả vào kho. ELT (extract, load, transform) nạp dữ liệu thô vào kho trước, rồi biến đổi ngay trong kho, nên tính lại bao nhiêu lần cũng được.
% Nguồn: https://aws.amazon.com/compare/the-difference-between-etl-and-elt/ (ETL biến đổi trên một máy xử lý riêng rồi mới nạp; ELT: "you can interact with and transform the raw data as many times as needed")

Repo là ELT. Bronze nạp từng version JSON vào \code{raw} gần như nguyên trạng, chỉ thêm các cột dòng dõi như \code{\_ingested\_\allowbreak{}at} và \code{\_run\_\allowbreak{}id}. Silver biến đổi bằng Spark SQL ngay trong lakehouse. ETL vẫn hợp trong vài trường hợp: AWS kể database cũ và luồng cảm biến IoT. Thêm một trường hợp nữa: cột dữ liệu cá nhân không được phép vào kho thì phải bỏ trước khi nạp. Module 6 bàn dữ liệu cá nhân kỹ hơn.
% Nguồn: lakehouse/docs/data-model.md dòng 22 tới 32 (cột dòng dõi)

\begin{ghichu}
\textbf{So với SWE.} ELT giống event sourcing. Martin Fowler mô tả event sourcing là ``Capture all changes to an application state as a sequence of events'', và nhờ đó có thể bỏ trạng thái đi rồi dựng lại bằng cách chạy lại các event. Bronze giữ mọi version, nên Silver dựng lại được bất kỳ lúc nào.
% Nguồn: https://martinfowler.com/eaaDev/EventSourcing.html

\textbf{Chỗ so sánh gãy.} Event của event sourcing do chính bạn thiết kế. Dữ liệu thô của ELT do team khác thiết kế, có thể bẩn và có thể chứa dữ liệu cá nhân. Giữ thô cũng tốn chỗ: \repo{raw.olist\_\allowbreak{}orders} có gần bốn lần số dòng của bảng đơn đã làm sạch.
\end{ghichu}

\subsection{Làm thử}

Khối dưới tính doanh thu theo ngày bằng hai cách. ETL tính bằng Python rồi chỉ nạp kết quả vào kho. ELT nạp hai file CSV vào kho rồi tính bằng SQL trong kho. Sau đó luật đổi: doanh thu cộng thêm phí vận chuyển.

% chạy: E
\begin{codeblock}
lakehouse/.venv/bin/python - <<'EOF'
import collections, csv, decimal, duckdb
D = 'lakehouse/source_sim/data/'
doc = lambda f: csv.DictReader(open(D + f, newline=''))
# ETL: biến đổi bằng Python trước khi nạp; kho chỉ nhận bảng kết quả
ngay = {r['order_id']: r['order_purchase_timestamp'][:10]
        for r in doc('olist_orders_dataset.csv')
        if r['order_status'] not in ('canceled', 'unavailable')}
tong = collections.defaultdict(decimal.Decimal)
for r in doc('olist_order_items_dataset.csv'):
    if r['order_id'] in ngay:
        tong[ngay[r['order_id']]] += decimal.Decimal(r['price'])
etl = duckdb.connect()
etl.sql('CREATE TABLE daily_revenue (order_date DATE, revenue DECIMAL(14,2))')
etl.executemany('INSERT INTO daily_revenue VALUES (?, ?)', list(tong.items()))
# ELT: nạp nguyên CSV vào kho, rồi biến đổi bằng SQL ngay trong kho
elt = duckdb.connect()
elt.sql(f"CREATE TABLE raw_orders AS SELECT * FROM '{D}olist_orders_dataset.csv'")
elt.sql(f"CREATE TABLE raw_items AS SELECT * FROM '{D}olist_order_items_dataset.csv'")
luat = """CREATE OR REPLACE TABLE daily_revenue AS
  SELECT CAST(o.order_purchase_timestamp AS DATE) AS order_date,
         sum(CAST({tien} AS DECIMAL(14,2))) AS revenue
  FROM raw_orders o JOIN raw_items i USING (order_id)
  WHERE o.order_status NOT IN ('canceled', 'unavailable') GROUP BY 1"""
elt.sql(luat.format(tien='i.price'))
bang = lambda con: [r[0] for r in con.sql('SHOW TABLES').fetchall()]
lay = lambda con: sorted(con.sql('SELECT * FROM daily_revenue').fetchall())
print('ETL, bảng trong kho:', bang(etl))
print('ELT, bảng trong kho:', bang(elt))
print(f'hai cách cho cùng {len(lay(etl))} dòng daily_revenue:', lay(etl) == lay(elt))
print('Đổi luật: doanh thu cộng thêm freight_value')
elt.sql(luat.format(tien='i.price + i.freight_value'))
print('  ELT, chạy lại một câu SQL:',
      elt.sql('SELECT sum(revenue) FROM daily_revenue').fetchall()[0][0])
try:
    etl.sql('SELECT sum(freight_value) FROM daily_revenue')
except duckdb.Error as e:
    print('  ETL:', type(e).__name__, 'vì kho không có cột freight_value')
EOF
\end{codeblock}
\textit{Kết quả:}
% kết quả: E
\begin{codeblock}
ETL, bảng trong kho: ['daily_revenue']
ELT, bảng trong kho: ['daily_revenue', 'raw_items', 'raw_orders']
hai cách cho cùng 614 dòng daily_revenue: True
Đổi luật: doanh thu cộng thêm freight_value
  ELT, chạy lại một câu SQL: 15740216.93
  ETL: BinderException vì kho không có cột freight_value
\end{codeblock}

Trước khi luật đổi, hai cách cho cùng 614 dòng. Sau khi luật đổi, ELT chỉ cần chạy lại một câu SQL trên dữ liệu thô đã có trong kho. Kho của ETL không còn cột \code{freight\_\allowbreak{}value}, nên phải đọc lại nguồn và viết lại bước biến đổi.

\begin{luuy}
\textbf{Bẫy thường gặp.}
\begin{itemize}
  \item \textbf{Nghĩ ELT là không kiểm gì khi nạp.} Bronze vẫn ghi cursor và giờ nạp cho từng dòng, để biết đã đọc tới đâu và lần ngược một dòng về lần chạy đã sinh ra nó.
  \item \textbf{Nghĩ ELT đã thay hẳn ETL.} Cột dữ liệu cá nhân không được phép vào kho thì phải bỏ trước khi nạp; khi đó một bước ETL nhỏ đứng trước ELT.
\end{itemize}
\end{luuy}

% =====================================================================
\section{Bài 6: Lướt nhanh vòng đời, kho, tầng và batch}

\textbf{Vấn đề.} Repo \repo{lakehouse/} đã chia các tầng. Bài này chỉ gắn tên chuẩn cho chúng, để bạn đọc được tài liệu và đề thi mà không phải dịch trong đầu.

\textbf{Giải thích.} Mỗi dòng dưới là một khái niệm, một câu lời thường, và chỗ của nó trong repo.

\begin{xltabular}{\textwidth}{@{}P{3.0cm} L P{4.6cm}@{}}
\toprule
\thd{khái niệm} & \thd{nói gọn} & \thd{trong repo} \\
\midrule
\endhead
\bottomrule
\endfoot
vòng đời DE & sinh dữ liệu, lưu trữ, nạp, biến đổi, phục vụ & \repo{source\_\allowbreak{}sim}; rustfs và Iceberg; Bronze; Silver; Gold, Trino, Metabase ở các module sau \\
undercurrent & sáu mối lo chạy qua mọi giai đoạn: bảo mật, quản lý dữ liệu, DataOps, kiến trúc dữ liệu, orchestration, kỹ thuật phần mềm & Airflow, test trong \repo{tests/}, quyền ở Module 6 \\
data warehouse & kho phân tích, dữ liệu có cấu trúc, phục vụ BI & repo không dùng \\
data lake & kho file rẻ, giữ mọi kiểu dữ liệu ở dạng gốc & bucket \code{lakehouse} trên rustfs \\
lakehouse & file mở trên kho đối tượng, cộng một table format có giao dịch & Parquet, Iceberg 1.8.1, REST catalog \\
medallion & bronze thô chưa kiểm, silver đã kiểm và làm sạch, gold đã mô hình cho phân tích & \code{raw}; \code{staging}, \code{semantic}, \code{core}; \code{analytics} và \code{marts} chưa làm \\
batch và streaming & xử lý theo lô theo lịch, hay xử lý từng sự kiện khi nó tới & batch: Bronze mỗi 15 phút, Silver mỗi 4 giờ; streaming ở Module 12 \\
Lambda và Kappa & hai đường batch và speed song song, hay một đường stream tính lại bằng cách phát lại log & repo không cần kiến trúc nào trong hai: độ tươi tính bằng giờ là đủ, nên chỉ cần batch \\
\end{xltabular}
% Nguồn: FDE, mục lục chương 2; https://learn.microsoft.com/en-us/azure/architecture/databases/guide/big-data-architectures (Lambda, Kappa, lakehouse; ms.date 2025-09-12)
% Nguồn: https://docs.databricks.com/aws/en/lakehouse/medallion (cập nhật 29/10/2026); lakehouse/docs/data-model.md dòng 9 tới 15; lakehouse/README.md dòng 13 tới 24 và 78

Ba ý đáng nhớ hơn tên gọi. Tầng sau luôn dựng lại được từ tầng trước: \code{raw} giữ cả 391.551 version, nên 99.441 dòng của staging có thể dựng lại bất kỳ lúc nào. Medallion chia tầng theo chất lượng, không thay cho mô hình hoá: Databricks gọi nó là ``a recommended best practice but not a requirement'', và nói gold vẫn cần dimensional model. Còn Lambda có nhược điểm mà Azure Architecture Center nêu rõ: logic xử lý nằm ở hai nơi, với hai framework.

\begin{ghichu}
\textbf{So với SWE.} Các tầng giống kiến trúc nhiều lớp: adapter nhận dữ liệu ngoài vào, domain làm sạch, read model phục vụ việc đọc. Lambda giống ghi đôi vào database và cache: một đường nhanh, một đường chắc.

\textbf{Chỗ so sánh gãy.} Trong app, sửa dữ liệu ở domain là bình thường; ở đây không ai sửa tay \code{raw}, vì mọi tầng sau dựng từ nó. Cache sai thì xoá đi là xong, còn hai đường của Lambda phải ra cùng một con số, nên mỗi luật phải viết và test hai lần.
\end{ghichu}

\begin{luuy}
\textbf{Bẫy thường gặp.}
\begin{itemize}
  \item \textbf{Chọn streaming khi độ tươi tính bằng giờ là đủ.} Thẻ yêu cầu ở Bài 1 quyết định chuyện này, không phải sở thích công nghệ.
  \item \textbf{Coi lakehouse là một sản phẩm mua về.} Nó là một cách ghép: file mở, table format, catalog, và nhiều engine cùng đọc.
\end{itemize}
\end{luuy}

% =====================================================================
\section{Bài 7: Data mesh và sản phẩm dữ liệu}

\textbf{Vấn đề.} Khi có hàng chục team cùng cần dữ liệu, một team dữ liệu trung tâm thành nút cổ chai.

\textbf{Giải thích.} Data mesh là một cách tổ chức team, không phải một công nghệ. Zhamak Dehghani nêu bốn nguyên tắc: mỗi domain sở hữu dữ liệu phân tích của mình; dữ liệu là sản phẩm, và ``the consumers of that data should be treated as customers''; một nền tảng tự phục vụ chung; và luật quản trị chung chạy tự động. Trong repo này, một \textit{sản phẩm dữ liệu} (data product) sẽ là \repo{core.orders} có team chủ, có thẻ yêu cầu như Lab 2.2, có hợp đồng và SLO độ tươi. Một nền tảng nhỏ với một team như repo chưa cần data mesh; nó chỉ đáng cân nhắc khi có nhiều domain và nhiều team dữ liệu.
% Nguồn: https://martinfowler.com/articles/data-mesh-principles.html (Zhamak Dehghani, 03/12/2020)

\begin{ghichu}
\textbf{So với SWE.} Sản phẩm dữ liệu giống một service có team sở hữu, có hợp đồng và người dùng.

\textbf{Chỗ so sánh gãy.} Người dùng của sản phẩm dữ liệu là người phân tích, và mọi domain phải dùng chung nền tảng và luật quản trị, nên không team nào được tự chọn hết như với service riêng.
\end{ghichu}

\begin{luuy}
\textbf{Bẫy thường gặp.} Đổi tên team thành ``domain'' mà không có nền tảng chung và luật quản trị chung. Khi đó bạn có nhiều kho dữ liệu rời rạc, không phải data mesh.
\end{luuy}

% =====================================================================
\section{Bài 8: Tự viết hay mua}

\textbf{Vấn đề.} Bronze của repo là code tự viết. Mỗi nguồn mới là một quyết định mới: viết tiếp, hay mua một connector có sẵn.

\textbf{Giải thích.} Phần code thường nhỏ hơn bạn nghĩ, còn phần giữ cho nó đúng thì lớn hơn. Bronze của repo gồm client API 89 dòng và job nạp 81 dòng. Ba file test của phần này dài 373 dòng, và hợp đồng API mà code phải giữ đúng dài 442 dòng: cursor, trang, nhãn thời gian, chạy lại sau khi bị kill.
% Kiểm: wc -l trong lakehouse ngày 06/11/2026: src/lakehouse/sources/olist_api.py 89, src/lakehouse/bronze/ingest_olist.py 81, tests/test_bronze.py 173, tests/test_olist_api.py 130, tests/stub_api.py 70, docs/source-api-contract.md 442.

Năm câu hỏi dưới giúp quyết định. Sách Fundamentals of Data Engineering có mục ``Build Versus Buy'' trong chương 4 bàn kỹ hơn.
% Nguồn: FDE, mục lục chương 4

\begin{xltabular}{\textwidth}{@{}P{4.6cm} L L@{}}
\toprule
\thd{câu hỏi} & \thd{nghiêng về tự viết} & \thd{nghiêng về mua} \\
\midrule
\endhead
\bottomrule
\endfoot
nguồn đã có connector sẵn chưa? & chưa, như API nội bộ của Olist & có, như một SaaS phổ biến \\
phần này có giúp công ty khác đối thủ không? & có & không; chép dữ liệu là việc ai cũng làm \\
ai trực khi nguồn đổi? & team có người trực và có test như \repo{tests/test\_\allowbreak{}bronze.py} & không có ai \\
dữ liệu được phép đi qua bên thứ ba không? & không & có, sau khi xét hợp đồng \\
tổng chi phí sở hữu là gì? & giờ kỹ sư để viết, sửa và trực & tiền thuê bao; hỏi rõ cách tính giá trước khi ký \\
\end{xltabular}

\begin{ghichu}
\textbf{So với SWE.} Đây là câu hỏi quen: dùng thư viện hay tự viết.

\textbf{Chỗ so sánh gãy.} Một thư viện bạn ghim phiên bản rồi để yên. Một connector chạy mỗi 15 phút trên nguồn do người khác đổi, nên nó hỏng vào lúc bạn không chọn, dù là code tự viết hay đồ mua về.
\end{ghichu}

\begin{luuy}
\textbf{Bẫy thường gặp.} Chỉ tính giờ viết code. Với Bronze của repo, test và hợp đồng dài gần gấp năm lần code nạp; đó mới là phần tốn công lâu dài.
\end{luuy}
% Kiểm: (373 + 442) / (89 + 81) = 4,8.

% =====================================================================
\section{Lab 2.1: Một câu hỏi, một đáp án}

Lab này làm lại câu hỏi ở phần Vấn đề thật cho tới khi chỉ còn một đáp án. Mọi file nằm trong \code{lab\_\allowbreak{}m02/}.

\textbf{Bước 1: tạo công cụ chạy SQL và kiểm múi giờ.} File \code{hoi.py} mở DuckDB, tạo view cho 9 file CSV, rồi chạy từng câu SQL trong file bạn đưa vào. Mỗi câu kết thúc bằng dấu chấm phẩy. \code{hoi.py} bỏ comment trước khi tách câu, nên comment chứa gì cũng được; chỉ chuỗi trong câu SQL là không được chứa ký hiệu mở comment. Công cụ đặt session ở UTC, giống session Spark của repo, để kết quả không phụ thuộc múi giờ máy bạn. File CSV không ghi múi giờ, nên hai câu trong \code{mui\_\allowbreak{}gio.sql} tìm bằng chứng ngay trong dữ liệu.
% Nguồn: lakehouse/src/lakehouse/spark.py dòng 12 (spark.sql.session.timeZone UTC)

% chạy: L1
\begin{codeblock}
mkdir -p lab_m02
cat > lab_m02/hoi.py <<'EOF'
"""Chạy từng câu SQL trong một file trên 9 CSV Olist và in kết quả."""
import pathlib, re, sys, duckdb
D = 'lakehouse/source_sim/data/'
con = duckdb.connect()
con.sql("SET TimeZone = 'UTC'")  # như session Spark của repo
for t in ['orders', 'order_items', 'order_payments', 'order_reviews',
          'customers', 'products', 'sellers', 'geolocation']:
    con.sql(f"CREATE VIEW {t} AS SELECT * FROM '{D}olist_{t}_dataset.csv'")
con.sql(f"CREATE VIEW category_translation AS "
        f"SELECT * FROM '{D}product_category_name_translation.csv'")
van_ban = pathlib.Path(sys.argv[1]).read_text()
# bỏ comment trước, để dấu ; trong comment không cắt nhầm câu
van_ban = re.sub(r'/\*.*?\*/|--[^\n]*', ' ', van_ban, flags=re.S)
for cau in con.extract_statements(van_ban):
    sql = cau.query.strip().rstrip(';')  # DuckDB 2.0 giữ dấu ; cuối câu
    if cau.type == duckdb.StatementType.SELECT:
        con.sql(sql).show()
    else:  # CREATE, SET: chạy, không in gì
        con.execute(sql)
EOF
cat > lab_m02/mui_gio.sql <<'EOF'
/* 1. Review tạo lúc 01:00 thay vì 00:00, đếm theo ngày */
SELECT CAST(review_creation_date AS DATE) AS ngay, count(*) AS so_review
FROM order_reviews WHERE hour(review_creation_date) = 1
GROUP BY 1 ORDER BY 1;
/* 2. Đồng hồ São Paulo quanh 03:00 UTC ngày 16/10/2016 */
SELECT t AS gio_utc, t AT TIME ZONE 'America/Sao_Paulo' AS gio_sao_paulo
FROM (VALUES (TIMESTAMPTZ '2016-10-16 02:59:59+00'),
             (TIMESTAMPTZ '2016-10-16 03:00:00+00')) v(t);
EOF
lakehouse/.venv/bin/python lab_m02/hoi.py lab_m02/mui_gio.sql
\end{codeblock}
\textit{Kết quả:}
% kết quả: L1
\begin{codeblock}
┌────────────┬───────────┐
│    ngay    │ so_review │
│    date    │   int64   │
├────────────┼───────────┤
│ 2016-10-16 │         1 │
│ 2017-10-15 │        84 │
└────────────┴───────────┘

┌──────────────────────────┬─────────────────────┐
│         gio_utc          │    gio_sao_paulo    │
│ timestamp with time zone │      timestamp      │
├──────────────────────────┼─────────────────────┤
│ 2016-10-16 02:59:59+00   │ 2016-10-15 23:59:59 │
│ 2016-10-16 03:00:00+00   │ 2016-10-16 01:00:00 │
└──────────────────────────┴─────────────────────┘
\end{codeblock}

Câu 2 cho thấy nửa đêm ngày 16/10/2016 không có ở São Paulo: đồng hồ nhảy từ 23:59:59 sang 01:00:00, vì Brazil bắt đầu giờ mùa hè. Review được tạo ``lúc nửa đêm'' của ngày đó và của ngày 15/10/2017 nên mang giờ 01:00, đúng như câu 1.

Hai dấu hiệu này khớp \repo{docs/data-model.md} mục 9 hàng 18, nơi còn ghi rằng ít đơn nhất được đặt lúc 4 tới 5 giờ sáng. Vậy mọi mốc là giờ địa phương Brazil, dù hợp đồng API coi chúng là nhãn không múi giờ và README của mock gọi là UTC.
% Nguồn: lakehouse/docs/data-model.md dòng 331 (mục 9 hàng 18: 85 review mang 01:00:00, 1 vào 2016-10-16 và 84 vào 2017-10-15; ít đơn nhất lúc 4 tới 5 giờ, 206 và 188 đơn)

\textit{Nếu lỗi:} \code{IO Error: No files found} nghĩa là bạn không chạy từ thư mục cha của \repo{lakehouse/}. Nếu cột \code{gio\_\allowbreak{}utc} hiện \code{+07} thay vì \code{+00}, file \code{hoi.py} của bạn thiếu dòng \code{SET TimeZone}.

\textbf{Bước 2: đếm theo từng cách hiểu.} Viết file \code{lab\_\allowbreak{}m02/cach\_\allowbreak{}hieu.sql}, một câu truy vấn trả về sáu dòng, mỗi dòng một cách hiểu, cột \code{cach\_\allowbreak{}hieu} và \code{so\_\allowbreak{}don}, sắp theo \code{cach\_\allowbreak{}hieu}:

\begin{enumerate}
  \item[a.] đơn mua trong tháng 2, mốc theo giờ São Paulo như dữ liệu;
  \item[b.] đơn mua trong tháng 2, nhưng hai mốc 00:00 ngày 01/02 và 00:00 ngày 01/03 tính theo giờ UTC;
  \item[c.] trong a, đơn giao tới khách từ 00:00 ngày 01/03;
  \item[d.] trong a, đơn giao tới khách từ 00:00 ngày 02/03;
  \item[e.] trong a, đơn chưa giao lúc 00:00 ngày 01/03, kể cả đơn không bao giờ giao;
  \item[f.] trong a, đơn giao từ 00:00 ngày 01/03 giờ UTC.
\end{enumerate}

Để đổi một mốc UTC sang giờ São Paulo, dùng \code{TIMESTAMPTZ '2017-03-01 00:00:00+00' AT TIME ZONE 'America/Sao\_\allowbreak{}Paulo'} như câu 2 ở bước 1. Muốn sáu dòng ra đúng thứ tự, đặt nhãn bắt đầu bằng chữ cái và thêm \code{ORDER BY 1}.

% chạy: L2 cần-lời-giải
\begin{codeblock}
lakehouse/.venv/bin/python lab_m02/hoi.py lab_m02/cach_hieu.sql
\end{codeblock}
\textit{Kết quả, với lời giải tham khảo:}
% kết quả: L2
\begin{codeblock}
┌──────────────────────────────────────────────┬────────┐
│                  cach_hieu                   │ so_don │
│                   varchar                    │ int64  │
├──────────────────────────────────────────────┼────────┤
│ a. mua trong tháng 2, giờ São Paulo          │   1780 │
│ b. mua trong tháng 2, giờ UTC                │   1771 │
│ c. trong a, giao từ 00:00 ngày 01/03         │    746 │
│ d. trong a, giao từ 00:00 ngày 02/03         │    673 │
│ e. trong a, chưa giao lúc 00:00 ngày 01/03   │    873 │
│ f. trong a, giao từ 00:00 ngày 01/03 giờ UTC │    746 │
└──────────────────────────────────────────────┴────────┘
\end{codeblock}

\textit{Nếu lệch:} dòng a khác 1.780 thì kiểm khoảng nửa mở; dòng e khác 873 thì kiểm cách bạn xử lý \code{NULL}, vì 127 đơn chưa bao giờ được giao.

Dòng b mất 9 đơn so với dòng a: múi giờ đổi cả mẫu số. Dòng f bằng dòng c chỉ là may: mốc 00:00 UTC ngày 01/03 là 21:00 ngày 28/02 ở São Paulo, và từ giờ đó tới hết ngày không có đơn nào được giao.
% Kiểm: hoi.py ngày 06/11/2026, DuckDB 1.5.6: đơn tháng 2 giao theo ngày 24/02: 80, 25/02: 14, 27/02: 1 (20:48:20), 01/03: 73, 02/03: 83; không có dòng cho 26/02 và 28/02. Mốc 00:00 giờ Hà Nội ngày 01/03 là 14:00 ngày 28/02 ở São Paulo và cũng cho 746; tháng 2 theo giờ Hà Nội có 1.776 đơn.

\textbf{Bước 3: viết lại câu hỏi.} BA trả lời: ``sau 01/03'' nghĩa là từ ngày 02/03 trở đi; mọi mốc theo giờ Brazil như dữ liệu; đơn chưa giao không tính; lấy số trên toàn bộ dữ liệu hiện có. Viết \code{lab\_\allowbreak{}m02/cau\_\allowbreak{}hoi\_\allowbreak{}chot.sql}: một comment ở đầu chứa câu hỏi đã viết lại thành một câu, rồi câu SQL trả về đúng một số.

% chạy: L3 cần-lời-giải
\begin{codeblock}
lakehouse/.venv/bin/python lab_m02/hoi.py lab_m02/cau_hoi_chot.sql
\end{codeblock}
\textit{Kết quả, với lời giải tham khảo:}
% kết quả: L3
\begin{codeblock}
┌────────┐
│ so_don │
│ int64  │
├────────┤
│    673 │
└────────┘
\end{codeblock}

\textit{Nếu lệch:} kết quả khác 673 thì kiểm lại mốc ``sau 01/03'' như dòng d ở bước 2, và kiểm rằng đơn chưa giao không được tính.

\hoanthanh{\code{cach\_\allowbreak{}hieu.sql} in đúng sáu số 1.780, 1.771, 746, 673, 873 và 746; \code{cau\_\allowbreak{}hoi\_\allowbreak{}chot.sql} in 673; câu hỏi viết lại nêu đủ bốn thứ: khoảng nửa mở của từng mốc, múi giờ, cách tính đơn chưa giao, và dữ liệu tới ngày nào.}

% =====================================================================
\section{Lab 2.2: Thẻ yêu cầu cho năm câu hỏi Gold}

Năm câu hỏi ở \repo{docs/data-model.md} mục 11 sẽ thành bảng Gold ở Module 3 và Module 7. Lab này viết thẻ yêu cầu cho cả năm, rồi kiểm ba thẻ bằng SQL. Một thẻ đạt khi người khác đọc nó và tính ra đúng số của bạn mà không phải hỏi lại.

\textbf{Bước 1: tạo file thẻ.} Câu 1 đã có thẻ mẫu, viết từ biên bản ở bước 2. Bốn câu còn lại để trống.

% chạy: L4
\begin{codeblock}
cat > lab_m02/yeu_cau.md <<'EOF'
# Thẻ yêu cầu cho năm câu hỏi Gold

## 1. Doanh thu theo ngày, danh mục, bang người bán
- Người dùng: team BA, dashboard xem lúc 08:00 mỗi sáng; chọn danh mục chạy khuyến mãi
- Grain: một dòng mỗi (ngày đặt, danh mục, bang người bán)
- Metric: tổng price của dòng hàng, BRL; không cộng freight_value
- Lọc: bỏ đơn canceled và unavailable
- Danh mục: tên tiếng Anh; không có bản dịch thì tên gốc; không có tên thì unknown
- Mốc cắt: ngày của order_purchase_timestamp, giờ São Paulo, [00:00, 00:00 hôm sau)
- Ngày chốt: số của một ngày chốt sau 14 ngày; trước đó là số tạm
- Độ tươi: lúc 08:00 có mọi đơn đặt trước 00:00; đơn chưa duyệt chưa có doanh thu
- Câu kiểm: 24/11/2017 trên toàn bộ dữ liệu: số dòng, tổng cả ngày, dòng lớn nhất

## 2. Giao hàng nhanh hay chậm
## 3. Khách quay lại
## 4. Điểm review theo người bán hay danh mục
## 5. Bản đồ khách và người bán
EOF
grep -c '^## ' lab_m02/yeu_cau.md
\end{codeblock}
\textit{Kết quả:}
% kết quả: L4
\begin{codeblock}
5
\end{codeblock}
% Kiểm: mô phỏng bằng source_sim.replay ngày 06/11/2026, Python 3.10.20, trạng thái đơn theo giờ mô phỏng: doanh thu ngày 24/11/2017, bỏ canceled và unavailable, thấy 130347.74 lúc 25/11 08:00, 146442.66 lúc 28/11 08:00, 152588.84 lúc 01/12 08:00, và 152653.74 từ 08/12 00:00, bằng số trên toàn bộ dữ liệu. Đơn ngày 24/11 được duyệt muộn nhất lúc 2017-12-07 09:50:08. Trên toàn bộ dữ liệu, sau khi bỏ canceled và unavailable, chỉ 7 đơn được duyệt sau hơn 14 ngày kể từ lúc mua.

\textit{Nếu lỗi:} số nhỏ hơn 5 thường là do khi chép từ PDF, đầu dòng tiêu đề bị thêm dấu cách. Mỗi tiêu đề thẻ phải bắt đầu bằng \code{\#\#} ngay ở đầu dòng.

\textbf{Bước 2: đọc biên bản phỏng vấn.} Đây là câu trả lời của BA cho ba câu hỏi đầu.

\begin{ghichu}
\begin{itemize}
  \item \textbf{Câu 1.} ``Doanh thu là tiền hàng, không tính tiền ship. Đơn huỷ hay hết hàng thì không tính. Ngày là ngày khách đặt, theo giờ Brazil. Tôi xem lúc 8 giờ sáng và cần đủ số của hôm qua.''
  \item \textbf{Câu 2.} ``Trễ là giao tới khách sau ngày hẹn, so theo ngày, không theo giờ. Đơn chưa giao thì để riêng, không tính là trễ hay đúng hẹn.''
  \item \textbf{Câu 3.} ``Khách là một người mua thật, không phải một mã khách của từng đơn. Quay lại là có từ hai đơn trở lên, đơn ở trạng thái nào cũng tính.''
\end{itemize}
\end{ghichu}

\textbf{Bước 3: viết thẻ cho câu 2 tới 5.} Mỗi thẻ có đủ các dòng như thẻ mẫu. Bảng dưới liệt kê điều mỗi thẻ phải chốt, kèm số trong repo cho thấy vì sao phải chốt.

\begin{xltabular}{\textwidth}{@{}P{2.8cm} L P{4.8cm}@{}}
\toprule
\thd{câu hỏi} & \thd{thẻ phải chốt} & \thd{số trong repo} \\
\midrule
\endhead
\bottomrule
\endfoot
2. giao nhanh hay chậm & so theo ngày hay theo giờ; mẫu số có tính đơn chưa giao không & ngày hẹn luôn là 00:00:00; 2.965 đơn chưa giao (mục 9 hàng 19) \\
3. khách quay lại & khách là \code{customer\_\allowbreak{}id} hay \code{customer\_\allowbreak{}unique\_\allowbreak{}id}; mấy đơn thì là quay lại & 99.441 \code{customer\_\allowbreak{}id} nhưng 96.096 người mua (mục 9 hàng 8) \\
4. review theo người bán hay danh mục & đơn nhiều review lấy review nào; review của đơn nhiều người bán tính cho ai & 547 đơn có nhiều review (mục 9 hàng 4); 1.278 đơn có hơn một người bán \\
5. bản đồ khách và người bán & toạ độ lấy theo zip thế nào; người không có toạ độ thì làm gì & 270 người mua và 8 người bán không có toạ độ (mục 9 hàng 13) \\
\end{xltabular}
% Kiểm: hoi.py ngày 06/11/2026: số order_id có hơn một seller_id khác nhau trong order_items là 1278; số đơn có hơn một review là 547, khớp mục 9 hàng 4.

\textbf{Bước 4: kiểm ba thẻ bằng SQL.} Viết \code{lab\_\allowbreak{}m02/q1.sql}, \code{q2.sql} và \code{q3.sql} chỉ từ thẻ của bạn, không nhìn lại biên bản. File \code{q1.sql} tạo bảng tạm theo grain của thẻ, rồi chạy câu kiểm ngày 24/11/2017 trả về một dòng: số dòng, tổng, và dòng có doanh thu lớn nhất. File \code{q2.sql} đếm đơn trễ, đơn đúng hẹn và đơn chưa giao. File \code{q3.sql} đếm người mua, người quay lại và phần trăm làm tròn 2 chữ số.

% chạy: L5 cần-lời-giải
\begin{codeblock}
for q in q1 q2 q3; do lakehouse/.venv/bin/python lab_m02/hoi.py lab_m02/$q.sql; done
\end{codeblock}
\textit{Kết quả, với lời giải tham khảo:}
% kết quả: L5
\begin{codeblock}
┌─────────┬───────────────┬───────────────────┬────────────────────┐
│ so_dong │     tong      │   dong_lon_nhat   │ doanh_thu_lon_nhat │
│  int64  │ decimal(38,2) │      varchar      │   decimal(38,2)    │
├─────────┼───────────────┼───────────────────┼────────────────────┤
│     146 │     152653.74 │ bed_bath_table SP │           14717.34 │
└─────────┴───────────────┴───────────────────┴────────────────────┘

┌───────┬──────────┬───────────┐
│  tre  │ dung_hen │ chua_giao │
│ int64 │  int64   │   int64   │
├───────┼──────────┼───────────┤
│  6535 │    89941 │      2965 │
└───────┴──────────┴───────────┘

┌───────────┬──────────┬───────────┐
│ nguoi_mua │ quay_lai │ phan_tram │
│   int64   │  int64   │  double   │
├───────────┼──────────┼───────────┤
│     96096 │     2997 │      3.12 │
└───────────┴──────────┴───────────┘
\end{codeblock}

Số của câu 2 và câu 3 khớp đúng \repo{docs/data-model.md} mục 9: 6.535 đơn trễ, 89.941 đơn đúng hẹn, 2.965 đơn chưa giao, và 2.997 người mua có hơn một \code{customer\_\allowbreak{}id}. Repo đã đi tới cùng định nghĩa đó cho cột \code{is\_\allowbreak{}delivered\_\allowbreak{}late}.
% Nguồn: lakehouse/docs/data-model.md dòng 321 (hàng 8: 2.997 người mua) và dòng 332 (hàng 19: 6.535, 89.941, 2.965)
% Kiểm: hoi.py ngày 06/11/2026, DuckDB 1.5.6: so order_delivered_customer_date > order_estimated_delivery_date theo timestamp ra 7827 đơn trễ; 99441 customer_id khác nhau cho 99441 đơn.

\textit{Nếu lệch:} câu 2 ra 7.827 đơn trễ thì bạn đang so theo giờ, vì ngày hẹn luôn là 00:00 nên đơn giao lúc 10:00 đúng ngày hẹn bị tính là trễ. Câu 3 ra 0 người quay lại thì bạn đang đếm theo \code{customer\_\allowbreak{}id}, vì mỗi đơn có một \code{customer\_\allowbreak{}id} mới.

\textbf{Bước 5: người đọc thứ hai.} Đưa thẻ câu 4 và câu 5 cho một người khác, hoặc tự đọc lại sau hai ngày. Người đó viết SQL chỉ từ thẻ, và hai bạn so số. Chỗ nào lệch là chỗ thẻ còn thiếu một quyết định; ghi quyết định đó vào thẻ rồi so lại.

\hoanthanh{\code{yeu\_\allowbreak{}cau.md} có năm thẻ, thẻ nào cũng đủ người dùng, grain, metric, lọc, mốc cắt có cả ngày chốt, độ tươi và câu kiểm; \code{q1.sql} in 146 dòng, tổng 152.653,74 BRL, và dòng lớn nhất là \code{bed\_\allowbreak{}bath\_\allowbreak{}table} của bang SP với 14.717,34 BRL; \code{q2.sql} in 6.535, 89.941 và 2.965; \code{q3.sql} in 96.096, 2.997 và 3,12; người đọc thứ hai ra cùng số với bạn ở câu 4 và câu 5.}

% =====================================================================
\section{Lab 2.3: Chọn kho cho năm việc của Olist}

Lab này chọn kho cho năm việc trên dữ liệu Olist. Hai việc đã có số đo từ Bài 2.

\textbf{Bước 1: tìm review theo chữ.} \code{LIKE}, dù chạy trên kho nào, vẫn quét mọi bình luận và chỉ khớp theo chuỗi ký tự. Khối dưới nạp 40.977 bình luận vào SQLite hai lần: một bảng thường để tìm bằng \code{LIKE}, và một bảng FTS5 có index đảo ngược.
% Kiểm: ngày 06/11/2026, DuckDB 1.5.6 một luồng, trung vị 7 lần đo sau một lần chạy nóng: LIKE '%não recebi%' trên cùng 40.977 bình luận mất 1,9 ms, SQLite LIKE 3,9 ms, FTS5 0,23 ms; LIKE của DuckDB phân biệt hoa thường nên chỉ thấy 1130 bình luận.

% chạy: L6
\begin{codeblock}
lakehouse/.venv/bin/python - <<'EOF'
import csv, sqlite3
F = 'lakehouse/source_sim/data/olist_order_reviews_dataset.csv'
binh_luan = [r['review_comment_message'] for r in csv.DictReader(open(F, newline=''))
             if r['review_comment_message']]
db = sqlite3.connect(':memory:')
db.execute('CREATE TABLE review (msg TEXT)')
db.execute('CREATE VIRTUAL TABLE review_fts USING fts5(msg)')  # index đảo ngược
for bang in ['review', 'review_fts']:
    db.executemany(f'INSERT INTO {bang} VALUES (?)', [(m,) for m in binh_luan])
dem = lambda sql, *p: db.execute(sql, p).fetchall()[0][0]
fts = 'SELECT count(*) FROM review_fts WHERE review_fts MATCH ?'
print('bình luận:', len(binh_luan))
like = "SELECT count(*) FROM review WHERE msg LIKE ?"
print("LIKE '%não recebi%':", dem(like, '%não recebi%'))
print('FTS5, cụm từ "nao recebi":', dem(fts, '"nao recebi"'))
print("LIKE '%atras%':", dem(like, '%atras%'))
print('FTS5, tiền tố atras*:', dem(fts, 'atras*'))
EOF
\end{codeblock}
\textit{Kết quả:}
% kết quả: L6
\begin{codeblock}
bình luận: 40977
LIKE '%não recebi%': 1696
FTS5, cụm từ "nao recebi": 2044
LIKE '%atras%': 481
FTS5, tiền tố atras*: 509
\end{codeblock}

\textit{Nếu lỗi:} \code{sqlite3.OperationalError: no such module: fts5} nghĩa là bạn không chạy bằng \repo{lakehouse/\allowbreak{}.venv/\allowbreak{}bin/\allowbreak{}python}.

FTS5 thấy 2.044 bình luận, còn \code{LIKE} thấy 1.696; 1.673 bình luận có mặt ở cả hai. \textit{Bộ tách từ} (tokenizer) mặc định của FTS5 là \code{unicode61}: nó không phân biệt hoa thường theo Unicode và bỏ dấu của chữ Latin. Trong 371 bình luận chỉ FTS5 thấy, 264 bình luận gõ ``nao recebi'' không dấu, và 95 bình luận viết hoa ``NÃO''. \code{LIKE} so từng ký tự và chỉ không phân biệt hoa thường với chữ ASCII: với nó, ``a'' và ``Ã'' đều khác ``ã'', nên nó bỏ sót cả hai nhóm. Ngược lại, FTS5 khớp theo từ nguyên vẹn, nên bỏ qua 23 bình luận như ``não recebido'', còn \code{LIKE} bắt được chúng vì nó khớp theo chuỗi con. Viết \code{"nao recebi"*} thì FTS5 thấy 2.081 bình luận, gồm cả 1.696 bình luận của \code{LIKE}.
% Nguồn: https://www.sqlite.org/fts5.html (unicode61 là tokenizer mặc định; remove_diacritics mặc định 1); https://www.sqlite.org/lang_expr.html (LIKE: "SQLite only understands upper/lower case for ASCII characters by default")
% Kiểm: ngày 06/11/2026, Python 3.10.20, SQLite 3.50.4: 1673 bình luận khớp cả hai; 371 bình luận chỉ FTS5 thấy (264 có "nao recebi" không dấu, 95 có "NÃO" viết hoa trong cụm từ, 12 dạng khác như "näo", "Ñão", "recebí"); 23 bình luận chỉ LIKE thấy ("não recebido" 12, "não recebir" 6, "não recebimento" 4, "não recebie" 1); FTS5 '"nao recebi"*' thấy 2081, chứa cả 1696 của LIKE và 2044 của cụm từ.

Bỏ dấu cũng có giá: tiền tố \code{atras*} khớp cả ``atrás'', nghĩa là ``trước đây'', chứ không chỉ ``atraso'' là trễ. Tìm theo chữ là một kiểu truy cập riêng, với bộ tách từ riêng, và index đảo ngược trả lời nhanh mà không quét. Trên máy viết bài, một câu FTS5 mất khoảng 0,2 ms, còn \code{LIKE} quét cả bảng mất khoảng 3,7 ms.
% Kiểm: ngày 06/11/2026, Python 3.10.20, SQLite 3.50.4: tập bình luận khớp atras* mà không khớp LIKE '%atras%' chỉ chứa các dạng của "atrás" (21 "atrás", 4 "atrás,", 2 "atrás.", 1 "Atrás"); trung vị 7 lần đo: FTS5 "nao recebi" 0,21 ms, LIKE '%não recebi%' 3,7 ms.

\textbf{Bước 2: điền bảng chọn kho.} Chép bảng dưới vào \code{lab\_\allowbreak{}m02/chon\_\allowbreak{}kho.md} và điền bốn dòng còn lại. Mỗi dòng cần một bằng chứng là con số: số đo ở Bài 2, số đếm ở bước 1, hoặc số dòng và nhịp chạy trong repo.

\begin{xltabular}{\textwidth}{@{}P{2.4cm} P{2.4cm} P{2.6cm} P{2.2cm} L@{}}
\toprule
\thd{việc} & \thd{ai đọc, nhịp đọc} & \thd{kiểu truy cập} & \thd{kho chọn} & \thd{bằng chứng} \\
\midrule
\endhead
\bottomrule
\endfoot
API tra đơn theo \code{order\_\allowbreak{}id} & app, nhiều lần mỗi giây & tra theo khoá & kho theo dòng có index, như PostgreSQL & Bài 2: SQLite tra nhanh hơn DuckDB khoảng 70 lần \\
dashboard doanh thu & BA, mỗi sáng & & & \\
tìm review theo chữ & đội chăm sóc khách, khi cần & & & \\
metric độ tươi & người trực pipeline, mỗi 15 phút & & & \\
tra toạ độ theo zip & dashboard bản đồ & & & \\
\end{xltabular}

Hai gợi ý. Metric độ tươi là một con số cho mỗi bảng sau mỗi lần Bronze chạy, tức 96 điểm mỗi bảng mỗi ngày. Tra toạ độ theo zip là tra theo khoá trên một bảng nhỏ: \repo{semantic.zip\_\allowbreak{}geo} có 19.010 zip. Hỏi thêm ai đọc nó: \repo{core.customers} đã có sẵn toạ độ cho dashboard, còn một API cho app thì cần kho khác.
% Nguồn: lakehouse/airflow/dags/olist_bronze_ingest.py dòng 35 (schedule "*/15 * * * *", 24 x 4 = 96 lần mỗi ngày); lakehouse/docs/data-model.md mục 10 (semantic.zip_geo 19.010 dòng) và mục 11 (customer_lat, customer_lng)

\textbf{Bước 3: viết lý do ngược.} Với mỗi dòng, thêm cột thứ sáu ``lý do ngược'' vào \code{chon\_\allowbreak{}kho.md} và ghi một câu: chuyện gì xảy ra nếu chọn kho của dòng bên cạnh. Ví dụ: đặt API tra đơn lên DuckDB thì mỗi lần tra tốn khoảng 160 µs thay vì 2 µs, đã gồm chi phí cố định của mỗi câu.

\hoanthanh{khối ở bước 1 in đúng 40.977, 1.696, 2.044, 481 và 509; \code{chon\_\allowbreak{}kho.md} có đủ năm dòng, mỗi dòng có kiểu truy cập, kho chọn, một con số làm bằng chứng và một câu lý do ngược.}

% =====================================================================
\section{Thử thách kỹ năng}

\begin{thuthach}
Không có hướng dẫn từng bước. Đội vận hành Olist muốn hai thứ: báo cáo doanh thu theo ngày cho ban giám đốc trước 08:00 mỗi sáng, và cảnh báo trong vòng 10 phút khi một đơn có nguy cơ giao trễ. Viết một bản ghi quyết định kiến trúc dài một tới hai trang cho cả hai nhu cầu.

Bản ghi gồm: thẻ yêu cầu đủ năm phần cho mỗi nhu cầu; kho chọn theo kiểu truy cập, kèm một số đo; batch hay streaming; ETL hay ELT; tự viết hay mua phần nạp dữ liệu. Nêu ít nhất hai đánh đổi và một rủi ro.
\hoanthanh{mỗi nhu cầu có đủ người dùng, grain, metric, mốc cắt và độ tươi; độ tươi được so với lịch thật của repo, và bản ghi nói lịch hiện tại có đạt không; mỗi lựa chọn có một lý do dẫn về một bài trong module hoặc một số đo; nếu chọn Lambda thì có cách tránh viết logic hai lần.}

\textbf{Mang vào capstone:} thẻ yêu cầu của Lab 2.2 là đầu vào của mốc capstone 1 ở cuối Module 3, nơi bạn dựng star schema Gold cho đúng năm câu hỏi đó.
\end{thuthach}
% Nguồn: mốc 1 của lộ trình: bus matrix, grain và star schema Gold trên DuckDB cho năm câu hỏi

% =====================================================================
\section{Tự kiểm tra}

Trả lời bằng lời của bạn trước, rồi mới đọc gợi ý.

\begin{enumerate}
  \item Năm phần của một thẻ yêu cầu là gì? Phần nào hay bị quên nhất?\par
  \textit{Gợi ý:} người dùng, grain, định nghĩa metric, mốc cắt, độ tươi. Hay quên ngày chốt trong mốc cắt, vì số đơn tháng 2/2017 đã giao còn tăng từ 907 lên 1.653.
  \item Vì sao ``đơn giao sau 01/03/2017'' có ít nhất ba đáp án?\par
  \textit{Gợi ý:} từ đầu hay từ cuối ngày 01/03 cho 746 hay 673; tính cả đơn chưa giao thì ra 873.
  \item Làm sao biết các mốc của Olist là giờ Brazil khi file không ghi múi giờ?\par
  \textit{Gợi ý:} 85 review mang 01:00 đúng hai ngày Brazil bắt đầu giờ mùa hè, và ít đơn nhất lúc 4 tới 5 giờ sáng.
  \item Vì sao DuckDB tra theo khoá chậm hơn SQLite mà gộp lại nhanh hơn?\par
  \textit{Gợi ý:} khi tra, SQLite có B-tree dẫn thẳng tới dòng, còn DuckDB quét cả cột \code{order\_\allowbreak{}id} và trả chi phí cố định mỗi câu. Khi gộp, DuckDB tính theo lô trên kiểu \code{TIMESTAMP}, còn SQLite xử lý từng dòng và đọc mốc thời gian từ chữ; đọc ít cột giúp ít khi dữ liệu nằm trong RAM.
  \item Độ tươi khác độ trễ của API ở đâu?\par
  \textit{Gợi ý:} độ tươi đo tuổi của dữ liệu mới nhất; pipeline trễ không sinh lỗi nào, dashboard vẫn nhanh nhưng số cũ.
  \item Khi luật doanh thu đổi, vì sao ELT chỉ cần chạy lại SQL còn ETL phải đọc lại nguồn?\par
  \textit{Gợi ý:} ELT giữ dữ liệu thô trong kho; kho của ETL chỉ có bảng kết quả, không còn \code{freight\_\allowbreak{}value}.
  \item Medallion có thay cho mô hình dữ liệu ở gold không?\par
  \textit{Gợi ý:} không; Databricks gọi medallion là thực hành được khuyên, không bắt buộc, và gold vẫn cần dimensional model.
\end{enumerate}

% =====================================================================
\section{Đọc thêm}

Đọc sau khi làm xong các lab, để gắn tên chuẩn cho những gì bạn đã tự làm.

\begin{itemize}
  \item \href{https://www.oreilly.com/library/view/fundamentals-of-data/9781098108298/}{Fundamentals of Data Engineering}, sách của Joe Reis và Matt Housley, bản in 2022: đọc chương 2 về vòng đời và undercurrent, và mục Build Versus Buy ở chương 4, khoảng 2 giờ. Đây là khung mà Bài 6 và Bài 8 dùng.
  % Nguồn: http://web.archive.org/web/20260330095128/https://www.oreilly.com/library/view/fundamentals-of-data/9781098108298/ (bản lưu của trang O'Reilly, trang gốc trả 403 cho công cụ tải)
  \item \href{https://martin.kleppmann.com/2026/03/24/designing-data-intensive-applications-2e.html}{Designing Data-Intensive Applications, bản 2}, sách của Martin Kleppmann và Chris Riccomini, 2026: ở chương 4, đọc mục ``Column-Oriented Storage'' và mục ``Query Execution: Compilation and Vectorization'', khoảng 1,5 giờ, để hiểu vì sao DuckDB gộp nhanh ở Bài 2.
  % Kiểm: curl ngày 06/11/2026 bản lưu web.archive.org/web/20260616232347 của trang O'Reilly 9781098119058: chương 4 Storage and Retrieval, phần Data Storage for Analytics, có Column-Oriented Storage và Query Execution: Compilation and Vectorization.
  \item \href{https://sre.google/workbook/data-processing/}{SRE Workbook, chương Data Processing Pipelines}, sách của Google SRE: đọc phần SLO độ tươi và độ đúng, khoảng 45 phút. Bản web không ghi ngày; đây là nguồn của cách viết độ tươi ở Bài 1.
  \item \href{https://martinfowler.com/articles/data-mesh-principles.html}{Data Mesh Principles and Logical Architecture}, bài viết của Zhamak Dehghani, 03/12/2020: đọc phần bốn nguyên tắc, khoảng 45 phút, nếu bạn muốn hiểu Bài 7 kỹ hơn một đoạn.
  \item \href{https://learn.microsoft.com/en-us/azure/architecture/databases/guide/big-data-architectures}{Azure Architecture Center: Big data architectures}, bản ngày 12/09/2025: đọc phần Lambda, Kappa và Lakehouse, khoảng 30 phút, để thấy vì sao Lambda phải viết logic hai lần còn Kappa thì không.
  \item \href{https://docs.databricks.com/aws/en/lakehouse/medallion}{Databricks: Medallion architecture}, bản cập nhật 29/10/2026: cả trang, khoảng 15 phút, để so ba tầng của Databricks với các namespace của repo.
  \item \href{https://www.getdbt.com/what-is-analytics-engineering}{What is analytics engineering?}, bài viết của Claire Carroll, 26/03/2019: cả bài, khoảng 15 phút, để thấy vai analytics engineer ra đời từ đâu.
\end{itemize}
% Kiểm: tổng giờ đọc thêm 2 + 1,5 + 0,75 + 0,75 + 0,5 + 0,25 + 0,25 = 6 giờ, khớp dòng Giờ của thẻ module.

% =====================================================================
\section{Kế hoạch theo tuần}

Module này có 13 giờ và học song song với Module 3 được, nên tuần 2 còn 5 giờ cho Module 3.

\begin{xltabular}{\textwidth}{@{}P{1.4cm} L P{3.6cm} P{1.2cm}@{}}
\toprule
\thd{tuần} & \thd{việc lõi} & \thd{thêm nếu còn giờ} & \thd{giờ} \\
\midrule
\endhead
\bottomrule
\endfoot
1 & Vấn đề thật và Bài 1, 2 giờ; Lab 2.1, 1 giờ; Lab 2.2, 2,5 giờ; Bài 2, 2 giờ; Lab 2.3, 1 giờ; Bài 3, 0,5 giờ & chương 2 của Fundamentals of Data Engineering & 9 \\
2 & Bài 4, 0,5 giờ; Bài 5, 0,75 giờ; Bài 6 và bài kiểm 15 phút, 0,5 giờ; Bài 7, 0,25 giờ; Bài 8, 0,5 giờ; Tự kiểm tra, 0,25 giờ; thử thách kỹ năng, 1,25 giờ & phần còn lại của Đọc thêm & 4 \\
\end{xltabular}


% =====================================================================
\modulechapter{3}{Mô hình dữ liệu phân tích: thiết kế Gold cho Olist}
\moduleinfo{Cơ bản tới trung cấp}{30 giờ}{Module 1}

\begin{xltabular}{\textwidth}{@{}P{3.4cm} L@{}}
\toprule
\thd{thẻ module} & \thd{nội dung} \\
\midrule
\endhead
\bottomrule
\endfoot
Giờ & 30 giờ lõi: ba tuần 9 giờ và một tuần 3 giờ, theo nhịp 9 giờ module mỗi tuần; thêm khoảng 2,5 giờ cho các bài trong Đọc thêm, và khoảng 30 giờ nếu học cả khoá 4 của DeepLearning.AI \\
Dùng trong repo & 9 file CSV Olist trong \repo{source\_\allowbreak{}sim/\allowbreak{}data/}; \repo{docs/\allowbreak{}data-model.md} mục 5, 6, 9, 10 và 11; server \repo{source\_\allowbreak{}sim/\allowbreak{}app.py} chạy như một process ở cổng 8010 trong Lab 3.4, không cần Docker; tên bảng và tên cột của \repo{core.orders}, \repo{core.customers}, \repo{core.order\_\allowbreak{}items}, \repo{core.products}, \repo{core.sellers} và \repo{semantic.zip\_\allowbreak{}geo} cho Lab 3.5 \\
Phiên bản & DuckDB 1.5.6, pyarrow 25.0.1, FastAPI 0.146.0 và uvicorn 0.54.0 trong \repo{lakehouse/.venv}, Python 3.10.20; code cũng chạy trên DuckDB 2.0.0. Docs: \href{https://duckdb.org/docs/1.5/sql/statements/merge_into}{DuckDB 1.5: MERGE INTO}, \href{https://duckdb.org/docs/2.0/sql/statements/merge_into.html}{DuckDB 2.0: MERGE INTO}, \href{https://arrow.apache.org/docs/25.0/python/generated/pyarrow.Table.html}{pyarrow 25.0: Table} \\
Lướt được & Không nên bỏ bài nào. Bài 1 và Bài 10 đọc nhanh được nếu bạn đã quen chuẩn hoá schema và TDD. Kiểm 15 phút: chạy khối Làm thử của Bài 1 và khối Làm thử của Bài 10, rồi giải thích vì sao tiền kiểu \code{DOUBLE} làm hàng trăm nhóm lệch, và vì sao test đối chiếu đỏ khi chỉ thêm hai dòng. Giải thích được cả hai thì đọc nhanh hai bài đó \\
\end{xltabular}
% Kiểm: lakehouse/.venv/bin/python -c "import duckdb, pyarrow, fastapi, uvicorn" in 1.5.6, 25.0.1, 0.146.0, 0.54.0; Python 3.10.20; 06/11/2026.
% Kiểm: trang docs DuckDB 1.5 và 2.0 cho MERGE INTO trả 200 ngày 06/11/2026; chỉ trang 2.0 có câu "cannot affect the same target row more than once".
% =====================================================================
%  Bản 2 của Module 3 theo kiểu tutorial.
%  Mỗi khối code có dòng "% chạy: <id>"; output thật nằm sau "% kết quả: <id>".
%  Output chạy lại ngày 06/11/2026 bằng lakehouse/.venv/bin/python, Python 3.10.20,
%  DuckDB 1.5.6, từ một thư mục chứa symlink lakehouse trỏ tới repo. Mọi khối cũng chạy
%  được trên DuckDB 2.0.0 với Python 3.11; output giống hệt, trừ số byte ở Bài 8.
%  Khối ghi "cần-lời-giải" chỉ chạy được sau khi người đọc viết code của mình; lời giải
%  tham khảo dùng để kiểm output nằm ở solutions/m03/.
% =====================================================================

Module 1 dạy bạn hỏi dữ liệu bằng SQL trên chính schema của app. Module này dạy bạn dựng lại dữ liệu đó thành một dạng khác, gọi là \textit{star schema}, để câu hỏi phân tích dễ viết và số liệu không sai.

Bạn học bằng cách làm trên bộ Olist. Mỗi khái niệm đi kèm một lỗi thật đo được trên dữ liệu, một đoạn code chạy được, và chỗ nó giống hoặc khác thứ bạn đã quen khi làm phần mềm. Sản phẩm cuối là một tầng Gold chạy thử cho Olist trên DuckDB: một star schema có bộ test tự động, một accumulating snapshot dựng từ API của repo, và năm mart cho năm câu hỏi của \repo{docs/\allowbreak{}data-model.md}, kèm phương trình đối soát. Đó là mốc capstone 1.

\begin{ghichu}
\textbf{Cách chạy code trong module này.}
\begin{itemize}
  \item Lời giải mẫu nằm trong \repo{solutions/m03/}. Đọc sau khi tự làm, để đối chiếu.
  \item Mọi khối code chạy từ thư mục cha của \repo{lakehouse/} bằng \repo{lakehouse/.venv/bin/python}, trong đó có DuckDB 1.5.6. DuckDB là database phân tích chạy ngay trong process Python, không cần server, nên module này không cần Docker.
  \item Lộ trình ghim DuckDB 2.0.0, nhưng bản này cần Python 3.11 trở lên. Venv của \repo{lakehouse} dùng Python 3.10 nên có DuckDB 1.5.6, và mọi output trong module là của 1.5.6. Theo thông báo ra mắt bản 2.0, DuckDB 1.5 được hỗ trợ tới tháng 3/2027. Trên 2.0.0, mọi khối cho cùng output, trừ số byte ở Bài 8 lệch một chút.
  % Nguồn: https://pypi.org/pypi/duckdb/2.0.0/json (requires_python >=3.11.0); https://duckdb.org/2026/10/21/announcing-duckdb-200
  \item Dữ liệu là 9 file CSV Olist trong \repo{lakehouse/source\_\allowbreak{}sim/data/}. Nếu bạn để CSV ở chỗ khác, ví dụ \code{olist/} như ở Module 1, hãy sửa biến \code{D} ở đầu mỗi khối.
  \item Lab 3.4 chạy server giả lập nguồn của repo, \repo{source\_\allowbreak{}sim}, như một process Python ở cổng 8010, không cần Docker. Mọi file mà các lab tạo ra nằm trong \code{lab\_\allowbreak{}m03/}; không khối nào ghi vào \repo{lakehouse/}.
  \item Phần ``Kết quả'' dưới mỗi khối là output thật, chạy lại ngày 06/11/2026. Nếu máy bạn in ra số khác, dừng lại tìm lý do trước khi đọc tiếp; số byte ở Bài 8 là ngoại lệ, như mục trên.
\end{itemize}
\end{ghichu}

\begin{ghichu}
\textbf{Cú pháp DuckDB mà Module 1 chưa dạy.} Module 1 đã dạy \code{NOT EXISTS}, \code{FILTER} và \code{QUALIFY}. Các khối code ở đây dùng thêm:
\begin{itemize}
  \item \code{A ANTI JOIN B USING (k)} giữ các dòng của \code{A} không có dòng khớp ở \code{B}, giống \code{WHERE NOT EXISTS}.
  \item \code{GROUP BY ALL} nhóm theo mọi cột của \code{SELECT} không nằm trong aggregate. \code{ORDER BY ALL} sắp theo mọi cột, từ trái sang phải.
  \item \code{A EXCEPT B} trả các dòng có ở \code{A} mà không có ở \code{B}, sau khi bỏ dòng trùng. \code{SELECT * EXCLUDE (c)} chọn mọi cột trừ cột \code{c}.
  \item \code{range(a, b, bước) r(d)} sinh một bảng tên \code{r} có cột \code{d}, chạy từ \code{a} tới trước \code{b}. Giống \code{range()} của Python, điểm cuối không được tính.
\end{itemize}
% Nguồn: https://duckdb.org/docs/2.0/sql/query_syntax/from (mục Semi and Anti Joins); https://duckdb.org/docs/2.0/sql/query_syntax/groupby; https://duckdb.org/docs/2.0/sql/query_syntax/orderby; https://duckdb.org/docs/2.0/sql/query_syntax/setops
% Nguồn: https://duckdb.org/docs/2.0/sql/functions/timestamp (range: half open range); https://docs.python.org/3/library/stdtypes.html (range: r[i] < stop)
\end{ghichu}

% =====================================================================
\section{Vấn đề thật: doanh thu tăng 26,9\% mà không ai bán thêm gì}

Bạn là kỹ sư phần mềm vừa chuyển sang team dữ liệu. Một BA, tức người phân tích nghiệp vụ, nhắn: ``Cho mình tổng tiền khách đã trả, chia theo danh mục sản phẩm.''

Tiền thanh toán nằm ở bảng \code{order\_\allowbreak{}payments}, còn danh mục gắn với sản phẩm của từng dòng hàng trong \code{order\_\allowbreak{}items}. Bạn làm việc quen tay: join hai bảng theo \code{order\_\allowbreak{}id} rồi cộng \code{payment\_\allowbreak{}value}. Trước khi chia theo danh mục, bạn thử xem con số tổng.

% chạy: A
\begin{codeblock}
lakehouse/.venv/bin/python - <<'EOF'
import duckdb
D = 'lakehouse/source_sim/data/'
con = duckdb.connect()
for t in ['order_items', 'order_payments']:
    con.sql(f"CREATE VIEW {t} AS SELECT * FROM '{D}olist_{t}_dataset.csv'")
one = lambda q: con.sql(q).fetchall()[0][0]
dung = one("SELECT sum(CAST(payment_value AS DECIMAL(14,2))) FROM order_payments")
sai = one("""SELECT sum(CAST(p.payment_value AS DECIMAL(14,2)))
             FROM order_payments p JOIN order_items i USING (order_id)""")
print('chỉ cộng thanh toán:', dung)
print('join rồi mới cộng  :', sai, f'(cao hơn {sai / dung - 1:.1%})')
con.sql("""SELECT p.payment_value, count(*) AS so_dong_sau_join,
                  sum(CAST(p.payment_value AS DECIMAL(14,2))) AS tong_sau_join
           FROM order_payments p JOIN order_items i USING (order_id)
           WHERE order_id = '8272b63d03f5f79c56e9e4120aec44ef'
           GROUP BY 1""").show()
EOF
\end{codeblock}
\textit{Kết quả:}
% kết quả: A
\begin{codeblock}
chỉ cộng thanh toán: 16008872.12
join rồi mới cộng  : 20308134.71 (cao hơn 26.9%)
┌───────────────┬──────────────────┬───────────────┐
│ payment_value │ so_dong_sau_join │ tong_sau_join │
│    double     │      int64       │ decimal(38,2) │
├───────────────┼──────────────────┼───────────────┤
│        196.11 │               21 │       4118.31 │
└───────────────┴──────────────────┴───────────────┘
\end{codeblock}

Con số sau join cao hơn con số đúng 4.299.262,59 BRL. Không có lỗi, không có cảnh báo nào: database làm đúng điều bạn viết.

Bảng nhỏ ở cuối cho thấy lý do. Đơn \code{8272b63d...} trả tiền một lần, 196,11 BRL, nhưng có 21 dòng hàng. Join lặp khoản tiền đó cho từng dòng hàng, nên nó bị cộng 21 lần thành 4.118,31 BRL.

Lỗi này có tên là \textit{fan-out}, và nó là một lỗi rất hay gặp khi đọc thẳng schema của app để phân tích. Module này dạy cách thiết kế bảng để lỗi này khó xảy ra, và cách viết test để bắt nó khi nó vẫn xảy ra. Câu hỏi của cả module là: thiết kế tầng Gold cho Olist thế nào để mỗi con số chỉ có một định nghĩa, và một test đỏ ngay khi tổng tiền lệch dù chỉ 0,01 BRL?

% =====================================================================
\section{Từ SWE sang DE}

Bảng dưới nối thứ bạn đã biết với tên gọi trong data engineering, theo thứ tự các bài học. Cột cuối quan trọng nhất: so sánh nào cũng có chỗ gãy, và mỗi bài học sẽ chỉ ra chỗ gãy đó trên dữ liệu thật.

\begin{xltabular}{\textwidth}{@{}P{3.2cm} P{2.8cm} L L@{}}
\toprule
\thd{bạn đã biết (SWE)} & \thd{trong DE gọi là} & \thd{giống ở đâu} & \thd{khác ở đâu} \\
\midrule
\endhead
\bottomrule
\endfoot
schema do ORM và migration tạo ra & schema OLTP, đã chuẩn hoá & mỗi sự thật nằm ở một chỗ, ghi nhanh và an toàn & câu hỏi phân tích phải join nhiều bảng và dễ nhân dòng \\
read model trong CQRS & star schema ở tầng gold & một mô hình riêng cho việc đọc, dựng từ dữ liệu ghi & dựng theo lô từ nhiều nguồn, trả lời cả câu hỏi chưa ai đặt ra \\
quyết định một resource của API là gì & grain & khoá suy ra từ định nghĩa & grain là một câu nghiệp vụ viết trước khi chọn cột, và trong lakehouse không ai ép khoá giúp bạn \\
lỗi gộp nhiều aggregation trong \code{annotate()} của Django & fan-out & cùng cơ chế: join nhân dòng rồi mới cộng & \code{Count(distinct=True)} chữa được phép đếm; với tổng tiền thì \code{sum(DISTINCT)} sai, phải gộp từng bảng về grain trước khi join \\
bảng \code{events} hay audit log chỉ thêm & bảng fact & mỗi dòng là một chuyện đã xảy ra & có measure để cộng và khoá tới dimension; loại accumulating snapshot lại được sửa \\
counter, gauge, tỉ lệ lỗi trong monitoring & measure cộng được, cộng được một phần, không cộng được & cùng luật cộng, và luật đó phải ghi trong đặc tả & counter là số cộng dồn, nên phải lấy phần tăng giữa hai lần đo rồi mới cộng; measure của fact đã là con số của từng sự kiện \\
enum dùng chung giữa các service, ma trận service và capability & conformed dimension, bus matrix & một định nghĩa cho mọi nơi; nhìn một bảng biết ai dùng gì & dimension là bảng có dữ liệu và lịch sử, đổi nó làm đổi số liệu quá khứ; bus matrix còn quyết định thứ tự xây, mỗi lần một quy trình \\
tách bảng cho khỏi lặp chuỗi, hay nhúng sẵn như một document & star, snowflake, One Big Table & cùng câu hỏi: tách ra hay chép vào & Parquet đã lưu mỗi giá trị một lần trong từ điển của cột, nên tách không tiết kiệm; chép hết vào một bảng thì đổi tên \code{bed\_\allowbreak{}bath\_\allowbreak{}table} phải ghi lại 11.115 dòng (Bài 8) \\
câu \code{UPDATE} ghi đè; S3 Versioning & SCD Type 1, SCD Type 2 & giá trị mới nhất; hoặc nhiều version với một version hiện hành & Type 1 làm sai số liệu quá khứ; Type 2 chỉ thấy thay đổi mà pipeline kịp nhìn, và fact phải trỏ đúng version \\
bảng nối của \code{ManyToManyField}, bảng lịch sử chỉ thêm & Data Vault: hub, link, satellite & tách định danh, quan hệ và thuộc tính ra các bảng riêng & vault phục vụ việc nạp và kiểm toán; người phân tích vẫn cần một star dựng bên trên (Bài 11) \\
state machine của đơn: cột \code{status} bị \code{UPDATE} qua \code{created}, \code{approved}, \code{shipped} & accumulating snapshot & một dòng mỗi đơn, đi dần qua các mốc & app chỉ giữ trạng thái hiện tại, snapshot giữ mỗi mốc một cột; engine nào cũng cần lọc version trùng trong lô và chặn version cũ; trạng thái trung gian của mock là giả lập (Lab 3.4) \\
unit test, checksum, \code{Collectors.\allowbreak{}partitioningBy} của Java & data test, phương trình đối soát & assert tự động; hai nửa của một phép chia cộng lại bằng tổng & chạy trên dữ liệu thật nên đỏ có thể do nguồn; phương trình chỉ chứng minh không mất đồng nào, không chứng minh bộ lọc đúng (Bài 10, Bài 12) \\
\end{xltabular}

\textbf{Thói quen SWE gây hại nhất ở module này là tin rằng câu SQL chạy không lỗi thì con số đúng.} Fan-out, join làm rơi dòng, version cũ ghi đè version mới đều chạy mà không báo gì; chỉ test đối chiếu với nguồn mới bắt được chúng.

% =====================================================================
\section{Bạn sẽ làm được}

\begin{itemize}
  \item Viết grain của mỗi fact trong một câu, suy ra khoá, chứng minh khoá không trùng bằng một câu SQL, và phân loại measure kèm lý do.
  \item Dựng star schema Olist trên DuckDB gồm bốn dimension và hai fact, với bộ test viết trước trả về 0 dòng lỗi, rồi chứng minh fan-out trên chính star đó.
  \item Cài SCD Type 2 bằng SQL với khoảng hiệu lực nửa mở, chạy lại không sinh thêm dòng, và gắn mỗi dòng hàng với đúng một version người bán.
  \item Dựng accumulating snapshot của đơn hàng từ 391.551 version mà \repo{source\_\allowbreak{}sim} phát qua API, sao cho áp lại hay áp lộn thứ tự đều không đổi dòng nào.
  \item Định nghĩa metric doanh thu trong mart cho năm câu hỏi của repo, và chứng minh bằng phương trình đối soát rằng không đồng nào bị mất hay bị cộng hai lần.
  \item Tự chọn grain cho fact review và viết phương trình đối soát cho nó, không cần hướng dẫn từng bước.
\end{itemize}

% =====================================================================
\section{Bài 1: Schema của app và schema cho phân tích}

\textbf{Vấn đề.} Bộ Olist có 9 bảng, đúng kiểu một app thương mại điện tử: đơn, dòng hàng, thanh toán, đánh giá, khách, sản phẩm, người bán, bảng dịch tên danh mục và toạ độ theo zip. Câu hỏi ``doanh thu theo danh mục tiếng Anh và tháng'' cần 3 join qua 4 bảng và hai luật: dịch tên danh mục, cắt thời điểm mua về đầu tháng. Mỗi người viết câu này sẽ tự viết hai luật đó, và số của hai người có thể lệch nhau.

\textbf{Giải thích.} Schema của app là \textit{OLTP}, viết tắt của online transaction processing. Đó là database phục vụ app đang chạy: nhận rất nhiều giao dịch nhỏ, ưu tiên ghi nhanh và ghi đúng.

Schema OLTP được \textit{chuẩn hoá}: dữ liệu được tách ra nhiều bảng để mỗi sự thật chỉ nằm ở một chỗ, sửa một chỗ là xong. Mức hay gặp là 3NF, trong đó mỗi cột chỉ phụ thuộc vào khoá của bảng chứa nó.

Phân tích cần thứ ngược lại: quét hàng triệu dòng rồi gộp, với ít join và các luật được viết sẵn một lần. \textit{Star schema} làm đúng việc đó. Nó có một bảng \textit{fact} ở giữa, mỗi dòng là một sự kiện kèm các con số đo được. Xung quanh là các bảng \textit{dimension}, mỗi bảng mô tả một thứ như khách, sản phẩm hay ngày.

Mỗi dimension cách fact đúng một join. Dimension còn được làm phẳng sẵn: \code{dim\_\allowbreak{}product} chứa luôn tên danh mục tiếng Anh, nên người đọc không cần biết bảng dịch tồn tại.

\begin{codeblock}
Schema của app (OLTP):  ──<  nghĩa là một dòng bên trái, nhiều dòng bên phải

customers ──< orders ──< order_items >── products >── category_translation
                │             │
                ├──< order_payments        sellers ──< order_items
                └──< order_reviews

Star schema:  fact ở giữa, mỗi dimension cách fact đúng một join

                       dim_date
                          │
      dim_customer ── fct_order_items ── dim_product
                          │
                      dim_seller
\end{codeblock}

\begin{ghichu}
\textbf{So với SWE.} Schema OLTP chính là thứ ORM và migration của bạn tạo ra. Tài liệu Django gọi migration là cách đưa các thay đổi trên model vào schema của database.
% Nguồn: https://docs.djangoproject.com/en/6.1/topics/migrations/

Star schema giống \textit{read model} trong CQRS: Martin Fowler mô tả CQRS là dùng một model để cập nhật và một model khác để đọc. Fowler còn có khái niệm \textit{reporting database}: một bản sao dữ liệu vận hành đặt trong schema khác, không cần chuẩn hoá vì chỉ để đọc, và thường được nạp bằng một job chạy đêm. Star schema là một reporting database có luật thiết kế rõ ràng.
% Nguồn: https://martinfowler.com/bliki/CQRS.html
% Nguồn: https://martinfowler.com/bliki/ReportingDatabase.html

\textbf{Chỗ so sánh gãy.} Read model của CQRS thường được cập nhật sau mỗi lệnh ghi, với độ trễ nhỏ. Star schema phục vụ cả những câu hỏi chưa ai đặt ra, được dựng theo lô từ nhiều nguồn, và giữ lại lịch sử mà app đã ghi đè.
\end{ghichu}

\subsection{Làm thử}

Khối dưới dựng một star tối giản gồm \code{dim\_\allowbreak{}product}, \code{dim\_\allowbreak{}date} và \code{fct\_\allowbreak{}order\_\allowbreak{}items}, rồi chạy cùng một câu hỏi trên hai schema. Nó chạy hai lần: lần đầu tiền có kiểu \code{DECIMAL}, lần sau có kiểu \code{DOUBLE}. Khoá của \code{dim\_\allowbreak{}date} ở đây tạm là chính giá trị ngày; Bài 6 chuyển sang khoá số dạng \code{yyyymmdd}.

% chạy: B
\begin{codeblock}
lakehouse/.venv/bin/python - <<'EOF'
import duckdb
D = 'lakehouse/source_sim/data/'
con = duckdb.connect()
for t in ['orders', 'order_items', 'products']:
    con.sql(f"CREATE VIEW {t} AS SELECT * FROM '{D}olist_{t}_dataset.csv'")
con.sql(f"CREATE VIEW category_translation AS "
        f"SELECT * FROM '{D}product_category_name_translation.csv'")
one = lambda q: con.sql(q).fetchall()[0][0]

for money in ['DECIMAL(14,2)', 'DOUBLE']:
    # Star tối giản: luật dịch tên và luật cắt tháng nằm trong dimension
    con.sql("""CREATE OR REPLACE TABLE dim_product AS
      SELECT p.product_id, coalesce(t.product_category_name_english,
                                    p.product_category_name, 'unknown') AS category
      FROM products p
      LEFT JOIN category_translation t USING (product_category_name)""")
    con.sql("""CREATE OR REPLACE TABLE dim_date AS
      SELECT CAST(d AS DATE) AS date_key,
             CAST(date_trunc('month', d) AS DATE) AS month
      FROM range(DATE '2016-01-01', DATE '2019-01-01', INTERVAL 1 DAY) r(d)""")
    con.sql(f"""CREATE OR REPLACE TABLE fct_order_items AS
      SELECT i.order_id, i.order_item_id, i.product_id,
             CAST(o.order_purchase_timestamp AS DATE) AS date_key,
             CAST(i.price AS {money}) AS price
      FROM order_items i JOIN orders o USING (order_id)""")
    # Cùng một câu hỏi: doanh thu theo danh mục tiếng Anh và tháng
    oltp = f"""SELECT coalesce(t.product_category_name_english,
                              p.product_category_name, 'unknown') AS category,
                     CAST(date_trunc('month', o.order_purchase_timestamp) AS DATE)
                       AS month,
                     sum(CAST(i.price AS {money})) AS revenue
              FROM order_items i
              JOIN orders o ON o.order_id = i.order_id
              JOIN products p ON p.product_id = i.product_id
              LEFT JOIN category_translation t
                ON t.product_category_name = p.product_category_name
              GROUP BY 1, 2"""
    star = """SELECT p.category, d.month, sum(f.price) AS revenue
              FROM fct_order_items f
              JOIN dim_product p USING (product_id)
              JOIN dim_date d USING (date_key)
              GROUP BY 1, 2"""
    rows = one(f"SELECT count(*) FROM ({star})")
    a_not_b = one(f"SELECT count(*) FROM (({oltp}) EXCEPT ({star}))")
    b_not_a = one(f"SELECT count(*) FROM (({star}) EXCEPT ({oltp}))")
    print(money, '| số dòng:', rows, '| có nhóm lệch:', a_not_b + b_not_a > 0)
con.sql("""SELECT p.category, d.month,
                  sum(CAST(f.price AS DECIMAL(14,2))) AS revenue
           FROM fct_order_items f JOIN dim_product p USING (product_id)
           JOIN dim_date d USING (date_key)
           GROUP BY 1, 2 ORDER BY revenue DESC LIMIT 3""").show()
EOF
\end{codeblock}
\textit{Kết quả:}
% kết quả: B
\begin{codeblock}
DECIMAL(14,2) | số dòng: 1283 | có nhóm lệch: False
DOUBLE | số dòng: 1283 | có nhóm lệch: True
┌───────────────┬────────────┬───────────────┐
│   category    │   month    │    revenue    │
│    varchar    │    date    │ decimal(38,2) │
├───────────────┼────────────┼───────────────┤
│ watches_gifts │ 2018-05-01 │     123872.66 │
│ health_beauty │ 2018-08-01 │     120803.94 │
│ health_beauty │ 2018-06-01 │     107908.82 │
└───────────────┴────────────┴───────────────┘
\end{codeblock}

Với \code{DECIMAL}, hai cách cho cùng 1.283 nhóm, và \code{EXCEPT} theo cả hai chiều không còn dòng nào: hai schema trả lời giống hệt nhau. Câu trên star schema chỉ có 2 join và không còn luật nào, vì luật đã nằm trong dimension.

Đó là nghĩa của chữ ``nhất quán'' trong phân tích: mọi BA nhận cùng một định nghĩa, vì định nghĩa chỉ được viết một lần.

\begin{luuy}
\textbf{Bẫy thường gặp.}
\begin{itemize}
  \item \textbf{Để tiền ở kiểu \code{DOUBLE}.} Khoảng 680 trên 1.283 nhóm lệch ở vài chữ số cuối, vì cộng số thực theo thứ tự khác nhau cho kết quả hơi khác. Luật ``không dùng float cho tiền'' của SWE vẫn đúng.
  % Kiểm: chạy lại khối B với tiền DOUBLE, 10 luồng, DuckDB 1.5.6, 06/11/2026: bốn lần chạy cho 689, 674, 689, 689 nhóm lệch; tổng star của watches_gifts 2018-05 là 123872.65999999979 ở ba lần, 123872.65999999955 ở một lần; OLTP luôn 123872.66000000009. Ngày 05/11/2026 với SET threads = 1: cả hai 123872.65999999955, lệch 0.
  \item \textbf{Tin số nhóm lệch là cố định.} Nó đổi theo số luồng và giữa các lần chạy, nên khối ở trên chỉ in có lệch hay không. Với \code{SET threads = 1}, thứ tự cộng cố định nên hết lệch.
  % Kiểm: khối B thêm SET threads = 1 ngay sau connect, chạy ba lần ngày 06/11/2026, DuckDB 1.5.6: dòng DOUBLE in "có nhóm lệch: False" cả ba lần.
  \item \textbf{Nghĩ star schema thay thế OLTP.} App vẫn ghi vào OLTP. Star schema được dựng lại từ dữ liệu đó, ở một tầng riêng cho phân tích.
  \item \textbf{Coi việc lặp tên danh mục trong \code{dim\_\allowbreak{}product} là lỗi.} Lặp là cố ý. Bài 8 đo xem nó tốn bao nhiêu byte.
\end{itemize}
\end{luuy}

% =====================================================================
\section{Bài 2: Quy trình nghiệp vụ và grain}

\textbf{Vấn đề.} Ai đó yêu cầu ``một bảng doanh thu''. Một dòng của bảng đó là gì: một đơn, một món hàng, hay một lần trả tiền? Nếu không trả lời câu này trước, cột nào cũng có vẻ hợp lý để thêm vào, và \code{payment\_\allowbreak{}value} sẽ nằm cạnh \code{price} trên cùng một dòng. Đó chính là lỗi ở phần mở đầu.

\textbf{Giải thích.} Phương pháp dùng trong module này là của Ralph Kimball, thường gọi tắt là Kimball. Kimball chia việc thiết kế thành bốn bước, làm đúng thứ tự:
% Nguồn: https://www.kimballgroup.com/wp-content/uploads/2013/08/2013.09-Kimball-Dimensional-Modeling-Techniques11.pdf (trang 1 và 2)
\begin{enumerate}
  \item Chọn \textit{quy trình nghiệp vụ}: một hoạt động có thật mà mỗi lần xảy ra lại sinh ra con số, như nhận một đơn hàng. Đừng chọn theo phòng ban hay theo báo cáo.
  \item Khai báo \textit{grain}: một câu nói chính xác một dòng fact là gì.
  \item Chọn dimension: ai, cái gì, ở đâu, khi nào.
  \item Chọn fact: các con số đo đúng ở grain đó.
\end{enumerate}

Grain là bước then chốt. Kimball gọi câu grain là một hợp đồng ràng buộc cả thiết kế: mọi dimension và mọi fact phải khớp với nó, nên nó phải được viết ra trước.

Hãy bắt đầu ở \textit{atomic grain}, tức mức chi tiết nhất mà quy trình ghi nhận. Bảng đã gộp sẵn, như doanh thu theo ngày, chỉ trả lời được những câu hỏi đã đoán trước. Mỗi grain là một bảng riêng; hai grain không bao giờ nằm chung một fact.

Trong Olist, module này chọn bốn quy trình. Mỗi quy trình cho một fact với grain riêng:

\begin{xltabular}{\textwidth}{@{}P{2.1cm} L P{2.6cm} r P{4.2cm}@{}}
\toprule
\thd{quy trình} & \thd{một sự kiện là} & \thd{nguồn} & \thd{số dòng} & \thd{fact và khoá} \\
\midrule
\endhead
\bottomrule
\endfoot
đặt hàng & một dòng hàng được mua & \code{order\_\allowbreak{}items} & 112.650 & \code{fct\_\allowbreak{}order\_\allowbreak{}items}: \code{order\_\allowbreak{}id}, \code{order\_\allowbreak{}item\_\allowbreak{}id} \\
thanh toán & một lần trả tiền cho một đơn & \code{order\_\allowbreak{}payments} & 103.886 & \code{fct\_\allowbreak{}payments}: \code{order\_\allowbreak{}id}, \code{payment\_\allowbreak{}sequential} \\
giao hàng & một đơn đi qua các mốc từ lúc đặt tới lúc giao & \code{orders} & 99.441 & \code{fct\_\allowbreak{}orders}: \code{order\_\allowbreak{}id} \\
đánh giá & một review của khách & \code{order\_\allowbreak{}reviews} & 99.224 & \code{fct\_\allowbreak{}reviews}: bạn tự xác định ở thử thách kỹ năng \\
\end{xltabular}

\begin{ghichu}
\textbf{So với SWE.} Grain giống lúc bạn thiết kế một resource của REST API. \code{GET /orders/\{id\}} trả về một đơn, còn \code{GET /orders/\{id\}/items/\{n\}} trả về một dòng hàng. Khi đã quyết định một resource là gì, id của nó tự suy ra. Grain cũng vậy: câu ``một dòng là một dòng hàng trong một đơn'' cho ra khoá gồm \code{order\_\allowbreak{}id} và \code{order\_\allowbreak{}item\_\allowbreak{}id}.

\textbf{Chỗ so sánh gãy.} Primary key chỉ nói dòng nào là duy nhất; grain nói dòng đó đo sự kiện gì. Ngoài ra, trong app thì database chặn dòng trùng khoá giúp bạn. Trong lakehouse thì không: spec của Iceberg ghi rõ tính duy nhất theo khoá định danh không được Iceberg bảo đảm hay đòi hỏi. Bạn phải tự kiểm bằng truy vấn, như khối dưới.
% Nguồn: https://iceberg.apache.org/spec/#identifier-field-ids
\end{ghichu}

\subsection{Làm thử}

Mỗi câu grain được viết thành một khoá. Nếu câu grain đúng, số dòng của bảng phải bằng số khoá khác nhau.

% chạy: C
\begin{codeblock}
lakehouse/.venv/bin/python - <<'EOF'
import duckdb
D = 'lakehouse/source_sim/data/'
con = duckdb.connect()
for t in ['orders', 'order_items', 'order_payments']:
    con.sql(f"CREATE VIEW {t} AS SELECT * FROM '{D}olist_{t}_dataset.csv'")
# grain viết thành khoá: số dòng phải bằng số khoá khác nhau
checks = [('order_items', 'order_id, order_item_id'),
          ('order_payments', 'order_id, payment_sequential'),
          ('orders', 'order_id'),
          ('order_items', 'order_id')]   # thử sai: grain "một dòng mỗi đơn"
for table, key in checks:
    rows, keys = con.sql(f"SELECT count(*), count(DISTINCT ({key})) "
                         f"FROM {table}").fetchall()[0]
    print(f'{table}, khoá ({key}): {rows} dòng, {keys} khoá')
# Dạng test: trả về các khoá bị trùng; 0 dòng là đạt
con.sql("""SELECT order_id, order_item_id, count(*) AS n
           FROM order_items GROUP BY ALL HAVING count(*) > 1""").show()
EOF
\end{codeblock}
\textit{Kết quả:}
% kết quả: C
\begin{codeblock}
order_items, khoá (order_id, order_item_id): 112650 dòng, 112650 khoá
order_payments, khoá (order_id, payment_sequential): 103886 dòng, 103886 khoá
orders, khoá (order_id): 99441 dòng, 99441 khoá
order_items, khoá (order_id): 112650 dòng, 98666 khoá
┌──────────┬───────────────┬───────┐
│ order_id │ order_item_id │   n   │
│ varchar  │     int64     │ int64 │
└──────────┴───────────────┴───────┘
               0 rows
\end{codeblock}

Ba câu grain đầu đúng trên dữ liệu nguồn. Dòng thứ tư là một grain sai cố ý: nếu bạn tuyên bố \code{order\_\allowbreak{}items} có ``một dòng mỗi đơn'', 112.650 dòng chỉ có 98.666 khoá khác nhau.

Câu \code{GROUP BY} kèm \code{HAVING count(*) > 1} ở cuối là dạng chuẩn của một \textit{test grain}. Nó trả về các khoá bị trùng, và 0 dòng nghĩa là đạt. Bạn sẽ gặp lại đúng dạng này ở Lab 3.2.

\begin{luuy}
\textbf{Bẫy thường gặp.}
\begin{itemize}
  \item \textbf{Thiết kế theo báo cáo.} Mỗi dashboard một bảng, mỗi bảng tính ``doanh thu'' một kiểu, rồi số giữa các dashboard lệch nhau. Kimball làm ngược lại: một fact gắn với một sự kiện có thật, không gắn với một báo cáo.
  \item \textbf{Nhầm grain với primary key.} Hai bảng có thể cùng dùng \code{order\_\allowbreak{}id} trong khoá mà đo hai thứ khác nhau: một bảng ghi các mốc giao của mỗi đơn, một bảng ghi trạng thái của mỗi đơn vào cuối mỗi ngày.
  \item \textbf{Gộp quá sớm.} Bảng doanh thu theo ngày chạy nhanh, nhưng không trả lời được câu ``doanh thu theo người bán''. Giữ fact ở atomic grain, gộp sau.
\end{itemize}
\end{luuy}

% =====================================================================
\section{Bài 3: Fan-out, lỗi nhân dòng}

\textbf{Vấn đề.} Phần mở đầu cho thấy join thanh toán với dòng hàng rồi cộng làm tổng tiền tăng 26,9\%. Bài này trả lời hai câu: sửa thế nào cho đúng, và vì sao vài cách sửa trông có vẻ đúng vẫn sai.

\textbf{Giải thích.} Khi join theo quan hệ một nhiều, mỗi dòng bên ``một'' bị lặp lại cho mỗi dòng khớp bên ``nhiều''. Nếu sau đó bạn cộng một cột của bên ``một'', mỗi giá trị bị cộng nhiều lần.

Luật sửa rất ngắn: không cộng qua một join làm đổi grain. Gộp từng bảng về cùng một grain trước, ví dụ một dòng mỗi đơn, rồi mới join các kết quả đã gộp.

\begin{ghichu}
\textbf{So với SWE.} Nếu bạn từng dùng Django ORM, bạn đã gặp lỗi này dưới một cái tên khác. Tài liệu Django ghi rằng gộp nhiều aggregation trong \code{annotate()} sẽ cho kết quả sai, vì ORM dùng join thay cho subquery. Ví dụ của họ là một cuốn sách có 2 tác giả và 3 cửa hàng. Gọi \code{annotate()} với hai \code{Count}, một trên \code{authors} và một trên \code{store}, thì cả hai số đều ra 6, vì 2 nhân 3 bằng 6. Đó chính là fan-out.
% Nguồn: https://docs.djangoproject.com/en/6.1/topics/db/aggregation/#combining-multiple-aggregations

\textbf{Chỗ so sánh gãy.} Django có cách chữa riêng cho phép đếm là tham số \code{distinct=True}, và chính tài liệu của họ viết rằng với phần lớn aggregate khác thì không có cách nào tránh vấn đề này trong \code{annotate()}. Django cũng không cảnh báo lúc chạy, lỗi chỉ được ghi trong tài liệu. Ở kho dữ liệu cũng vậy, và cách chữa kiểu distinct không dùng được cho tổng tiền: khối dưới cho thấy \code{sum(DISTINCT)} ra 7.700.442,27 BRL.
\end{ghichu}

\subsection{Làm thử}

Khối dưới thử cách sửa đúng và hai cách sửa trông có vẻ đúng.

% chạy: D
\begin{codeblock}
lakehouse/.venv/bin/python - <<'EOF'
import duckdb
D = 'lakehouse/source_sim/data/'
con = duckdb.connect()
for t in ['order_items', 'order_payments']:
    con.sql(f"CREATE VIEW {t} AS SELECT * FROM '{D}olist_{t}_dataset.csv'")
# Gộp từng bảng về grain "một dòng mỗi đơn" trước, rồi mới join
con.sql("""CREATE TABLE p AS SELECT order_id,
             sum(CAST(payment_value AS DECIMAL(14,2))) AS paid
           FROM order_payments GROUP BY 1""")
con.sql("""CREATE TABLE i AS SELECT order_id, count(*) AS n_items
           FROM order_items GROUP BY 1""")
one = lambda q: con.sql(q).fetchall()[0][0]
print('gộp trước, LEFT JOIN :',
      one("SELECT sum(paid) FROM p LEFT JOIN i USING (order_id)"))
print('gộp trước, INNER JOIN:',
      one("SELECT sum(paid) FROM p JOIN i USING (order_id)"))
print('đơn có tiền mà không có hàng:',
      one("SELECT count(*) FROM p ANTI JOIN i USING (order_id)"))
print('sum(DISTINCT) sau join:', one("""
  SELECT sum(DISTINCT CAST(p.payment_value AS DECIMAL(14,2)))
  FROM order_payments p JOIN order_items i USING (order_id)"""))
EOF
\end{codeblock}
\textit{Kết quả:}
% kết quả: D
\begin{codeblock}
gộp trước, LEFT JOIN : 16008872.12
gộp trước, INNER JOIN: 15846280.17
đơn có tiền mà không có hàng: 775
sum(DISTINCT) sau join: 7700442.27
\end{codeblock}

\begin{itemize}
  \item Gộp về grain đơn trước rồi \code{LEFT JOIN} cho đúng 16.008.872,12 BRL.
  \item Cùng cách đó nhưng dùng \code{INNER JOIN} cho 15.846.280,17 BRL, tức mất 162.591,95 BRL của 775 đơn có thanh toán mà không có dòng hàng nào. Đây là nửa còn lại của lỗi join: join không chỉ nhân dòng mà còn có thể làm rơi dòng, và cũng không báo gì.
  \item \code{sum(DISTINCT ...)} sau join cho 7.700.442,27 BRL, chưa tới một nửa số đúng. \code{DISTINCT} bỏ luôn những khoản tiền của các đơn khác nhau nhưng trùng số tiền.
\end{itemize}

\begin{luuy}
\textbf{Bẫy thường gặp.}
\begin{itemize}
  \item \textbf{Thêm \code{DISTINCT} cho ``hết trùng''.} Nó che lỗi ở số đếm và làm sai số tổng theo một cách khác.
  \item \textbf{Dùng \code{INNER JOIN} khi bên kia có thể thiếu dòng.} Hãy hỏi trước: có đơn nào không có dòng hàng không? Ở Olist có 775 đơn như vậy.
  \item \textbf{Nghĩ fan-out chỉ làm sai tổng.} Nó làm sai cả trung bình, vì dòng bị lặp được tính thêm trọng số. Bài 7 cho thấy điểm review trung bình của bang SP thành 4,13 thay vì 4,17 sau một join như vậy.
\end{itemize}
\end{luuy}

% =====================================================================
\section{Bài 4: Bảng fact, loại fact và loại measure}

\textbf{Vấn đề.} Không phải con số nào cũng cộng được. Tỉ lệ giao trễ của cả nước có phải là trung bình tỉ lệ của 27 bang không? Số đơn đang chờ giao vào cuối mỗi ngày cộng qua một tuần có nghĩa gì không? Và với câu hỏi ``mỗi bước giao hàng mất bao lâu'', một dòng fact nên là gì?

\textbf{Giải thích.} Bảng fact chứa các con số đo được, gọi là \textit{measure}, cùng với khoá trỏ sang các dimension. Kimball chia fact thành các loại theo cách một dòng được sinh ra:
% Nguồn: Kimball techniques PDF, trang 4 và 5

\begin{xltabular}{\textwidth}{@{}P{2.7cm} L P{2.6cm} L@{}}
\toprule
\thd{loại fact} & \thd{một dòng là} & \thd{dòng có bị sửa} & \thd{ví dụ Olist} \\
\midrule
\endhead
\bottomrule
\endfoot
transaction & một sự kiện tại một thời điểm & không, chỉ thêm & \code{fct\_\allowbreak{}order\_\allowbreak{}items}, \code{fct\_\allowbreak{}payments} \\
periodic snapshot & trạng thái của một thứ vào cuối mỗi kỳ, như mỗi ngày & không, mỗi kỳ thêm một lượt dòng & số đơn đang chờ giao của mỗi người bán vào cuối mỗi ngày \\
accumulating snapshot & một quy trình có các mốc biết trước, mỗi mốc một cột ngày & có, sửa khi quy trình qua từng mốc & \code{fct\_\allowbreak{}orders}: đặt, duyệt, giao cho hãng vận chuyển, giao tới khách, ngày dự kiến \\
factless & các dimension gặp nhau tại một thời điểm, không có con số & không & Olist không có loại này; ví dụ của Kimball là một học sinh có mặt ở một lớp vào một ngày \\
\end{xltabular}

Measure cũng có ba loại, theo việc nó cộng được theo những chiều nào:

\begin{xltabular}{\textwidth}{@{}P{3.3cm} P{3.0cm} L L@{}}
\toprule
\thd{loại measure} & \thd{cộng được theo} & \thd{ví dụ Olist} & \thd{cách gộp đúng} \\
\midrule
\endhead
\bottomrule
\endfoot
cộng được, \textit{additive} & mọi dimension & \code{price}, \code{freight\_\allowbreak{}value} & \code{sum} \\
cộng được một phần, \textit{semi-additive} & mọi dimension trừ thời gian & số đơn đang chờ giao vào cuối ngày & lấy giá trị cuối kỳ, hoặc trung bình theo thời gian \\
không cộng được, \textit{non-additive} & không chiều nào & \code{review\_\allowbreak{}score}, tỉ lệ giao trễ & cộng tử số và mẫu số riêng, chia ở bước cuối \\
\end{xltabular}

\begin{ghichu}
\textbf{So với SWE.} Fact transaction giống bảng \code{events} hay audit log chỉ thêm mà bạn từng viết: mỗi dòng là một chuyện đã xảy ra, và không ai sửa nó.

Ba loại measure giống ba loại số liệu bạn gặp khi làm monitoring. Measure cộng được giống từng lần tăng của một counter: cộng theo service hay theo giờ đều có nghĩa. Measure cộng được một phần giống một gauge, loại metric mà Prometheus định nghĩa là một giá trị có thể tăng giảm tuỳ ý. Cộng bộ nhớ của mọi pod tại một thời điểm thì có nghĩa, còn cộng bộ nhớ của một pod qua 24 giờ thì không. Measure không cộng được giống tỉ lệ lỗi: tỉ lệ lỗi của cả hệ thống là tổng số lỗi chia tổng số request, không phải trung bình tỉ lệ lỗi của từng service.
% Nguồn: https://prometheus.io/docs/concepts/metric_types/

\textbf{Chỗ so sánh gãy.} Accumulating snapshot không chỉ thêm: dòng của một đơn được sửa mỗi khi đơn qua một mốc, và Kimball nói đây là loại fact duy nhất được cập nhật như vậy. Ngoài ra, bảng events của app hay chứa một payload JSON; fact thì chứa measure có kiểu rõ ràng và khoá tới dimension.
\end{ghichu}

\subsection{Làm thử}

Khối đầu dựng \code{fct\_\allowbreak{}orders} dạng accumulating snapshot: một dòng mỗi đơn, đủ năm mốc ở bảng trên, mỗi mốc một cột ngày. Bảng còn có một cột độ trễ và điểm của review mới nhất. Khối này trả lời câu hỏi thứ nhất và thứ ba ở phần vấn đề; khối thứ hai trả lời câu hỏi thứ hai.

% chạy: E
\begin{codeblock}
lakehouse/.venv/bin/python - <<'EOF'
import duckdb
D = 'lakehouse/source_sim/data/'
con = duckdb.connect()
for t in ['orders', 'order_reviews', 'customers']:
    con.sql(f"CREATE VIEW {t} AS SELECT * FROM '{D}olist_{t}_dataset.csv'")
# Accumulating snapshot: một dòng mỗi đơn, mỗi mốc một cột ngày
con.sql("""CREATE TABLE fct_orders AS
  SELECT o.order_id, c.customer_state,
         CAST(o.order_purchase_timestamp AS DATE)      AS purchase_date,
         CAST(o.order_approved_at AS DATE)             AS approved_date,
         CAST(o.order_delivered_carrier_date AS DATE)  AS carrier_date,
         CAST(o.order_delivered_customer_date AS DATE) AS delivered_date,
         CAST(o.order_estimated_delivery_date AS DATE) AS estimated_date,
         CAST(o.order_delivered_customer_date AS DATE)
           - CAST(o.order_purchase_timestamp AS DATE)  AS days_to_customer,
         o.order_delivered_customer_date > o.order_estimated_delivery_date
           AS late_by_timestamp,
         r.review_score
  FROM orders o JOIN customers c USING (customer_id)
  LEFT JOIN (SELECT order_id, review_score FROM order_reviews
             QUALIFY row_number() OVER (PARTITION BY order_id
               ORDER BY review_creation_date DESC, review_answer_timestamp DESC,
                        review_id DESC) = 1) r USING (order_id)""")
con.sql("""SELECT count(*) AS da_giao,
                  count(*) FILTER (WHERE late_by_timestamp) AS tre_theo_gio,
                  count(*) FILTER (WHERE delivered_date > estimated_date)
                    AS tre_theo_ngay
           FROM fct_orders WHERE delivered_date IS NOT NULL""").show()
con.sql("""SELECT delivered_date > estimated_date AS tre, count(*) AS so_don,
                  round(avg(review_score), 2) AS diem_tb
           FROM fct_orders WHERE delivered_date IS NOT NULL
           GROUP BY 1 ORDER BY 1""").show()
# Tỉ lệ là measure không cộng được: tính từ tổng, đừng lấy trung bình tỉ lệ con
con.sql("""WITH s AS (
             SELECT customer_state,
                    count(*) FILTER (WHERE delivered_date > estimated_date) AS late,
                    count(*) AS delivered
             FROM fct_orders WHERE delivered_date IS NOT NULL GROUP BY 1)
           SELECT round(sum(late) / sum(delivered), 4) AS ti_le_tu_tong,
                  round(avg(late / delivered), 4)      AS trung_binh_27_bang
           FROM s""").show()
# Mỗi bước mất bao lâu: các mốc nằm trên cùng một dòng, nên chỉ cần trừ hai cột
con.sql("""SELECT
             round(avg(approved_date - purchase_date), 1) AS dat_toi_duyet,
             round(avg(carrier_date - approved_date), 1)  AS duyet_toi_van_chuyen,
             round(avg(delivered_date - carrier_date), 1) AS van_chuyen_toi_khach,
             round(avg(days_to_customer), 1)              AS ca_quy_trinh
           FROM fct_orders WHERE delivered_date IS NOT NULL""").show()
EOF
\end{codeblock}
\textit{Kết quả:}
% kết quả: E
\begin{codeblock}
┌─────────┬──────────────┬───────────────┐
│ da_giao │ tre_theo_gio │ tre_theo_ngay │
│  int64  │    int64     │     int64     │
├─────────┼──────────────┼───────────────┤
│   96476 │         7827 │          6535 │
└─────────┴──────────────┴───────────────┘

┌─────────┬────────┬─────────┐
│   tre   │ so_don │ diem_tb │
│ boolean │ int64  │ double  │
├─────────┼────────┼─────────┤
│ false   │  89941 │    4.29 │
│ true    │   6535 │    2.27 │
└─────────┴────────┴─────────┘

┌───────────────┬────────────────────┐
│ ti_le_tu_tong │ trung_binh_27_bang │
│    double     │       double       │
├───────────────┼────────────────────┤
│        0.0677 │             0.0915 │
└───────────────┴────────────────────┘

┌───────────────┬──────────────────────┬──────────────────────┬──────────────┐
│ dat_toi_duyet │ duyet_toi_van_chuyen │ van_chuyen_toi_khach │ ca_quy_trinh │
│    double     │        double        │        double        │    double    │
├───────────────┼──────────────────────┼──────────────────────┼──────────────┤
│           0.5 │                  2.7 │                  9.3 │         12.5 │
└───────────────┴──────────────────────┴──────────────────────┴──────────────┘
\end{codeblock}

\begin{itemize}
  \item So theo giờ cho 7.827 đơn trễ; so theo ngày cho 6.535 đơn. Chênh lệch là 1.292 đơn giao đúng ngày dự kiến nhưng sau 00:00, vì \code{order\_\allowbreak{}estimated\_\allowbreak{}delivery\_\allowbreak{}date} của Olist luôn là 00:00:00. Repo \repo{lakehouse} cũng so theo ngày vì lý do này, ở mục 6 của \repo{docs/data-model.md}.
  % Kiểm: count(*) FILTER (WHERE CAST(order_estimated_delivery_date AS TIME) <> '00:00:00') = 0 trên 99.441 đơn.
  \item Đơn giao trễ được 2,27 điểm review trung bình, đơn đúng hẹn được 4,29 điểm. Accumulating snapshot làm câu hỏi này dễ, vì mọi mốc của một đơn nằm trên cùng một dòng.
  \item Tỉ lệ trễ tính từ tổng là 6,77\%; trung bình tỉ lệ của 27 bang là 9,15\%. Trung bình của các tỉ lệ cho mỗi bang cùng một trọng số, nên bang AL, với 397 đơn đã giao và tỉ lệ trễ 21,4\%, kéo số chung lên như thể AL lớn bằng SP.
  \item Mỗi bước là phép trừ hai cột trên cùng một dòng: trung bình 0,5 ngày từ lúc đặt tới lúc duyệt, 2,7 ngày từ lúc duyệt tới lúc giao cho hãng vận chuyển, và 9,3 ngày từ đó tới tay khách. Nếu mỗi mốc là một dòng trong bảng events, bạn phải ghép các dòng của cùng một đơn trước khi trừ.
  \item Cột \code{days\_\allowbreak{}to\_\allowbreak{}customer} là một \textit{lag fact}. Kimball gợi ý lưu cho mỗi mốc một cột độ trễ tính từ mốc đầu tiên; khi đó độ trễ giữa hai mốc bất kỳ chỉ là hiệu của hai cột. Cả quy trình mất trung bình 12,5 ngày.
  % Nguồn: Kimball techniques PDF, trang 15, mục Lag/Duration Facts
\end{itemize}

Khối đầu của phần Làm thử dựng bảng một lần, từ trạng thái cuối của mọi đơn trong CSV. Lab 3.4 dựng chính bảng này theo cách pipeline làm: từ từng version của đơn, sửa dòng mỗi khi đơn qua một mốc.

Khối thứ hai dựng một periodic snapshot nhỏ cho cả sàn: mỗi ngày từ 07/05 tới 13/05/2018 một dòng, đếm số đơn đã duyệt mà khách chưa nhận vào cuối ngày đó. Đơn không có ngày giao tới khách không được tính, vì không biết nó hết chờ lúc nào.

% chạy: E2
\begin{codeblock}
lakehouse/.venv/bin/python - <<'EOF'
import duckdb
D = 'lakehouse/source_sim/data/'
con = duckdb.connect()
con.sql(f"CREATE VIEW orders AS SELECT * FROM '{D}olist_orders_dataset.csv'")
# Periodic snapshot: một dòng mỗi ngày. Đơn đang chờ vào cuối ngày d là đơn
# đã duyệt trong hoặc trước ngày d, và tới tay khách sau ngày d
con.sql("""CREATE TABLE fct_pending_daily AS
  SELECT CAST(r.d AS DATE) AS snapshot_date, count(o.order_id) AS pending_orders
  FROM range(DATE '2018-05-07', DATE '2018-05-14', INTERVAL 1 DAY) r(d)
  LEFT JOIN orders o
    ON CAST(o.order_approved_at AS DATE) <= r.d
   AND CAST(o.order_delivered_customer_date AS DATE) > r.d
  GROUP BY 1""")
con.sql("SELECT * FROM fct_pending_daily ORDER BY 1").show()
# arg_max(x, ngày): giá trị x ở ngày lớn nhất, tức giá trị cuối kỳ
con.sql("""SELECT sum(pending_orders) AS cong_7_ngay,
                  round(avg(pending_orders), 1) AS trung_binh_ngay,
                  arg_max(pending_orders, snapshot_date) AS cuoi_tuan
           FROM fct_pending_daily""").show()
EOF
\end{codeblock}
\textit{Kết quả:}
% kết quả: E2
\begin{codeblock}
┌───────────────┬────────────────┐
│ snapshot_date │ pending_orders │
│     date      │     int64      │
├───────────────┼────────────────┤
│ 2018-05-07    │           2746 │
│ 2018-05-08    │           2791 │
│ 2018-05-09    │           2782 │
│ 2018-05-10    │           2722 │
│ 2018-05-11    │           2626 │
│ 2018-05-12    │           2730 │
│ 2018-05-13    │           2901 │
└───────────────┴────────────────┘

┌─────────────┬─────────────────┬───────────┐
│ cong_7_ngay │ trung_binh_ngay │ cuoi_tuan │
│   int128    │     double      │   int64   │
├─────────────┼─────────────────┼───────────┤
│       19298 │          2756.9 │      2901 │
└─────────────┴─────────────────┴───────────┘
\end{codeblock}

\begin{itemize}
  \item Cộng 7 ngày được 19.298, nhưng ngày cao nhất của tuần chỉ có 2.901 đơn đang chờ. Một đơn chờ cả tuần bị đếm 7 lần, giống cộng bộ nhớ của một pod qua 24 giờ ở ghi chú trên.
  \item Measure cộng được một phần phải gộp theo thời gian bằng giá trị cuối kỳ, 2.901 đơn vào cuối ngày 13/05/2018, hoặc bằng trung bình, 2.756,9 đơn mỗi ngày. Ví dụ quen thuộc của Kimball là số dư tài khoản: cộng được theo mọi dimension trừ thời gian.
  % Nguồn: Kimball techniques PDF, trang 4, mục Additive, Semi-Additive, and Non-Additive Facts
\end{itemize}

\begin{luuy}
\textbf{Bẫy thường gặp.}
\begin{itemize}
  \item \textbf{So một timestamp với một ngày lưu dạng 00:00:00.} Đọc dữ liệu trước khi viết điều kiện: 1.292 đơn bị tính trễ oan.
  \item \textbf{Lấy trung bình của các tỉ lệ con.} Lưu tử số và mẫu số, cộng chúng, rồi mới chia.
  \item \textbf{Cộng periodic snapshot theo thời gian.} Số đơn đang chờ giao của 7 ngày cộng lại không phải số đơn đang chờ của tuần đó.
  \item \textbf{Để mốc chưa xảy ra thành NULL trong khoá.} Dùng dòng khoá 0 ``chưa biết'' của \code{dim\_\allowbreak{}date}, ở Bài 6.
\end{itemize}
\end{luuy}

% =====================================================================
\section{Bài 5: Dimension và khoá}

\textbf{Vấn đề.} Bạn được hỏi ``bao nhiêu phần trăm khách mua lại lần hai''. Bảng \code{customers} có cột \code{customer\_\allowbreak{}id}, nên bạn đếm số khách có hơn một đơn theo cột đó và được 0. Không ai mua lại sao? Thật ra Olist cấp một \code{customer\_\allowbreak{}id} mới cho mỗi đơn. Người mua thật được định danh bằng \code{customer\_\allowbreak{}unique\_\allowbreak{}id}.

\textbf{Giải thích.} Dimension là bảng mô tả, trả lời các câu ai, cái gì, ở đâu, khi nào. Cột của nó được dùng trong \code{WHERE} và \code{GROUP BY}. Bảng dimension ngắn và rộng: mỗi thực thể một dòng, nhiều cột chữ dễ đọc. Một dimension có ba loại khoá:
% Nguồn: Kimball techniques PDF, trang 4, 7 và 8
\begin{itemize}
  \item \textbf{Khoá tự nhiên}, \textit{natural key}: khoá đến từ hệ nguồn, như \code{product\_\allowbreak{}id}.
  \item \textbf{Khoá thay thế}, \textit{surrogate key}: số nguyên không mang nghĩa do kho dữ liệu tự cấp, bắt đầu từ 1. Kimball khuyên mọi dimension dùng loại khoá này làm khoá chính, vì khoá tự nhiên có thể đổi, có thể trùng giữa hai hệ nguồn, và sẽ lặp lại khi dimension giữ lịch sử ở Bài 9. Dimension ngày là ngoại lệ: nó dùng được khoá dạng \code{yyyymmdd}.
  \item \textbf{Khoá bền}, \textit{durable key}: khoá không đổi của một thực thể qua mọi version của nó. Ở Olist, \code{customer\_\allowbreak{}unique\_\allowbreak{}id} đóng vai này.
\end{itemize}

Mỗi dimension còn có một dòng đặc biệt, thường mang khoá 0, cho giá trị ``chưa biết''. Kimball viết rằng khoá ngoại NULL trên fact vi phạm toàn vẹn tham chiếu; trong thực tế nó còn làm \code{INNER JOIN} rơi dòng. Thay cho khoá NULL, dimension phải có một dòng mặc định cho trường hợp chưa biết hoặc không áp dụng.

Trong star schema chuẩn, các dimension không trỏ tới nhau. Quan hệ giữa hai dimension là một sự kiện, nên nó nằm trong fact. Câu ``người mua \code{2e875ea5} mua của người bán \code{7299e27e}'' chỉ có nghĩa khi gắn với một đơn, nên nó là một dòng của \code{fct\_\allowbreak{}order\_\allowbreak{}items}. Có ba ngoại lệ có tên riêng. Kimball khuyên tránh snowflake, dùng outrigger hạn chế, và dùng bridge khi một dimension có nhiều giá trị cho một dòng fact:
% Nguồn: Kimball techniques PDF, trang 9 và 18
\begin{itemize}
  \item \textbf{Snowflake}: dimension tách tiếp thành dimension con, như \code{dim\_\allowbreak{}product} trỏ sang \code{dim\_\allowbreak{}category}. Bài 8 đo nó.
  \item \textbf{Outrigger}: dimension trỏ sang một dimension khác, như ngày mua đầu tiên của người mua trỏ sang \code{dim\_\allowbreak{}date}.
  \item \textbf{Bridge}: bảng nối cho quan hệ nhiều nhiều, như một đơn trả bằng nhiều kiểu thanh toán.
\end{itemize}

Đây là cách đọc một đơn thật qua fact và dimension. Đơn \code{002f98c0f7efd42638ed6100ca699b42} có 2 dòng hàng; các id được cắt còn 8 ký tự. Fact chỉ chứa khoá và con số:

\begin{xltabular}{\textwidth}{@{}l l l l L r r@{}}
\toprule
\thd{dòng} & \thd{ngày đặt} & \thd{người mua} & \thd{sản phẩm} & \thd{người bán} & \thd{price} & \thd{freight} \\
\midrule
\endhead
\bottomrule
\endfoot
1 & 2017-08-04 & \code{2e875ea5} & \code{d41dc2f2} & \code{7299e27e} & 8,99 & 32,57 \\
2 & 2017-08-04 & \code{2e875ea5} & \code{880be32f} & \code{fa40cc5b} & 44,90 & 7,16 \\
\end{xltabular}

Nghĩa của các khoá nằm ở dimension. Sản phẩm \code{d41dc2f2} thuộc danh mục \code{consoles\_\allowbreak{}games}, nặng 450 g, có 1 ảnh; sản phẩm \code{880be32f} thuộc \code{toys}, nặng 200 g, có 4 ảnh. Người bán \code{7299e27e} ở araguari, bang MG; người bán \code{fa40cc5b} ở sao paulo, bang SP. Người mua \code{2e875ea5} ở porto alegre, bang RS.

Ghép lại, dòng 1 nghĩa là: một người mua ở porto alegre, bang RS, mua một món \code{consoles\_\allowbreak{}games} của người bán ở araguari, bang MG, ngày 2017-08-04, giá 8,99 BRL, phí vận chuyển 32,57 BRL.

\begin{ghichu}
\textbf{So với SWE.} Khoá thay thế có cùng lý do với việc bạn không lấy email làm primary key của bảng \code{users}: email có thể đổi. Cặp \code{customer\_\allowbreak{}id} và \code{customer\_\allowbreak{}unique\_\allowbreak{}id} giống một id gắn với từng đơn so với user id thật.

Dòng ``unknown'' là Null Object pattern: thay vì trả về NULL rồi buộc mọi nơi phải kiểm, bạn trả về một đối tượng mặc định có hành vi an toàn. Ở đây join không rơi dòng, và báo cáo hiện chữ ``unknown'' thay vì một ô trống.

\textbf{Chỗ so sánh gãy.} Trong app, một user có một id suốt đời. Trong dimension có lịch sử, một người mua có nhiều khoá thay thế, mỗi version một khoá, và chỉ khoá bền là không đổi. Bài 9 làm rõ điều này.

Null Object cũng gãy ở một chỗ: dòng unknown có thể nuốt lỗi âm thầm. Nếu bước tra khoá hỏng, fact vẫn nạp đủ dòng, chỉ là trỏ về khoá 0, nên cần thêm một test theo dõi tỉ lệ dòng mang khoá 0.
\end{ghichu}

\subsection{Làm thử}

% chạy: F
\begin{codeblock}
lakehouse/.venv/bin/python - <<'EOF'
import duckdb
D = 'lakehouse/source_sim/data/'
con = duckdb.connect()
for t in ['orders', 'customers']:
    con.sql(f"CREATE VIEW {t} AS SELECT * FROM '{D}olist_{t}_dataset.csv'")
# Ai là "một người mua"? Đếm số đơn theo từng loại id
con.sql("""WITH n AS (
             SELECT c.customer_id, c.customer_unique_id, count(*) OVER
                      (PARTITION BY c.customer_unique_id) AS orders_per_buyer
             FROM orders o JOIN customers c USING (customer_id))
           SELECT count(DISTINCT customer_id)        AS customer_id,
                  count(DISTINCT customer_unique_id) AS customer_unique_id,
                  count(DISTINCT customer_unique_id)
                    FILTER (WHERE orders_per_buyer > 1) AS mua_tu_2_lan
           FROM n""").show()
# dim_customer: khoá thay thế từ 1, dòng "unknown" mang khoá 0
con.sql("""CREATE TABLE dim_customer AS
  SELECT 0 AS customer_key, 'unknown' AS customer_unique_id,
         'unknown' AS customer_city, 'unknown' AS customer_state
  UNION ALL
  SELECT row_number() OVER (ORDER BY customer_unique_id), customer_unique_id,
         customer_city, customer_state
  FROM (SELECT c.customer_unique_id, c.customer_city, c.customer_state
        FROM orders o JOIN customers c USING (customer_id)
        QUALIFY row_number() OVER (PARTITION BY c.customer_unique_id
                ORDER BY o.order_purchase_timestamp DESC,
                         c.customer_id DESC) = 1)""")
con.sql("""SELECT count(*) AS so_dong, count(DISTINCT customer_key) AS so_khoa,
                  min(customer_key) AS khoa_min, max(customer_key) AS khoa_max
           FROM dim_customer""").show()
EOF
\end{codeblock}
\textit{Kết quả:}
% kết quả: F
\begin{codeblock}
┌─────────────┬────────────────────┬──────────────┐
│ customer_id │ customer_unique_id │ mua_tu_2_lan │
│    int64    │       int64        │    int64     │
├─────────────┼────────────────────┼──────────────┤
│       99441 │              96096 │         2997 │
└─────────────┴────────────────────┴──────────────┘

┌─────────┬─────────┬──────────┬──────────┐
│ so_dong │ so_khoa │ khoa_min │ khoa_max │
│  int64  │  int64  │  int64   │  int64   │
├─────────┼─────────┼──────────┼──────────┤
│   96097 │   96097 │        0 │    96096 │
└─────────┴─────────┴──────────┴──────────┘
\end{codeblock}

\begin{itemize}
  \item Có 99.441 \code{customer\_\allowbreak{}id} cho 99.441 đơn, nên mỗi \code{customer\_\allowbreak{}id} chỉ có một đơn, và tỉ lệ mua lại tính theo cột này luôn là 0.
  \item Theo \code{customer\_\allowbreak{}unique\_\allowbreak{}id} có 96.096 người mua, trong đó 2.997 người mua từ 2 lần trở lên, khoảng 3,1\%.
  \item \code{dim\_\allowbreak{}customer} có 96.097 dòng: 96.096 người mua cộng dòng khoá 0. Địa chỉ được lấy ở đơn mới nhất; Bài 9 cho thấy cách này làm sai thành phố của 130 đơn và sai bang của 43 đơn.
\end{itemize}

\begin{luuy}
\textbf{Bẫy thường gặp.}
\begin{itemize}
  \item \textbf{Dùng \code{customer\_\allowbreak{}id} của Olist làm người mua.} Kiểm bằng cách đếm giá trị khác nhau của hai cột trước khi chọn khoá.
  \item \textbf{Dùng khoá tự nhiên làm khoá chính của dimension có lịch sử.} Khi một người mua có hai version, khoá tự nhiên xuất hiện hai lần.
  \item \textbf{Để NULL trong khoá của fact.} Join với dimension sẽ làm rơi dòng đó khỏi mọi báo cáo. Dùng khoá 0.
  \item \textbf{Cho các dimension trỏ sang nhau ``cho tiện''.} Mỗi lần như vậy, người đọc phải join thêm một bảng, và quan hệ đáng lẽ nằm trong fact bị giấu đi.
\end{itemize}
\end{luuy}

% =====================================================================
\section{Bài 6: Dimension ngày, vai, degenerate và junk}

\textbf{Vấn đề.} Vài dimension không hề có trong app. App tính thứ trong tuần bằng code mỗi khi cần, còn phân tích lại cần nhóm theo tuần, tháng, ngày lễ ở hầu hết câu hỏi. Một đơn có năm mốc thời gian, và mốc nào cũng là một ngày. \code{order\_\allowbreak{}id} không còn thuộc tính nào sau khi khách, sản phẩm và người bán đã thành dimension riêng. Đơn còn vài cờ nhỏ như trạng thái, kiểu thanh toán, giao đúng hẹn hay không, có review hay không. Làm mỗi cờ một dimension thì fact sẽ có cả chục khoá.

\textbf{Giải thích.}
% Nguồn: Kimball techniques PDF, trang 7 tới 9
\begin{itemize}
  \item \textbf{Dimension ngày.} Mỗi ngày một dòng, với các cột như năm, tháng, thứ, cờ ngày lễ. Kimball viết rằng bạn không bao giờ muốn tính ngày lễ Phục sinh bằng SQL, mà tra nó trong dimension ngày. Khoá có thể là số \code{yyyymmdd}. Bảng cần một dòng khoá 0 cho ngày chưa biết, như ngày giao của đơn chưa giao.
  \item \textbf{Role-playing.} Một bảng \code{dim\_\allowbreak{}date} được fact tham chiếu nhiều lần, mỗi lần với một vai: ngày đặt, ngày giao, ngày dự kiến. Kimball khuyên mỗi vai là một view riêng với tên cột riêng, để hai vai dùng được trong cùng một câu truy vấn.
  \item \textbf{Degenerate dimension.} \code{order\_\allowbreak{}id} nằm thẳng trên fact và không có bảng \code{dim\_\allowbreak{}order}, vì sau khi tách hết thuộc tính thì đơn không còn gì để mô tả. Kimball dùng ví dụ số hoá đơn trên các dòng của hoá đơn.
  \item \textbf{Junk dimension.} Gom các cờ nhỏ vào một dimension, mỗi dòng là một tổ hợp cờ, và fact chỉ giữ một khoá. Kimball viết rằng bảng này chỉ nên chứa các tổ hợp thật sự xuất hiện trong dữ liệu. Hướng dẫn Power BI của Microsoft lại nói bảng này thường là tích Descartes của mọi giá trị, tức kết quả \code{CROSS JOIN} của tất cả các cờ. Bảng \code{CROSS JOIN} không bao giờ thiếu tổ hợp; cách của Kimball nhỏ hơn, nhưng pipeline phải thêm dòng khi gặp tổ hợp mới.
  % Nguồn: https://learn.microsoft.com/en-us/power-bi/guidance/star-schema (mục Junk dimensions)
\end{itemize}

\begin{ghichu}
\textbf{So với SWE.} Dimension ngày giống một hàm tiện ích về lịch mà cả team dùng chung, nhưng được lưu thành dữ liệu để mọi công cụ SQL nhận cùng một câu trả lời. Role-playing giống một class \code{Date} dùng cho nhiều field như \code{created\_\allowbreak{}at}, \code{shipped\_\allowbreak{}at}, \code{delivered\_\allowbreak{}at}.

Degenerate dimension giống \code{trace\_\allowbreak{}id} trong log: bạn nhóm theo nó để gom các dòng của cùng một request, nhưng không có bảng \code{traces} nào chứa thuộc tính của nó. Junk dimension giống flyweight pattern hay string interning: rất nhiều đối tượng dùng chung vài tổ hợp giá trị, nên mỗi tổ hợp được lưu một lần và đối tượng chỉ giữ một tham chiếu nhỏ.

\textbf{Chỗ so sánh gãy.} Hàm tiện ích chạy lúc được gọi; dimension ngày phải được nạp trước và có hạn. \code{dim\_\allowbreak{}date} của module này dừng ở 2018-12-31, nên một đơn đặt sau ngày đó sẽ không tìm thấy ngày của nó. Junk dimension theo cách của Kimball cũng vậy: tổ hợp mới phải được thêm vào trước khi fact dùng tới.

Với \code{trace\_\allowbreak{}id}, chỗ gãy nằm ở lý do không có bảng. Đơn hàng vốn có bảng \code{orders} ở nguồn; star bỏ \code{dim\_\allowbreak{}order} vì ngày đặt, người mua và các thuộc tính khác của đơn đã thành khoá tới các dimension khác trên từng dòng fact.
% Nguồn: Kimball techniques PDF, trang 7, mục Degenerate Dimensions
\end{ghichu}

\subsection{Làm thử}

Khối đầu dựng \code{dim\_\allowbreak{}date} có dòng khoá 0, tạo hai vai ngày đặt và ngày giao, rồi dùng cả hai vai trong một câu: đơn đặt trong tháng 12/2017 được giao vào tháng nào?

% chạy: G
\begin{codeblock}
lakehouse/.venv/bin/python - <<'EOF'
import duckdb
D = 'lakehouse/source_sim/data/'
con = duckdb.connect()
con.sql(f"CREATE VIEW orders AS SELECT * FROM '{D}olist_orders_dataset.csv'")
# dim_date: khoá yyyymmdd, thêm dòng khoá 0 cho ngày chưa biết
con.sql("""CREATE TABLE dim_date AS
  SELECT CAST(strftime(d, '%Y%m%d') AS INTEGER) AS date_key, CAST(d AS DATE) AS day,
         year(d) AS year, month(d) AS month, dayname(d) AS weekday
  FROM range(DATE '2016-01-01', DATE '2019-01-01', INTERVAL 1 DAY) r(d)
  UNION ALL SELECT 0, NULL, NULL, NULL, 'unknown'""")
# Mỗi vai là một view riêng, tên cột riêng
con.sql("""CREATE VIEW dim_purchase_date AS SELECT date_key AS purchase_date_key,
             year AS purchase_year, month AS purchase_month FROM dim_date""")
con.sql("""CREATE VIEW dim_delivered_date AS SELECT date_key AS delivered_date_key,
             year AS delivered_year, month AS delivered_month FROM dim_date""")
con.sql("""CREATE TABLE fct_orders AS
  SELECT order_id,
         CAST(strftime(order_purchase_timestamp, '%Y%m%d') AS INTEGER)
           AS purchase_date_key,
         coalesce(CAST(strftime(order_delivered_customer_date, '%Y%m%d')
                       AS INTEGER), 0) AS delivered_date_key
  FROM orders""")
print('dim_date:', con.sql("SELECT count(*) FROM dim_date").fetchall()[0][0], 'dòng')
# Đơn đặt trong tháng 12/2017 được giao vào tháng nào?
con.sql("""SELECT d.delivered_year, d.delivered_month, count(*) AS so_don
           FROM fct_orders f
           JOIN dim_purchase_date p USING (purchase_date_key)
           JOIN dim_delivered_date d USING (delivered_date_key)
           WHERE p.purchase_year = 2017 AND p.purchase_month = 12
           GROUP BY ALL ORDER BY ALL""").show()
EOF
\end{codeblock}
\textit{Kết quả:}
% kết quả: G
\begin{codeblock}
dim_date: 1097 dòng
┌────────────────┬─────────────────┬────────┐
│ delivered_year │ delivered_month │ so_don │
│     int64      │      int64      │ int64  │
├────────────────┼─────────────────┼────────┤
│           2017 │              12 │   3332 │
│           2018 │               1 │   2112 │
│           2018 │               2 │     58 │
│           2018 │               3 │      8 │
│           2018 │               4 │      3 │
│           NULL │            NULL │    160 │
└────────────────┴─────────────────┴────────┘
\end{codeblock}

\code{dim\_\allowbreak{}date} có 1.097 dòng: 1.096 ngày từ đầu 2016 tới hết 2018, cộng dòng khoá 0. Phần lớn đơn đặt tháng 12/2017 được giao trong tháng 12 và tháng 1. 160 đơn chưa giao nhận khoá 0, nên hiện NULL ở cột năm và tháng thay vì biến mất khỏi kết quả.

Khối thứ hai đo degenerate dimension và dựng junk dimension cho các cờ của đơn.

% chạy: G2
\begin{codeblock}
lakehouse/.venv/bin/python - <<'EOF'
import duckdb
D = 'lakehouse/source_sim/data/'
con = duckdb.connect()
for t in ['orders', 'order_items', 'order_payments', 'order_reviews']:
    con.sql(f"CREATE VIEW {t} AS SELECT * FROM '{D}olist_{t}_dataset.csv'")
# Degenerate dimension: order_id nằm trên fact, không có bảng dim_order
n_rows, n_orders = con.sql("SELECT count(*), count(DISTINCT order_id) "
                           "FROM order_items").fetchall()[0]
print('dòng hàng:', n_rows, '| đơn khác nhau:', n_orders)
# Junk dimension: gom các cờ nhỏ của đơn, chỉ giữ tổ hợp thật sự xuất hiện
con.sql("""CREATE TABLE flags AS
  SELECT o.order_status,
         coalesce(p.pay_types, 'none') AS pay_types,
         CASE WHEN o.order_delivered_customer_date IS NULL THEN 'chưa giao'
              WHEN CAST(o.order_delivered_customer_date AS DATE)
                   <= CAST(o.order_estimated_delivery_date AS DATE) THEN 'đúng hẹn'
              ELSE 'trễ' END AS delivery,
         r.order_id IS NOT NULL AS has_review
  FROM orders o
  LEFT JOIN (SELECT order_id, string_agg(DISTINCT payment_type, '+'
                    ORDER BY payment_type) AS pay_types
             FROM order_payments GROUP BY 1) p USING (order_id)
  LEFT JOIN (SELECT DISTINCT order_id FROM order_reviews) r USING (order_id)""")
con.sql("""SELECT count(DISTINCT order_status) AS n_status,
                  count(DISTINCT pay_types)    AS n_pay,
                  count(DISTINCT delivery)     AS n_delivery,
                  count(DISTINCT has_review)   AS n_review,
                  n_status * n_pay * n_delivery * n_review AS co_the,
                  (SELECT count(*) FROM (SELECT DISTINCT * FROM flags)) AS xuat_hien
           FROM flags""").show()
con.sql("""CREATE TABLE dim_order_profile AS
  SELECT row_number() OVER (ORDER BY order_status, pay_types, delivery, has_review)
           AS profile_key, *
  FROM (SELECT DISTINCT * FROM flags)""")
con.sql("""SELECT p.*, count(*) AS so_don FROM flags f
           JOIN dim_order_profile p
             USING (order_status, pay_types, delivery, has_review)
           GROUP BY ALL ORDER BY so_don DESC LIMIT 3""").show()
EOF
\end{codeblock}
\textit{Kết quả:}
% kết quả: G2
\begin{codeblock}
dòng hàng: 112650 | đơn khác nhau: 98666
┌──────────┬───────┬────────────┬──────────┬────────┬───────────┐
│ n_status │ n_pay │ n_delivery │ n_review │ co_the │ xuat_hien │
│  int64   │ int64 │   int64    │  int64   │ int64  │   int64   │
├──────────┼───────┼────────────┼──────────┼────────┼───────────┤
│        8 │     8 │          3 │        2 │    384 │        70 │
└──────────┴───────┴────────────┴──────────┴────────┴───────────┘

┌─────────────┬──────────────┬─────────────┬──────────┬────────────┬────────┐
│ profile_key │ order_status │  pay_types  │ delivery │ has_review │ so_don │
│    int64    │   varchar    │   varchar   │ varchar  │  boolean   │ int64  │
├─────────────┼──────────────┼─────────────┼──────────┼────────────┼────────┤
│          24 │ delivered    │ credit_card │ đúng hẹn │ true       │  66922 │
│          19 │ delivered    │ boleto      │ đúng hẹn │ true       │  17690 │
│          22 │ delivered    │ credit_card │ trễ      │ true       │   4714 │
└─────────────┴──────────────┴─────────────┴──────────┴────────────┴────────┘
\end{codeblock}

\begin{itemize}
  \item 112.650 dòng hàng thuộc 98.666 đơn. \code{order\_\allowbreak{}id} nhóm các dòng hàng của cùng một đơn mà không cần bảng riêng.
  \item Các cờ có 8 trạng thái, 8 tổ hợp kiểu thanh toán, 3 trạng thái giao và 2 giá trị có review, tức 384 tổ hợp có thể. Chỉ 70 tổ hợp thật sự xuất hiện trên 99.441 đơn, và ba tổ hợp lớn nhất đã chiếm 89.326 đơn.
  \item Kiểu thanh toán được gộp thành một chuỗi như \code{credit\_\allowbreak{}card+voucher}, vì 2.246 đơn dùng hơn một kiểu. Nếu không gộp, cờ này không có một giá trị duy nhất cho mỗi đơn.
  % Kiểm: count(*) FROM (SELECT order_id FROM order_payments GROUP BY 1 HAVING count(DISTINCT payment_type) > 1) = 2246.
\end{itemize}

\begin{luuy}
\textbf{Bẫy thường gặp.}
\begin{itemize}
  \item \textbf{Lọc theo khoá \code{yyyymmdd} thay vì theo cột năm, tháng.} Kimball khuyên lọc và nhóm theo thuộc tính của dimension ngày, không theo khoá.
  \item \textbf{Mỗi cờ một khoá ngoại.} Fact phình ra thành một bảng có cả chục khoá, mỗi khoá trỏ tới một dimension chỉ vài dòng.
  \item \textbf{Quên dòng khoá 0.} Đơn chưa giao rơi khỏi mọi báo cáo theo ngày giao mà không ai biết.
  \item \textbf{Hai vai dùng chung tên cột.} Nếu cả hai view đều có cột \code{year}, câu truy vấn và biểu đồ sẽ không cho biết đó là năm đặt hay năm giao.
\end{itemize}
\end{luuy}

% =====================================================================
\section{Bài 7: Conformed dimension, bus matrix và drill-across}

\textbf{Vấn đề.} Giám đốc muốn một bảng: doanh thu và điểm review trung bình theo bang của khách. Doanh thu nằm ở \code{fct\_\allowbreak{}order\_\allowbreak{}items}, điểm nằm ở \code{fct\_\allowbreak{}reviews}, và hai fact khác grain: một đơn có nhiều dòng hàng, và có đơn có nhiều review. Join thẳng hai fact rồi gộp là gặp lại fan-out, lần này ở cả tổng lẫn trung bình.

\textbf{Giải thích.}
% Nguồn: Kimball techniques PDF, trang 10 và 11
\begin{itemize}
  \item \textbf{Conformed dimension}: một dimension được nhiều fact dùng chung, với cùng tên cột và cùng miền giá trị. Nhờ đó ``bang SP'' có cùng một nghĩa ở mọi fact.
  \item \textbf{Drill-across}: truy vấn riêng từng fact, gộp theo cùng một thuộc tính của conformed dimension, rồi ghép các kết quả theo thuộc tính đó. Hai fact không bao giờ được join với nhau ở mức dòng.
  \item \textbf{Bus matrix}: bảng có dòng là quy trình nghiệp vụ, cột là dimension, và đánh dấu ô nào áp dụng. Đọc theo dòng để kiểm một dimension có đúng một giá trị cho mỗi sự kiện không. Đọc theo cột để thấy dimension nào phải conform giữa nhiều quy trình. Kimball khuyên xây từng dòng một.
\end{itemize}

Bus matrix của Olist trong module này:

\begin{xltabular}{\textwidth}{@{}P{2.3cm} L P{1.9cm} P{1.9cm} P{1.9cm} P{2.3cm}@{}}
\toprule
\thd{quy trình} & \thd{\code{dim\_\allowbreak{}date}} & \thd{\code{dim\_\allowbreak{}customer}} & \thd{\code{dim\_\allowbreak{}product}} & \thd{\code{dim\_\allowbreak{}seller}} & \thd{\code{dim\_\allowbreak{}order\_\allowbreak{}profile}} \\
\midrule
\endhead
\bottomrule
\endfoot
đặt hàng & ngày đặt, ngày giao & x & x & x & x \\
thanh toán & ngày đặt & x & & & \\
giao hàng & năm mốc của đơn & x & & & x \\
đánh giá & ngày tạo review & x & & & \\
\end{xltabular}

Ô người bán của dòng giao hàng bỏ trống có lý do: ở grain một dòng mỗi đơn, 1.278 đơn có từ hai người bán trở lên, nên người bán không có một giá trị duy nhất. Muốn phân tích giao hàng theo người bán, grain phải là một dòng mỗi đơn và mỗi người bán.
% Kiểm: count(*) FROM (SELECT order_id FROM order_items GROUP BY 1 HAVING count(DISTINCT seller_id) > 1) = 1278.

Một hệ OLTP thường sinh ra nhiều star, mỗi quy trình một star, dùng chung các conformed dimension; tập các star đó đôi khi được gọi là \textit{fact constellation}.
% Nguồn: https://en.wikipedia.org/wiki/Fact_constellation
Chiều ngược lại cũng đúng: một star có thể lấy dữ liệu từ nhiều service, như thanh toán từ service thanh toán và dòng hàng từ service đơn hàng.

\begin{ghichu}
\textbf{So với SWE.} Conformed dimension giống một enum hay một kiểu dùng chung mà mọi service cùng import, như một file protobuf định nghĩa mã bang. Nếu mỗi service tự đặt mã bang của mình, không ai ghép được số liệu giữa các service. Bus matrix giống ma trận service và capability mà team kiến trúc hay vẽ: nhìn một bảng là biết ai dùng gì, và đổi một cột thì ảnh hưởng tới những dòng nào.

\textbf{Chỗ so sánh gãy.} Enum dùng chung là hằng số trong code, đổi bằng một lần deploy. Conformed dimension là một bảng có dữ liệu và lịch sử, có pipeline và người chịu trách nhiệm. Đổi nó có thể làm đổi số liệu quá khứ của mọi fact dùng nó.

Bus matrix còn là kế hoạch xây: Kimball dùng nó để xếp thứ tự ưu tiên với phía nghiệp vụ, và mỗi lần chỉ làm một dòng, tức một quy trình.
% Nguồn: Kimball techniques PDF, trang 11, mục Enterprise Data Warehouse Bus Matrix
\end{ghichu}

\subsection{Làm thử}

Trong khối dưới, \code{dim\_\allowbreak{}customer} tạm dùng bảng \code{customers} của Olist, mỗi dòng là địa chỉ của khách tại một đơn. Bài 9 dựng bản đúng bằng SCD Type 2 và cho ra cùng các con số này.

% chạy: H
\begin{codeblock}
lakehouse/.venv/bin/python - <<'EOF'
import duckdb
D = 'lakehouse/source_sim/data/'
con = duckdb.connect()
for t in ['orders', 'order_items', 'order_reviews', 'customers']:
    con.sql(f"CREATE VIEW {t} AS SELECT * FROM '{D}olist_{t}_dataset.csv'")
# Hai fact khác grain, cùng dùng một dimension khách hàng
con.sql("""CREATE TABLE fct_order_items AS
  SELECT i.order_id, o.customer_id, CAST(i.price AS DECIMAL(12,2)) AS price
  FROM order_items i JOIN orders o USING (order_id)""")
con.sql("""CREATE TABLE fct_reviews AS
  SELECT r.review_id, r.order_id, o.customer_id, r.review_score
  FROM order_reviews r JOIN orders o USING (order_id)""")
con.sql("""CREATE TABLE dim_customer AS
  SELECT customer_id, customer_state FROM customers""")
print('Sai: join hai fact với nhau rồi mới gộp')
con.sql("""SELECT c.customer_state, sum(f.price) AS revenue,
                  round(avg(r.review_score), 2) AS avg_review
           FROM fct_order_items f JOIN dim_customer c USING (customer_id)
           LEFT JOIN fct_reviews r USING (order_id)
           GROUP BY 1 ORDER BY revenue DESC LIMIT 3""").show()
print('Đúng: drill-across, mỗi fact gộp riêng theo customer_state rồi ghép')
con.sql("""WITH rev AS (
             SELECT c.customer_state, sum(f.price) AS revenue
             FROM fct_order_items f JOIN dim_customer c USING (customer_id)
             GROUP BY 1),
           rv AS (
             SELECT c.customer_state, round(avg(r.review_score), 2) AS avg_review
             FROM fct_reviews r JOIN dim_customer c USING (customer_id)
             GROUP BY 1)
           SELECT * FROM rev JOIN rv USING (customer_state)
           ORDER BY revenue DESC LIMIT 3""").show()
EOF
\end{codeblock}
\textit{Kết quả:}
% kết quả: H
\begin{codeblock}
Sai: join hai fact với nhau rồi mới gộp
┌────────────────┬───────────────┬────────────┐
│ customer_state │    revenue    │ avg_review │
│    varchar     │ decimal(38,2) │   double   │
├────────────────┼───────────────┼────────────┤
│ SP             │    5228869.17 │       4.13 │
│ RJ             │    1831678.85 │       3.81 │
│ MG             │    1591518.47 │       4.09 │
└────────────────┴───────────────┴────────────┘

Đúng: drill-across, mỗi fact gộp riêng theo customer_state rồi ghép
┌────────────────┬───────────────┬────────────┐
│ customer_state │    revenue    │ avg_review │
│    varchar     │ decimal(38,2) │   double   │
├────────────────┼───────────────┼────────────┤
│ SP             │    5202955.05 │       4.17 │
│ RJ             │    1824092.67 │       3.87 │
│ MG             │    1585308.03 │       4.14 │
└────────────────┴───────────────┴────────────┘
\end{codeblock}

\begin{itemize}
  \item Join thẳng hai fact làm doanh thu của SP thành 5.228.869,17 BRL thay vì 5.202.955,05 BRL, vì đơn có nhiều review làm dòng hàng bị lặp. Điểm trung bình thành 4,13 thay vì 4,17, vì mỗi review bị đếm một lần cho mỗi dòng hàng của đơn.
  \item Drill-across cho số đúng: SP được 5.202.955,05 BRL và 4,17 điểm, RJ được 1.824.092,67 BRL và 3,87 điểm, MG được 1.585.308,03 BRL và 4,14 điểm.
\end{itemize}

\begin{luuy}
\textbf{Bẫy thường gặp.}
\begin{itemize}
  \item \textbf{Join hai fact ở mức dòng.} Luôn gộp riêng từng fact về thuộc tính chung trước.
  \item \textbf{Ghép theo hai vai khác nhau của một thuộc tính.} Bang của khách và bang của người bán đều là ``bang'', nhưng ghép doanh thu theo bang người bán với điểm review theo bang khách là sai nghĩa.
  \item \textbf{Quên rằng trung bình phụ thuộc grain.} Bang AM được 4,18 điểm khi lấy trung bình 147 dòng review, nhưng được 4,21 điểm khi mỗi đơn chỉ giữ điểm của review mới nhất, trên 146 đơn. Ghi rõ định nghĩa ngay cạnh con số.
  % Kiểm: AM, 615 điểm trên 147 dòng review = 4,1837; review mới nhất theo thứ tự ở khối E, 614 điểm trên 146 đơn = 4,2055. AM lệch nhiều nhất trong 27 bang; RJ chỉ lệch từ 3,8750 lên 3,8771.
\end{itemize}
\end{luuy}

% =====================================================================
\section{Bài 8: Star, snowflake hay One Big Table}

\textbf{Vấn đề.} Phản xạ SWE của bạn nhìn \code{dim\_\allowbreak{}product} và thấy tên danh mục bị lặp trên 32.951 dòng, trong khi chỉ có 74 danh mục. Bạn muốn chuẩn hoá: tách ra \code{dim\_\allowbreak{}category} và để \code{dim\_\allowbreak{}product} chỉ giữ khoá. Một đồng nghiệp lại muốn đi hướng ngược lại: chép mọi thuộc tính của dimension lên từng dòng fact, để BA khỏi phải join. Ai đúng?

\textbf{Giải thích.} Tách một dimension thành các bảng con như vậy gọi là \textit{snowflake}. Kimball khuyên tránh snowflake: nó khó đọc với người dùng nghiệp vụ, có thể làm truy vấn chậm hơn, và một dimension phẳng chứa đúng cùng thông tin.
% Nguồn: Kimball techniques PDF, trang 9

Còn lý do tiết kiệm chỗ thì sao? Trên file dạng cột như Parquet, một cột có ít giá trị khác nhau gần như không tốn chỗ, nhờ \textit{dictionary encoding}. Mỗi giá trị khác nhau được lưu một lần trong một trang từ điển của cột, còn mỗi dòng chỉ lưu một số nguyên nhỏ trỏ vào từ điển. Module 4 sẽ giải thích kỹ cách mã hoá này.
% Nguồn: https://parquet.apache.org/docs/file-format/data-pages/encodings/

Hướng ngược lại có tên là \textit{One Big Table}, viết tắt OBT: một bảng rộng ở grain của fact, mang theo mọi thuộc tính của các dimension. Năm 2022, Michael Kaminsky so hai cách trên Redshift, Snowflake và BigQuery với dữ liệu TPC-DS. Các câu kiểu dashboard trên OBT chạy nhanh hơn star khoảng 25 tới 50\%, nhưng trên Redshift OBT chiếm khoảng gấp đôi chỗ của bảng fact. Ông gợi ý dựng star trước, rồi join lại thành OBT cho người dùng cuối.
% Nguồn: https://www.fivetran.com/blog/star-schema-vs-obt (bài viết của Michael Kaminsky, 16/08/2022): "25%-50% faster depending on the warehouse"; one_big_table khoảng 60.250 MB so với store_sales 29.778 MB; "staging your ELT process such that the data all get transformed into something like a star schema before everything gets re-joined back together for end-user querying".

\begin{ghichu}
\textbf{So với SWE.} Dictionary encoding giống string interning, hay giống việc bạn lưu một enum thay vì chuỗi: mỗi giá trị có một bản, phần còn lại chỉ là tham chiếu. Nói cách khác, Parquet đã tự ``chuẩn hoá'' cột đó ở tầng lưu trữ, nên bạn không cần chuẩn hoá ở tầng mô hình.

OBT giống một document nhúng sẵn dữ liệu liên quan. Tài liệu MongoDB khuyên nhúng khi hai thực thể có quan hệ ``chứa'', như document \code{contacts} chứa \code{address}, để đọc mọi thứ trong một lần truy cập.
% Nguồn: https://www.mongodb.com/docs/manual/data-modeling/embedding/

\textbf{Chỗ so sánh gãy.} Khác string interning, từ điển của Parquet chỉ nằm trong từng column chunk, và bị bỏ khi cột có quá nhiều giá trị khác nhau, như cột \code{product\_\allowbreak{}id} ở khối dưới.

Lý do chuẩn hoá trong OLTP không phải là chỗ lưu, mà là tránh sửa sót: sửa tên danh mục ở một bảng rồi quên bảng khác. Dimension do code của pipeline cập nhật, không ai sửa tay, nên lý do đó cũng không còn.

OBT thì làm lý do đó quay lại. Một dòng hàng không ``chứa'' danh mục của sản phẩm; OBT chỉ chép tên danh mục lên mọi dòng hàng của nó, nên đổi tên một danh mục là phải ghi lại hàng nghìn dòng, như khối thứ hai cho thấy.
\end{ghichu}

\subsection{Làm thử}

Khối đầu ghi bản star và bản snowflake ra Parquet trong một thư mục tạm, với nén mặc định của DuckDB, rồi so kích thước.

% chạy: I
\begin{codeblock}
lakehouse/.venv/bin/python - <<'EOF'
import duckdb, os, tempfile
D = 'lakehouse/source_sim/data/'
con = duckdb.connect()
con.sql(f"CREATE VIEW products AS SELECT * FROM '{D}olist_products_dataset.csv'")
con.sql(f"CREATE VIEW category_translation AS "
        f"SELECT * FROM '{D}product_category_name_translation.csv'")
# Star: tên danh mục nằm thẳng trong dim_product
con.sql("""CREATE TABLE dim_product AS
  SELECT p.product_id, coalesce(t.product_category_name_english,
                                p.product_category_name, 'unknown') AS category
  FROM products p LEFT JOIN category_translation t USING (product_category_name)""")
# Snowflake: tách danh mục ra bảng riêng, dim_product chỉ giữ khoá
con.sql("""CREATE TABLE dim_category AS
  SELECT row_number() OVER (ORDER BY category) AS category_id, category
  FROM (SELECT DISTINCT category FROM dim_product)""")
con.sql("""CREATE TABLE dim_product_sf AS
  SELECT p.product_id, c.category_id
  FROM dim_product p JOIN dim_category c USING (category)""")
with tempfile.TemporaryDirectory() as tmp:
    for t in ['dim_product', 'dim_product_sf', 'dim_category']:
        f = f'{tmp}/{t}.parquet'
        con.sql(f"COPY {t} TO '{f}' (FORMAT parquet)")
        n = con.sql(f"SELECT count(*) FROM {t}").fetchall()[0][0]
        print(f'{t:15} {n:6} dòng {os.path.getsize(f):10,} byte')
    con.sql(f"""SELECT path_in_schema AS cot, encodings,
                       total_compressed_size AS byte
                FROM parquet_metadata('{tmp}/dim_product.parquet')""").show()
EOF
\end{codeblock}
\textit{Kết quả:}
% kết quả: I
\begin{codeblock}
dim_product      32951 dòng  1,092,443 byte
dim_product_sf   32951 dòng  1,091,566 byte
dim_category        74 dòng      1,820 byte
┌────────────┬──────────────────┬─────────┐
│    cot     │    encodings     │  byte   │
│  varchar   │     varchar      │  int64  │
├────────────┼──────────────────┼─────────┤
│ product_id │ PLAIN            │ 1061589 │
│ category   │ PLAIN_DICTIONARY │   30231 │
└────────────┴──────────────────┴─────────┘
\end{codeblock}

\begin{itemize}
  \item Bản star: \code{dim\_\allowbreak{}product} nặng 1.092.443 byte.
  \item Bản snowflake: \code{dim\_\allowbreak{}product\_\allowbreak{}sf} nặng 1.091.566 byte, cộng \code{dim\_\allowbreak{}category} nặng 1.820 byte, tổng 1.093.386 byte. Bản snowflake nặng hơn 943 byte.
  \item Cột \code{category} chỉ chiếm 30.231 byte, khoảng 2,8\% file, vì nó được ghi bằng dictionary encoding. DuckDB in ra \code{PLAIN\_\allowbreak{}DICTIONARY}, giá trị cũ mà spec Parquet đã đánh dấu deprecated; Module 4 dùng giá trị hiện tại là \code{RLE\_\allowbreak{}DICTIONARY}. Phần lớn file là cột \code{product\_\allowbreak{}id}, 1.061.589 byte.
\end{itemize}

Snowflake không tiết kiệm được gì ở đây, lại thêm một join cho mọi câu hỏi theo danh mục.

Khối thứ hai dựng OBT cho dòng hàng, rồi đo hai cái giá của nó: chỗ lưu, và số dòng phải sửa khi một dimension đổi.

% chạy: I2
\begin{codeblock}
lakehouse/.venv/bin/python - <<'EOF'
import duckdb, tempfile
D = 'lakehouse/source_sim/data/'
con = duckdb.connect()
for t in ['orders', 'order_items', 'customers', 'products', 'sellers']:
    con.sql(f"CREATE VIEW {t} AS SELECT * FROM '{D}olist_{t}_dataset.csv'")
con.sql(f"CREATE VIEW category_translation AS "
        f"SELECT * FROM '{D}product_category_name_translation.csv'")
con.sql("""CREATE TABLE dim_product AS
  SELECT p.product_id, coalesce(t.product_category_name_english,
                                p.product_category_name, 'unknown') AS category
  FROM products p LEFT JOIN category_translation t USING (product_category_name)""")
# OBT: thuộc tính của dimension chép lên từng dòng hàng
con.sql("""CREATE TABLE obt_order_items AS
  SELECT i.order_id, i.order_item_id, i.product_id, i.seller_id, o.customer_id,
         CAST(o.order_purchase_timestamp AS DATE) AS purchase_date,
         CAST(i.price AS DECIMAL(12,2)) AS price, p.category,
         s.seller_city, s.seller_state, c.customer_city, c.customer_state
  FROM order_items i JOIN orders o USING (order_id)
  JOIN dim_product p USING (product_id) JOIN sellers s USING (seller_id)
  JOIN customers c USING (customer_id) ORDER BY i.order_id, i.order_item_id""")
with tempfile.TemporaryDirectory() as tmp:
    con.sql(f"COPY obt_order_items TO '{tmp}/obt.parquet' (FORMAT parquet)")
    total, copied = con.sql(f"""SELECT sum(total_compressed_size),
        sum(total_compressed_size) FILTER (WHERE path_in_schema IN ('category',
          'seller_city', 'seller_state', 'customer_city', 'customer_state'))
        FROM parquet_metadata('{tmp}/obt.parquet')""").fetchall()[0]
print(f'OBT trên Parquet: {total:,} byte; 5 cột chép từ dimension: '
      f'{copied:,} byte ({copied / total:.1%})')
# Đổi tên một danh mục: star sửa dimension, OBT sửa mọi dòng hàng của danh mục
print('dòng phải sửa khi đổi tên bed_bath_table, ở star và ở OBT:', con.sql("""
  SELECT (SELECT count(*) FROM dim_product WHERE category = 'bed_bath_table'),
         (SELECT count(*) FROM obt_order_items
          WHERE category = 'bed_bath_table')""").fetchall()[0])
EOF
\end{codeblock}
\textit{Kết quả:}
% kết quả: I2
\begin{codeblock}
OBT trên Parquet: 10,467,714 byte; 5 cột chép từ dimension: 585,411 byte (5.6%)
dòng phải sửa khi đổi tên bed_bath_table, ở star và ở OBT: (3029, 11115)
\end{codeblock}

\begin{itemize}
  \item Năm cột chép từ dimension chỉ tốn 585.411 byte, khoảng 5,6\% file, cũng nhờ dictionary encoding. Trên Parquet, chép chữ lặp lại gần như không tốn chỗ.
  \item Kaminsky dựng bảng bằng một công cụ không đặt được kiểu mã hoá cột trên Redshift, và chính ông ghi rằng số chỗ lưu đo được có lẽ là trường hợp xấu nhất.
  % Nguồn: https://www.fivetran.com/blog/star-schema-vs-obt: "Because dbt doesn't have the ability to specify column compression or encoding style in Redshift, this is probably the worst-possible-case in terms of disk storage size."
  \item Cái giá thật nằm ở cập nhật. Đổi tên danh mục \code{bed\_\allowbreak{}bath\_\allowbreak{}table} là sửa 3.029 dòng của \code{dim\_\allowbreak{}product}, hoặc 1 dòng của \code{dim\_\allowbreak{}category} nếu dùng snowflake, nhưng là ghi lại 11.115 dòng của OBT.
\end{itemize}

Câu hỏi theo danh mục và bang trên OBT không cần join nào, vì mọi cột đã nằm trên cùng một dòng. Kết luận cho Olist: giữ star làm mô hình chính; nếu một dashboard cần OBT, dựng nó từ star như một mart, và dựng lại nó mỗi khi dimension đổi.

\begin{luuy}
\textbf{Bẫy thường gặp.}
\begin{itemize}
  \item \textbf{Tối ưu trước khi đo.} Khối trên chạy chưa tới một giây trên máy này. Hãy đo trước khi tách bảng.
  \item \textbf{Snowflake vì nguồn đã chuẩn hoá.} Nguồn chuẩn hoá vì nó phục vụ việc ghi; star schema phục vụ việc đọc.
  \item \textbf{Đem kết luận này sang dimension rất lớn mà không đo.} Dimension lớn nhất ở đây chỉ có 99.441 dòng. Với dimension lớn hơn nhiều, hãy đo như trên rồi mới quyết định có tách hay không.
  \item \textbf{Nhét cột của grain khác vào OBT.} Thêm \code{payment\_\allowbreak{}value} vào OBT ở grain dòng hàng là lặp lại đúng lỗi ở phần mở đầu: tổng tiền tăng 26,9\%.
  \item \textbf{Coi OBT là nơi giữ dữ liệu gốc.} OBT là bản sao để đọc. Sửa trên OBT mà không sửa star thì lần dựng lại sau sẽ xoá mất chỗ sửa đó.
\end{itemize}
\end{luuy}

% =====================================================================
\section{Bài 9: Slowly changing dimension}

\textbf{Vấn đề.} Người mua \code{37266d5c6bf5cb0700372d16599a9b21} đặt đơn đầu tiên vào ngày 2017-10-08 ở timoteo, zip 35180, rồi đặt đơn thứ hai ngày 2018-06-17 ở coronel fabriciano, zip 35170. Nếu \code{dim\_\allowbreak{}customer} chỉ giữ địa chỉ mới nhất, như \repo{core.customers} của repo \repo{lakehouse} hiện nay, đơn năm 2017 bị tính cho coronel fabriciano. Trên cả Olist, 250 người mua có hơn một zip.
% Nguồn: lakehouse/docs/data-model.md, mục 6, core.customers lấy địa chỉ ở đơn mới nhất.

\textbf{Giải thích.} \textit{Slowly changing dimension}, viết tắt SCD, là cách một dimension xử lý khi thuộc tính của một thực thể đổi theo thời gian. Kimball đánh số các cách làm từ Type 0 tới Type 7. Module này dùng năm cách dưới; Type 4, 5 và 7 nằm ngoài phạm vi module.
% Nguồn: Kimball techniques PDF, trang 12 và 13

\begin{xltabular}{\textwidth}{@{}P{1.6cm} L P{4.6cm}@{}}
\toprule
\thd{kiểu} & \thd{làm gì} & \thd{đơn 2017 của người mua trên được tính vào} \\
\midrule
\endhead
\bottomrule
\endfoot
Type 0 & giữ giá trị gốc, không bao giờ đổi & timoteo \\
Type 1 & ghi đè bằng giá trị mới, mất lịch sử & coronel fabriciano \\
Type 2 & thêm một dòng cho mỗi version, có ngày bắt đầu hiệu lực, ngày hết hiệu lực và cờ dòng hiện hành & timoteo \\
Type 3 & ghi đè cột chính, thêm một cột giữ giá trị ngay trước đó & coronel fabriciano theo cột chính, timoteo theo cột giá trị trước; chỉ nhớ một lần đổi \\
Type 6 & Type 2, cộng thêm các cột giá trị hiện hành được ghi đè trên mọi version & cả hai: timoteo theo version lúc đặt, coronel fabriciano theo giá trị hiện tại \\
\end{xltabular}

Type 2 cần khoá thay thế cho mỗi version, ví dụ \code{customer\_\allowbreak{}sk}, vì \code{customer\_\allowbreak{}unique\_\allowbreak{}id} giờ lặp lại trên nhiều dòng. Fact trỏ tới version còn hiệu lực lúc sự kiện xảy ra.

Khoảng hiệu lực là khoảng nửa mở \code{[valid\_\allowbreak{}from, valid\_\allowbreak{}to)}: tính cả điểm đầu, không tính điểm cuối. Hai version liền nhau dùng chung một mốc ranh giới nhưng không bao giờ chồng nhau, nên mỗi thời điểm thuộc đúng một version. Dòng hiện hành để \code{valid\_\allowbreak{}to} là NULL, hoặc là một ngày rất xa như 9999-12-31; Microsoft nhắc tới cả hai cách.
% Nguồn: https://learn.microsoft.com/en-us/power-bi/guidance/star-schema (mục Type 2 SCD)

Một trường hợp liên quan là \textit{late arriving dimension}: fact tới trước dimension của nó. Kimball tạo một dòng dimension giữ chỗ chứa khoá tự nhiên và các thuộc tính ``chưa biết'', rồi ghi đè kiểu Type 1 khi thông tin tới. Repo \repo{lakehouse} chọn cách khác: \repo{core.orders} để \code{customer\_\allowbreak{}key} là NULL khi customer chưa tới, rồi tính lại đơn khi customer tới.
% Nguồn: Kimball techniques PDF, trang 20; lakehouse/docs/data-model.md, mục 6, core.orders.

Lịch sử Type 2 và các data test ở Bài 10 cũng viết được bằng dbt, công cụ mà nhánh tuỳ chọn T1 dạy; module này tự viết bằng SQL để bạn thấy từng bước.
% Nguồn: https://docs.getdbt.com/docs/build/snapshots; https://docs.getdbt.com/docs/build/data-tests

\begin{ghichu}
\textbf{So với SWE.} Type 1 là câu \code{UPDATE customers SET city = ...} mà app của bạn vẫn chạy mỗi ngày.

Type 2 giống S3 Versioning. Bucket bật versioning giữ một version hiện hành cùng không hoặc nhiều version cũ của mỗi object; mỗi version có version ID riêng, còn key của object giữ nguyên. Ở đây key là khoá bền \code{customer\_\allowbreak{}unique\_\allowbreak{}id}, còn version ID là khoá thay thế \code{customer\_\allowbreak{}sk}.
% Nguồn: https://docs.aws.amazon.com/AmazonS3/latest/userguide/Versioning.html

Khoảng nửa mở giống lịch đặt phòng: hai lượt \code{[14:30, 15:30)} và \code{[15:30, 16:30)} nối nhau mà không chồng. Hàm tạo range của PostgreSQL mặc định dùng đúng quy ước này. PostgreSQL 18 còn có ràng buộc \code{UNIQUE (id, valid\_\allowbreak{}at WITHOUT OVERLAPS)}, đúng luật ``một thực thể không có hai version chồng thời gian'' của Type 2. Cột \code{valid\_\allowbreak{}at} phải có kiểu range như \code{tstzrange}, và cột \code{id}, nếu là số hay chuỗi, cần extension \code{btree\_\allowbreak{}gist}.
% Nguồn: https://www.postgresql.org/docs/18/rangetypes.html; https://www.postgresql.org/docs/18/sql-createtable.html

\textbf{Chỗ so sánh gãy.} S3 tạo một version cho mỗi lệnh ghi. Type 2 chỉ thấy các thay đổi mà pipeline kịp nhìn thấy giữa hai lần nạp: nếu địa chỉ đổi hai lần trong một ngày mà pipeline chạy mỗi ngày một lần, version ở giữa mất vĩnh viễn. Fact cũng phải giữ đúng khoá của version, hoặc tra theo thời gian lúc nạp; S3 không có khái niệm này.
\end{ghichu}

\subsection{Làm thử}

Bảng \code{customers} của Olist có một dòng cho mỗi đơn, nên nó chụp lại địa chỉ của người mua tại lúc đặt. Khối dưới dùng các ảnh chụp đó để dựng Type 1 và Type 2, rồi đếm sai số. Thay đổi được theo dõi theo zip.

% chạy: J
\begin{codeblock}
lakehouse/.venv/bin/python - <<'EOF'
import duckdb
D = 'lakehouse/source_sim/data/'
con = duckdb.connect()
for t in ['orders', 'order_items', 'customers']:
    con.sql(f"CREATE VIEW {t} AS SELECT * FROM '{D}olist_{t}_dataset.csv'")
one = lambda q: con.sql(q).fetchall()[0]
# Mỗi đơn chụp lại địa chỉ của người mua tại lúc đặt
con.sql("""CREATE TABLE addr AS
  SELECT o.order_id, o.order_purchase_timestamp AS ts, c.customer_id,
         c.customer_unique_id AS buyer, c.customer_zip_code_prefix AS zip,
         c.customer_city AS city, c.customer_state AS state
  FROM orders o JOIN customers c USING (customer_id)""")

# Type 1: chỉ giữ địa chỉ mới nhất, giống một câu UPDATE ghi đè
con.sql("""CREATE TABLE dim_customer_t1 AS
  SELECT buyer, zip, city, state FROM addr
  QUALIFY row_number() OVER (PARTITION BY buyer
                             ORDER BY ts DESC, customer_id DESC) = 1""")
print('Type 1, số đơn sai thành phố, sai bang:', one("""
  SELECT count(*) FILTER (WHERE a.city <> t.city),
         count(*) FILTER (WHERE a.state <> t.state)
  FROM addr a JOIN dim_customer_t1 t USING (buyer)"""))
print('Type 1, tiền hàng bị tính sang bang khác:', one("""
  SELECT sum(CAST(i.price AS DECIMAL(14,2)))
  FROM order_items i JOIN addr a USING (order_id)
  JOIN dim_customer_t1 t USING (buyer)
  WHERE a.state <> t.state""")[0])

# Type 2: mỗi lần zip đổi là một version, hiệu lực trong [valid_from, valid_to)
con.sql("""CREATE TABLE dim_customer_t2 AS
  WITH marked AS (
    SELECT *, zip IS DISTINCT FROM lag(zip) OVER (PARTITION BY buyer ORDER BY ts)
                AS is_new_version
    FROM addr)
  SELECT row_number() OVER (ORDER BY buyer, ts) AS customer_sk,
         buyer, zip, city, state,
         ts AS valid_from,
         lead(ts) OVER (PARTITION BY buyer ORDER BY ts) AS valid_to,
         lead(ts) OVER (PARTITION BY buyer ORDER BY ts) IS NULL AS is_current
  FROM marked WHERE is_new_version""")
print('Type 2, số version, số người mua, số dòng hiện hành:', one("""
  SELECT count(*), count(DISTINCT buyer), count(*) FILTER (WHERE is_current)
  FROM dim_customer_t2"""))

# Mỗi đơn phải khớp đúng một version
match = """SELECT a.order_id, count(t.customer_sk) AS n
           FROM addr a LEFT JOIN dim_customer_t2 t ON t.buyer = a.buyer AND {cond}
           GROUP BY 1"""
half_open = "a.ts >= t.valid_from AND (t.valid_to IS NULL OR a.ts < t.valid_to)"
closed = ("a.ts BETWEEN t.valid_from "
          "AND coalesce(t.valid_to, TIMESTAMP '9999-12-31')")
for name, cond in [('nửa mở [from, to)', half_open),
                   ('đóng, dùng BETWEEN', closed)]:
    print(name, '(số version khớp, số đơn):', con.sql(
        f"SELECT n, count(*) FROM ({match.format(cond=cond)}) GROUP BY 1 ORDER BY 1"
    ).fetchall())

example = con.sql("""
  SELECT customer_sk, zip, city, valid_from, valid_to, is_current
  FROM dim_customer_t2 WHERE buyer = '37266d5c6bf5cb0700372d16599a9b21'""")
for row in example.fetchall():
    print(*row, sep=' | ')
EOF
\end{codeblock}
\textit{Kết quả:}
% kết quả: J
\begin{codeblock}
Type 1, số đơn sai thành phố, sai bang: (130, 43)
Type 1, tiền hàng bị tính sang bang khác: 6164.99
Type 2, số version, số người mua, số dòng hiện hành: (96352, 96096, 96096)
nửa mở [from, to) (số version khớp, số đơn): [(1, 99441)]
đóng, dùng BETWEEN (số version khớp, số đơn): [(1, 99185), (2, 256)]
20825 | 35180 | timoteo | 2017-10-08 21:38:21 | 2018-06-17 23:51:02 | False
20826 | 35170 | coronel fabriciano | 2018-06-17 23:51:02 | None | True
\end{codeblock}

\begin{itemize}
  \item Type 1 tính sai thành phố của 130 đơn và sai bang của 43 đơn; 6.164,99 BRL tiền hàng bị tính sang bang khác. Doanh thu của SP theo Type 1 là 5.205.024,43 BRL thay vì 5.202.955,05 BRL.
  % Kiểm: tổng price theo bang của bản Type 1, so với bang lúc đặt; chạy ngày 05/11/2026.
  \item Type 2 có 96.352 version cho 96.096 người mua, và mỗi người mua có đúng một dòng hiện hành. Với khoảng nửa mở, cả 99.441 đơn khớp đúng một version. Doanh thu theo bang qua bản Type 2 bằng đúng các con số ở Bài 7.
  % Kiểm: bang theo version Type 2 khác bang lúc đặt ở 0 đơn.
  \item Với \code{BETWEEN}, tức khoảng đóng ở cả hai đầu, 256 đơn khớp hai version. Đó là 256 đơn nằm đúng ranh giới, mỗi lần đổi zip một đơn. Fact join theo cách này sẽ bị fan-out như ở Bài 3.
  \item Hai dòng cuối là hai version của người mua ở phần vấn đề, mỗi version một \code{customer\_\allowbreak{}sk}.
\end{itemize}

\begin{luuy}
\textbf{Bẫy thường gặp.}
\begin{itemize}
  \item \textbf{Dùng \code{BETWEEN} cho khoảng hiệu lực.} Dùng \code{>=} ở đầu và \code{<} ở cuối.
  \item \textbf{Quên dòng hiện hành có \code{valid\_\allowbreak{}to} là NULL.} Điều kiện \code{ts < valid\_\allowbreak{}to} loại mọi dòng hiện hành; hãy thêm \code{valid\_\allowbreak{}to IS NULL OR}, hoặc dùng một ngày rất xa.
  \item \textbf{Theo dõi cột nhiễu.} Nếu so cả tên thành phố, bạn được 96.355 version thay vì 96.352, vì 2 người mua có cùng zip nhưng tên thành phố ghi khác nhau, như planaltina và planaltina de goias. Chỉ theo dõi cột mà thay đổi của nó có nghĩa với nghiệp vụ.
  % Kiểm: dựng Type 2 theo (zip, city, state) cho 96.355 version.
  \item \textbf{Dựng Type 2 khi không ai cần lịch sử.} Type 2 làm bảng to hơn và join khó hơn. Hỏi BA trước: báo cáo cần địa chỉ lúc đặt hay địa chỉ hiện tại?
\end{itemize}
\end{luuy}

% =====================================================================
\section{Bài 10: Mô hình trong lakehouse và test dữ liệu}

\textbf{Vấn đề.} Trên PostgreSQL ở Module 1, ràng buộc \code{PRIMARY KEY} và \code{FOREIGN KEY} chặn dòng sai ngay lúc ghi. Trong lakehouse, không ai chặn giúp bạn. Spec của Iceberg ghi rằng tính duy nhất theo khoá định danh không được Iceberg bảo đảm hay đòi hỏi, và việc đó thuộc trách nhiệm của engine hoặc bên cung cấp dữ liệu. Databricks có cú pháp khoá chính và khoá ngoại, nhưng tài liệu của họ viết rằng các ràng buộc này chỉ để thông tin và không được thực thi. Vậy ai giữ cho khoá của star schema đúng?
% Nguồn: https://iceberg.apache.org/spec/#identifier-field-ids; https://docs.databricks.com/aws/en/tables/constraints

\textbf{Giải thích.} Star schema nằm ở tầng gold. Databricks mô tả gold là nơi mô hình hoá dữ liệu bằng dimensional model, thiết lập quan hệ và định nghĩa measure. Star schema được dựng từ tầng silver đã làm sạch, không đọc thẳng từ OLTP.
% Nguồn: https://docs.databricks.com/aws/en/lakehouse/medallion

Trong repo \repo{lakehouse}, tầng silver \repo{core} đã gần như một star: dimension \code{customers}, \code{sellers}, \code{products}, fact \code{orders}, \code{order\_\allowbreak{}items}, và fact không chép cột của dimension. Tầng gold, gồm \code{analytics} và \code{marts}, còn nằm trong kế hoạch. Chi tiết ở mục 6 và mục 11 của \repo{docs/data-model.md}; Lab 3.5 dựng thử \code{marts} trên DuckDB.

Thứ giữ cho khoá đúng là \textit{data test}: một câu \code{SELECT} đi tìm các dòng ``lỗi'', tức các dòng bác bỏ điều bạn khẳng định; trả về 0 dòng nghĩa là đạt. Một star schema cần ít nhất các loại sau:
% Nguồn: https://docs.getdbt.com/docs/build/data-tests

\begin{xltabular}{\textwidth}{@{}P{2.4cm} L L@{}}
\toprule
\thd{loại test} & \thd{khẳng định} & \thd{bắt được lỗi gì} \\
\midrule
\endhead
\bottomrule
\endfoot
grain & khoá của fact và dimension không trùng & nạp một lô hai lần \\
quan hệ & mọi khoá trên fact có dòng ở dimension, không có NULL & dimension thiếu dòng, khoá quên đổi về 0 \\
đối chiếu & số dòng và tổng tiền ở fact bằng ở nguồn & fan-out, join làm rơi dòng \\
SCD Type 2 & mỗi khoá bền có đúng một dòng hiện hành, các version không chồng nhau & lần nạp đóng thiếu hoặc mở thừa một version \\
\end{xltabular}

\begin{ghichu}
\textbf{So với SWE.} Data test giống câu assert trong unit test cộng với checksum sau khi chép file. Test grain và test quan hệ là các assert về cấu trúc. Test đối chiếu là checksum: tổng \code{price} ở fact bằng tổng ở nguồn nghĩa là không dòng nào bị nhân lên hay rơi mất trên đường đi.

\textbf{Chỗ so sánh gãy.} Unit test chạy trên fixture cố định, và đỏ nghĩa là code sai. Data test chạy trên dữ liệu thật ở mỗi lần build, nên đỏ có thể do nguồn đổi chứ không do code của bạn. Bạn phải quyết định dừng pipeline hay chỉ cảnh báo, và đó là chủ đề của Module 10.
\end{ghichu}

\subsection{Làm thử}

Khối dưới chạy ba test trên một fact nhỏ, cố ý làm hỏng dữ liệu, rồi chạy lại. Đây là vòng đỏ và xanh quen thuộc của TDD, chỉ khác là thứ bị kiểm là dữ liệu.

% chạy: K
\begin{codeblock}
lakehouse/.venv/bin/python - <<'EOF'
import duckdb
D = 'lakehouse/source_sim/data/'
con = duckdb.connect()
for t in ['order_items', 'products']:
    con.sql(f"CREATE VIEW {t} AS SELECT * FROM '{D}olist_{t}_dataset.csv'")
con.sql("CREATE TABLE dim_product AS SELECT product_id FROM products")
con.sql("""CREATE TABLE fct_order_items AS
  SELECT order_id, order_item_id, product_id, CAST(price AS DECIMAL(12,2)) AS price
  FROM order_items""")
# Mỗi test là một câu SELECT trả về các dòng lỗi; 0 dòng nghĩa là đạt
tests = {
  'grain: (order_id, order_item_id) không trùng': """
     SELECT order_id, order_item_id FROM fct_order_items
     GROUP BY ALL HAVING count(*) > 1""",
  'khoá ngoại: product_id nào cũng có trong dim_product': """
     SELECT f.product_id FROM fct_order_items f
     ANTI JOIN dim_product p USING (product_id)""",
  'đối chiếu: tổng price bằng nguồn': """
     SELECT 1 WHERE (SELECT sum(price) FROM fct_order_items)
       IS DISTINCT FROM
       (SELECT sum(CAST(price AS DECIMAL(14,2))) FROM order_items)""",
}
def run_tests(label):
    print(label)
    for name, sql in tests.items():
        n = len(con.sql(sql).fetchall())
        print(f"  {'ĐẠT ' if n == 0 else 'LỖI '} {name} ({n} dòng lỗi)")
run_tests('Lần 1, dữ liệu sạch:')
# Cố ý làm hỏng: chép lại một dòng, và thêm một dòng trỏ tới sản phẩm không tồn tại
con.sql("INSERT INTO fct_order_items SELECT * FROM fct_order_items LIMIT 1")
con.sql("INSERT INTO fct_order_items VALUES ('x', 1, 'khong_co_san_pham', 10.00)")
run_tests('Lần 2, sau khi làm hỏng:')
EOF
\end{codeblock}
\textit{Kết quả:}
% kết quả: K
\begin{codeblock}
Lần 1, dữ liệu sạch:
  ĐẠT  grain: (order_id, order_item_id) không trùng (0 dòng lỗi)
  ĐẠT  khoá ngoại: product_id nào cũng có trong dim_product (0 dòng lỗi)
  ĐẠT  đối chiếu: tổng price bằng nguồn (0 dòng lỗi)
Lần 2, sau khi làm hỏng:
  LỖI  grain: (order_id, order_item_id) không trùng (1 dòng lỗi)
  LỖI  khoá ngoại: product_id nào cũng có trong dim_product (1 dòng lỗi)
  LỖI  đối chiếu: tổng price bằng nguồn (1 dòng lỗi)
\end{codeblock}

Lần đầu cả ba test đạt. Lần sau, dòng bị chép làm test grain đỏ, dòng trỏ tới sản phẩm không tồn tại làm test quan hệ đỏ, và cả hai dòng thừa làm tổng tiền không còn khớp nên test đối chiếu cũng đỏ. Mỗi test chỉ là một câu \code{SELECT}, rẻ hơn nhiều so với việc dò lỗi trên một dashboard đã sai.

\begin{luuy}
\textbf{Bẫy thường gặp.}
\begin{itemize}
  \item \textbf{Tin ràng buộc khoá trên bảng lakehouse.} Iceberg không ép tính duy nhất, và khoá của Databricks chỉ để thông tin.
  \item \textbf{Chỉ test cấu trúc, không test đối chiếu.} Test grain vẫn đạt khi một \code{INNER JOIN} làm rơi hẳn một nhóm dòng, như 775 đơn không có dòng hàng ở Bài 3. Chỉ test đối chiếu số dòng và tổng tiền mới bắt được.
  \item \textbf{Đối chiếu fact với chính nó.} Test đối chiếu phải so với nguồn, không so với một bảng dựng từ cùng câu SQL.
\end{itemize}
\end{luuy}

% =====================================================================
\section{Bài 11: Kimball, Inmon và Data Vault}

\textbf{Vấn đề.} Khi đọc tài liệu DE hay đi phỏng vấn, bạn sẽ gặp hai cái tên khác ngoài Kimball: Inmon và Data Vault. Một team khác có thể đã chọn một trong hai. Bạn cần biết mỗi cách đặt dữ liệu chi tiết ở đâu và người phân tích đọc bảng nào, để không tranh luận chỉ vì tên gọi.

\textbf{Giải thích.} Ba cách khác nhau chủ yếu ở chỗ dữ liệu chi tiết nhất được lưu dưới dạng nào.
\begin{itemize}
  \item \textbf{Kimball}, cách của module này: dữ liệu chi tiết được lưu ngay ở dạng dimension và fact, mỗi quy trình một star, các star ghép với nhau qua conformed dimension.
  \item \textbf{Inmon}, qua mô hình Corporate Information Factory, viết tắt CIF: dữ liệu chi tiết vào một kho trung tâm chuẩn hoá theo 3NF, và data mart của từng phòng ban chỉ được nạp từ kho đó.
  % Nguồn: https://williaminmon.substack.com/p/a-tale-of-two-architectures-kimball (bài viết của William Inmon, 23/12/2025): "the data in the corporate information factory data warehouse is stored in a normalized relational format"; "the only source of data for the data marts is the data warehouse".
  \item Margy Ross của Kimball Group đặt khác biệt cốt lõi ở cùng chỗ: Kimball đòi dữ liệu chi tiết có dạng dimension, còn CIF đặt nó trong kho chuẩn hoá; mart của CIF chứa dữ liệu đã gộp, dạng dimension.
  % Nguồn: https://decisionworks.com/2004/03/differences-of-opinion/ (bài viết của Margy Ross, 03/03/2004): "a third normal form (3NF) relational database containing atomic data"; "marts are tailored by business department/function with dimensionally structured summary data"; "The Kimball approach says atomic data must be dimensionally structured." Bài không nhắc tên Inmon, chỉ nói CIF.
  \item \textbf{Data Vault}, của Dan Linstedt, tách dữ liệu thành ba loại bảng:
  \begin{itemize}
    \item \textit{Hub} giữ danh sách khoá nghiệp vụ không trùng, kèm lúc khoá tới kho và hệ nguồn của nó.
    \item \textit{Link} giữ quan hệ nhiều nhiều giữa các hub.
    \item \textit{Satellite} giữ thuộc tính mô tả của một hub hoặc một link. Mỗi lần thuộc tính đổi, một dòng mới được thêm vào, kèm giờ nạp và hệ nguồn.
  \end{itemize}
  % Nguồn: https://dataversity.net/data-vault-series-1-data-vault-overview/5054 (bài viết của Dan Linstedt, TDAN.com, 07/2002): "a unique list of business keys"; "Load Date Time Stamp, recording when the key itself first arrived in the warehouse"; "a physical representation of a many-to-many 3NF relationship"; "Satellites may be connected to Hubs or Links"; "the new row is always inserted"; hub, link và satellite đều có "Record Source". Kiểm lại ngày 06/11/2026.
\end{itemize}

Hướng dẫn của Databricks cho Data Vault đặt mô hình dimension ở tầng gold, phía trên vault, và nhắc rằng Data Vault khuyên dùng mã băm của khoá nghiệp vụ làm khoá chính.
% Nguồn: https://www.databricks.com/blog/2022/06/24/prescriptive-guidance-for-implementing-a-data-vault-model-on-the-databricks-lakehouse-platform.html (24/06/2022): "Data Vault modeling recommends using a hash of business keys as the primary keys"; "loading a dimensional model Data Warehouse in Gold layer".

Repo \repo{lakehouse} đi gần Kimball hơn: \repo{core} đã có dạng dimension và fact, và Gold sẽ đọc thẳng \repo{core}. Nhưng nó có một chi tiết giống Data Vault: \code{customer\_\allowbreak{}key} của \repo{semantic.customer\_\allowbreak{}identity} là mã băm SHA-256 của \code{customer\_\allowbreak{}unique\_\allowbreak{}id}, như khoá băm của một hub.
% Nguồn: lakehouse/docs/data-model.md, mục 5, semantic.customer_identity: customer_key = sha2(customer_unique_id, 256).

Phác thảo Data Vault cho đơn hàng và người mua của Olist. Như sơ đồ ở Bài 1, dấu \code{<} hay \code{>} nằm ở phía có nhiều dòng:

\begin{codeblock}[samepage=true]
hub_customer               link_order_customer        hub_order
┌────────────────────┐     ┌────────────────────┐     ┌────────────────────┐
│ customer_hk        │────<│ customer_hk        │     │ order_hk           │
│ customer_unique_id │     │ order_hk           │>────│ order_id           │
│ load_ts            │     │ load_ts            │     │ load_ts            │
│ record_source      │     │ record_source      │     │ record_source      │
└─────────┬──────────┘     └────────────────────┘     └─────────┬──────────┘
          │ một hub, nhiều dòng satellite                       │
┌─────────┴──────────┐                                ┌─────────┴──────────┐
│ customer_hk        │                                │ order_hk           │
│ load_ts            │                                │ load_ts            │
│ record_source      │                                │ record_source      │
│ zip, city, state   │                                │ trạng thái, mốc    │
└────────────────────┘                                └────────────────────┘
sat_customer_address                                  sat_order_status
\end{codeblock}

\begin{ghichu}
\textbf{So với SWE.} Ba loại bảng của Data Vault đều quen với bạn. Hub giống một bảng chỉ có id và khoá nghiệp vụ. Link giống bảng nối mà Django tự tạo cho một \code{ManyToManyField}: tài liệu Django ghi rằng phía sau, Django tạo một bảng trung gian để biểu diễn quan hệ nhiều nhiều. Satellite giống một bảng lịch sử chỉ thêm, như audit log của một bảng.
% Nguồn: https://docs.djangoproject.com/en/6.1/ref/models/fields/ (ManyToManyField, Database Representation): "Behind the scenes, Django creates an intermediary join table to represent the many-to-many relationship."

\textbf{Chỗ so sánh gãy.} Trong app, bảng nối và bảng lịch sử là chi tiết nằm sau ORM; trong vault, chúng là toàn bộ mô hình. Một câu đơn giản như ``doanh thu theo bang của người mua'' phải đi qua hub đơn, link, hub người mua, rồi chọn đúng dòng satellite theo thời gian, trong khi star chỉ cần một join. Vault được thiết kế để nạp và kiểm toán, không để người phân tích đọc trực tiếp.
\end{ghichu}

Một bản dựng thử từ CSV, lấy giờ nạp là giờ đặt đơn như cách \repo{source\_\allowbreak{}sim} phát customer, có 96.096 dòng \code{hub\_\allowbreak{}customer}, đúng bằng số người mua, và 99.441 dòng cho \code{hub\_\allowbreak{}order} cũng như cho link. Satellite địa chỉ có 96.355 dòng, đúng bằng số version Type 2 khi theo dõi cả tên thành phố ở phần bẫy của Bài 9. Chính Linstedt so satellite với dimension Type 2 của Kimball.
% Kiểm: dựng hub_customer (sha256(customer_unique_id), min giờ đặt), hub_order, link_order_customer và sat_customer_address (thêm dòng khi (zip, city, state) IS DISTINCT FROM lag theo giờ đặt) từ orders JOIN customers; DuckDB 1.5.6, 06/11/2026: 96.096, 99.441, 99.441, 96.355 dòng.
% Nguồn: lakehouse/source_sim/README.md dòng 39: customers có updated_at bằng order_purchase_timestamp của đơn.
% Nguồn: Linstedt 2002 như trên: "The Satellite is most closely related to a Type 2 dimension as defined by Ralph Kimball."

\begin{luuy}
\textbf{Bẫy thường gặp.}
\begin{itemize}
  \item \textbf{Chọn Data Vault cho một nguồn nhỏ.} Mọi bảng của vault ghi hệ nguồn và giờ nạp, thứ có ích khi nhiều hệ nguồn cùng đổ về một kho. Olist chỉ có một nguồn, nên vault chỉ thêm join.
  \item \textbf{Cho người phân tích truy vấn thẳng vault.} Dựng star hoặc mart bên trên, như Gold của repo đọc \repo{core}.
  \item \textbf{Tranh luận Kimball hay Inmon bằng tên gọi.} Hỏi hai câu cụ thể: dữ liệu chi tiết nằm ở dạng nào, và người phân tích đọc bảng nào.
\end{itemize}
\end{luuy}

% =====================================================================
\section{Bài 12: Metric trong mart và phương trình đối soát}

\textbf{Vấn đề.} Giám đốc hỏi doanh thu của Olist là bao nhiêu. Bốn người trả lời bốn con số, từ 13,22 tới 16,01 triệu BRL, và không ai viết sai SQL. Mỗi người dùng một định nghĩa ``doanh thu'' riêng. Khi các con số đã lên dashboard, không ai còn biết dashboard nào dùng định nghĩa nào.

\textbf{Giải thích.} Một \textit{metric} là một con số có tên mà nhiều người cùng dùng. Định nghĩa của nó phải trả lời đủ năm câu: cộng cột nào, ở grain nào, lọc bỏ những gì, theo mốc thời gian nào, và tính bằng đơn vị gì. Định nghĩa được viết một lần, thành một cột trong mart, thay vì mỗi dashboard tự viết câu \code{WHERE} của mình. Bài 1 của Module 2 viết định nghĩa metric vào thẻ yêu cầu, giống tiêu chí chấp nhận của một user story; bài này thêm phần kiểm bằng số.
% Nguồn: roadmap/modules/m02_vong_doi.tex dòng 90 (cầu nối "user story có tiêu chí chấp nhận" với "định nghĩa metric") và dòng 145 (thẻ yêu cầu giống user story kèm tiêu chí chấp nhận), grep ngày 06/11/2026.

Đây là định nghĩa doanh thu mà Lab 3.5 dùng:

\begin{xltabular}{\textwidth}{@{}P{3.2cm} L@{}}
\toprule
\thd{trường} & \thd{doanh thu, cột \code{revenue} của \code{marts.revenue\_\allowbreak{}daily}} \\
\midrule
\endhead
\bottomrule
\endfoot
một câu & tổng giá hàng của các dòng hàng thuộc đơn không bị huỷ và không unavailable \\
cộng cột nào & \code{price} của \repo{core.order\_\allowbreak{}items} \\
không gồm & phí vận chuyển \code{freight\_\allowbreak{}value}, nằm ở cột \code{freight} riêng \\
bộ lọc & \code{order\_\allowbreak{}status NOT IN ('canceled', 'unavailable')} \\
mốc thời gian & ngày của \code{order\_\allowbreak{}purchase\_\allowbreak{}timestamp}, giờ địa phương Brazil \\
grain & một dòng mỗi ngày đặt, danh mục và bang người bán \\
đơn vị & BRL, kiểu \code{DECIMAL}, đúng tới 0,01 \\
phần bị loại & \code{marts.revenue\_\allowbreak{}excluded}, theo ngày và trạng thái đơn \\
\end{xltabular}
% Nguồn: lakehouse/docs/data-model.md mục 9 hàng 18 (giờ Brazil) và mục 11 (core.orders chứa cả đơn canceled và unavailable, lọc theo order_status).

Bộ lọc nào cũng bỏ đi một phần dữ liệu. \textit{Phương trình đối soát} (reconciliation) chứng minh phần bị bỏ là có chủ đích: phần giữ lại cộng phần bị loại phải bằng đúng tổng ở nguồn. Muốn kiểm được, phần bị loại phải được lưu thành một bảng có lý do, chứ không biến mất trong một câu \code{WHERE}.

\begin{ghichu}
\textbf{So với SWE.} Java có \code{Collectors.partitioningBy}: thay vì \code{filter()} bỏ đi phần không khớp, nó trả về cả hai nửa trong một \code{Map<Boolean, List<T>>}, và tài liệu Java ghi rằng map đó luôn có đủ hai khoá \code{false} và \code{true}. Mart nên làm giống vậy. Khi có đủ hai nửa, bạn kiểm được bất biến quen thuộc: hai nửa cộng lại bằng cái ban đầu.
% Nguồn: https://docs.oracle.com/en/java/javase/21/docs/api/java.base/java/util/stream/Collectors.html (partitioningBy): "The returned Map always contains mappings for both false and true keys."

\textbf{Chỗ so sánh gãy.} Với list, bất biến chỉ là đếm phần tử. Với tiền, phải so tổng bằng \code{DECIMAL} đúng tới 0,01 BRL, vì \code{DOUBLE} lệch ở chữ số cuối như Bài 1. Hai nửa cũng phải cùng grain với metric: lọc theo trạng thái của đơn, nhưng cộng theo dòng hàng. Cuối cùng, phương trình đúng chỉ chứng minh không đồng nào bị mất hay bị cộng hai lần; nó không chứng minh bỏ đơn huỷ là lựa chọn đúng. Lựa chọn đó là quyết định nghiệp vụ, và phải nằm trong định nghĩa.
\end{ghichu}

\subsection{Làm thử}

Khối dưới tính bốn con số cùng tên ``doanh thu'', dựng một mart theo định nghĩa trên cùng bảng phần bị loại, kiểm phương trình, rồi cố ý làm sai bảng phần bị loại.

% chạy: R
\begin{codeblock}
lakehouse/.venv/bin/python - <<'EOF'
import duckdb
D = 'lakehouse/source_sim/data/'
con = duckdb.connect()
for t in ['orders', 'order_items', 'order_payments']:
    con.sql(f"CREATE VIEW {t} AS SELECT * FROM '{D}olist_{t}_dataset.csv'")
con.sql("""CREATE VIEW items AS
  SELECT CAST(o.order_purchase_timestamp AS DATE) AS purchase_date, o.order_status,
         CAST(i.price AS DECIMAL(14,2)) AS price
  FROM order_items i JOIN orders o USING (order_id)""")
one = lambda q: con.sql(q).fetchall()[0][0]
print('Bốn con số cùng tên "doanh thu":')
print('  mọi dòng hàng                  :', one("SELECT sum(price) FROM items"))
print('  bỏ đơn canceled và unavailable :', one("""SELECT sum(price) FROM items
                                  WHERE order_status NOT IN ('canceled', 'unavailable')"""))
print('  chỉ đơn delivered              :', one("""SELECT sum(price) FROM items
                                  WHERE order_status = 'delivered'"""))
print('  tiền khách trả, payment_value  :', one("""SELECT sum(CAST(payment_value
                                  AS DECIMAL(14,2))) FROM order_payments"""))
# Mart: metric viết một lần; phần bị loại cũng thành một bảng, kèm lý do
con.sql("""CREATE TABLE mart_revenue_daily AS
  SELECT purchase_date, sum(price) AS revenue FROM items
  WHERE order_status NOT IN ('canceled', 'unavailable') GROUP BY 1""")
con.sql("""CREATE TABLE mart_revenue_excluded AS
  SELECT purchase_date, order_status AS reason, sum(price) AS revenue FROM items
  WHERE order_status IN ('canceled', 'unavailable') GROUP BY ALL""")
check = """SELECT k, x, k + x AS cong, s, k + x = s AS khop FROM (SELECT
  (SELECT sum(revenue) FROM mart_revenue_daily) AS k,
  (SELECT sum(revenue) FROM mart_revenue_excluded) AS x,
  (SELECT sum(CAST(price AS DECIMAL(14,2))) FROM order_items) AS s)"""
def doi_soat(label):
    k, x, c, s, ok = con.sql(check).fetchall()[0]
    print(f'{label}: {k} + {x} = {c}; nguồn {s}; khớp: {ok}')
doi_soat('giữ lại + bị loại')
# Lỗi cố ý: người viết bảng bị loại quên trạng thái unavailable
con.sql("DELETE FROM mart_revenue_excluded WHERE reason = 'unavailable'")
doi_soat('quên unavailable ')
EOF
\end{codeblock}
\textit{Kết quả:}
% kết quả: R
\begin{codeblock}
Bốn con số cùng tên "doanh thu":
  mọi dòng hàng                  : 13591643.70
  bỏ đơn canceled và unavailable : 13494400.74
  chỉ đơn delivered              : 13221498.11
  tiền khách trả, payment_value  : 16008872.12
giữ lại + bị loại: 13494400.74 + 97242.96 = 13591643.70; nguồn 13591643.70; khớp: True
quên unavailable : 13494400.74 + 95235.27 = 13589636.01; nguồn 13591643.70; khớp: False
\end{codeblock}

\begin{itemize}
  \item Mỗi con số đúng với định nghĩa của nó. Tiền khách trả còn gồm phí vận chuyển và tiền của 775 đơn không có dòng hàng ở Bài 3, nên không dùng được cho doanh thu theo danh mục.
  % Kiểm: 98.299 trên 98.665 đơn có cả dòng hàng và thanh toán có tổng payment_value bằng tổng price cộng freight_value, lệch không quá 0,01; DuckDB 1.5.6, 06/11/2026.
  \item Với định nghĩa của bài, phần giữ lại là 13.494.400,74 BRL và phần bị loại là 97.242,96 BRL. Cộng lại đúng bằng 13.591.643,70 BRL, tổng \code{price} của file CSV mà mục 10 của \repo{docs/data-model.md} dùng để kiểm chéo.
  \item Khi bảng phần bị loại quên trạng thái \code{unavailable}, phương trình lệch 2.007,69 BRL, tức khoảng 0,015\% tổng tiền. Một dashboard không bao giờ cho bạn thấy độ lệch nhỏ như vậy; phương trình thì thấy ngay.
  % Kiểm: 2.007,69 BRL là tổng price của 7 dòng hàng thuộc đơn unavailable; đếm bằng order_items JOIN orders GROUP BY order_status, DuckDB 1.5.6, 06/11/2026.
  \item Nếu chỉ tính đơn \code{delivered}, phần giữ lại là 13.221.498,11 BRL, phần bị loại là 370.145,59 BRL, và phương trình vẫn đúng. Đây là lý do phương trình không thay được định nghĩa.
  % Kiểm: sum(price) theo order_status <> 'delivered' = 370145.59, cùng lệnh, 06/11/2026.
\end{itemize}

\begin{luuy}
\textbf{Bẫy thường gặp.}
\begin{itemize}
  \item \textbf{Lọc trong dashboard thay vì trong mart.} Mỗi dashboard sẽ có một bộ lọc riêng, và bạn quay lại bốn con số ở phần vấn đề.
  \item \textbf{Chỉ đối soát tổng cả kỳ.} Tổng khớp không chứng minh tiền nằm đúng ngày. Đối soát thêm theo ngày hoặc theo tháng, nhất là khi đổi mốc thời gian hay múi giờ.
  \item \textbf{Đặt hai tên cho một con số, hoặc một tên cho hai con số.} Nếu \code{revenue} ở hai mart có hai định nghĩa, đổi tên một trong hai.
\end{itemize}
\end{luuy}

% =====================================================================
\section{Lab 3.1: Bus matrix và grain}

Lab này viết đặc tả trước khi viết SQL, giống việc viết interface trước khi cài đặt. Bạn cần giấy, một file markdown và mẫu truy vấn ở Bài 2.

\begin{enumerate}
  \item Với mỗi quy trình đặt hàng, thanh toán, giao hàng và đánh giá, viết một câu grain theo mẫu ``Một dòng là ...''.
  \item Vẽ bus matrix: dòng là quy trình, cột là \code{dim\_\allowbreak{}date} kèm các vai, \code{dim\_\allowbreak{}customer}, \code{dim\_\allowbreak{}product}, \code{dim\_\allowbreak{}seller} và \code{dim\_\allowbreak{}order\_\allowbreak{}profile}. Với mỗi ô bỏ trống, viết một câu lý do, như lý do người bán ở dòng giao hàng.
  \item Từ mỗi câu grain, viết ra khoá, rồi kiểm trên CSV bằng mẫu ở Bài 2: số dòng phải bằng số khoá khác nhau.
  \item Phân loại \code{price}, \code{freight\_\allowbreak{}value}, \code{payment\_\allowbreak{}value}, \code{payment\_\allowbreak{}installments}, \code{review\_\allowbreak{}score} thành cộng được, cộng được một phần hay không cộng được, kèm một câu lý do cho mỗi cột.
\end{enumerate}

\hoanthanh{bus matrix có đủ bốn quy trình; mỗi fact có một câu grain và một khoá; câu kiểm khoá của \code{fct\_\allowbreak{}order\_\allowbreak{}items} và \code{fct\_\allowbreak{}payments} cho số dòng bằng số khoá; năm measure đều có phân loại kèm lý do.}

% =====================================================================
\section{Lab 3.2: Dựng star schema Olist, test viết trước}

Lab này làm theo kiểu TDD: bạn nhận bộ test trước, chạy và thấy nó đỏ, rồi viết model cho tới khi mọi test xanh. Mọi thứ nằm trong thư mục \code{lab\_\allowbreak{}m03/}, và kết quả là file \code{lab\_\allowbreak{}m03/olist\_\allowbreak{}star.duckdb}.

\textbf{Bước 1: tạo runner.} Runner chạy các file \code{models/NN\_\allowbreak{}ten\_\allowbreak{}bang.sql} theo thứ tự tên, in số dòng của từng bảng, rồi chạy mọi test trong \code{tests/}. Nó trả mã lỗi 1 khi có test đỏ, nên sau này gắn vào CI được ngay.

% chạy: L1
\begin{codeblock}
mkdir -p lab_m03/models lab_m03/tests
cat > lab_m03/build.py <<'EOF'
import duckdb, pathlib, sys
D = 'lakehouse/source_sim/data/'
LAB = pathlib.Path('lab_m03')
con = duckdb.connect(str(LAB / 'olist_star.duckdb'))
# Nguồn là CSV Olist, đọc qua view tạm nên không bị chép vào file DuckDB
for t in ['orders', 'order_items', 'order_payments', 'order_reviews',
          'customers', 'products', 'sellers']:
    con.sql(f"CREATE OR REPLACE TEMP VIEW {t} AS "
            f"SELECT * FROM '{D}olist_{t}_dataset.csv'")
con.sql(f"CREATE OR REPLACE TEMP VIEW category_translation AS "
        f"SELECT * FROM '{D}product_category_name_translation.csv'")
# models/NN_ten_bang.sql tạo bảng ten_bang; chạy theo thứ tự tên file
for f in sorted((LAB / 'models').glob('*.sql')):
    con.execute(f.read_text())
    table = f.stem.split('_', 1)[1]
    rows = con.sql(f'SELECT count(*) FROM {table}').fetchall()[0][0]
    print(f'{f.name:26} {rows:>7} dòng')
# tests/*.sql: mỗi file là một câu SELECT trả về dòng lỗi; 0 dòng là đạt
tests = sorted((LAB / 'tests').glob('*.sql'))
failed = 0
for f in tests:
    try:
        n = len(con.sql(f.read_text()).fetchall())
        msg = f'{n} dòng lỗi'
    except duckdb.Error as e:
        n, msg = 1, str(e).splitlines()[0].split(': ', 1)[-1]
    failed += n > 0
    print(f"{'ĐẠT' if n == 0 else 'LỖI'}  {f.stem}: {msg}")
print(f'{len(tests)} test, {failed} lỗi')
sys.exit(1 if failed else 0)
EOF
\end{codeblock}

\textbf{Bước 2: model đầu tiên.} Mỗi file model bắt đầu bằng câu grain trong comment. \code{dim\_\allowbreak{}date} được cho sẵn làm mẫu.

% chạy: L2
\begin{codeblock}
cat > lab_m03/models/01_dim_date.sql <<'EOF'
/* Grain: một dòng mỗi ngày từ 2016-01-01 tới 2018-12-31,
   cộng một dòng khoá 0 cho ngày chưa biết hoặc chưa xảy ra. */
CREATE OR REPLACE TABLE dim_date AS
SELECT CAST(strftime(d, '%Y%m%d') AS INTEGER) AS date_key,
       CAST(d AS DATE) AS day, year(d) AS year, month(d) AS month,
       dayname(d) AS weekday, dayofweek(d) IN (0, 6) AS is_weekend
FROM range(DATE '2016-01-01', DATE '2019-01-01', INTERVAL 1 DAY) r(d)
UNION ALL
SELECT 0, NULL, NULL, NULL, 'unknown', NULL;
EOF
\end{codeblock}

\textbf{Bước 3: bộ test.} Tám test dưới là đặc tả của lab. Mỗi file là một câu \code{SELECT} trả về dòng lỗi. Tạo chúng rồi chạy runner.

% chạy: L3
\begin{codeblock}
cat > lab_m03/tests/dim_customer_grain.sql <<'EOF'
/* Mỗi người mua một dòng, mỗi customer_key một dòng */
SELECT customer_unique_id AS k FROM dim_customer GROUP BY 1 HAVING count(*) > 1
UNION ALL
SELECT CAST(customer_key AS VARCHAR) FROM dim_customer
GROUP BY 1 HAVING count(*) > 1;
EOF
cat > lab_m03/tests/dim_unknown_row.sql <<'EOF'
/* Dimension nào cũng có đúng một dòng khoá 0 cho giá trị chưa biết */
SELECT 'dim_date' AS dim FROM dim_date WHERE date_key = 0 HAVING count(*) <> 1
UNION ALL SELECT 'dim_customer' FROM dim_customer
          WHERE customer_key = 0 HAVING count(*) <> 1
UNION ALL SELECT 'dim_product' FROM dim_product
          WHERE product_key = 0 HAVING count(*) <> 1
UNION ALL SELECT 'dim_seller' FROM dim_seller
          WHERE seller_key = 0 HAVING count(*) <> 1;
EOF
cat > lab_m03/tests/fct_order_items_grain.sql <<'EOF'
/* Grain: một dòng mỗi (order_id, order_item_id) */
SELECT order_id, order_item_id FROM fct_order_items
GROUP BY ALL HAVING count(*) > 1;
EOF
cat > lab_m03/tests/fct_order_items_khoa_ngoai.sql <<'EOF'
/* Mọi khoá trên fact phải có dòng ở dimension; NULL cũng tính là lỗi */
SELECT 'customer_key' AS col, f.customer_key AS k FROM fct_order_items f
  LEFT JOIN dim_customer d USING (customer_key) WHERE d.customer_key IS NULL
UNION ALL SELECT 'product_key', f.product_key FROM fct_order_items f
  LEFT JOIN dim_product d USING (product_key) WHERE d.product_key IS NULL
UNION ALL SELECT 'seller_key', f.seller_key FROM fct_order_items f
  LEFT JOIN dim_seller d USING (seller_key) WHERE d.seller_key IS NULL
UNION ALL SELECT 'purchase_date_key', f.purchase_date_key FROM fct_order_items f
  LEFT JOIN dim_date d ON d.date_key = f.purchase_date_key WHERE d.date_key IS NULL
UNION ALL SELECT 'delivered_date_key', f.delivered_date_key FROM fct_order_items f
  LEFT JOIN dim_date d ON d.date_key = f.delivered_date_key
  WHERE d.date_key IS NULL;
EOF
cat > lab_m03/tests/fct_order_items_doi_chieu.sql <<'EOF'
/* Số dòng, tổng price, tổng freight_value phải bằng nguồn */
SELECT * FROM
  (SELECT count(*), sum(price), sum(freight_value) FROM fct_order_items)
EXCEPT
SELECT * FROM
  (SELECT count(*), sum(CAST(price AS DECIMAL(12,2))),
          sum(CAST(freight_value AS DECIMAL(12,2))) FROM order_items);
EOF
cat > lab_m03/tests/fct_payments_grain.sql <<'EOF'
/* Grain: một dòng mỗi (order_id, payment_sequential) */
SELECT order_id, payment_sequential FROM fct_payments
GROUP BY ALL HAVING count(*) > 1;
EOF
cat > lab_m03/tests/fct_payments_doi_chieu.sql <<'EOF'
/* Số dòng và tổng payment_value phải bằng nguồn */
SELECT * FROM (SELECT count(*), sum(payment_value) FROM fct_payments)
EXCEPT
SELECT * FROM (SELECT count(*), sum(CAST(payment_value AS DECIMAL(12,2)))
               FROM order_payments);
EOF
cat > lab_m03/tests/fct_payments_khoa_ngoai.sql <<'EOF'
/* Khoá người mua và khoá ngày đặt của thanh toán phải có ở dimension */
SELECT 'customer_key' AS col, f.customer_key AS k FROM fct_payments f
  LEFT JOIN dim_customer d USING (customer_key) WHERE d.customer_key IS NULL
UNION ALL SELECT 'purchase_date_key', f.purchase_date_key FROM fct_payments f
  LEFT JOIN dim_date d ON d.date_key = f.purchase_date_key WHERE d.date_key IS NULL;
EOF

lakehouse/.venv/bin/python lab_m03/build.py
\end{codeblock}
\textit{Kết quả:}
% kết quả: L3
\begin{codeblock}
01_dim_date.sql               1097 dòng
LỖI  dim_customer_grain: Table with name dim_customer does not exist!
LỖI  dim_unknown_row: Table with name dim_customer does not exist!
LỖI  fct_order_items_doi_chieu: Table with name fct_order_items does not exist!
LỖI  fct_order_items_grain: Table with name fct_order_items does not exist!
LỖI  fct_order_items_khoa_ngoai: Table with name fct_order_items does not exist!
LỖI  fct_payments_doi_chieu: Table with name fct_payments does not exist!
LỖI  fct_payments_grain: Table with name fct_payments does not exist!
LỖI  fct_payments_khoa_ngoai: Table with name fct_payments does not exist!
8 test, 8 lỗi
\end{codeblock}

Mọi test đỏ vì các bảng chưa tồn tại. Đây là điểm xuất phát bạn muốn.

\textbf{Bước 4: viết năm model.} Viết mỗi bảng thành một file trong \code{lab\_\allowbreak{}m03/models/}, theo đặc tả dưới. Bài 1, Bài 5 và Bài 6 đã có gần đủ các mảnh SQL bạn cần.

\begin{xltabular}{\textwidth}{@{}P{3.4cm} L L@{}}
\toprule
\thd{file và grain} & \thd{cột} & \thd{luật} \\
\midrule
\endhead
\bottomrule
\endfoot
\code{02\_\allowbreak{}dim\_\allowbreak{}customer.sql}: một dòng mỗi \code{customer\_\allowbreak{}unique\_\allowbreak{}id}, cộng dòng khoá 0 & \code{customer\_\allowbreak{}key}, \code{customer\_\allowbreak{}unique\_\allowbreak{}id}, \code{customer\_\allowbreak{}city}, \code{customer\_\allowbreak{}state} & khoá thay thế bằng \code{row\_\allowbreak{}number()} từ 1; địa chỉ lấy ở đơn mới nhất, hai đơn cùng lúc thì lấy \code{customer\_\allowbreak{}id} lớn hơn, như \repo{core.customers} \\
\code{03\_\allowbreak{}dim\_\allowbreak{}product.sql}: một dòng mỗi \code{product\_\allowbreak{}id}, cộng dòng khoá 0 & \code{product\_\allowbreak{}key}, \code{product\_\allowbreak{}id}, \code{category}, \code{product\_\allowbreak{}weight\_\allowbreak{}g}, \code{product\_\allowbreak{}photos\_\allowbreak{}qty} & \code{category} là tên tiếng Anh; không có bản dịch thì dùng tên gốc; không có tên gốc thì ghi \code{unknown} \\
\code{04\_\allowbreak{}dim\_\allowbreak{}seller.sql}: một dòng mỗi \code{seller\_\allowbreak{}id}, cộng dòng khoá 0 & \code{seller\_\allowbreak{}key}, \code{seller\_\allowbreak{}id}, \code{seller\_\allowbreak{}city}, \code{seller\_\allowbreak{}state} & khoá thay thế bằng \code{row\_\allowbreak{}number()} từ 1 \\
\code{05\_\allowbreak{}fct\_\allowbreak{}order\_\allowbreak{}items.sql}: một dòng mỗi \code{order\_\allowbreak{}id} và \code{order\_\allowbreak{}item\_\allowbreak{}id} & \code{order\_\allowbreak{}id}, \code{order\_\allowbreak{}item\_\allowbreak{}id}, \code{customer\_\allowbreak{}key}, \code{product\_\allowbreak{}key}, \code{seller\_\allowbreak{}key}, \code{purchase\_\allowbreak{}date\_\allowbreak{}key}, \code{delivered\_\allowbreak{}date\_\allowbreak{}key}, \code{order\_\allowbreak{}status}, \code{price}, \code{freight\_\allowbreak{}value} & tra người mua qua \code{orders.customer\_\allowbreak{}id} rồi \code{customers.customer\_\allowbreak{}unique\_\allowbreak{}id}; khoá không tìm thấy thì ghi 0; đơn chưa giao có \code{delivered\_\allowbreak{}date\_\allowbreak{}key} bằng 0; tiền có kiểu \code{DECIMAL(12,2)}; \code{order\_\allowbreak{}id} là degenerate dimension; \code{order\_\allowbreak{}status} tạm nằm thẳng trên fact thay cho khoá của \code{dim\_\allowbreak{}order\_\allowbreak{}profile} \\
\code{06\_\allowbreak{}fct\_\allowbreak{}payments.sql}: một dòng mỗi \code{order\_\allowbreak{}id} và \code{payment\_\allowbreak{}sequential} & \code{order\_\allowbreak{}id}, \code{payment\_\allowbreak{}sequential}, \code{customer\_\allowbreak{}key}, \code{purchase\_\allowbreak{}date\_\allowbreak{}key}, \code{payment\_\allowbreak{}type}, \code{payment\_\allowbreak{}installments}, \code{payment\_\allowbreak{}value} & tra khoá như \code{fct\_\allowbreak{}order\_\allowbreak{}items}; tiền có kiểu \code{DECIMAL(12,2)} \\
\end{xltabular}

\textbf{Bước 5: chạy lại cho tới khi xanh.} Khi cả năm file đúng, runner in ra:

% chạy: L5 cần-lời-giải
\begin{codeblock}
lakehouse/.venv/bin/python lab_m03/build.py
\end{codeblock}
\textit{Kết quả:}
% kết quả: L5
\begin{codeblock}
01_dim_date.sql               1097 dòng
02_dim_customer.sql          96097 dòng
03_dim_product.sql           32952 dòng
04_dim_seller.sql             3096 dòng
05_fct_order_items.sql      112650 dòng
06_fct_payments.sql         103886 dòng
ĐẠT  dim_customer_grain: 0 dòng lỗi
ĐẠT  dim_unknown_row: 0 dòng lỗi
ĐẠT  fct_order_items_doi_chieu: 0 dòng lỗi
ĐẠT  fct_order_items_grain: 0 dòng lỗi
ĐẠT  fct_order_items_khoa_ngoai: 0 dòng lỗi
ĐẠT  fct_payments_doi_chieu: 0 dòng lỗi
ĐẠT  fct_payments_grain: 0 dòng lỗi
ĐẠT  fct_payments_khoa_ngoai: 0 dòng lỗi
8 test, 0 lỗi
\end{codeblock}

\textbf{Bước 6: chứng minh fan-out trên star của bạn.} Viết hai câu truy vấn trên \code{olist\_\allowbreak{}star.duckdb}. Câu thứ nhất join \code{fct\_\allowbreak{}payments} với \code{fct\_\allowbreak{}order\_\allowbreak{}items} theo \code{order\_\allowbreak{}id} rồi cộng \code{payment\_\allowbreak{}value}. Câu thứ hai gộp mỗi fact về grain đơn trước, rồi mới join, và giữ cả đơn có thanh toán mà không có dòng hàng.

\hoanthanh{runner in \code{8 test, 0 lỗi} với đúng số dòng như trên; mỗi file model bắt đầu bằng câu grain; câu thứ nhất ở bước 6 cho 20.308.134,71 BRL và câu thứ hai cho 16.008.872,12 BRL.}

% =====================================================================
\section{Lab 3.3: SCD Type 2 cho người bán}

Olist chỉ có một ảnh chụp bảng \code{sellers}, nên bạn tự tạo thay đổi để luyện. Bài 9 dựng lịch sử một lần từ quá khứ. Lab này làm việc mà pipeline phải làm mỗi ngày: nhận thay đổi mới, đóng version cũ, mở version mới, và chạy lại thì không sinh thêm dòng.

\textbf{Bước 1: tạo dimension, dữ liệu thay đổi và câu SQL áp thay đổi.} Ba người bán có nhiều dòng hàng nhất chuyển thành phố vào 2018-01-01, và có thêm một người bán mới. Câu áp thay đổi chạy trong một transaction: bước đầu đóng version hiện hành của người bán có thay đổi, bước sau thêm version mới cho người bán không còn dòng hiện hành.

% chạy: M1
\begin{codeblock}
mkdir -p lab_m03/scd2
cat > lab_m03/scd2/00_init.sql <<'EOF'
/* Nạp lần đầu: mỗi người bán một version, hiệu lực từ 2016-01-01 */
CREATE OR REPLACE TABLE dim_seller_scd2 AS
SELECT row_number() OVER (ORDER BY seller_id) AS seller_key,
       seller_id, seller_city, seller_state,
       DATE '2016-01-01' AS valid_from, CAST(NULL AS DATE) AS valid_to,
       true AS is_current
FROM sellers;
/* Thay đổi giả lập: ba người bán chuyển thành phố, một người bán mới */
CREATE OR REPLACE TABLE seller_changes AS
SELECT * FROM (VALUES
  ('6560211a19b47992c3666cc44a7e94c0', 'campinas',       'SP', DATE '2018-01-01'),
  ('4a3ca9315b744ce9f8e9374361493884', 'curitiba',       'PR', DATE '2018-01-01'),
  ('1f50f920176fa81dab994f9023523100', 'belo horizonte', 'MG', DATE '2018-01-01'),
  ('new_seller_0001',                  'recife',         'PE', DATE '2018-01-01'))
  t(seller_id, seller_city, seller_state, change_date);
EOF
cat > lab_m03/scd2/01_apply.sql <<'EOF'
BEGIN;
/* Bước 1: đóng version hiện hành của người bán có thuộc tính đổi */
UPDATE dim_seller_scd2 d
SET valid_to = c.change_date, is_current = false
FROM seller_changes c
WHERE d.seller_id = c.seller_id AND d.is_current
  AND (d.seller_city IS DISTINCT FROM c.seller_city
       OR d.seller_state IS DISTINCT FROM c.seller_state);
/* Bước 2: thêm version mới cho người bán không còn dòng hiện hành */
INSERT INTO dim_seller_scd2
SELECT (SELECT max(seller_key) FROM dim_seller_scd2)
         + row_number() OVER (ORDER BY c.seller_id),
       c.seller_id, c.seller_city, c.seller_state,
       c.change_date, NULL, true
FROM seller_changes c
WHERE NOT EXISTS (SELECT 1 FROM dim_seller_scd2 d
                  WHERE d.seller_id = c.seller_id AND d.is_current);
COMMIT;
EOF
\end{codeblock}

\textbf{Bước 2: nạp lần đầu, rồi áp thay đổi hai lần.} Lần áp thứ hai phải không đổi gì; tính chất này gọi là \textit{idempotent}, giống một API \code{PUT} gọi lại hai lần vẫn cho cùng trạng thái.
% Nguồn: https://www.rfc-editor.org/rfc/rfc9110.txt (mục 9.2.2)

% chạy: M2
\begin{codeblock}
lakehouse/.venv/bin/python - <<'EOF'
import duckdb
D = 'lakehouse/source_sim/data/'
con = duckdb.connect('lab_m03/olist_star.duckdb')
con.sql(f"CREATE TEMP VIEW sellers AS SELECT * FROM '{D}olist_sellers_dataset.csv'")
con.execute(open('lab_m03/scd2/00_init.sql').read())
count = "SELECT count(*), count(*) FILTER (WHERE is_current) FROM dim_seller_scd2"
print('sau khi nạp lần đầu, (số dòng, số dòng hiện hành):',
      con.sql(count).fetchall()[0])
for lan in [1, 2]:
    con.execute(open('lab_m03/scd2/01_apply.sql').read())
    print(f'sau lần áp thay đổi thứ {lan}:', con.sql(count).fetchall()[0])
con.sql("""SELECT seller_key, seller_city, seller_state,
                  valid_from, valid_to, is_current
           FROM dim_seller_scd2
           WHERE seller_id = '4a3ca9315b744ce9f8e9374361493884'
           ORDER BY valid_from""").show()
EOF
\end{codeblock}
\textit{Kết quả:}
% kết quả: M2
\begin{codeblock}
sau khi nạp lần đầu, (số dòng, số dòng hiện hành): (3095, 3095)
sau lần áp thay đổi thứ 1: (3099, 3096)
sau lần áp thay đổi thứ 2: (3099, 3096)
┌────────────┬─────────────┬──────────────┬────────────┬────────────┬────────────┐
│ seller_key │ seller_city │ seller_state │ valid_from │  valid_to  │ is_current │
│   int64    │   varchar   │   varchar    │    date    │    date    │  boolean   │
├────────────┼─────────────┼──────────────┼────────────┼────────────┼────────────┤
│        882 │ ibitinga    │ SP           │ 2016-01-01 │ 2018-01-01 │ false      │
│       3097 │ curitiba    │ PR           │ 2018-01-01 │ NULL       │ true       │
└────────────┴─────────────┴──────────────┴────────────┴────────────┴────────────┘
\end{codeblock}

\textbf{Bước 3: test.} Thêm hai test vào \code{lab\_\allowbreak{}m03/tests/}: mỗi \code{seller\_\allowbreak{}id} có đúng một dòng \code{is\_\allowbreak{}current}, và không có hai version nào của cùng một người bán chồng khoảng \code{[valid\_\allowbreak{}from, valid\_\allowbreak{}to)}. Chạy lại runner.

\textbf{Bước 4: tra version cho từng dòng hàng.} Viết câu truy vấn gắn mỗi dòng của \code{fct\_\allowbreak{}order\_\allowbreak{}items} với version người bán còn hiệu lực vào ngày đặt, dùng khoảng nửa mở như Bài 9.

\hoanthanh{lần áp thứ hai vẫn giữ 3.099 dòng và 3.096 dòng hiện hành; runner in \code{10 test, 0 lỗi}; cả 112.650 dòng hàng khớp đúng một version người bán; với ba người bán đã chuyển, 3.260 dòng hàng trỏ tới version cũ và 2.691 dòng hàng trỏ tới version mới.}

% =====================================================================
\section{Lab 3.4: Accumulating snapshot từ API của repo}

Ở Bài 4, \code{fct\_\allowbreak{}orders} được dựng một lần từ CSV, tức từ trạng thái cuối của mọi đơn. Pipeline thật không có trạng thái cuối: nó nhận từng version của đơn theo thời gian, và phải sửa dòng của đơn trong snapshot mỗi khi đơn qua một mốc. Lab này làm đúng việc đó với \repo{source\_\allowbreak{}sim}, server giả lập nguồn REST của repo. Module 8 sẽ làm phần kéo dữ liệu này cho Bronze thật; ở đây mọi thứ nằm trong \code{lab\_\allowbreak{}m03/snapshot/}.

Server chạy như một process Python ở cổng 8010, vì cổng 8000 thuộc service \code{source-sim} của docker compose. Nó nạp mọi version vào DuckDB trong RAM, chỉ ghi file tạm của hệ điều hành, và không ghi gì vào \repo{lakehouse/}.
% Nguồn: lakehouse/docker-compose.yml:149 (127.0.0.1:8000:8000); lakehouse/source_sim/app.py:39-49 (DuckDB in-memory, NamedTemporaryFile).

\begin{luuy}
\textbf{Thật và giả lập trong dữ liệu của source\_sim.}
\begin{itemize}
  \item \textbf{Thật:} mọi mốc thời gian, vì chúng chép từ CSV, và trạng thái ở version cuối, vì version cuối đúng bằng dòng CSV.
  \item \textbf{Giả lập:} trạng thái của mọi version trung gian. Mock gán \code{created}, \code{approved}, \code{shipped}, \code{delivered} theo mốc xa nhất đã tới, và không cho trạng thái lùi.
  \item \textbf{Giả lập:} lúc đơn chuyển sang \code{canceled} hay \code{unavailable}. Olist không ghi mốc này, nên trạng thái thật hiện ra ở mốc cuối cùng của đơn.
  \item \textbf{Cần nhớ:} mock coi mọi mốc là UTC, nhưng chúng là giờ địa phương Brazil.
  \item \textbf{Lạ nhưng thật:} 166 đơn có giờ giao cho hãng vận chuyển sớm hơn giờ đặt ngay trong CSV, nên version đầu của chúng là \code{shipped} mà chưa có giờ đặt.
\end{itemize}
% Nguồn: lakehouse/source_sim/README.md dòng 29 và dòng 42 tới 44; lakehouse/source_sim/replay.py:59-73; lakehouse/docs/data-model.md mục 9 hàng 17 và hàng 18.
\end{luuy}

\textbf{Bước 1: chạy server.} \code{SIM\_\allowbreak{}FIXED\_\allowbreak{}NOW} đóng băng đồng hồ mô phỏng. Đặt nó sau version cuối của dữ liệu thì mọi version hiện ra ngay; README của mock dùng đúng giá trị 2018-12-31T00:00:00 cho việc này. \code{PYTHONDONTWRITEBYTECODE=1} giữ cho Python không ghi file \code{.pyc} vào repo.
% Nguồn: lakehouse/source_sim/README.md dòng 54.

% chạy: N1
\begin{codeblock}
mkdir -p lab_m03/snapshot
PYTHONDONTWRITEBYTECODE=1 SIM_FIXED_NOW=2018-12-31T00:00:00 \
  lakehouse/.venv/bin/python -m uvicorn --app-dir lakehouse source_sim.app:app \
  --port 8010 > lab_m03/snapshot/sim.log 2>&1 < /dev/null &
echo $! > lab_m03/snapshot/sim.pid
# Server cần vài giây để nạp mọi version vào RAM trước khi trả lời
for i in $(seq 60); do curl -s localhost:8010/health && break; sleep 1; done
echo
curl -s 'localhost:8010/api/v1/orders?page_size=1'; echo
\end{codeblock}
\textit{Kết quả:}
% kết quả: N1
\begin{codeblock}
{"status":"ok","sim_time":"2018-12-31T00:00:00"}
{"entity":"orders","page_size":1,"has_more":true,"sim_time":"2018-12-31T00:00:00","data":[{"order_id":"2e7a8482f6fb09756ca50c10d7bfc047","customer_id":"08c5351a6aca1c1589a38f244edeee9d","order_status":"created","order_purchase_timestamp":"2016-09-04 21:15:19","order_approved_at":null,"order_delivered_carrier_date":null,"order_delivered_customer_date":null,"order_estimated_delivery_date":"2016-10-20 00:00:00","updated_at":"2016-09-04T21:15:19","cursor":"2016-09-04T21:15:19|000000001"}]}
\end{codeblock}

Dòng đầu là \code{/health}, sau đó là trang đầu của \code{orders} với một version. Nếu vòng chờ hết 60 giây mà không in gì, mở \code{lab\_\allowbreak{}m03/snapshot/sim.log}. Dòng \code{address already in use} nghĩa là cổng 8010 đang bận: tìm process bằng \code{lsof -nP -iTCP:8010} rồi dừng nó.

\textbf{Bước 2: kéo mọi version vào bảng thô.} Mỗi version có \code{cursor} dạng \code{<updated\_\allowbreak{}at>|<seq>}. Trang sau gửi \code{after} bằng cursor của version cuối trang trước, và dừng khi \code{has\_\allowbreak{}more} là \code{false}. Client chỉ so cursor như chuỗi, không tách nó ra. Bảng \code{raw\_\allowbreak{}order\_\allowbreak{}versions} giữ mọi version kèm số trang, giống Bronze chỉ thêm và không sửa.
% Nguồn: lakehouse/source_sim/README.md dòng 13 tới 27 (cursor, after, has_more, page_size tối đa 5000).

% chạy: N2
\begin{codeblock}
lakehouse/.venv/bin/python - <<'EOF'
import duckdb, json, urllib.parse, urllib.request
import pyarrow as pa
COLS = ['order_id', 'customer_id', 'order_status', 'order_purchase_timestamp',
        'order_approved_at', 'order_delivered_carrier_date',
        'order_delivered_customer_date', 'order_estimated_delivery_date',
        'updated_at', 'cursor']
con = duckdb.connect('lab_m03/snapshot/orders.duckdb')
con.sql("CREATE OR REPLACE TABLE raw_order_versions (page INTEGER, "
        + ', '.join(f'{c} VARCHAR' for c in COLS) + ")")
page, after = 0, None
while True:
    url = 'http://localhost:8010/api/v1/orders?page_size=5000'
    if after:  # trang sau bắt đầu ngay sau cursor của version cuối trang trước
        url += '&after=' + urllib.parse.quote(after)
    body = json.load(urllib.request.urlopen(url))
    page += 1
    batch = pa.Table.from_pylist(body['data'],
                                 schema=pa.schema([(c, pa.string()) for c in COLS]))
    con.sql(f"INSERT INTO raw_order_versions SELECT {page}, * FROM batch")
    after = body['data'][-1]['cursor']
    if not body['has_more']:
        break
print('số trang:', page)
print('số version, số đơn, số cursor, version cuối:', con.sql("""SELECT count(*),
  count(DISTINCT order_id), count(DISTINCT cursor), max(updated_at)
  FROM raw_order_versions""").fetchall()[0])
print('(số version của một đơn, số đơn):', con.sql("""SELECT n, count(*) FROM
  (SELECT count(*) AS n FROM raw_order_versions GROUP BY order_id)
  GROUP BY 1 ORDER BY 1""").fetchall())
EOF
\end{codeblock}
\textit{Kết quả:}
% kết quả: N2
\begin{codeblock}
số trang: 79
số version, số đơn, số cursor, version cuối: (391551, 99441, 391551, '2018-10-17T17:30:18')
(số version của một đơn, số đơn): [(1, 177), (2, 1626), (3, 2430), (4, 95208)]
\end{codeblock}

79 trang chứa 391.551 version của 99.441 đơn, mỗi version một cursor riêng, và phân bố số version mỗi đơn khớp đúng hàng 17 ở mục 9 của \repo{docs/data-model.md}. Nếu bạn nhận ít version hơn, \code{sim\_\allowbreak{}time} ở bước 1 đang sớm hơn 2018-10-17T17:30:18, giờ của version cuối.

\textbf{Bước 3: xem version của hai đơn.} Khối dưới in các version của một đơn đã giao và một đơn bị huỷ, kèm các mốc đã có ở mỗi version.

% chạy: N3
\begin{codeblock}
lakehouse/.venv/bin/python - <<'EOF'
import duckdb
con = duckdb.connect('lab_m03/snapshot/orders.duckdb', read_only=True)
MOC = {'order_purchase_timestamp': 'đặt', 'order_approved_at': 'duyệt',
       'order_delivered_carrier_date': 'giao vận chuyển',
       'order_delivered_customer_date': 'tới khách'}
for oid in ['00010242fe8c5a6d1ba2dd792cb16214', '0166cc8756d58b4a30fb49de83527120']:
    print('đơn', oid[:8])
    rows = con.sql(f"""SELECT updated_at, order_status, {', '.join(MOC)}
                       FROM raw_order_versions WHERE order_id = '{oid}'
                       ORDER BY cursor""").fetchall()
    for r in rows:
        co = [ten for ten, v in zip(MOC.values(), r[2:]) if v]
        print(f'  {r[0]}  {r[1]:9}  mốc đã có: {", ".join(co)}')
EOF
\end{codeblock}
\textit{Kết quả:}
% kết quả: N3
\begin{codeblock}
đơn 00010242
  2017-09-13T08:59:02  created    mốc đã có: đặt
  2017-09-13T09:45:35  approved   mốc đã có: đặt, duyệt
  2017-09-19T18:34:16  shipped    mốc đã có: đặt, duyệt, giao vận chuyển
  2017-09-20T23:43:48  delivered  mốc đã có: đặt, duyệt, giao vận chuyển, tới khách
đơn 0166cc87
  2018-02-08T21:09:18  created    mốc đã có: đặt
  2018-02-08T21:27:56  approved   mốc đã có: đặt, duyệt
  2018-02-09T23:12:30  canceled   mốc đã có: đặt, duyệt, giao vận chuyển
\end{codeblock}

Ở đơn \code{00010242}, bốn giờ ở cột đầu là thật, còn ba trạng thái đầu là mock suy ra từ mốc đã tới. Đơn \code{0166cc87} thành \code{canceled} đúng vào giờ giao cho hãng vận chuyển: đó là giờ mock chọn, không phải giờ đơn bị huỷ, vì Olist không ghi giờ huỷ.

\textbf{Bước 4: áp từng trang vào snapshot bằng \code{MERGE}.} Snapshot có grain một dòng mỗi đơn. Câu \code{MERGE} thêm dòng cho đơn mới và sửa dòng của đơn đã có, như upsert ở Module 1, Bài 10. Hai chi tiết giữ cho bảng đúng. Thứ nhất, một trang có thể chứa nhiều version của cùng một đơn, nên câu nguồn chỉ giữ version có cursor lớn nhất của mỗi đơn. Thứ hai, điều kiện \code{s.last\_\allowbreak{}cursor > t.last\_\allowbreak{}cursor} chỉ cho version mới hơn ghi đè, nên áp lại một trang hay áp lộn thứ tự đều vô hại. \code{UPDATE BY NAME} và \code{INSERT BY NAME} ghép cột theo tên.

% chạy: N4
\begin{codeblock}
cat > lab_m03/snapshot/00_create.sql <<'EOF'
/* Grain: một dòng mỗi đơn. Mỗi mốc một cột, được điền dần khi đơn qua mốc đó. */
CREATE OR REPLACE TABLE fct_order_snapshot (
  order_id VARCHAR, customer_id VARCHAR, order_status VARCHAR,
  purchase_ts TIMESTAMP, approved_ts TIMESTAMP, carrier_ts TIMESTAMP,
  delivered_ts TIMESTAMP, estimated_date DATE, last_cursor VARCHAR);
EOF
cat > lab_m03/snapshot/01_apply_page.sql <<'EOF'
/* Áp các version của một trang vào snapshot; $page là số trang */
MERGE INTO fct_order_snapshot t
USING (
  SELECT order_id, customer_id, order_status,
         CAST(order_purchase_timestamp AS TIMESTAMP)      AS purchase_ts,
         CAST(order_approved_at AS TIMESTAMP)             AS approved_ts,
         CAST(order_delivered_carrier_date AS TIMESTAMP)  AS carrier_ts,
         CAST(order_delivered_customer_date AS TIMESTAMP) AS delivered_ts,
         CAST(order_estimated_delivery_date AS DATE)      AS estimated_date,
         cursor AS last_cursor
  FROM raw_order_versions WHERE page = $page
  /* Một trang có thể chứa nhiều version của cùng đơn: giữ version mới nhất */
  QUALIFY row_number() OVER (PARTITION BY order_id ORDER BY cursor DESC) = 1) s
ON t.order_id = s.order_id
/* Chỉ ghi đè bằng version mới hơn: áp lại hay áp lộn thứ tự đều không làm lùi */
WHEN MATCHED AND s.last_cursor > t.last_cursor THEN UPDATE BY NAME
WHEN NOT MATCHED THEN INSERT BY NAME;
EOF
lakehouse/.venv/bin/python - <<'EOF'
import duckdb
con = duckdb.connect('lab_m03/snapshot/orders.duckdb')
con.execute(open('lab_m03/snapshot/00_create.sql').read())
apply = open('lab_m03/snapshot/01_apply_page.sql').read()
def run(pages):  # số dòng snapshot mà MERGE đã thêm hoặc sửa
    return sum(con.execute(apply, {'page': p}).fetchall()[0][0] for p in pages)
print('áp trang 1 tới 40, số dòng đổi:', run(range(1, 41)))
# Đơn chưa tới version cuối của nó đang mang trạng thái do mock giả lập
con.sql("""SELECT s.order_status, count(*) FILTER (WHERE s.last_cursor = f.c) AS that,
                  count(*) FILTER (WHERE s.last_cursor < f.c) AS gia_lap
           FROM fct_order_snapshot s JOIN (SELECT order_id, max(cursor) AS c
                FROM raw_order_versions GROUP BY 1) f USING (order_id)
           GROUP BY 1 ORDER BY count(*) DESC""").show()
print('áp trang 41 tới 79, số dòng đổi:', run(range(41, 80)))
print('áp lại cả 79 trang, số dòng đổi:', run(range(1, 80)))
print('áp trang 79 rồi trang 1, số dòng đổi:', run([79, 1]))
csv = """SELECT order_id, customer_id, order_status, order_purchase_timestamp,
           order_approved_at, order_delivered_carrier_date,
           order_delivered_customer_date, CAST(order_estimated_delivery_date AS DATE)
         FROM 'lakehouse/source_sim/data/olist_orders_dataset.csv'"""
snap = "SELECT * EXCLUDE (last_cursor) FROM fct_order_snapshot"
print('số dòng, lệch với CSV hai chiều:', con.sql(f"""SELECT
  (SELECT count(*) FROM fct_order_snapshot),
  (SELECT count(*) FROM ({snap} EXCEPT {csv})),
  (SELECT count(*) FROM ({csv} EXCEPT {snap}))""").fetchall()[0])
EOF
\end{codeblock}
\textit{Kết quả:}
% kết quả: N4
\begin{codeblock}
áp trang 1 tới 40, số dòng đổi: 105619
┌──────────────┬───────┬─────────┐
│ order_status │ that  │ gia_lap │
│   varchar    │ int64 │  int64  │
├──────────────┼───────┼─────────┤
│ delivered    │ 47136 │       0 │
│ shipped      │   601 │    2262 │
│ approved     │     2 │     754 │
│ unavailable  │   509 │       0 │
│ canceled     │   308 │       0 │
│ processing   │   271 │       0 │
│ invoiced     │   208 │       0 │
│ created      │     4 │     122 │
└──────────────┴───────┴─────────┘

áp trang 41 tới 79, số dòng đổi: 111064
áp lại cả 79 trang, số dòng đổi: 0
áp trang 79 rồi trang 1, số dòng đổi: 0
số dòng, lệch với CSV hai chiều: (99441, 0, 0)
\end{codeblock}

\begin{itemize}
  \item Sau 40 trang, snapshot có 52.177 đơn, và 3.138 đơn đang mang trạng thái giả lập: 2.262 \code{shipped}, 754 \code{approved}, 122 \code{created}. Đếm ``số đơn đang chờ duyệt'' trên bảng này là đếm số do mock tạo ra.
  \item Áp lại cả 79 trang, rồi áp trang 79 trước trang 1, đều đổi 0 dòng. Đó là tính idempotent bạn đã gặp ở Lab 3.3.
  \item Snapshot cuối có 99.441 dòng và trùng khớp với CSV ở cả tám cột, theo cả hai chiều \code{EXCEPT}.
\end{itemize}

Nếu snapshot có hơn 99.441 dòng, câu nguồn thiếu bước giữ version mới nhất; nếu lần áp lại đổi dòng nào, điều kiện so cursor đang thiếu hoặc sai chiều.

\begin{luuy}
\textbf{Bẫy thường gặp.}
\begin{itemize}
  \item \textbf{Bỏ bước giữ version mới nhất trong trang.} Engine nào cũng cần bước này; chỉ cách hỏng là khác nhau:
  \begin{itemize}
    \item DuckDB 1.5.6 không báo lỗi: trang 1 thêm 5.000 dòng cho 1.562 đơn, và snapshot cuối có 257.365 dòng cho 99.441 đơn.
    \item DuckDB 2.0.0 dừng ở trang 2 với lỗi \code{MERGE INTO command cannot affect the same target row more than once}.
    \item PostgreSQL từ chối một lệnh sửa cùng một dòng hai lần, với lỗi \code{MERGE command cannot affect row a second time} mà bạn đã gặp ở Module 1, Bài 10.
  \end{itemize}
  % Kiểm: chạy 01_apply_page.sql không có dòng QUALIFY trên bản chép của raw_order_versions, DuckDB 1.5.6 và 2.0.0, 06/11/2026. PostgreSQL: roadmap/modules/m01_sql.tex, Bài 10, output "ERROR:  MERGE command cannot affect row a second time"; https://www.postgresql.org/docs/18/sql-merge.html ("a repeated UPDATE or DELETE will cause a cardinality violation").
  \item \textbf{Viết \code{UPDATE} hay \code{INSERT} không có \code{BY NAME}.} DuckDB khi đó ghép cột theo vị trí, cả với \code{UPDATE SET *} và \code{INSERT *}; bảng nguồn đổi thứ tự cột là dữ liệu vào sai cột.
  % Kiểm: MERGE với nguồn có cột (b, k, a) vào bảng (k, a, b): UPDATE và INSERT * báo Conversion Error do ghép theo vị trí, BY NAME ghép đúng; DuckDB 1.5.6 và 2.0.0, 06/11/2026.
  \item \textbf{Tin trạng thái trung gian.} Với dữ liệu thật, trạng thái là thứ hệ nguồn ghi; với mock, chỉ các mốc thời gian là thật. Metric về thời gian xử lý nên tính từ các cột mốc, như Bài 4.
\end{itemize}
\end{luuy}

\textbf{Bước 5: dựng lại snapshot mà không dùng \code{MERGE}.} Viết file \code{lab\_\allowbreak{}m03/snapshot/02\_\allowbreak{}as\_\allowbreak{}of.sql} tạo temp view \code{order\_\allowbreak{}snapshot\_\allowbreak{}as\_\allowbreak{}of}: trạng thái mọi đơn sau trang 40, dựng thẳng từ \code{raw\_\allowbreak{}order\_\allowbreak{}versions} bằng cách giữ version có cursor lớn nhất của mỗi đơn, với đúng chín cột của snapshot. Khối dưới dựng lại snapshot bằng \code{MERGE} tới trang 40 rồi so hai cách. Cách dựng thẳng chậm hơn khi dữ liệu lớn, nhưng là thước đo để kiểm cách tăng dần.

% chạy: N5 cần-lời-giải
\begin{codeblock}
lakehouse/.venv/bin/python - <<'EOF'
import duckdb
con = duckdb.connect()  # làm trong RAM, không đụng snapshot của bước 4
con.sql("ATTACH 'lab_m03/snapshot/orders.duckdb' AS lab (READ_ONLY)")
con.sql("CREATE TEMP VIEW raw_order_versions AS SELECT * FROM lab.raw_order_versions")
con.execute(open('lab_m03/snapshot/00_create.sql').read())
apply = open('lab_m03/snapshot/01_apply_page.sql').read()
for p in range(1, 41):
    con.execute(apply, {'page': p})
con.execute(open('lab_m03/snapshot/02_as_of.sql').read())  # file bạn viết
a, b = 'fct_order_snapshot', 'order_snapshot_as_of'
print('số dòng của hai cách:', con.sql(f"""SELECT (SELECT count(*) FROM {a}),
  (SELECT count(*) FROM {b})""").fetchall()[0])
print('lệch hai chiều:', con.sql(f"""SELECT
  (SELECT count(*) FROM (SELECT * FROM {a} EXCEPT SELECT * FROM {b})),
  (SELECT count(*) FROM (SELECT * FROM {b} EXCEPT SELECT * FROM {a}))
  """).fetchall()[0])
EOF
\end{codeblock}
\textit{Kết quả:}
% kết quả: N5
\begin{codeblock}
số dòng của hai cách: (52177, 52177)
lệch hai chiều: (0, 0)
\end{codeblock}

Hai cách cùng ra 52.177 dòng và không lệch dòng nào. Nếu lệch, so tên và kiểu của chín cột trước, rồi kiểm điều kiện \code{page <= 40}.

\textbf{Bước 6: dừng server.}

% chạy: N6
\begin{codeblock}
kill "$(cat lab_m03/snapshot/sim.pid)"
while kill -0 "$(cat lab_m03/snapshot/sim.pid)" 2>/dev/null; do sleep 0.2; done
grep -c 'Application shutdown complete' lab_m03/snapshot/sim.log
curl -s localhost:8010/health || echo "không còn server ở cổng 8010, curl trả mã $?"
\end{codeblock}
\textit{Kết quả:}
% kết quả: N6
\begin{codeblock}
1
không còn server ở cổng 8010, curl trả mã 7
\end{codeblock}

Số 1 là số dòng \code{Application shutdown complete} trong log, tức server đã tắt hẳn; mã 7 của \code{curl} nghĩa là không kết nối được.

\hoanthanh{khối ở bước 4 in \code{(99441, 0, 0)} ở dòng cuối, và hai lần áp lại đều đổi 0 dòng; khối ở bước 5 in \code{(52177, 52177)} và \code{(0, 0)}; sau bước 6, \code{curl -s localhost:8010/health} trả mã 7.}

% =====================================================================
\section{Lab 3.5: Mart cho năm câu hỏi của repo}

Mục 11 của \repo{docs/data-model.md} liệt kê năm câu hỏi mà Gold của repo phải trả lời, cùng các bảng \repo{core} cần đọc. Lab này thiết kế và dựng năm mart đó trên DuckDB, từ một bản \repo{core} thu nhỏ dựng từ CSV, với đúng tên bảng và tên cột của repo. Đây là bản chạy thử. Ở M7, bạn viết lại các mart này bằng Spark SQL trên \repo{core} thật, và phương trình đối soát của lab này là tiêu chí hoàn thành ở đó.

\textbf{Bước 1: dựng core thu nhỏ và runner.} File \code{00\_\allowbreak{}core.sql} được cho sẵn, vì core là việc của tầng Silver, mà repo đã làm bằng Spark trong \repo{src/\allowbreak{}lakehouse/\allowbreak{}silver/\allowbreak{}core.py}. Runner chạy mọi file \code{.sql} trong \code{lab\_\allowbreak{}m03/gold/} theo thứ tự tên, in số dòng của từng bảng, rồi chạy năm test. Mỗi test so mart với CSV chứ không so với core, để bắt được cả lỗi của core lẫn lỗi của mart.

% chạy: P1
\begin{codeblock}
mkdir -p lab_m03/gold
cat > lab_m03/gold/00_core.sql <<'EOF'
/* Bản thu nhỏ của semantic và core trong repo, dựng từ CSV; tên bảng, tên cột
   theo docs/data-model.md mục 5 và 6, chỉ giữ cột mà năm câu hỏi ở mục 11 cần. */
CREATE SCHEMA IF NOT EXISTS semantic;
CREATE SCHEMA IF NOT EXISTS core;
/* Toạ độ mỗi zip: trung bình các điểm khác nhau nằm trong khung Brazil */
CREATE OR REPLACE TABLE semantic.zip_geo AS
SELECT zip_code_prefix, avg(lat) AS lat, avg(lng) AS lng FROM (
  SELECT DISTINCT lpad(CAST(geolocation_zip_code_prefix AS VARCHAR), 5, '0')
           AS zip_code_prefix, geolocation_lat AS lat, geolocation_lng AS lng,
         geolocation_city, geolocation_state
  FROM geolocation)
WHERE lat BETWEEN -33.75 AND 5.27 AND lng BETWEEN -73.99 AND -34.79
GROUP BY 1;
CREATE OR REPLACE TABLE core.products AS
SELECT p.product_id, coalesce(t.product_category_name_english,
         p.product_category_name, 'unknown') AS product_category_name_english
FROM products p LEFT JOIN category_translation t USING (product_category_name);
CREATE OR REPLACE TABLE core.sellers AS
SELECT s.seller_id, s.zip AS seller_zip_code_prefix, s.seller_state,
       g.lat AS seller_lat, g.lng AS seller_lng
FROM (SELECT *, lpad(CAST(seller_zip_code_prefix AS VARCHAR), 5, '0') AS zip
      FROM sellers) s
LEFT JOIN semantic.zip_geo g ON g.zip_code_prefix = s.zip;
CREATE OR REPLACE TABLE core.customers AS
WITH b AS (  /* địa chỉ ở đơn mới nhất; đặt cùng lúc thì lấy customer_id lớn hơn */
  SELECT c.customer_unique_id,
         arg_max(lpad(CAST(c.customer_zip_code_prefix AS VARCHAR), 5, '0'),
                 (o.order_purchase_timestamp, c.customer_id)) AS zip,
         arg_max(c.customer_state, (o.order_purchase_timestamp, c.customer_id))
           AS customer_state,
         min(o.order_purchase_timestamp) AS first_order_purchase_timestamp,
         max(o.order_purchase_timestamp) AS last_order_purchase_timestamp,
         count(DISTINCT c.customer_id) AS customer_id_count
  FROM customers c JOIN orders o USING (customer_id) GROUP BY 1)
SELECT sha256(customer_unique_id) AS customer_key, customer_unique_id,
       zip AS customer_zip_code_prefix, customer_state,
       g.lat AS customer_lat, g.lng AS customer_lng,
       first_order_purchase_timestamp, last_order_purchase_timestamp,
       customer_id_count
FROM b LEFT JOIN semantic.zip_geo g ON g.zip_code_prefix = b.zip;
CREATE OR REPLACE TABLE core.orders AS
SELECT o.order_id, sha256(c.customer_unique_id) AS customer_key, o.order_status,
       o.order_purchase_timestamp, o.order_delivered_customer_date, r.review_score,
       CAST(o.order_delivered_customer_date AS DATE)
         - CAST(o.order_purchase_timestamp AS DATE) AS delivery_days,
       CAST(o.order_delivered_customer_date AS DATE)
         > CAST(o.order_estimated_delivery_date AS DATE) AS is_delivered_late
FROM orders o JOIN customers c USING (customer_id)
LEFT JOIN (SELECT order_id, review_score FROM order_reviews  /* review mới nhất */
           QUALIFY row_number() OVER (PARTITION BY order_id ORDER BY
             review_creation_date DESC, review_answer_timestamp DESC,
             review_id DESC) = 1) r USING (order_id);
CREATE OR REPLACE TABLE core.order_items AS
SELECT order_id, order_item_id, product_id, seller_id,
       CAST(price AS DECIMAL(12,2)) AS price,
       CAST(freight_value AS DECIMAL(12,2)) AS freight_value
FROM order_items;
EOF
cat > lab_m03/gold/build_gold.py <<'EOF'
import duckdb, pathlib, sys
D = 'lakehouse/source_sim/data/'
G = pathlib.Path('lab_m03/gold')
con = duckdb.connect(str(G / 'gold.duckdb'))
for t in ['orders', 'order_items', 'order_reviews', 'customers', 'products',
          'sellers', 'geolocation']:
    con.sql(f"CREATE OR REPLACE TEMP VIEW {t} AS "
            f"SELECT * FROM '{D}olist_{t}_dataset.csv'")
con.sql(f"CREATE OR REPLACE TEMP VIEW category_translation AS "
        f"SELECT * FROM '{D}product_category_name_translation.csv'")
for f in sorted(G.glob('*.sql')):  # 00_core.sql, rồi 10_marts.sql khi bạn đã viết
    con.execute(f.read_text())
for s, t in con.sql("SELECT schema_name, table_name FROM duckdb_tables() "
                    "WHERE database_name = 'gold' ORDER BY ALL").fetchall():
    n = con.sql(f'SELECT count(*) FROM {s}.{t}').fetchall()[0][0]
    print(f'{s + "." + t:28} {n:>7} dòng')
GRAIN = {'revenue_daily': 'purchase_date, category, seller_state',
         'revenue_excluded': 'purchase_date, order_status',
         'delivery_state_month': 'purchase_month, customer_state',
         'repeat_buyers_month': 'first_purchase_month',
         'review_seller': 'seller_id', 'map_zip': 'zip_code_prefix'}
# Mỗi test trả về dòng lỗi, 0 dòng là đạt; so mart với CSV, không so với core
TESTS = {
  'grain sáu mart': ' UNION ALL '.join(
      f"SELECT '{m}' FROM marts.{m} GROUP BY {k} HAVING count(*) > 1"
      for m, k in GRAIN.items()),
  'doanh thu: giữ + loại = nguồn': """SELECT 1 WHERE
     (SELECT sum(revenue) FROM marts.revenue_daily)
     + (SELECT sum(revenue) FROM marts.revenue_excluded)
     IS DISTINCT FROM (SELECT sum(CAST(price AS DECIMAL(12,2))) FROM order_items)""",
  'giao hàng: đã tới + chưa tới = số đơn': """SELECT 1 WHERE
     (SELECT sum(delivered_orders) FROM marts.delivery_state_month)
     + (SELECT count(*) FROM orders WHERE order_delivered_customer_date IS NULL)
     IS DISTINCT FROM (SELECT count(*) FROM orders)""",
  'người mua: hai mart = nguồn': """SELECT n FROM (
     SELECT sum(buyers) AS n FROM marts.repeat_buyers_month
     UNION ALL SELECT sum(buyers) FROM marts.map_zip)
     WHERE n IS DISTINCT FROM (SELECT count(DISTINCT customer_unique_id)
                               FROM customers)""",
  'review: mỗi cặp đơn, người bán một lần': """SELECT 1 WHERE
     (SELECT sum(reviewed_orders) FROM marts.review_seller)
     IS DISTINCT FROM (SELECT count(*) FROM (SELECT DISTINCT order_id, seller_id
       FROM order_items WHERE order_id IN (SELECT order_id FROM order_reviews)))""",
}
failed = 0
for name, sql in TESTS.items():
    try:
        n = len(con.sql(sql).fetchall())
        msg = f'{n} dòng lỗi'
    except duckdb.Error as e:
        n, msg = 1, str(e).splitlines()[0].split(': ', 1)[-1]
    failed += n > 0
    print(f"{'ĐẠT' if n == 0 else 'LỖI'}  {name}: {msg}")
if not failed:
    k, x = con.sql("SELECT (SELECT sum(revenue) FROM marts.revenue_daily), "
                   "(SELECT sum(revenue) FROM marts.revenue_excluded)").fetchall()[0]
    print(f'đối soát doanh thu: {k} + {x} = {k + x}')
sys.exit(1 if failed else 0)
EOF
lakehouse/.venv/bin/python lab_m03/gold/build_gold.py
\end{codeblock}
\textit{Kết quả:}
% kết quả: P1
\begin{codeblock}
core.customers                 96096 dòng
core.order_items              112650 dòng
core.orders                    99441 dòng
core.products                  32951 dòng
core.sellers                    3095 dòng
semantic.zip_geo               19010 dòng
LỖI  grain sáu mart: Table with name "marts.revenue_daily" does not exist because schema "marts" does not exist.
LỖI  doanh thu: giữ + loại = nguồn: Table with name "marts.revenue_daily" does not exist because schema "marts" does not exist.
LỖI  giao hàng: đã tới + chưa tới = số đơn: Table with name "marts.delivery_state_month" does not exist because schema "marts" does not exist.
LỖI  người mua: hai mart = nguồn: Table with name "marts.repeat_buyers_month" does not exist because schema "marts" does not exist.
LỖI  review: mỗi cặp đơn, người bán một lần: Table with name "marts.review_seller" does not exist because schema "marts" does not exist.
\end{codeblock}

Số dòng của sáu bảng khớp mục 10 của \repo{docs/data-model.md}: 19.010 zip, 96.096 người mua, 3.095 người bán, 32.951 sản phẩm, 99.441 đơn và 112.650 dòng hàng. Năm test đỏ vì schema \code{marts} chưa có. Nếu số dòng khác, kiểm biến \code{D} trỏ đúng thư mục CSV.

\textbf{Bước 2: đọc đặc tả.} Mỗi mart có grain, cột và luật dưới đây; nếu đã làm Lab 2.2, so bảng này với thẻ yêu cầu của bạn. Mỗi \code{CREATE TABLE} trong file của bạn phải mở đầu bằng comment ghi câu grain và định nghĩa một câu của từng measure, theo mẫu định nghĩa metric ở Bài 12.

\begin{xltabular}{\textwidth}{@{}P{3.8cm} P{4.2cm} L@{}}
\toprule
\thd{mart và grain} & \thd{cột} & \thd{luật} \\
\midrule
\endhead
\bottomrule
\endfoot
\code{revenue\_\allowbreak{}daily}: một dòng mỗi ngày đặt, danh mục, bang người bán & \code{purchase\_\allowbreak{}date}, \code{category}, \code{seller\_\allowbreak{}state}, \code{item\_\allowbreak{}count}, \code{revenue}, \code{freight} & định nghĩa doanh thu ở Bài 12; \code{category} là \code{product\_\allowbreak{}category\_\allowbreak{}name\_\allowbreak{}english} \\
\code{revenue\_\allowbreak{}excluded}: một dòng mỗi ngày đặt và trạng thái đơn & \code{purchase\_\allowbreak{}date}, \code{order\_\allowbreak{}status}, \code{item\_\allowbreak{}count}, \code{revenue} & chỉ đơn \code{canceled} và \code{unavailable} \\
\code{delivery\_\allowbreak{}state\_\allowbreak{}month}: một dòng mỗi tháng đặt và bang người mua & \code{purchase\_\allowbreak{}month}, \code{customer\_\allowbreak{}state}, \code{delivered\_\allowbreak{}orders}, \code{late\_\allowbreak{}orders}, \code{delivery\_\allowbreak{}days\_\allowbreak{}sum} & chỉ đơn đã tới tay khách; lưu tử số và mẫu số, không lưu tỉ lệ hay trung bình \\
\code{repeat\_\allowbreak{}buyers\_\allowbreak{}month}: một dòng mỗi tháng mua lần đầu & \code{first\_\allowbreak{}purchase\_\allowbreak{}month}, \code{buyers}, \code{repeat\_\allowbreak{}buyers} & người mua quay lại có \code{customer\_\allowbreak{}id\_\allowbreak{}count} lớn hơn 1 \\
\code{review\_\allowbreak{}seller}: một dòng mỗi người bán & \code{seller\_\allowbreak{}id}, \code{reviewed\_\allowbreak{}orders}, \code{review\_\allowbreak{}score\_\allowbreak{}sum} & mỗi cặp đơn và người bán đếm một lần, dù đơn có nhiều dòng hàng của người bán đó \\
\code{map\_\allowbreak{}zip}: một dòng mỗi zip của người mua hoặc người bán & \code{zip\_\allowbreak{}code\_\allowbreak{}prefix}, \code{lat}, \code{lng}, \code{buyers}, \code{sellers} & toạ độ từ \code{semantic.zip\_\allowbreak{}geo}; zip không có toạ độ vẫn giữ, với \code{lat} là NULL \\
\end{xltabular}

\textbf{Bước 3: viết \code{lab\_\allowbreak{}m03/gold/10\_\allowbreak{}marts.sql}.} File bắt đầu bằng \code{CREATE SCHEMA IF NOT EXISTS marts;} và chỉ đọc \code{core} cùng \code{semantic}. Bài 3, Bài 4, Bài 7 và Bài 12 có đủ các mảnh bạn cần.

\textbf{Bước 4: chạy lại runner cho tới khi xanh.} Khi cả sáu bảng đúng, runner in ra:

% chạy: P2 cần-lời-giải
\begin{codeblock}
lakehouse/.venv/bin/python lab_m03/gold/build_gold.py
\end{codeblock}
\textit{Kết quả:}
% kết quả: P2
\begin{codeblock}
core.customers                 96096 dòng
core.order_items              112650 dòng
core.orders                    99441 dòng
core.products                  32951 dòng
core.sellers                    3095 dòng
marts.delivery_state_month       556 dòng
marts.map_zip                  15068 dòng
marts.repeat_buyers_month         25 dòng
marts.revenue_daily            35548 dòng
marts.revenue_excluded           286 dòng
marts.review_seller             3090 dòng
semantic.zip_geo               19010 dòng
ĐẠT  grain sáu mart: 0 dòng lỗi
ĐẠT  doanh thu: giữ + loại = nguồn: 0 dòng lỗi
ĐẠT  giao hàng: đã tới + chưa tới = số đơn: 0 dòng lỗi
ĐẠT  người mua: hai mart = nguồn: 0 dòng lỗi
ĐẠT  review: mỗi cặp đơn, người bán một lần: 0 dòng lỗi
đối soát doanh thu: 13494400.74 + 97242.96 = 13591643.70
\end{codeblock}

Nếu một test đỏ, tìm tên test trong danh sách dưới:
\begin{itemize}
  \item \textbf{Grain:} chạy riêng câu \code{GROUP BY} của từng mart trong \code{GRAIN} để tìm mart có khoá trùng. Ví dụ, nhóm theo \code{order\_\allowbreak{}purchase\_\allowbreak{}timestamp} thay vì theo ngày sẽ cho nhiều dòng cùng một ngày.
  \item \textbf{Doanh thu:} cộng \code{revenue} của hai mart. Thiếu đúng 2.007,69 BRL so với 13.591.643,70 nghĩa là bạn quên \code{unavailable} trong \code{revenue\_\allowbreak{}excluded}, như Bài 12.
  \item \textbf{Giao hàng:} tính tổng \code{delivered\_\allowbreak{}orders}. 96.478 thay vì 96.476 nghĩa là bạn lọc theo trạng thái \code{delivered} thay vì theo \code{order\_\allowbreak{}delivered\_\allowbreak{}customer\_\allowbreak{}date IS NOT NULL}.
  \item \textbf{Người mua:} tính tổng \code{buyers} của \code{map\_\allowbreak{}zip}. 95.826 thay vì 96.096 nghĩa là bạn dùng \code{INNER JOIN} với \code{zip\_\allowbreak{}geo} và làm rơi 270 người mua không có toạ độ.
  \item \textbf{Review:} tính tổng \code{reviewed\_\allowbreak{}orders}. 111.708 thay vì 99.247 nghĩa là bạn đang đếm mỗi dòng hàng thay vì mỗi cặp đơn và người bán.
\end{itemize}
% Kiểm: lời giải tham khảo với INNER JOIN zip_geo cho sum(buyers) = 95.826; đếm dòng hàng có review cho 111.708; DuckDB 1.5.6, 06/11/2026.
% Kiểm: olist_orders_dataset.csv, DuckDB 1.5.6, 06/11/2026: order_status = 'delivered' cho 96.478 đơn, order_delivered_customer_date IS NOT NULL cho 96.476 đơn, NULL cho 2.965 đơn; 2.007,69 BRL là phần lệch của khối R ở Bài 12.
% Kiểm: runner với lời giải tham khảo bị sửa từng chỗ, DuckDB 1.5.6, 06/11/2026: lọc order_status = 'delivered' cho tổng delivered_orders 96.478 và chỉ test giao hàng đỏ; revenue_excluded chỉ có canceled cho tổng hai mart 13.589.636,01 và chỉ test doanh thu đỏ; revenue_daily GROUP BY order_purchase_timestamp cho test grain 16.147 dòng lỗi.

\hoanthanh{runner in năm dòng \code{ĐẠT} và dòng \code{đối soát doanh thu: 13494400.74 + 97242.96 = 13591643.70}; mỗi \code{CREATE TABLE} có câu grain và định nghĩa measure trong comment.}

Thêm nếu còn giờ: đổi định nghĩa doanh thu sang chỉ đơn \code{delivered}. Phương trình vẫn đúng, với 13.221.498,11 BRL giữ lại và 370.145,59 BRL bị loại.

% =====================================================================
\section{Thử thách kỹ năng}

\begin{thuthach}
Không có hướng dẫn từng bước. Thêm quy trình đánh giá vào Gold của Lab 3.5.

Thêm vào \code{lab\_\allowbreak{}m03/gold/} một bảng \code{fct\_\allowbreak{}reviews} với grain do bạn tự xác định, sau khi kiểm xem \code{review\_\allowbreak{}id} có duy nhất trong dữ liệu không. Từ đó dựng \code{marts.review\_\allowbreak{}category\_\allowbreak{}month}: điểm review theo danh mục và tháng tạo review, lưu tử số và mẫu số. Tự quyết và ghi lại ba điều: đơn có nhiều review thì tính review nào, đơn có hàng của hai danh mục thì điểm tính cho danh mục nào, và review của đơn bị huỷ có được tính không. Runner chỉ kiểm grain của bảng trong schema \code{marts}, nên đặt \code{fct\_\allowbreak{}reviews} ở đó, hoặc sửa runner để \code{GRAIN} nhận tên có schema.
\hoanthanh{bảng và mart mới có câu grain trong comment và có mặt trong \code{GRAIN} của runner; phương trình đối soát cho review đúng: số review được tính cộng số review bị loại, mỗi phần có lý do, bằng 99.224; runner in 0 dòng \code{LỖI}.}

\textbf{Mang vào capstone.} Mốc capstone 1 là bus matrix, grain và tầng Gold trên DuckDB cho năm câu hỏi ở mục 11 của \repo{docs/\allowbreak{}data-model.md}, có phương trình đối soát. Đem bus matrix của Lab 3.1, star của Lab 3.2, snapshot của Lab 3.4, mart của Lab 3.5 và phần review này vào repo capstone, kèm lệnh chạy runner. Dùng thẻ yêu cầu của Lab 2.2 làm đặc tả cho năm mart. Chạy lại năm câu của thử thách Module 1 trên star của Lab 3.2 và snapshot của Lab 3.4, rồi so với \code{lab\_\allowbreak{}m01/\allowbreak{}thu\_\allowbreak{}thach.sql}.
\end{thuthach}
% Nguồn: lakehouse/docs/data-model.md mục 9 hàng 3 (99.224 review, 98.410 review_id); mốc 1 của lộ trình.
% Nguồn: roadmap/modules/m01_sql.tex, thử thách kỹ năng, dòng Mang vào capstone (mốc 1 khớp lab_m01/thu_thach.sql); roadmap/modules/m02_vong_doi.tex, Lab 2.2 và dòng Mang vào capstone (thẻ yêu cầu là đầu vào của mốc 1); grep ngày 06/11/2026.

% =====================================================================
\section{Tự kiểm tra}

Trả lời bằng lời của bạn trước, rồi mới đọc gợi ý.

\begin{enumerate}
  \item Vì sao grain phải được viết trước khi chọn cột? Join thanh toán với dòng hàng rồi cộng \code{payment\_\allowbreak{}value} cho 20.308.134,71 BRL; sửa thế nào để ra 16.008.872,12 BRL, và vì sao \code{sum(DISTINCT ...)} không phải là cách sửa?\par
  \textit{Gợi ý:} grain là hợp đồng mà mọi cột phải khớp; một dòng hàng và một lần trả tiền là hai grain. Gộp mỗi bảng về grain đơn rồi \code{LEFT JOIN}; \code{DISTINCT} bỏ cả những khoản tiền trùng số của các đơn khác nhau.
  \item Loại fact nào hợp với câu hỏi ``mỗi bước giao hàng mất bao lâu''? Vì sao tỉ lệ giao trễ của cả nước không phải trung bình tỉ lệ của 27 bang?\par
  \textit{Gợi ý:} accumulating snapshot, vì mọi mốc của một đơn nằm trên cùng một dòng. Tỉ lệ không cộng được: cộng tử số và mẫu số rồi mới chia, ra 6,77\% thay vì 9,15\%.
  \item Vì sao tỉ lệ khách quay lại tính theo \code{customer\_\allowbreak{}id} của Olist luôn bằng 0? Khoá tự nhiên, khoá thay thế và khoá bền khác nhau thế nào trong một dimension Type 2?\par
  \textit{Gợi ý:} mỗi đơn có một \code{customer\_\allowbreak{}id} mới; người mua thật là \code{customer\_\allowbreak{}unique\_\allowbreak{}id}. Nghĩ tới S3 Versioning: khoá bền giống key của object, khoá thay thế giống version ID, còn khoá tự nhiên là thứ hệ nguồn đưa tới và có thể đổi.
  \item Một người bán chuyển từ SP sang RJ. Viết các dòng dimension sau thay đổi theo Type 1, Type 2 và Type 3. Vì sao khoảng hiệu lực của Type 2 là nửa mở?\par
  \textit{Gợi ý:} Type 1 còn một dòng RJ; Type 2 có hai dòng nối nhau tại ngày đổi; Type 3 có một dòng với bang hiện tại là RJ và bang trước đó là SP. Với \code{BETWEEN}, 256 đơn Olist khớp hai version và bị nhân dòng.
  \item Vì sao tách \code{dim\_\allowbreak{}category} không tiết kiệm chỗ trên Parquet? Cái giá thật của OBT là gì?\par
  \textit{Gợi ý:} dictionary encoding lưu mỗi tên danh mục một lần, cả cột \code{category} chỉ tốn 30.231 byte; năm cột chép vào OBT chỉ chiếm 5,6\% file. Cái giá là cập nhật: đổi tên \code{bed\_\allowbreak{}bath\_\allowbreak{}table} phải ghi lại 11.115 dòng, và dễ lỡ tay trộn hai grain trong một bảng.
  \item Trong Kimball, Inmon và Data Vault, dữ liệu chi tiết được lưu ở dạng nào, và người phân tích đọc bảng nào?\par
  \textit{Gợi ý:} Kimball lưu ngay ở dạng dimension và fact, người phân tích đọc star. Inmon lưu trong kho trung tâm 3NF, người phân tích đọc data mart dạng dimension nạp từ kho. Data Vault lưu ở hub, link và satellite, người phân tích đọc star dựng bên trên.
  \item Khi áp các trang version của \repo{source\_\allowbreak{}sim} vào snapshot bằng \code{MERGE}, hai chi tiết nào giữ cho bảng đúng, và bỏ mỗi chi tiết thì hỏng thế nào? Phần nào của các version là thật?\par
  \textit{Gợi ý:} giữ version mới nhất của mỗi đơn trong trang, nếu không DuckDB 1.5.6 cho 257.365 dòng thay vì 99.441; chỉ ghi đè bằng cursor lớn hơn, nếu không áp lại hay áp lộn thứ tự sẽ làm dòng lùi về version cũ. Các mốc thời gian và trạng thái ở version cuối là thật; trạng thái trung gian là mock giả lập.
  \item Viết phương trình đối soát của doanh thu Olist theo định nghĩa ở Bài 12. Nó chứng minh điều gì, và không chứng minh điều gì?\par
  \textit{Gợi ý:} 13.494.400,74 cộng 97.242,96 bằng 13.591.643,70. Nó chứng minh không đồng nào bị mất hay bị cộng hai lần, nhưng không chứng minh bỏ đơn huỷ là đúng: định nghĩa chỉ giữ đơn \code{delivered} cũng thoả phương trình.
\end{enumerate}

% =====================================================================
\section{Đọc thêm}

Đọc sau khi đã làm xong các lab, để gắn tên chuẩn cho những gì bạn đã tự làm. Thời gian ở mỗi mục là ước lượng.

\begin{itemize}
  \item \href{https://www.kimballgroup.com/wp-content/uploads/2013/08/2013.09-Kimball-Dimensional-Modeling-Techniques11.pdf}{Kimball Dimensional Modeling Techniques}: tài liệu PDF 24 trang của Kimball Group, bản 2013. Đọc trang 1 tới 15, khoảng 60 phút: mỗi kỹ thuật trong module có ở đó một đoạn ngắn, kèm tên chuẩn.
  \item \href{https://dataversity.net/data-vault-series-1-data-vault-overview/5054}{Dan Linstedt: Data Vault Series 1}: bài viết năm 2002 của người tạo ra Data Vault. Đọc phần định nghĩa hub, link và satellite, khoảng 20 phút; đây là bản đầu, chưa có khoá băm của bản 2.0.
  \item \href{https://www.fivetran.com/blog/star-schema-vs-obt}{Michael Kaminsky: Star Schema vs. OBT}: bài viết ngày 16/08/2022, đo hai cách trên Redshift, Snowflake và BigQuery. Đọc phần kết quả và phần chỗ lưu, khoảng 15 phút; số chỗ lưu là của Redshift, không phải của Parquet.
  \item \href{https://decisionworks.com/2004/03/differences-of-opinion/}{Margy Ross: Differences of Opinion}: bài viết năm 2004 so cách của Kimball với Corporate Information Factory của Inmon, nhìn từ phía Kimball. Đọc cả bài, khoảng 15 phút.
  \item \href{https://github.com/DataTalksClub/data-engineering-zoomcamp/tree/main/04-analytics-engineering}{DE Zoomcamp, Module 4 Analytics Engineering}: miễn phí, nhánh main của repo, xem ngày 06/11/2026. Unit 1 tóm tắt fact và dimension theo Kimball, khoảng 30 phút; các unit sau thuộc nhánh T1.
  \item \href{https://www.coursera.org/learn/data-modeling-transformation-serving}{DeepLearning.AI, khoá 4: Data Modeling, Transformation, and Serving}: ghi danh miễn phí. Xem tuần 1, khoảng 10 giờ theo trang khoá: Inmon so với Kimball, Data Vault và OBT, ứng với Bài 8 và Bài 11.
\end{itemize}
% Kiểm: trang Coursera khoá 4 xem ngày 06/11/2026: "3 weeks to complete", "10 hours a week"; tuần 1 có video "Inmon VS Kimball", "Data Vault", "One Big Table".
% Kiểm: raw.githubusercontent.com/DataTalksClub/data-engineering-zoomcamp/main/04-analytics-engineering/01-analytics-engineering-basics.md có mục "Dimensional Modeling" theo Kimball, 06/11/2026.

% =====================================================================
\section{Kế hoạch theo tuần}

Mỗi tuần 9 giờ cho module; giờ thứ mười dành cho chương nghề nghiệp. Tuần 4 chỉ cần 3 giờ, nên 6 giờ module còn lại của tuần đó dành cho Module 4.

\begin{xltabular}{\textwidth}{@{}P{1.0cm} L P{4.4cm} P{0.8cm}@{}}
\toprule
\thd{tuần} & \thd{việc lõi} & \thd{thêm nếu còn giờ} & \thd{giờ} \\
\midrule
\endhead
\bottomrule
\endfoot
1 & Bài 1 tới Bài 7, chạy mọi khối code và đọc kỹ từng output; Lab 3.1 & chạy lại khối Làm thử của Bài 1 với \code{SET threads = 1} & 9 \\
2 & Bài 8 tới Bài 10; Lab 3.2 và Lab 3.3 & bài viết của Kaminsky & 9 \\
3 & Bài 11 và Bài 12; Lab 3.4 và Lab 3.5; Tự kiểm tra; bắt đầu thử thách kỹ năng & phần thêm nếu còn giờ của Lab 3.5; bài viết của Linstedt và của Margy Ross & 9 \\
4 & hết thử thách kỹ năng; mốc capstone 1 & DE Zoomcamp Module 4 unit 1; khoá 4 của DeepLearning.AI & 3 \\
\end{xltabular}
% Ước lượng ngày 06/11/2026, chưa đo trên người học. Tuần 1 khoảng 8 giờ: 23 trang và 10 khối code từ đầu module tới hết Bài 7, cộng khoảng 2 giờ cho Lab 3.1. Tuần 2 khoảng 9 giờ: Bài 8 tới Bài 10 khoảng 2,5 giờ, Lab 3.2 khoảng 4 giờ, Lab 3.3 khoảng 2,5 giờ. Tuần 3 khoảng 9 giờ: Bài 11 và Bài 12 khoảng 1,25 giờ, Lab 3.4 khoảng 2 giờ, Lab 3.5 khoảng 4 giờ, Tự kiểm tra 0,5 giờ, phần còn lại cho thử thách. Tuần 4 khoảng 3 giờ. Tổng khoảng 29 tới 31 giờ. Số trang lấy từ bản dựng tectonic cùng ngày.

Nếu thiếu giờ, đọc nhanh các phần sau thay vì bỏ hẳn:
\begin{itemize}
  \item Bài 1 và Bài 10, khi bạn qua được bài kiểm 15 phút ở thẻ module.
  \item Bài 11: đọc phần Giải thích và sơ đồ; phần còn lại để dành cho lúc phải chọn mô hình cho một kho mới.
  \item Khối thứ hai của Bài 6 và bước 3 của Lab 3.4: chỉ đọc output in sẵn, không chạy.
\end{itemize}



\partpage{II}{Lưu trữ và lakehouse}{Object storage, định dạng file cột, Apache Iceberg từ bên trong, catalog và quản trị dữ liệu. Kết thúc bằng mốc capstone 2: tầng raw và staging trên Iceberg.}
% ---------------- m04_06 ----------------
% ===================================================================
% Module 4, 5, 6
% ===================================================================

\modulechapter{4}{Lưu trữ và định dạng file}
\moduleinfo{Trung cấp}{20 giờ}{Module 1, 2}

Mọi lakehouse cuối cùng đều là các file nằm trên object storage. Module này đi xuống tầng thấp nhất của vòng đời dữ liệu: dữ liệu được xếp trong file ra sao (Parquet, Avro, Arrow), và object storage tương thích S3 bảo đảm gì khi bạn ghi, liệt kê và đổi tên các file đó. Bạn sẽ mở footer của một file Parquet thật, đo tác động của nén, sắp xếp và chia partition, rồi kiểm các bảo đảm của S3 bằng AWS CLI. Học xong, bạn hiểu vì sao cần một table format như Iceberg (Module 5) đặt lên trên tầng này.

\section{Bạn sẽ học được}
\begin{itemize}
  \item Giải thích khác biệt giữa lưu theo dòng và lưu theo cột, và chọn Parquet, Avro hay Arrow cho đúng việc.
  \item Đọc cấu trúc một file Parquet: magic number, footer, row group, column chunk, page, thống kê min/max và page index.
  \item Chọn encoding, codec nén, kích thước row group và thứ tự sắp xếp để engine bỏ qua được nhiều dữ liệu.
  \item Áp dụng quy tắc schema resolution của Avro khi schema bên ghi và bên đọc khác nhau.
  \item Nêu chính xác S3 bảo đảm gì (strong read-after-write, ghi có điều kiện) và không bảo đảm gì (cập nhật nguyên tử trên nhiều key, rename).
  \item Đo và giải thích vấn đề nhiều file nhỏ, rồi chọn cách chia partition hợp lý.
\end{itemize}

\section{Nội dung chính}

\subsection{Định dạng file}
\begin{enumerate}
  \item \textbf{Lưu theo dòng và lưu theo cột.} Định dạng theo dòng (CSV, Avro) ghi trọn một bản ghi rồi mới tới bản ghi sau, hợp với ghi và đọc từng bản ghi. Định dạng theo cột (Parquet) gom giá trị của cùng một cột lại với nhau, nên truy vấn phân tích chỉ đọc những cột cần dùng và nén tốt hơn vì các giá trị giống kiểu đứng cạnh nhau. Parquet vẫn gom dòng: dữ liệu chia theo chiều ngang thành row group, trong mỗi row group mỗi cột là một column chunk (\href{https://parquet.apache.org/docs/concepts/}{Parquet: Concepts}).
  \item \textbf{Bố cục một file Parquet.} File bắt đầu bằng 4 byte magic \code{PAR1}; cuối file là file metadata, rồi 4 byte độ dài metadata (little endian), rồi lại \code{PAR1}. Metadata ghi sau dữ liệu để bên ghi chỉ cần đi một lượt; bên đọc đọc footer trước để biết vị trí mọi column chunk (\href{https://parquet.apache.org/docs/file-format/}{Parquet: File Format}). Các struct metadata là Thrift, tuần tự hoá bằng \code{TCompactProtocol} (\href{https://parquet.apache.org/docs/file-format/metadata/}{Parquet: Metadata}).
  \item \textbf{Row group, column chunk, page.} Row group là phân vùng ngang theo dòng; column chunk là dữ liệu của một cột trong một row group, nằm liền nhau trong file; page là đơn vị không chia nhỏ hơn về nén và encoding. Đơn vị song song: file hoặc row group cho việc chia job, column chunk cho I/O, page cho encoding và nén (\href{https://parquet.apache.org/docs/concepts/}{Parquet: Concepts}).
  \item \textbf{Thống kê, page index và bloom filter.} Metadata của mỗi column chunk mang thống kê như min, max, null count; \code{pyarrow} in chúng qua \code{metadata.\allowbreak row\_group(i).\allowbreak column(j).\allowbreak statistics} (\href{https://arrow.apache.org/docs/python/parquet.html}{pyarrow: Parquet}). Page index (ColumnIndex và OffsetIndex) là metadata tuỳ chọn, nằm gần footer, cho phép bỏ qua từng page; docs ghi nó giúp point lookup và range scan trên cột đã sắp xếp hiệu quả hơn về I/O (\href{https://parquet.apache.org/docs/file-format/pageindex/}{Parquet: Page Index}). DuckDB có \code{parquet\_bloom\_probe} để xem row group nào bỏ qua được nhờ bloom filter (\href{https://duckdb.org/docs/2.0/data/parquet/metadata}{DuckDB: Parquet Metadata}).
  \item \textbf{Encoding.} \code{PLAIN} dùng được cho mọi kiểu. Dictionary encoding hiện dùng \code{RLE\_DICTIONARY}; \code{PLAIN\_DICTIONARY} đã deprecated. Dictionary page được ghi trước các data page của column chunk; khi dictionary quá lớn, bên ghi quay về \code{PLAIN}. \code{DELTA\_BINARY\_PACKED} dành cho \code{INT32}, \code{INT64}. \code{BYTE\_STREAM\_SPLIT} tự nó không làm dữ liệu nhỏ đi; nó tách các byte của mỗi giá trị thành các luồng riêng để bước nén sau đó tốt hơn. Encoding \code{ALP} cho số thực đang ở trạng thái Preview (\href{https://parquet.apache.org/docs/file-format/data-pages/encodings/}{Parquet: Encodings}).
  \item \textbf{Nén.} Nén áp lên từng page. \code{pyarrow} mặc định dùng Snappy và hỗ trợ Brotli, Gzip, ZSTD, LZ4, không nén (\href{https://arrow.apache.org/docs/python/parquet.html}{pyarrow: Parquet}); Iceberg mặc định codec \code{zstd} (property \code{write.parquet.\allowbreak compression-codec}) (\href{https://iceberg.apache.org/docs/latest/configuration/}{Iceberg: Configuration}). Danh sách codec trong đặc tả ở \href{https://parquet.apache.org/docs/file-format/data-pages/compression/}{Parquet: Compression}.
  \item \textbf{Kích thước row group và page.} Parquet docs khuyên row group lớn (512 MB tới 1 GB) và page 8 KB (\href{https://parquet.apache.org/docs/file-format/configurations/}{Parquet: Configurations}); khái niệm block trong docs là block HDFS (\href{https://parquet.apache.org/docs/concepts/}{Concepts}). Iceberg mặc định row group 128 MB và file đích 512 MB (\code{write.parquet.\allowbreak row-group-size-bytes}, \code{write.target-\allowbreak file-size-bytes}, \href{https://iceberg.apache.org/docs/latest/configuration/}{Iceberg: Configuration}). Hai nguồn chính thức khác nhau, nên hãy tự đo trên dữ liệu và storage của bạn.
  \item \textbf{Dữ liệu lồng.} Parquet biểu diễn struct và list bằng definition level và repetition level; đọc mục Nested Encoding và Nulls trong \href{https://parquet.apache.org/docs/file-format/}{Parquet File Format}.
  \item \textbf{Avro và schema resolution.} Avro là định dạng theo dòng. Header của object container file bắt buộc có \code{avro.schema}, là schema JSON của dữ liệu trong file. Khi đọc bằng một schema khác, trường được khớp theo tên; trường bên ghi có mà bên đọc không có thì bị bỏ qua; trường bên đọc có mà bên ghi không có thì dùng giá trị mặc định, không có mặc định thì báo lỗi. Có các phép nâng kiểu như \code{int} lên \code{long}, \code{float}, \code{double}, và \code{string} với \code{bytes} đổi qua lại. Alias là tuỳ chọn của từng implementation (\href{https://avro.apache.org/docs/1.11.1/specification/}{Avro Specification}; trang có menu chọn bản mới hơn). Iceberg dùng Avro cho manifest (Module 5).
  \item \textbf{Arrow.} Arrow định nghĩa bố cục cột trong bộ nhớ để engine và thư viện trao đổi dữ liệu mà không phải chuyển đổi; \code{pyarrow} đọc Parquet thành bảng Arrow (\href{https://arrow.apache.org/docs/format/Columnar.html}{Arrow Columnar Format}). Arrow không thay Parquet làm định dạng lưu trữ lâu dài.
  \item \textbf{Schema khi ghi và khi đọc.} Parquet và Avro mang schema trong file; CSV thì không, nên engine phải đoán hoặc bạn phải khai báo. Đây là lý do tầng raw thường chuyển CSV sang Parquet sớm.
\end{enumerate}

\subsection{Object storage tương thích S3}
\begin{enumerate}[resume]
  \item \textbf{Key phẳng.} S3 có bucket và object, không có cây thư mục; ``folder'' trên console chỉ là prefix của key cộng dấu phân cách \code{/} (\href{https://docs.aws.amazon.com/AmazonS3/latest/userguide/object-keys.html}{S3: Object keys}).
  \item \textbf{Consistency.} S3 có strong read-after-write cho PUT và DELETE ở mọi region; GET hoặc LIST sau một PUT thành công sẽ thấy dữ liệu vừa ghi. Nhưng cập nhật tính theo từng key, không có cách cập nhật nguyên tử trên nhiều key; hai PUT cùng lúc vào một key thì request có timestamp mới hơn thắng, và S3 không có object locking cho writer đồng thời (\href{https://docs.aws.amazon.com/AmazonS3/latest/userguide/Welcome.html\#ConsistencyModel}{S3: data consistency model}).
  \item \textbf{Không có rename nguyên tử.} Với general purpose bucket, đổi tên là copy sang key mới rồi xoá (thêm delete marker cho) bản cũ; ``move'' một folder là nhiều thao tác trên từng object (\href{https://docs.aws.amazon.com/AmazonS3/latest/userguide/copy-object.html}{S3: Copying, moving, and renaming objects}). API \code{RenameObject} chỉ dành cho directory bucket (\href{https://docs.aws.amazon.com/AmazonS3/latest/userguide/directory-buckets-objects-rename.html}{S3: Renaming objects in directory buckets}).
  \item \textbf{Ghi có điều kiện.} \code{If-None-Match: *} chỉ ghi khi key chưa tồn tại; \code{If-Match: <ETag>} chỉ ghi khi ETag hiện tại khớp. Hai header dùng được với \code{PutObject}, \code{Complete\allowbreak MultipartUpload}, \code{CopyObject}; điều kiện sai trả \code{412 Precondition Failed}. Khi hai lần ghi có điều kiện đua nhau, lần hoàn thành trước thắng, các lần sau nhận 412; nếu có lệnh xoá chen vào thì có thể nhận \code{409 Conflict}. Bucket policy có thể bắt buộc ghi có điều kiện (\href{https://docs.aws.amazon.com/AmazonS3/latest/userguide/conditional-writes.html}{S3: Conditional writes}). Storage tương thích S3 khác có thể chưa hỗ trợ đủ; hãy kiểm bằng lab.
  \item \textbf{Liệt kê.} \code{ListObjectsV2} trả tối đa 1.000 key mỗi request, phân trang bằng \code{Continuation\allowbreak Token}; với general purpose bucket, key được trả theo thứ tự từ điển; \code{Prefix} cộng \code{Delimiter} gộp key thành \code{Common\allowbreak Prefixes} (\href{https://docs.aws.amazon.com/AmazonS3/latest/API/API_ListObjectsV2.html}{ListObjectsV2}).
  \item \textbf{Multipart upload.} AWS khuyên dùng multipart cho object từ 100 MB; số part từ 1 tới 10.000; object chỉ được dựng khi gọi complete; part của upload dở dang vẫn tính phí cho tới khi complete hoặc abort, nên đặt lifecycle rule \code{AbortIncompleteMultipartUpload} (\href{https://docs.aws.amazon.com/AmazonS3/latest/userguide/mpuoverview.html}{S3: Multipart upload}).
  \item \textbf{Partition và vấn đề nhiều file nhỏ.} Chia partition kiểu Hive (\code{col=value/}) giúp bỏ qua thư mục khi lọc, nhưng cột có quá nhiều giá trị sinh ra rất nhiều file nhỏ. Mỗi file tốn ít nhất một mục khi liệt kê (1.000 key mỗi request), một lần đọc footer và một lần mở file; Iceberg docs ghi rằng file nhỏ sinh ra lượng metadata không cần thiết và truy vấn kém hiệu quả vì chi phí mở file, và có thủ tục compaction để gộp chúng (\href{https://iceberg.apache.org/docs/latest/maintenance/}{Iceberg: Maintenance}, mục Compact data files). \code{pyarrow.parquet.\allowbreak write\_to\_dataset} với \code{partition\_cols} là cách nhanh để tạo dataset chia partition và thấy vấn đề này (\href{https://arrow.apache.org/docs/python/parquet.html}{pyarrow: Parquet}).
\end{enumerate}

\section{Bắt đầu}
\begin{itemize}
  \item \href{https://arrow.apache.org/docs/python/parquet.html}{pyarrow: Reading and Writing the Apache Parquet Format}: ghi một bảng, đọc metadata của file (\code{FileMetaData}), rồi xem \code{row\_group(0).\allowbreak column(0)}.
  \item \href{https://duckdb.org/docs/2.0/data/parquet/metadata}{DuckDB: Querying Parquet Metadata}: chạy \code{parquet\_metadata}, \code{parquet\_schema} trên cùng file.
  \item \href{https://docs.aws.amazon.com/AmazonS3/latest/userguide/conditional-writes.html}{S3: Conditional writes}: chạy ví dụ \code{put-object --if-none-match "*"} của trang, trỏ \code{--endpoint-url} tới RustFS.
\end{itemize}

\section{Tài liệu chính thức}
\begin{xltabular}{\textwidth}{@{}P{4.2cm} L P{1.8cm}@{}}
\toprule
\thd{tài liệu} & \thd{nội dung cần đọc} & \thd{thời lượng} \\
\midrule
\endhead
\bottomrule
\endfoot
\href{https://parquet.apache.org/docs/concepts/}{Parquet: Concepts} & block, file, row group, column chunk, page, đơn vị song song & 0,5 giờ \\
\href{https://parquet.apache.org/docs/file-format/}{Parquet: File Format} & bố cục file; các trang con Metadata, Nested Encoding, Nulls, Bloom Filter & 2 giờ \\
\href{https://parquet.apache.org/docs/file-format/data-pages/encodings/}{Parquet: Encodings} & \code{PLAIN}, \code{RLE\_DICTIONARY}, \code{DELTA\_BINARY\_PACKED}, \code{BYTE\_STREAM\_SPLIT}, fallback về \code{PLAIN} & 1 giờ \\
\href{https://parquet.apache.org/docs/file-format/pageindex/}{Parquet: Page Index} & ColumnIndex, OffsetIndex, vì sao đặt gần footer & 0,5 giờ \\
\href{https://parquet.apache.org/docs/file-format/configurations/}{Parquet: Configurations} & khuyến nghị row group và page size & 0,25 giờ \\
\href{https://arrow.apache.org/docs/python/parquet.html}{pyarrow: Parquet} & \code{write\_table}, \code{read\_metadata}, \code{ParquetFile}, \code{write\_to\_dataset} & 1,5 giờ \\
\href{https://duckdb.org/docs/2.0/data/parquet/metadata}{DuckDB: Parquet Metadata} & \code{parquet\_metadata}, \code{parquet\_schema}, \code{parquet\_file\_metadata}, \code{parquet\_kv\_metadata}, \code{parquet\_bloom\_probe} & 0,5 giờ \\
\href{https://avro.apache.org/docs/1.11.1/specification/}{Avro Specification} & Object Container Files, Schema Resolution, Aliases & 1,5 giờ \\
\href{https://arrow.apache.org/docs/format/Columnar.html}{Arrow Columnar Format} & buffer, validity bitmap & 1 giờ \\
\href{https://docs.aws.amazon.com/AmazonS3/latest/userguide/Welcome.html\#ConsistencyModel}{S3: data consistency model} & strong read-after-write, không cập nhật nguyên tử trên nhiều key & 0,5 giờ \\
\href{https://docs.aws.amazon.com/AmazonS3/latest/userguide/conditional-writes.html}{S3: Conditional writes} & \code{If-None-Match}, \code{If-Match}, 412, 409, giới hạn với multipart & 1 giờ \\
\href{https://docs.aws.amazon.com/AmazonS3/latest/userguide/object-keys.html}{S3: Object keys} và \href{https://docs.aws.amazon.com/AmazonS3/latest/userguide/copy-object.html}{Copy, move, rename} & mô hình phẳng; rename bằng copy rồi xoá & 0,5 giờ \\
\href{https://docs.aws.amazon.com/AmazonS3/latest/API/API_ListObjectsV2.html}{ListObjectsV2} và \href{https://docs.aws.amazon.com/AmazonS3/latest/userguide/mpuoverview.html}{Multipart upload} & phân trang, \code{CommonPrefixes}; ngưỡng 100 MB, part 1--10.000 & 1 giờ \\
\end{xltabular}

\section{Khoá học miễn phí}
\begin{itemize}
  \item \href{https://15445.courses.cs.cmu.edu/fall2025/schedule.html}{CMU 15-445/645 Database Systems, Fall 2025}: miễn phí slide và video; xem bài \#03 Database Storage I, \#05 Database Storage II, \#06 Storage Models \& Compression.
\end{itemize}

\section{Lab thực hành}
Chuẩn bị: tải bộ Olist từ Kaggle, giải nén vào thư mục \code{olist/}; cài \code{pyarrow} và DuckDB CLI; chạy RustFS trong Docker Compose của stack tham chiếu và có AWS CLI v2.

\subsection{Lab 4.1: CSV so với Parquet và các codec}
\begin{enumerate}
  \item Trong DuckDB, bật \code{.timer on} rồi nạp bảng lớn nhất của Olist (geolocation, khoảng một triệu dòng):
\begin{codeblock}
CREATE TABLE geo AS SELECT * FROM read_csv('olist/olist_geolocation_dataset.csv');
COPY geo TO 'geo_snappy.parquet' (FORMAT parquet, COMPRESSION snappy);
COPY geo TO 'geo_zstd.parquet'   (FORMAT parquet, COMPRESSION zstd);
COPY geo TO 'geo_none.parquet'   (FORMAT parquet, COMPRESSION uncompressed);
\end{codeblock}
  \item Ghi kích thước bốn file (\code{ls -l}). Chạy cùng truy vấn \code{SELECT geolocation\_state, avg(geolocation\_lat) FROM ... GROUP BY 1} trên CSV và ba file Parquet, ghi thời gian.
  \item Dùng \code{parquet\_metadata('geo\_zstd.parquet')} xem cột \code{compression} và \code{encodings} của từng cột; cột nào được dictionary encoding?
\end{enumerate}
\hoanthanh{bạn có bảng 4 dòng (định dạng, kích thước, thời gian truy vấn) và giải thích được bằng metadata vì sao cột \code{geolocation\_state} nhỏ hơn nhiều so với cột số thực.}

\subsection{Lab 4.2: Mở footer, sắp xếp và bỏ qua row group}
\begin{enumerate}
  \item Ghi bảng orders hai lần, một lần sắp theo thời gian mua, một lần xáo trộn, với row group nhỏ để dễ quan sát:
\begin{codeblock}
CREATE TABLE orders AS SELECT * FROM read_csv('olist/olist_orders_dataset.csv');
COPY (SELECT * FROM orders ORDER BY order_purchase_timestamp)
  TO 'orders_sorted.parquet' (FORMAT parquet, ROW_GROUP_SIZE 10000);
COPY (SELECT * FROM orders ORDER BY random())
  TO 'orders_random.parquet' (FORMAT parquet, ROW_GROUP_SIZE 10000);
\end{codeblock}
  \item Xem 4 byte đầu và 8 byte cuối: \code{head -c 4 orders\_sorted.parquet | xxd}, \code{tail -c 8 orders\_sorted.parquet | xxd}. Đổi 4 byte trước \code{PAR1} (little endian) ra số thập phân: đó là độ dài footer.
  \item Đọc footer bằng \code{pyarrow}:
\begin{codeblock}
import pyarrow.parquet as pq
md = pq.read_metadata("orders_sorted.parquet")
print(md.num_rows, md.num_row_groups, md.serialized_size)
for i in range(md.num_row_groups):
    c = md.row_group(i).column(3)   # order_purchase_timestamp
    print(i, c.path_in_schema, c.compression, c.encodings, c.statistics.min, c.statistics.max)
\end{codeblock}
  \item So với file xáo trộn:
\begin{codeblock}
SELECT row_group_id, stats_min, stats_max
FROM parquet_metadata('orders_random.parquet')
WHERE path_in_schema = 'order_purchase_timestamp';
\end{codeblock}
  \item Chạy truy vấn sau rồi lặp lại với file xáo trộn, so số dòng thật sự được quét:
\begin{codeblock}
EXPLAIN ANALYZE SELECT count(*) FROM 'orders_sorted.parquet'
WHERE order_purchase_timestamp BETWEEN '2017-11-24' AND '2017-11-25';
\end{codeblock}
\end{enumerate}
\hoanthanh{độ dài footer bạn tính từ 8 byte cuối khớp với \code{md.serialized\_size}, và bạn chỉ ra bằng \code{stats\_min}, \code{stats\_max} vì sao file đã sắp chỉ cần đọc một hai row group còn file xáo trộn phải đọc gần hết.}

\subsection{Lab 4.3: Kiểm các bảo đảm của S3 trên RustFS}
\begin{enumerate}
  \item Đặt biến môi trường credential của RustFS, rồi dùng alias cho gọn (thay endpoint bằng cổng S3 API trong compose của bạn):
\begin{codeblock}
alias s3='aws --endpoint-url http://localhost:9000 s3'
alias s3api='aws --endpoint-url http://localhost:9000 s3api'
s3 mb s3://lab
s3 cp orders_sorted.parquet s3://lab/raw/orders/year=2017/part-0.parquet
s3 cp orders_random.parquet s3://lab/raw/orders/year=2018/part-0.parquet
s3api list-objects-v2 --bucket lab --prefix raw/orders/ --delimiter /
\end{codeblock}
  \item Quan sát \code{CommonPrefixes}: ``thư mục'' chỉ là prefix. Chạy \code{s3 mv s3://lab/raw/orders/ s3://lab/raw/orders\_v2/ --recursive} và đọc output: mỗi object được chuyển riêng lẻ.
  \item Ghi có điều kiện hai lần vào cùng một key:
\begin{codeblock}
echo v1 > flag.txt
s3api put-object --bucket lab --key flag.txt --body flag.txt --if-none-match "*"
s3api put-object --bucket lab --key flag.txt --body flag.txt --if-none-match "*"
\end{codeblock}
  Trên Amazon S3, lần hai phải nhận 412. Ghi lại RustFS phiên bản bạn dùng trả gì.
  \item Làm một multipart upload bằng các lệnh \code{create-multipart-upload}, \code{upload-part} (ít nhất 2 part), \code{list-multipart-uploads}, \code{complete-multipart-upload} theo trang Multipart upload; kiểm rằng object chỉ xuất hiện trong \code{list-objects-v2} sau bước complete.
\end{enumerate}
\hoanthanh{bạn có một ghi chú ngắn: kết quả của lần PUT có điều kiện thứ hai, thời điểm object multipart xuất hiện, và giải thích vì sao \code{mv} một prefix không phải là thao tác nguyên tử.}

\subsection{Lab 4.4: Vấn đề nhiều file nhỏ}
\begin{enumerate}
  \item Ghi bảng customers thành dataset chia partition theo cột có rất nhiều giá trị, và theo cột có ít giá trị:
\begin{codeblock}
import pyarrow.csv as pv, pyarrow.parquet as pq
t = pv.read_csv("olist/olist_customers_dataset.csv")
pq.write_to_dataset(t, root_path="cust_by_zip",   partition_cols=["customer_zip_code_prefix"])
pq.write_to_dataset(t, root_path="cust_by_state", partition_cols=["customer_state"])
pq.write_table(t, "cust_one.parquet")
\end{codeblock}
  \item Đếm file bằng \code{find cust\_by\_zip -name '*.parquet' | wc -l}. Tính số request \code{ListObjectsV2} tối thiểu để liệt kê dataset đó trên S3.
  \item Trong DuckDB, đo truy vấn sau, rồi chạy cùng truy vấn trên \code{cust\_by\_state} và \code{cust\_one.parquet}:
\begin{codeblock}
SELECT customer_state, count(*)
FROM read_parquet('cust_by_zip/**/*.parquet', hive_partitioning = true)
GROUP BY 1;
\end{codeblock}
  \item Gộp \code{cust\_by\_zip} thành một file bằng \code{COPY (SELECT * FROM read\_parquet(...)) TO ...} rồi đo lại.
\end{enumerate}
\hoanthanh{bạn có bảng (cách ghi, số file, tổng dung lượng, thời gian truy vấn) và nêu được ba chi phí của file nhỏ: liệt kê, đọc footer và mở file, nén kém vì row group nhỏ.}

\section{Kế hoạch theo tuần}
\begin{xltabular}{\textwidth}{@{}P{1.4cm} L P{1.4cm}@{}}
\toprule
\thd{tuần} & \thd{việc} & \thd{giờ} \\
\midrule
\endhead
\bottomrule
\endfoot
1 & Đọc Parquet Concepts, File Format, Metadata, Encodings, Page Index, Configurations; pyarrow và DuckDB metadata; CMU bài \#06; làm Lab 4.1, 4.2 & 10 \\
2 & Đọc Avro Specification, Arrow Columnar Format, các trang S3; làm Lab 4.3, 4.4; thử thách kỹ năng; câu hỏi tự kiểm tra & 10 \\
thêm & CMU bài \#03, \#05; Parquet Nested Encoding và Bloom Filter; ORC specification để so sánh & 5 \\
\end{xltabular}

\begin{thuthach}
Đội phân tích cần hai truy vấn chạy nhiều lần mỗi ngày trên dữ liệu Olist: doanh thu theo tháng và theo bang của khách hàng, và tra cứu mọi dòng hàng của một \code{order\_id}. Hãy thiết kế và tạo trên RustFS một bố cục Parquet cho bảng nối \code{order\_items}, \code{orders}, \code{customers}: chọn cột chia partition (hoặc không chia), thứ tự sắp xếp, kích thước row group, codec. Không có hướng dẫn từng bước; tự chọn công cụ trong stack tham chiếu.
\hoanthanh{bạn trình bày được, bằng output của \code{parquet\_metadata} và \code{EXPLAIN ANALYZE}, rằng truy vấn một tháng chỉ đọc một phần nhỏ row group, truy vấn một \code{order\_id} bỏ qua được phần lớn dữ liệu, và dataset không có file nào nhỏ hơn mức bạn tự đặt ra kèm lý do.}
\end{thuthach}

\section{Câu hỏi tự kiểm tra}
\begin{enumerate}
  \item Vì sao bên đọc Parquet nhảy tới cuối file trước, và 8 byte cuối chứa gì?
  \item Row group, column chunk và page khác nhau thế nào? Cái nào là đơn vị của I/O, cái nào là đơn vị của nén?
  \item Min và max nằm ở đâu trong file, và vì sao sắp xếp trước khi ghi giúp bỏ qua nhiều dữ liệu hơn? Page index thêm được gì?
  \item Khi nào dictionary encoding hết có lợi, và bên ghi làm gì khi đó?
  \item \code{BYTE\_STREAM\_SPLIT} có tự làm file nhỏ đi không? Vì sao vẫn có ích?
  \item Parquet docs và Iceberg đưa ra kích thước row group khác nhau. Bạn sẽ quyết định thế nào cho dữ liệu trên object storage?
  \item Trong Avro, bên đọc có trường mà bên ghi không ghi và trường đó không có mặc định thì sao? Nếu có mặc định thì sao?
  \item Arrow giải quyết vấn đề gì mà Parquet không, và vì sao engine đọc Parquet thành Arrow?
  \item Kể hai bảo đảm S3 đưa ra và hai điều S3 không bảo đảm.
  \item \code{If-None-Match: *} giúp gì cho việc commit, và hai lần ghi đua nhau thì kết quả ra sao?
  \item Vì sao ``đổi tên thư mục'' không dùng làm commit nguyên tử trên S3 general purpose bucket?
  \item Bảng chia partition theo cột có hàng chục nghìn giá trị, mỗi file vài KB: kể ba cái giá phải trả.
\end{enumerate}

\section{Điểm dễ nhầm}
\begin{xltabular}{\textwidth}{@{}P{4.2cm} L P{3.2cm}@{}}
\toprule
\thd{điểm dễ nhầm} & \thd{thực tế} & \thd{docs chính thức} \\
\midrule
\endhead
\bottomrule
\endfoot
Nhiều người nghĩ S3 có thư mục & Mô hình dữ liệu phẳng; folder chỉ là prefix của key & \href{https://docs.aws.amazon.com/AmazonS3/latest/userguide/object-keys.html}{S3 Object keys} \\
``Đổi tên thư mục trên S3 là nguyên tử'' & Rename là copy rồi xoá từng object; \code{RenameObject} chỉ có ở directory bucket & \href{https://docs.aws.amazon.com/AmazonS3/latest/userguide/copy-object.html}{S3 Copy, move, rename} \\
``LIST trên S3 có thể thiếu file vừa ghi'' & GET và LIST sau một PUT thành công thấy dữ liệu mới & \href{https://docs.aws.amazon.com/AmazonS3/latest/userguide/Welcome.html\#ConsistencyModel}{S3 consistency} \\
``Strong consistency nghĩa là S3 có transaction'' & Không có cập nhật nguyên tử trên nhiều key; hai PUT cùng key thì bản mới hơn thắng & \href{https://docs.aws.amazon.com/AmazonS3/latest/userguide/Welcome.html\#ConsistencyModel}{S3 consistency} \\
``Lưu theo cột nghĩa là không gom dòng'' & Parquet gom dòng thành row group, rồi chia theo cột thành column chunk và page & \href{https://parquet.apache.org/docs/concepts/}{Parquet Concepts} \\
``Row group càng lớn càng tốt ở mọi nơi'' & Khuyến nghị 512 MB--1 GB của Parquet docs gắn với block HDFS; Iceberg mặc định 128 MB & \href{https://parquet.apache.org/docs/file-format/configurations/}{Parquet Configurations}, \href{https://iceberg.apache.org/docs/latest/configuration/}{Iceberg Configuration} \\
``\code{BYTE\_STREAM\_SPLIT} là một kiểu nén'' & Nó chỉ sắp lại byte để codec phía sau nén tốt hơn & \href{https://parquet.apache.org/docs/file-format/data-pages/encodings/}{Parquet Encodings} \\
``Arrow là định dạng lưu trữ cạnh tranh với Parquet'' & Arrow là bố cục cột trong bộ nhớ & \href{https://arrow.apache.org/docs/format/Columnar.html}{Arrow Columnar Format} \\
``Avro không cần schema khi đọc'' & File container mang schema bên ghi; bên đọc đối chiếu với schema của mình theo luật resolution & \href{https://avro.apache.org/docs/1.11.1/specification/}{Avro Specification} \\
``Càng nhiều partition càng nhanh'' & Nhiều file nhỏ làm tăng metadata và chi phí mở file; cần compaction & \href{https://iceberg.apache.org/docs/latest/maintenance/}{Iceberg Maintenance} \\
\end{xltabular}


% ===================================================================

\modulechapter{5}{Open table format: Apache Iceberg}
\moduleinfo{Trung cấp tới nâng cao}{35 giờ}{Module 4}

Module 4 cho thấy object storage không có cập nhật nguyên tử trên nhiều key và không có rename nguyên tử. Apache Iceberg bù vào bằng một cây metadata bất biến và một thao tác duy nhất: đổi nguyên tử con trỏ tới file metadata hiện tại. Module này đi từ đặc tả (metadata file, manifest list, manifest, sequence number, commit) tới vận hành (time travel, schema và partition evolution, xoá theo dòng, bảo trì, branch và tag). Học xong, bạn ghi được dữ liệu Olist vào Iceberg bằng PyIceberg và Spark, đọc được mọi file metadata nó tạo ra, và giữ bảng khoẻ mạnh bằng các thủ tục bảo trì.

\section{Bạn sẽ học được}
\begin{itemize}
  \item Vẽ cây metadata của Iceberg và nói được ai ghi file nào khi dùng REST catalog và khi dùng catalog kiểu JDBC hoặc Hive.
  \item Giải thích optimistic concurrency: đổi nguyên tử con trỏ, luật retry, vai trò của sequence number, và ý nghĩa của mã lỗi 409 so với 5xx trong REST catalog.
  \item Dùng time travel, schema evolution, hidden partitioning và partition evolution trên Spark và PyIceberg.
  \item Chọn copy-on-write hay merge-on-read; phân biệt position delete, equality delete và deletion vector của v3.
  \item Chạy và giải thích bốn thủ tục bảo trì: \code{expire\_snapshots}, \code{rewrite\_data\_files}, \code{rewrite\_manifests}, \code{remove\_orphan\_files}.
  \item Đọc metadata table (\code{snapshots}, \code{manifests}, \code{entries}, \code{files}) để chẩn đoán bảng.
  \item Dùng branch và tag cho write-audit-publish, và so sánh ngắn Iceberg với Delta Lake, Apache Hudi.
\end{itemize}

\section{Nội dung chính}

\subsection{Cấu trúc bảng}
\begin{enumerate}
  \item \textbf{Vì sao cần table format.} Bảng kiểu Hive giữ partition trong metastore và file trong file system, nên không đổi được nội dung bảng một cách nguyên tử. Iceberg theo dõi danh sách đầy đủ data file của bảng trong từng snapshot (\href{https://iceberg.apache.org/docs/latest/reliability/}{Iceberg: Reliability}).
  \item \textbf{Cây metadata.} Catalog giữ con trỏ tới file \code{metadata.json} hiện tại. File này (JSON) chứa schema, partition spec, sort order, danh sách snapshot, \code{snapshot-log}, ref (branch, tag) và table property. Mỗi snapshot trỏ tới một manifest list; manifest list liệt kê các manifest; manifest là file Avro bất biến liệt kê data file hoặc delete file kèm partition và thống kê cột (\href{https://iceberg.apache.org/spec/}{Iceberg spec}). Data file là Parquet, Avro hoặc ORC.
  \item \textbf{Thống kê cột trong manifest.} Iceberg ghi bound của cột vào manifest theo \code{write.metadata.\allowbreak metrics.default}, mặc định \code{truncate(16)}: cột chuỗi chỉ giữ 16 ký tự đầu làm bound; mặc định chỉ thu metrics cho 100 cột đầu (\code{write.metadata.\allowbreak metrics.max-inferred-\allowbreak column-defaults}, \href{https://iceberg.apache.org/docs/latest/configuration/}{Configuration}). Vì vậy min/max trong footer Parquet (Module 4) và bound trong manifest là hai chỗ khác nhau, có thể khác giá trị.
  \item \textbf{Format version.} Theo đặc tả, v1, v2, v3 đã hoàn chỉnh và được cộng đồng thông qua; v4 đang phát triển và chưa được thông qua. v3 thêm kiểu dữ liệu mới, giá trị mặc định cho cột, transform nhiều tham số, row lineage, deletion vector nhị phân và khoá mã hoá bảng (\href{https://iceberg.apache.org/spec/}{Iceberg spec}). Bảng mới mặc định \code{format-version = 2} từ Iceberg 1.4.0 (\href{https://iceberg.apache.org/docs/latest/configuration/}{Configuration}). Engine có thể chưa làm đủ tính năng v3; xem docs của engine trước khi nâng.
\end{enumerate}

\subsection{Commit và concurrency}
\begin{enumerate}[resume]
  \item \textbf{Optimistic concurrency.} Đặc tả viết: việc đổi nguyên tử một file metadata lấy một file khác là nền của serializable isolation. Writer tạo metadata mới dựa trên phiên bản hiện tại; nếu phiên bản đó đã cũ khi commit, writer phải thử lại dựa trên phiên bản mới. Một số thao tác áp lại được: append áp lên snapshot mới được, rewrite áp lại được nếu mọi file bị thay vẫn còn trong bảng (\href{https://iceberg.apache.org/spec/}{Iceberg spec}, mục Optimistic Concurrency). Số lần thử lại mặc định \code{commit.retry.num-retries = 4} (\href{https://iceberg.apache.org/docs/latest/configuration/}{Configuration}).
  \item \textbf{Sequence number.} Manifest mới thừa kế sequence number từ metadata của manifest list, nên manifest được ghi một lần và dùng lại khi retry; muốn đổi sequence number cho lần thử mới thì chỉ cần ghi lại manifest list (\href{https://iceberg.apache.org/spec/}{Iceberg spec}). Sequence number cũng quyết định delete file nào áp lên data file nào.
  \item \textbf{Mức isolation cho ghi theo dòng.} \code{write.delete.\allowbreak isolation-level}, \code{write.update.\allowbreak isolation-level}, \code{write.merge.\allowbreak isolation-level} mặc định \code{serializable}, có thể đặt \code{snapshot}; các property này áp cho lệnh delete, update, merge (\href{https://iceberg.apache.org/docs/latest/configuration/}{Configuration}).
  \item \textbf{Commit với catalog kiểu metastore.} Với JDBC, Hive và các catalog tương tự, chính client ghi file metadata mới (đặc tả mô tả tên \code{<V+1>-<uuid>.metadata.json}) rồi nhờ catalog đổi con trỏ bằng check-and-put. Bảng kiểu file system (Hadoop catalog) commit bằng rename; đặc tả ghi cách này không an toàn trên object store và file system cục bộ, và sẽ bỏ ở v4 (\href{https://iceberg.apache.org/spec/}{Iceberg spec}).
  \item \textbf{Commit qua REST catalog.} Client gửi \code{CommitTableRequest} gồm định danh bảng, \code{requirements} (ví dụ \code{assert-ref-snapshot-id}) và \code{updates} (ví dụ \code{add-snapshot}, \code{set-snapshot-ref}) tới endpoint \code{updateTable}. Server kiểm requirement, áp update lên metadata hiện tại, tự ghi \code{metadata.json} mới, đổi con trỏ, rồi trả metadata mới (\href{https://iceberg.apache.org/rest-catalog-spec/}{REST Catalog Spec}, \href{https://github.com/apache/iceberg/blob/main/open-api/rest-catalog-open-api.yaml}{OpenAPI}). PyIceberg ánh xạ 409 thành \code{CommitFailed\allowbreak Exception} (requirement không thoả hoặc xung đột), còn 500, 502, 504 thành \code{CommitState\allowbreak Unknown\allowbreak Exception} (\href{https://py.iceberg.apache.org/reference/pyiceberg/catalog/rest/}{PyIceberg: rest}). Khi trạng thái commit chưa biết, client không được xoá các file vừa ghi, vì commit có thể đã thành công.
\end{enumerate}

\textbf{Ai ghi file nào trong một lần append.} Sơ đồ dưới áp cho Spark; với PyIceberg, ``client'' là tiến trình Python làm cả việc của driver và executor.

\begin{codeblock}
tầng                    định dạng      REST catalog                catalog JDBC / Hive
con trỏ tới metadata    bản ghi DB     catalog server đổi nguyên tử catalog đổi bằng check-and-put
        |                              theo yêu cầu của client
metadata.json           JSON           catalog server              Spark driver
        |
manifest list           Avro           Spark driver                Spark driver
        |
manifest                Avro           Spark driver                Spark driver
        |
data file, delete file  Parquet/ORC/.. Spark executor              Spark executor
\end{codeblock}

\subsection{Đọc, tiến hoá và xoá}
\begin{enumerate}[resume]
  \item \textbf{Time travel.} Theo đặc tả, mỗi lần \code{current-snapshot-id} thay đổi thì \code{snapshot-log} có thêm một mục; đây không phải chuỗi \code{parent-snapshot-id} (\href{https://iceberg.apache.org/spec/}{Iceberg spec}). Trong Spark: \code{VERSION AS OF <snapshot\_id>} và \code{TIMESTAMP AS OF '...'} (\href{https://iceberg.apache.org/docs/latest/spark-queries/}{Spark Queries}); trong PyIceberg: \code{table.\allowbreak scan(snapshot\_id=...)} (\href{https://py.iceberg.apache.org/api/}{PyIceberg API}).
  \item \textbf{Schema evolution.} Cột được theo dõi bằng field ID, nên thêm, xoá, đổi tên, đổi thứ tự cột không phải ghi lại data file (\href{https://iceberg.apache.org/docs/latest/evolution/}{Evolution}). PyIceberg: \code{with table.\allowbreak update\_schema() as u:} với \code{add\_column}, \code{rename\_column}; thay đổi không tương thích cần \code{allow\_\allowbreak incompatible\_\allowbreak changes=True} (\href{https://py.iceberg.apache.org/api/}{PyIceberg API}).
  \item \textbf{Hidden partitioning và partition evolution.} Giá trị partition được suy ra từ cột gốc bằng transform (\code{year}, \code{month}, \code{day}, \code{hour}, \code{bucket}, \code{truncate}, \code{identity}); người dùng lọc trên cột gốc và Iceberg tự dùng partition. Đổi partition spec không ghi lại dữ liệu cũ; mỗi file giữ spec mà nó được ghi (\href{https://iceberg.apache.org/docs/latest/partitioning/}{Partitioning}, \href{https://iceberg.apache.org/docs/latest/evolution/}{Evolution}).
  \item \textbf{Copy-on-write và merge-on-read.} Ba property \code{write.\allowbreak delete.mode}, \code{write.\allowbreak update.mode}, \code{write.\allowbreak merge.mode} mặc định \code{copy-on-write}: file chứa dòng bị đổi được ghi lại toàn bộ. \code{merge-on-read} (v2 trở lên) ghi delete file và để bên đọc ghép lại, nên ghi rẻ hơn, đọc đắt hơn và cần compaction thường xuyên hơn (\href{https://iceberg.apache.org/docs/latest/configuration/}{Configuration}).
  \item \textbf{Delete file.} Position delete được mã hoá trong position delete file (v2) hoặc deletion vector (v3 trở lên); deletion vector lưu trong file Puffin, và phải chứa mọi xoá từ các position delete file hiện có của data file đó. Equality delete xoá theo giá trị cột (\href{https://iceberg.apache.org/spec/}{Iceberg spec}, mục Delete Formats; \href{https://github.com/apache/iceberg/blob/main/format/puffin-spec.md}{Puffin spec}).
\end{enumerate}

\subsection{Vận hành}
\begin{enumerate}[resume]
  \item \textbf{Metadata table.} Spark đọc được \code{<bảng>.snapshots}, \code{.history}, \code{.manifests}, \code{.entries}, \code{.files}, \code{.refs}, \code{.metadata\_\allowbreak log\_entries} (\href{https://iceberg.apache.org/docs/latest/spark-queries/}{Spark Queries}). PyIceberg có \code{table.\allowbreak inspect.\allowbreak snapshots()}, \code{manifests()}, \code{entries()}, \code{files()}, \code{history()}, \code{refs()}, \code{partitions()}, \code{data\_files()}, \code{delete\_files()} (\href{https://py.iceberg.apache.org/api/}{PyIceberg API}).
  \item \textbf{Bảo trì được khuyên.} Expire snapshot để bỏ snapshot cũ và file không còn dùng; mặc định giữ snapshot trong 5 ngày và ít nhất 1 snapshot (\code{history.\allowbreak expire.\allowbreak max-snapshot-\allowbreak age-ms}, \code{history.\allowbreak expire.\allowbreak min-snapshots-\allowbreak to-keep}). Xoá metadata cũ bằng \code{write.\allowbreak metadata.\allowbreak delete-after-\allowbreak commit.enabled} (mặc định \code{false}) và \code{write.\allowbreak metadata.\allowbreak previous-\allowbreak versions-max} (mặc định 100). Xoá file mồ côi do job hỏng để lại (\href{https://iceberg.apache.org/docs/latest/maintenance/}{Maintenance}, \href{https://iceberg.apache.org/docs/latest/configuration/}{Configuration}). Đừng chạy \code{remove\_orphan\_files} với ngưỡng tuổi ngắn hơn thời gian của job ghi dài nhất: file của một commit chưa xong trông giống file mồ côi.
  \item \textbf{Bảo trì tuỳ chọn.} Compaction (\code{rewrite\_data\_files}) gộp file nhỏ và áp delete file; \code{rewrite\_manifests} gom lại manifest để lập kế hoạch đọc nhanh hơn. Spark gọi chúng bằng \code{CALL <catalog>.system.<thủ tục>(...)} (\href{https://iceberg.apache.org/docs/latest/spark-procedures/}{Spark Procedures}); PyIceberg có \code{table.maintenance.\allowbreak expire\_snapshots()} (\href{https://py.iceberg.apache.org/api/}{PyIceberg API}).
  \item \textbf{Branch, tag và write-audit-publish.} Branch và tag là ref có tên trỏ tới snapshot, có thể đặt thời gian giữ riêng; nhánh \code{main} không bao giờ hết hạn (\href{https://iceberg.apache.org/docs/latest/branching/}{Branching and Tagging}, \href{https://iceberg.apache.org/spec/}{spec}). Write-audit-publish: ghi vào branch \code{audit}, kiểm tra, rồi \code{fast\_forward} \code{main} tới \code{audit} (\href{https://iceberg.apache.org/docs/latest/spark-procedures/}{Spark Procedures}).
  \item \textbf{Delta Lake và Apache Hudi.} Delta Lake commit bằng cách thêm file JSON có số thứ tự vào thư mục \code{\_delta\_log} và định kỳ ghi checkpoint (\href{https://github.com/delta-io/delta/blob/master/PROTOCOL.md}{Delta PROTOCOL.md}). Hudi có hai loại bảng copy-on-write và merge-on-read (\href{https://hudi.apache.org/docs/table_types/}{Hudi: Table \& Query Types}). Cả hai giải cùng bài toán commit trên object storage với cấu trúc metadata khác Iceberg.
\end{enumerate}

\section{Bắt đầu}
\begin{itemize}
  \item \href{https://iceberg.apache.org/spark-quickstart/}{Spark and Iceberg Quickstart}: chạy Docker Compose của trang, tạo bảng, insert, đọc metadata table.
  \item \href{https://py.iceberg.apache.org/api/}{PyIceberg: Python API}: \code{load\_catalog}, \code{create\_table}, \code{append}, \code{inspect}.
  \item \href{https://iceberg.apache.org/docs/latest/spark-getting-started/}{Iceberg: Spark Getting Started}: cấu hình catalog cho Spark.
\end{itemize}

\section{Tài liệu chính thức}
Đặc tả ở \code{iceberg.apache.org/spec/} không gắn phiên bản thư viện. Docs thư viện có phiên bản: \code{docs/latest/} hiện là 1.12.0; khi lab ghim phiên bản khác, đọc \code{docs/<phiên bản>/}.

\begin{xltabular}{\textwidth}{@{}P{4.2cm} L P{1.8cm}@{}}
\toprule
\thd{tài liệu} & \thd{nội dung cần đọc} & \thd{thời lượng} \\
\midrule
\endhead
\bottomrule
\endfoot
\href{https://iceberg.apache.org/spec/}{Iceberg Table Spec} & Overview, Format Versioning, Optimistic Concurrency, Sequence Numbers, Manifests, Manifest Lists, Table Metadata, Delete Formats, Commit Conflict Resolution, phụ lục Implementation Notes & 8 giờ \\
\href{https://iceberg.apache.org/docs/latest/reliability/}{Reliability} và \href{https://iceberg.apache.org/docs/latest/partitioning/}{Partitioning} & Hive thiếu gì; hidden partitioning & 1 giờ \\
\href{https://iceberg.apache.org/docs/latest/evolution/}{Evolution} & schema evolution, partition evolution, sort order & 1 giờ \\
\href{https://iceberg.apache.org/docs/latest/configuration/}{Configuration} & write properties, table behavior properties & 1 giờ \\
\href{https://iceberg.apache.org/docs/latest/maintenance/}{Maintenance} & recommended và optional maintenance & 1 giờ \\
\href{https://iceberg.apache.org/docs/latest/branching/}{Branching and Tagging} & ref, retention, WAP & 1 giờ \\
\href{https://iceberg.apache.org/docs/latest/spark-procedures/}{Spark Procedures} & \code{expire\_snapshots}, \code{rewrite\_data\_files}, \code{rewrite\_manifests}, \code{remove\_orphan\_files}, \code{rollback\_to\_snapshot}, \code{fast\_forward} & 1,5 giờ \\
\href{https://iceberg.apache.org/docs/latest/spark-queries/}{Spark Queries} & time travel, metadata table & 1 giờ \\
\href{https://iceberg.apache.org/rest-catalog-spec/}{REST Catalog Spec} và \href{https://github.com/apache/iceberg/blob/main/open-api/rest-catalog-open-api.yaml}{OpenAPI} & \code{updateTable}, \code{CommitTableRequest}, \code{TableRequirement}, \code{TableUpdate}, mã lỗi & 2 giờ \\
\href{https://py.iceberg.apache.org/api/}{PyIceberg API} & create, append, overwrite, delete, upsert, inspect, update\_schema, maintenance & 1,5 giờ \\
\href{https://github.com/delta-io/delta/blob/master/PROTOCOL.md}{Delta PROTOCOL.md} và \href{https://hudi.apache.org/docs/table_types/}{Hudi Table Types} & để so sánh: Delta log, checkpoint; COW và MOR của Hudi & 2 giờ \\
\end{xltabular}

\section{Khoá học miễn phí}
\begin{itemize}
  \item \href{https://15445.courses.cs.cmu.edu/fall2025/schedule.html}{CMU 15-445/645 Fall 2025}: miễn phí slide và video; bài \#17 Concurrency Control Theory và \#20 Multi-Version Concurrency Control giúp hiểu serializable isolation và snapshot isolation trước khi đọc mục Optimistic Concurrency của đặc tả.
\end{itemize}

\section{Lab thực hành}
Chuẩn bị: stack tham chiếu với RustFS, một REST catalog (\code{apache/iceberg-rest-fixture} để học, hoặc Polaris), Spark có Iceberg runtime và PyIceberg. Cấu hình Spark tối thiểu (thay host, cổng, bucket theo compose của bạn):
\begin{codeblock}
spark.sql.extensions            org.apache.iceberg.spark.extensions.IcebergSparkSessionExtensions
spark.sql.catalog.lab           org.apache.iceberg.spark.SparkCatalog
spark.sql.catalog.lab.type      rest
spark.sql.catalog.lab.uri       http://rest:8181
spark.sql.catalog.lab.warehouse s3://warehouse/
spark.sql.catalog.lab.io-impl   org.apache.iceberg.aws.s3.S3FileIO
spark.sql.catalog.lab.s3.endpoint          http://rustfs:9000
spark.sql.catalog.lab.s3.path-style-access true
\end{codeblock}

\subsection{Lab 5.1: Ghi Olist orders bằng PyIceberg và đọc cây metadata}
\begin{enumerate}
  \item Nạp orders với kiểu timestamp rõ ràng và ghi thành hai snapshot:
\begin{codeblock}
import pyarrow as pa, pyarrow.csv as pv, pyarrow.compute as pc
from pyiceberg.catalog import load_catalog

ts = ["order_purchase_timestamp", "order_approved_at",
      "order_delivered_carrier_date", "order_delivered_customer_date",
      "order_estimated_delivery_date"]
opts = pv.ConvertOptions(column_types={c: pa.timestamp("us") for c in ts})
orders = pv.read_csv("olist/olist_orders_dataset.csv", convert_options=opts)

cat = load_catalog("lab", **{
    "type": "rest", "uri": "http://localhost:8181", "warehouse": "s3://warehouse/",
    "s3.endpoint": "http://localhost:9000",
    "s3.access-key-id": "...", "s3.secret-access-key": "..."})
cat.create_namespace("olist")
tbl = cat.create_table("olist.orders", schema=orders.schema)

year = pc.year(orders["order_purchase_timestamp"])
tbl.append(orders.filter(pc.less(year, 2018)))
tbl.append(orders.filter(pc.greater_equal(year, 2018)))
print(tbl.metadata_location)
print(tbl.inspect.snapshots().to_pandas()[["snapshot_id", "parent_id", "operation", "manifest_list"]])
\end{codeblock}
  \item Tải file metadata hiện tại về và đọc bằng \code{jq}:
\begin{codeblock}
aws --endpoint-url http://localhost:9000 s3 cp <metadata_location> - \
  | jq '."current-snapshot-id", (.snapshots[] | ."manifest-list"), ."snapshot-log"'
\end{codeblock}
  \item Liệt kê toàn bộ object của bảng trên RustFS (\code{s3 ls --recursive}). Gán mỗi object vào một tầng của cây: \code{metadata.json}, manifest list, manifest, data file.
  \item Dùng \code{tbl.inspect.manifests()} và \code{tbl.inspect.entries()} để xem manifest nào thuộc snapshot nào và các bound của cột \code{order\_id} (chuỗi 32 ký tự bị cắt còn 16).
\end{enumerate}
\hoanthanh{bạn vẽ được cây của bảng sau hai lần append, với đúng tên file thật ở mỗi tầng, và chỉ ra file nào do catalog server ghi, file nào do PyIceberg ghi.}

\subsection{Lab 5.2: Hidden partitioning, time travel, schema evolution (Spark)}
\begin{enumerate}
  \item Tạo bảng chia partition theo tháng và nạp từ bảng ở Lab 5.1:
\begin{codeblock}
CREATE TABLE lab.olist.orders_p (
  order_id string, customer_id string, order_status string,
  order_purchase_timestamp timestamp, order_delivered_customer_date timestamp)
USING iceberg PARTITIONED BY (months(order_purchase_timestamp));

INSERT INTO lab.olist.orders_p
SELECT order_id, customer_id, order_status, order_purchase_timestamp,
       order_delivered_customer_date FROM lab.olist.orders;

SELECT partition, record_count, file_path FROM lab.olist.orders_p.files;
\end{codeblock}
  \item Chạy truy vấn lọc trên cột gốc \code{order\_purchase\_timestamp} cho tháng 11/2017 và xem số file được quét trong Spark UI (hoặc plan). Không có cột partition nào xuất hiện trong câu lệnh.
  \item Ghi lại \code{snapshot\_id} hiện tại từ \code{.snapshots}. Đổi schema và partition spec:
\begin{codeblock}
ALTER TABLE lab.olist.orders_p ADD COLUMN delivery_days int;
ALTER TABLE lab.olist.orders_p RENAME COLUMN order_status TO status;
ALTER TABLE lab.olist.orders_p
  REPLACE PARTITION FIELD order_purchase_timestamp_month WITH days(order_purchase_timestamp);
INSERT INTO lab.olist.orders_p SELECT ... ;   -- một ít dòng mới
SELECT spec_id, partition, file_path FROM lab.olist.orders_p.files;
SELECT count(*) FROM lab.olist.orders_p VERSION AS OF <snapshot_id cũ>;
\end{codeblock}
\end{enumerate}
\hoanthanh{bạn cho thấy file cũ và file mới có \code{spec\_id} khác nhau trong cùng một bảng, truy vấn time travel trả đúng số dòng cũ, và không data file cũ nào bị ghi lại khi đổi tên cột (so danh sách \code{file\_path} trước và sau).}

\subsection{Lab 5.3: Copy-on-write so với merge-on-read}
\begin{enumerate}
  \item Tạo hai bản sao:
\begin{codeblock}
CREATE TABLE lab.olist.orders_cow USING iceberg
  TBLPROPERTIES ('format-version'='2') AS SELECT * FROM lab.olist.orders;
CREATE TABLE lab.olist.orders_mor USING iceberg
  TBLPROPERTIES ('format-version'='2', 'write.delete.mode'='merge-on-read')
  AS SELECT * FROM lab.olist.orders;
DELETE FROM lab.olist.orders_cow WHERE order_status = 'canceled';
DELETE FROM lab.olist.orders_mor WHERE order_status = 'canceled';
\end{codeblock}
  \item So \code{SELECT operation, summary FROM <bảng>.snapshots} của hai bảng (số file thêm, file bỏ, delete file). Xem \code{SELECT content, file\_format, record\_count, file\_path FROM lab.olist.orders\_mor.files}: \code{content} bằng 0 là data file, 1 là position delete, 2 là equality delete.
  \item Tuỳ chọn, nếu phiên bản Spark và Iceberg của bạn hỗ trợ v3: lặp lại với \code{'format-version'=\allowbreak '3'} và tìm delete file định dạng Puffin (deletion vector).
  \item Chạy \code{CALL lab.system.rewrite\_data\_files(table => 'olist.orders\_mor')} và xem delete file còn được tham chiếu trong snapshot mới không.
\end{enumerate}
\hoanthanh{bạn có bảng so sánh hai chế độ (số byte ghi khi DELETE, số file mà lần đọc sau phải mở) và giải thích được vì sao merge-on-read cần compaction.}

\subsection{Lab 5.4: Bảo trì và write-audit-publish}
\begin{enumerate}
  \item Tạo bảng nhiều file nhỏ: bằng PyIceberg, append 30 lần, mỗi lần 1.000 dòng orders vào \code{olist.orders\_small}. Đếm \code{SELECT count(*) FROM lab.olist.orders\_small.files} và \code{.snapshots}.
  \item Chạy lần lượt và ghi lại số file, số manifest, số snapshot sau mỗi bước:
\begin{codeblock}
CALL lab.system.rewrite_data_files(table => 'olist.orders_small');
CALL lab.system.rewrite_manifests('olist.orders_small');
CALL lab.system.expire_snapshots(table => 'olist.orders_small',
     older_than => TIMESTAMP '<thời điểm hiện tại>', retain_last => 1);
CALL lab.system.remove_orphan_files(table => 'olist.orders_small', dry_run => true);
\end{codeblock}
  \item Write-audit-publish trên \code{orders\_p}:
\begin{codeblock}
ALTER TABLE lab.olist.orders_p CREATE TAG before_wap;
ALTER TABLE lab.olist.orders_p CREATE BRANCH audit;
INSERT INTO lab.olist.orders_p.branch_audit SELECT ... ;   -- dữ liệu mới
SELECT count(*) FROM lab.olist.orders_p.branch_audit WHERE order_id IS NULL;
SELECT count(*) FROM lab.olist.orders_p;                    -- main chưa đổi
CALL lab.system.fast_forward('olist.orders_p', 'main', 'audit');
\end{codeblock}
\end{enumerate}
\hoanthanh{bạn chứng minh bằng metadata table rằng compaction giảm số file mà không đổi số dòng, expire làm mất khả năng time travel về snapshot cũ, và \code{main} chỉ thấy dữ liệu mới sau \code{fast\_forward}.}

\section{Kế hoạch theo tuần}
\begin{xltabular}{\textwidth}{@{}P{1.4cm} L P{1.4cm}@{}}
\toprule
\thd{tuần} & \thd{việc} & \thd{giờ} \\
\midrule
\endhead
\bottomrule
\endfoot
1 & Reliability, Partitioning; spec: Overview, Manifests, Manifest Lists, Table Metadata; Spark Quickstart; Lab 5.1 & 10 \\
2 & spec: Optimistic Concurrency, Sequence Numbers, Commit Conflict Resolution; REST Catalog Spec \code{updateTable}; CMU bài \#17, \#20; Lab 5.2 & 10 \\
3 & Configuration, Evolution, spec Delete Formats; Maintenance, Spark Procedures, Branching; Lab 5.3, 5.4 & 10 \\
4 (nửa tuần) & Delta PROTOCOL.md, Hudi Table Types; thử thách kỹ năng; câu hỏi tự kiểm tra & 5 \\
thêm & spec Format Versioning chi tiết v3; Puffin spec; PyIceberg \code{upsert} và \code{overwrite} với filter & 6 \\
\end{xltabular}

\begin{thuthach}
Bảng \code{olist.orders} cần cập nhật trạng thái đơn hằng ngày. Hãy giả lập ba ngày thay đổi (đổi \code{order\_status}, thêm đơn mới), rồi xây quy trình: ghi thay đổi vào một branch, kiểm tra chất lượng, publish lên \code{main}, và chạy bảo trì để bảng không tích file nhỏ hay snapshot thừa. Thêm một tình huống hai writer cùng sửa các dòng chồng nhau và quan sát bên nào thất bại. Không có hướng dẫn từng bước; chọn Spark hoặc PyIceberg (\code{upsert} theo identifier field) tuỳ bạn.
\hoanthanh{bạn nộp được một bảng lấy từ metadata table cho mỗi ngày (snapshot, operation, số data file, số delete file), một lần commit thất bại kèm thông báo lỗi và giải thích theo mục Optimistic Concurrency của đặc tả, và trạng thái cuối có số file không vượt ngưỡng bạn tự đặt.}
\end{thuthach}

\section{Câu hỏi tự kiểm tra}
\begin{enumerate}
  \item Vẽ cây metadata từ con trỏ trong catalog tới data file; file nào là JSON, file nào là Avro?
  \item Trong một lần append bằng Spark qua REST catalog, liệt kê theo thứ tự các file được ghi, ai ghi, và thao tác nào làm commit hiện ra với người đọc.
  \item Vì sao retry chỉ cần ghi lại manifest list mà dùng lại được manifest?
  \item Hai writer cùng append: ai thắng, bên thua làm gì? Hai writer xoá các dòng chồng nhau thì sao?
  \item REST catalog trả 409 và trả 504 khác nhau thế nào, và client phải làm gì với các file đã ghi trong mỗi trường hợp?
  \item \code{snapshot-log} khác chuỗi cha của snapshot ở đâu?
  \item Hidden partitioning sửa được lỗi gì của cột ngày kiểu Hive mà người dùng phải tự lọc?
  \item Vì sao đổi tên cột không cần ghi lại dữ liệu?
  \item Khi nào chọn merge-on-read, và khi đó cần thêm việc bảo trì nào?
  \item Deletion vector của v3 khác position delete file của v2 ở đâu?
  \item Chuyện gì có thể xảy ra nếu \code{remove\_orphan\_files} chạy với ngưỡng tuổi rất ngắn trong lúc một job đang ghi?
  \item Kể bốn bước write-audit-publish bằng branch, và so cách commit của Delta Lake với Iceberg.
\end{enumerate}

\section{Điểm dễ nhầm}
\begin{xltabular}{\textwidth}{@{}P{4.2cm} L P{3.2cm}@{}}
\toprule
\thd{điểm dễ nhầm} & \thd{thực tế} & \thd{docs chính thức} \\
\midrule
\endhead
\bottomrule
\endfoot
Nhiều người nghĩ Iceberg là định dạng file như Parquet & Iceberg là table format quản lý data file Parquet, Avro hoặc ORC & \href{https://iceberg.apache.org/spec/}{Iceberg spec} \\
``Iceberg cần rename nguyên tử'' & Chỉ bảng kiểu file system dùng rename; đặc tả ghi cách này không an toàn trên object store & \href{https://iceberg.apache.org/spec/}{Iceberg spec} \\
``Catalog lưu metadata của bảng'' & Catalog giữ con trỏ tới \code{metadata.json} hiện tại; metadata và data nằm trên storage & \href{https://iceberg.apache.org/spec/}{Iceberg spec} \\
``Với REST catalog, Spark ghi \code{metadata.json}'' & Client gửi \code{requirements} và \code{updates}; server áp update, ghi metadata mới và đổi con trỏ & \href{https://iceberg.apache.org/rest-catalog-spec/}{REST Catalog Spec} \\
``REST trả 5xx nghĩa là commit hỏng, dọn file đi'' & 500, 502, 504 là trạng thái chưa biết; chỉ 409 là commit thất bại & \href{https://py.iceberg.apache.org/reference/pyiceberg/catalog/rest/}{PyIceberg rest} \\
``Iceberg mặc định snapshot isolation'' & Delete, update, merge mặc định \code{serializable} & \href{https://iceberg.apache.org/docs/latest/configuration/}{Configuration} \\
``Merge-on-read là mặc định và luôn nhanh hơn'' & Mặc định là copy-on-write; MOR đổi chi phí ghi lấy chi phí đọc & \href{https://iceberg.apache.org/docs/latest/configuration/}{Configuration} \\
``Hidden partitioning nghĩa là bảng không chia partition'' & Partition được suy ra bằng transform từ cột gốc & \href{https://iceberg.apache.org/docs/latest/partitioning/}{Partitioning} \\
``Time travel giữ mãi'' & Expire snapshot là bảo trì được khuyên; mặc định giữ 5 ngày và ít nhất 1 snapshot & \href{https://iceberg.apache.org/docs/latest/maintenance/}{Maintenance}, \href{https://iceberg.apache.org/docs/latest/configuration/}{Configuration} \\
``Min/max trong manifest bằng min/max trong footer Parquet'' & Manifest giữ bound theo metrics mode, mặc định \code{truncate(16)} & \href{https://iceberg.apache.org/docs/latest/configuration/}{Configuration} \\
``Mọi engine đã làm đủ v3, v4'' & v4 chưa được thông qua; engine có thể chưa làm đủ v3 & \href{https://iceberg.apache.org/spec/}{Iceberg spec} \\
\end{xltabular}


% ===================================================================

\modulechapter{6}{Catalog, metadata và governance}
\moduleinfo{Trung cấp}{15 giờ}{Module 5}

Module 5 cho thấy catalog Iceberg chỉ giữ một con trỏ, nhưng con trỏ đó quyết định ai thấy phiên bản nào của bảng. Module này hỏi tiếp: ai được đổi con trỏ, ai được đọc file trên storage, dữ liệu đi từ đâu tới đâu (lineage), và khi một người yêu cầu xoá dữ liệu cá nhân thì dữ liệu thật sự biến mất lúc nào. Bạn sẽ gọi REST catalog API bằng \code{curl}, dựng RBAC trên Apache Polaris, phát sự kiện OpenLineage vào Marquez, và diễn tập xoá dữ liệu trên bảng Iceberg.

\section{Bạn sẽ học được}
\begin{itemize}
  \item Phân biệt technical catalog (giữ con trỏ bảng, cho engine commit) với data catalog (để tìm, mô tả và quản trị dữ liệu).
  \item Gọi các endpoint chính của Iceberg REST catalog và đọc \code{metadata-location} của một bảng.
  \item Giải thích mô hình RBAC của Polaris và dựng một principal chỉ có quyền đọc.
  \item So sánh ba cách cấp quyền đọc file: key tĩnh, vended credential, remote signing.
  \item Mô tả row filter, column mask và giới hạn của RBAC ở mức catalog.
  \item Phát sự kiện OpenLineage cho một job và xem đồ thị lineage trong Marquez.
  \item Viết quy trình xoá dữ liệu cá nhân trên Iceberg đủ để dữ liệu biến mất khỏi storage.
\end{itemize}

\section{Nội dung chính}
\begin{enumerate}
  \item \textbf{Hai nghĩa của ``catalog''.} Technical catalog của Iceberg giữ con trỏ tới file metadata hiện tại và đổi nó nguyên tử khi commit (\href{https://iceberg.apache.org/spec/}{Iceberg spec}). Data catalog hay nền tảng metadata giúp người dùng tìm bảng, đọc mô tả, chủ sở hữu, lineage; nó không thay technical catalog trong đường commit.
  \item \textbf{Iceberg REST catalog.} Đặc tả OpenAPI định nghĩa: \code{GET /v1/config} (server trả cấu hình mặc định, cấu hình ghi đè như \code{prefix}, và danh sách endpoint), namespace, table (load, create, \code{updateTable}), commit nhiều bảng qua \code{/v1/\{prefix\}/\allowbreak transactions/commit}, và lập kế hoạch scan phía server (\href{https://iceberg.apache.org/rest-catalog-spec/}{REST Catalog Spec}, \href{https://github.com/apache/iceberg/blob/main/open-api/rest-catalog-open-api.yaml}{OpenAPI}). Load table trả \code{metadata-location}, metadata và \code{config} cho client.
  \item \textbf{Xác thực.} Server thường dùng OAuth2. Endpoint \code{/v1/oauth/tokens} trong đặc tả đã bị đánh dấu deprecated để gỡ bỏ; server có thể dùng identity provider bên ngoài. Polaris vẫn cấp token cho client qua \code{client\_id}, \code{client\_secret} và \code{scope = PRINCIPAL\_ROLE:ALL} (\href{https://polaris.apache.org/releases/1.1.0/getting-started/using-polaris/}{Polaris: Using Polaris}).
  \item \textbf{Cấp quyền đọc file.} Có ba cách: (a) client giữ key S3 tĩnh, đi vòng được qua catalog; (b) vended credential: client gửi \code{X-Iceberg-Access-\allowbreak Delegation: vended-credentials} và catalog trả credential ngắn hạn theo quyền trên bảng; (c) remote signing: catalog ký từng request S3. Lakekeeper ghi rằng remote signing chỉ giữ client trong phạm vi bảng khi object store kiểm mọi phần của request đã ký, và khuyên vended credential cho một số storage như MinIO, SeaweedFS (\href{https://docs.lakekeeper.io/docs/nightly/storage/}{Lakekeeper: Storage}).
  \item \textbf{Các catalog triển khai REST.}
  \begin{itemize}
    \item Apache Polaris: top-level project của Apache từ 02/2026. Quickstart chạy Polaris với metastore trong bộ nhớ và RustFS, không dành cho production (\href{https://polaris.apache.org/guides/quickstart/}{Polaris Quickstart}).
    \item Lakekeeper: chọn authorizer OpenFGA hoặc Cedar (\href{https://docs.lakekeeper.io/docs/nightly/authorization/}{Lakekeeper: Authorization}).
    \item AWS Glue Data Catalog: Iceberg REST endpoint dạng \code{https://glue.\allowbreak <region>.\allowbreak amazonaws.com/\allowbreak iceberg}, xác thực SigV4 với signing name \code{glue}, tham số \code{warehouse} là catalog ID (\href{https://docs.aws.amazon.com/glue/latest/dg/connect-glu-iceberg-rest.html}{AWS Glue: Iceberg REST endpoint}).
    \item \code{apache/iceberg-rest-fixture}: chỉ để học và test.
  \end{itemize}
  \item \textbf{RBAC của Polaris.} Principal (người hoặc dịch vụ) được gán principal role; principal role được gán catalog role; privilege (ví dụ \code{CATALOG\_MANAGE\_CONTENT}, \code{TABLE\_READ\_DATA}, \code{TABLE\_WRITE\_DATA}) cấp cho catalog role trên một securable (catalog, namespace, table). Không cấp privilege thẳng cho principal role. \code{TABLE\_READ\_DATA} cho phép nhận credential đọc ngắn hạn. Polaris dặn không để password, token hay access key trong property của bảng hoặc view, vì ai đọc được property sẽ thấy (\href{https://polaris.apache.org/in-dev/unreleased/managing-security/access-control/}{Polaris: Role-Based Access Control}).
  \item \textbf{Row filter và column mask.} Đây là kiểm soát ở mức dòng và cột, do engine hoặc nền tảng truy vấn thực thi. Ví dụ, Unity Catalog trên Databricks dùng SQL UDF: row filter trả boolean cho từng dòng lúc truy vấn, column mask trả giá trị gốc hoặc giá trị đã che; gắn bằng \code{ALTER TABLE ... SET ROW FILTER} và \code{ALTER COLUMN ... SET MASK} (\href{https://docs.databricks.com/aws/en/data-governance/unity-catalog/filters-and-masks/}{Databricks: Row filters and column masks}). Lakekeeper ghi rằng OpenFGA và Cedar không biểu diễn được row filter hay column mask (\href{https://docs.lakekeeper.io/docs/nightly/authorization/}{Lakekeeper: Authorization}).
  \item \textbf{Lineage với OpenLineage.} Mô hình gồm run, job, dataset và facet. Một run đi qua các trạng thái như \code{START}, \code{RUNNING}, \code{COMPLETE}, \code{FAIL}, \code{ABORT}; các sự kiện của cùng một run nối với nhau bằng \code{runId} (\href{https://openlineage.io/docs/spec/run-cycle/}{OpenLineage: Run Cycle}, \href{https://openlineage.io/getting-started/}{Getting Started}). Lineage cần bên phát (Spark listener, Airflow provider) và bên nhận như Marquez (\href{https://marquezproject.ai/docs/quickstart/}{Marquez Quickstart}).
  \item \textbf{Xoá dữ liệu cá nhân trên Iceberg.} \code{DELETE} tạo snapshot mới, nhưng snapshot cũ, tag và branch vẫn trỏ tới file cũ. Muốn dữ liệu biến mất khỏi storage: (1) với merge-on-read, chạy \code{rewrite\_data\_files} để file gốc không còn được dùng; (2) xoá hoặc rút ngắn tuổi các tag, branch đang giữ snapshot cũ (nhánh \code{main} không hết hạn, các ref khác theo \code{max-ref-age-ms}); (3) \code{expire\_snapshots}; (4) \code{remove\_orphan\_files} cho file của job hỏng (\href{https://iceberg.apache.org/docs/latest/branching/}{Branching}, \href{https://iceberg.apache.org/docs/latest/spark-procedures/}{Spark Procedures}, \href{https://iceberg.apache.org/docs/latest/maintenance/}{Maintenance}). Lưu ý bound của cột trong manifest có thể chứa tiền tố của giá trị cá nhân (mặc định \code{truncate(16)}); với cột nhạy cảm, đặt metrics mode \code{none} qua \code{write.metadata.\allowbreak metrics.column.<cột>} (\href{https://iceberg.apache.org/docs/latest/configuration/}{Configuration}).
\end{enumerate}

\begin{ghichu}
\textbf{Ghi chú pháp lý ở Việt Nam.} Luật Bảo vệ dữ liệu cá nhân số 91/2025/QH15 có hiệu lực từ 01/01/2026 và cho chủ thể dữ liệu quyền yêu cầu xoá (điều 4). Nghị định 356/2025/NĐ-CP hướng dẫn chi tiết, cũng có hiệu lực từ 01/01/2026; điều 5 khoản 4 yêu cầu phản hồi yêu cầu xoá trong 02 ngày làm việc và thực hiện trong 20 ngày (30 ngày nếu cần bên xử lý hoặc bên thứ ba, gia hạn một lần không quá 20 ngày). Xem văn bản trên \href{https://english.luatvietnam.vn/dan-su/law-on-personal-data-protection-law-no-91-2025-qh15-405135-d1.html}{LuatVietnam: Luật 91/2025/QH15} và \href{https://luatvietnam.vn/dan-su/nghi-dinh-356-2025-nd-cp-quy-dinh-chi-tiet-luat-bao-ve-du-lieu-ca-nhan-422896-d1.html}{LuatVietnam: Nghị định 356/2025}; đây là cơ sở dữ liệu pháp luật, khi cần trích dẫn chính thức hãy đối chiếu văn bản gốc trên Công báo. Quy trình xoá ở mục trên phải chạy xong trong các hạn này.
\end{ghichu}

\section{Bắt đầu}
\begin{itemize}
  \item \href{https://polaris.apache.org/guides/quickstart/}{Polaris Quickstart}: chạy Polaris cùng RustFS bằng Docker Compose, lấy credential từ \code{docker compose -p polaris-quickstart logs}.
  \item \href{https://openlineage.io/getting-started/}{OpenLineage Getting Started}: dựng Marquez, gửi sự kiện \code{START} và \code{COMPLETE}, xem đồ thị.
  \item \href{https://iceberg.apache.org/rest-catalog-spec/}{Iceberg REST Catalog Spec}: đọc danh sách endpoint trước khi gọi bằng \code{curl}.
\end{itemize}

\section{Tài liệu chính thức}
\begin{xltabular}{\textwidth}{@{}P{4.2cm} L P{1.8cm}@{}}
\toprule
\thd{tài liệu} & \thd{nội dung cần đọc} & \thd{thời lượng} \\
\midrule
\endhead
\bottomrule
\endfoot
\href{https://iceberg.apache.org/rest-catalog-spec/}{REST Catalog Spec} và \href{https://github.com/apache/iceberg/blob/main/open-api/rest-catalog-open-api.yaml}{OpenAPI} & \code{/v1/config}, namespace, table, \code{LoadTableResult}, \code{/transactions/commit}, header \code{X-Iceberg-Access-Delegation} & 2 giờ \\
\href{https://polaris.apache.org/guides/quickstart/}{Polaris Quickstart} & dựng Polaris với RustFS & 1 giờ \\
\href{https://polaris.apache.org/in-dev/unreleased/managing-security/access-control/}{Polaris: Role-Based Access Control} & RBAC model, privileges, ví dụ & 1,5 giờ \\
\href{https://polaris.apache.org/releases/1.1.0/getting-started/using-polaris/}{Polaris: Using Polaris} và \href{https://polaris.apache.org/releases/1.4.1/command-line-interface/}{CLI} & tạo catalog, principal, role, grant; kết nối Spark và PyIceberg & 1 giờ \\
\href{https://docs.lakekeeper.io/docs/nightly/authorization/}{Lakekeeper: Authorization} và \href{https://docs.lakekeeper.io/docs/nightly/storage/}{Storage} & authorizer, grant; vended credential và remote signing & 1,5 giờ \\
\href{https://docs.aws.amazon.com/glue/latest/dg/connect-glu-iceberg-rest.html}{AWS Glue: Iceberg REST endpoint} & endpoint, SigV4, ánh xạ \code{prefix} & 0,5 giờ \\
\href{https://docs.databricks.com/aws/en/data-governance/unity-catalog/filters-and-masks/}{Databricks: Row filters and column masks} & khái niệm row filter, column mask, ABAC & 0,5 giờ \\
\href{https://openlineage.io/getting-started/}{OpenLineage Getting Started} và \href{https://openlineage.io/docs/spec/run-cycle/}{Run Cycle} & run, job, dataset, trạng thái run & 1,5 giờ \\
\href{https://marquezproject.ai/docs/quickstart/}{Marquez Quickstart} & chạy Marquez, API lineage, UI & 0,5 giờ \\
\href{https://iceberg.apache.org/docs/latest/branching/}{Iceberg Branching} và \href{https://iceberg.apache.org/docs/latest/spark-procedures/}{Spark Procedures} & retention của ref; \code{expire\_snapshots}, \code{remove\_orphan\_files} & 1 giờ \\
\end{xltabular}

\section{Lab thực hành}
Chuẩn bị: Polaris Quickstart đang chạy (nó có sẵn catalog \code{quickstart\_catalog} trên RustFS). Nạp bảng \code{olist.orders} và \code{olist.customers} vào catalog này bằng mã của Lab 5.1, với \code{uri = http://localhost:8181/api/catalog}, \code{warehouse = quickstart\_catalog}, \code{credential = <client\_id>:<client\_secret>}, \code{scope = PRINCIPAL\_ROLE:ALL}.

\subsection{Lab 6.1: Gọi REST catalog bằng curl}
\begin{enumerate}
  \item Lấy token và cấu hình:
\begin{codeblock}
BASE=http://localhost:8181/api/catalog
TOKEN=$(curl -s $BASE/v1/oauth/tokens \
  -d grant_type=client_credentials -d client_id=$CLIENT_ID \
  -d client_secret=$CLIENT_SECRET -d scope=PRINCIPAL_ROLE:ALL | jq -r .access_token)
curl -s -H "Authorization: Bearer $TOKEN" "$BASE/v1/config?warehouse=quickstart_catalog" | jq
\end{codeblock}
  \item Đọc giá trị \code{prefix} trong \code{overrides}, rồi gọi tiếp:
\begin{codeblock}
curl -s -H "Authorization: Bearer $TOKEN" $BASE/v1/<prefix>/namespaces | jq
curl -s -H "Authorization: Bearer $TOKEN" $BASE/v1/<prefix>/namespaces/olist/tables | jq
curl -s -H "Authorization: Bearer $TOKEN" \
  $BASE/v1/<prefix>/namespaces/olist/tables/orders | jq '."metadata-location", .config'
curl -s -H "Authorization: Bearer $TOKEN" \
  -H "X-Iceberg-Access-Delegation: vended-credentials" \
  $BASE/v1/<prefix>/namespaces/olist/tables/orders | jq '.config, ."storage-credentials"'
\end{codeblock}
  \item So hai lần load table: lần nào có credential trong phản hồi? Mở console RustFS (cổng 9001) và tìm đúng file \code{metadata-location}.
\end{enumerate}
\hoanthanh{bạn có sơ đồ mọi đường giao tiếp trong lab (client tới Polaris, Polaris tới RustFS, client tới RustFS), mỗi đường ghi giao thức và loại credential, và chỉ ra được credential ngắn hạn trong phản hồi có header \code{vended-credentials}.}

\subsection{Lab 6.2: Principal chỉ đọc trên Polaris}
\begin{enumerate}
  \item Dùng CLI \code{polaris} với credential gốc (\code{--client-id}, \code{--client-secret}):
\begin{codeblock}
polaris principals create olist_reader          # ghi lại client id/secret trả về
polaris principal-roles create analyst
polaris principal-roles grant --principal olist_reader analyst
polaris catalog-roles create --catalog quickstart_catalog reader
polaris catalog-roles grant --catalog quickstart_catalog --principal-role analyst reader
polaris privileges catalog grant --catalog quickstart_catalog --catalog-role reader TABLE_READ_DATA
\end{codeblock}
  \item Kết nối PyIceberg bằng credential của \code{olist\_reader}, với header \code{X-Iceberg-Access-\allowbreak Delegation: vended-credentials}. Đọc \code{olist.orders} rồi thử \code{append} một dòng. Nếu không list được namespace, đọc trang RBAC để chọn thêm privilege cần thiết và ghi lại.
  \item Thử cấp privilege thẳng cho principal role và ghi lại vì sao không làm được.
\end{enumerate}
\hoanthanh{\code{olist\_reader} đọc được bảng, lần \code{append} bị từ chối với lỗi phân quyền, và bạn vẽ được chuỗi principal $\rightarrow$ principal role $\rightarrow$ catalog role $\rightarrow$ privilege của lab.}

\subsection{Lab 6.3: Lineage với OpenLineage và Marquez}
\begin{enumerate}
  \item Chạy Marquez theo Getting Started (\code{./docker/up.sh}); nếu cổng API trùng với dịch vụ khác của stack, đổi bằng \code{--api-port}. UI ở cổng 3000.
  \item Gửi hai sự kiện cho job nạp orders, cùng một \code{runId} (sinh bằng \code{uuidgen}):
\begin{codeblock}
{"eventType": "START", "eventTime": "2026-11-02T10:00:00Z",
 "run": {"runId": "<uuid>"},
 "job": {"namespace": "olist-lab", "name": "load_orders"},
 "inputs":  [{"namespace": "s3://lab", "name": "raw/olist_orders_dataset.csv"}],
 "outputs": [{"namespace": "iceberg://quickstart_catalog", "name": "olist.orders"}],
 "producer": "https://example.com/olist-lab",
 "schemaURL": "https://openlineage.io/spec/2-0-2/OpenLineage.json#/definitions/RunEvent"}
\end{codeblock}
  Lưu thành \code{start.json} và gửi, rồi gửi bản \code{COMPLETE} với \code{eventTime} muộn hơn:
\begin{codeblock}
curl -X POST http://localhost:<api-port>/api/v1/lineage \
  -H 'Content-Type: application/json' -d @start.json
\end{codeblock}
  \item Thêm job thứ hai \code{build\_customer\_orders} đọc \code{olist.orders} và \code{olist.customers}, ghi \code{olist.customer\_orders}. Xem đồ thị trong UI.
  \item Tuỳ chọn: bật Spark listener theo \href{https://openlineage.io/docs/integrations/spark/configuration/usage}{OpenLineage Spark integration} và chạy lại một truy vấn của Lab 5.2.
\end{enumerate}
\hoanthanh{Marquez hiển thị đồ thị CSV $\rightarrow$ \code{olist.orders} $\rightarrow$ \code{olist.customer\_orders}, và run của \code{load\_orders} có trạng thái hoàn thành với hai sự kiện cùng \code{runId}.}

\subsection{Lab 6.4: Diễn tập xoá dữ liệu cá nhân}
\begin{enumerate}
  \item Chọn một \code{customer\_id}. Tạo tag \code{before\_erasure} trên \code{olist.orders}, rồi \code{DELETE FROM ... WHERE customer\_id = '...'}.
  \item Chứng minh dữ liệu vẫn đọc được: truy vấn \code{VERSION AS OF} snapshot cũ và qua tag.
  \item Xoá tag (\code{ALTER TABLE ... DROP TAG before\_\allowbreak erasure}), chạy \code{rewrite\_data\_\allowbreak files} nếu bảng dùng merge-on-read, rồi \code{expire\_\allowbreak snapshots} với tham số \code{retain\_last => 1}, sau cùng \code{remove\_orphan\_\allowbreak files} (chạy \code{dry\_run} trước).
  \item Liệt kê lại object của bảng trên RustFS, kiểm rằng data file cũ chứa dòng đó đã mất và time travel về snapshot cũ báo lỗi.
\end{enumerate}
\hoanthanh{bạn có một runbook xoá dữ liệu gồm các bước, lệnh và bằng chứng (danh sách file trước và sau), và nêu được bước nào phải thêm nếu bảng có branch khác hoặc cột chứa dữ liệu cá nhân có metrics trong manifest.}

\section{Kế hoạch theo tuần}
\begin{xltabular}{\textwidth}{@{}P{1.4cm} L P{1.4cm}@{}}
\toprule
\thd{tuần} & \thd{việc} & \thd{giờ} \\
\midrule
\endhead
\bottomrule
\endfoot
1 & REST Catalog Spec; Polaris Quickstart, RBAC, Using Polaris; Lakekeeper Storage và Authorization; Lab 6.1, 6.2 & 10 \\
2 (nửa tuần) & OpenLineage Getting Started, Run Cycle; Marquez; Lab 6.3, 6.4; thử thách kỹ năng; câu hỏi tự kiểm tra & 5 \\
thêm & AWS Glue Iceberg REST endpoint; row filter và column mask; OpenLineage cho Airflow (\href{https://airflow.apache.org/docs/apache-airflow-providers-openlineage/stable/configurations-ref.html}{provider configuration}) & 5 \\
\end{xltabular}

\begin{thuthach}
Một nhóm có ba vai trò: analyst chỉ đọc bảng marts, data engineer ghi bảng staging và marts, và một dịch vụ báo cáo chỉ đọc một bảng duy nhất. Hãy dựng mô hình quyền này trên Polaris cho các bảng Olist, cấp quyền đọc file bằng vended credential (không phát key RustFS tĩnh cho ai ngoài admin), và phát lineage cho pipeline staging $\rightarrow$ marts. Không có hướng dẫn từng bước.
\hoanthanh{bạn trình bày được bảng principal, role, privilege; mỗi vai trò thử một thao tác được phép và một thao tác bị chặn với log lỗi; đồ thị lineage trong Marquez khớp với pipeline; và bạn chỉ ra được điều gì mô hình này vẫn không làm được (ví dụ che cột) cùng công cụ có thể bù.}
\end{thuthach}

\section{Câu hỏi tự kiểm tra}
\begin{enumerate}
  \item Catalog Iceberg lưu chính xác cái gì, và không lưu cái gì?
  \item Technical catalog và data catalog khác nhau thế nào? Công cụ nào nằm trên đường commit?
  \item Client gọi endpoint REST nào đầu tiên, và server trả gì ở đó để đổi hành vi client?
  \item Trong Polaris, vì sao không cấp \code{TABLE\_READ\_DATA} thẳng cho principal role? Vẽ chuỗi cấp quyền.
  \item Key S3 tĩnh dùng chung gây rủi ro gì mà vended credential gỡ được? Rủi ro nào còn lại?
  \item Theo Lakekeeper, khi nào remote signing không giữ được client trong phạm vi bảng?
  \item Row filter và column mask do thành phần nào thực thi, và vì sao RBAC của catalog không đủ?
  \item Trong OpenLineage, cái gì nối sự kiện \code{COMPLETE} với sự kiện \code{START}? Lineage tự có hay cần bên phát?
  \item Sau \code{DELETE} một khách hàng, liệt kê mọi nơi trong bảng Iceberg còn có thể giữ dữ liệu của họ.
  \item Vì sao tag tạo trước khi xoá có thể làm quy trình xoá thất bại, và bạn kiểm điều đó thế nào?
\end{enumerate}

\section{Điểm dễ nhầm}
\begin{xltabular}{\textwidth}{@{}P{4.2cm} L P{3.2cm}@{}}
\toprule
\thd{điểm dễ nhầm} & \thd{thực tế} & \thd{docs chính thức} \\
\midrule
\endhead
\bottomrule
\endfoot
Nhiều người nghĩ catalog giữ data hoặc file schema & Catalog giữ con trỏ tới file metadata hiện tại; metadata và data nằm trên storage & \href{https://iceberg.apache.org/spec/}{Iceberg spec} \\
``REST fixture là catalog production nhẹ'' & \code{iceberg-rest-fixture} chỉ để học và test; quickstart của Polaris cũng ghi không dành cho production & \href{https://polaris.apache.org/guides/quickstart/}{Polaris Quickstart} \\
``RBAC trên catalog cũng bảo vệ file'' & Ai giữ key S3 tĩnh đi vòng được qua catalog; vended credential được cấp theo quyền trên bảng & \href{https://polaris.apache.org/in-dev/unreleased/managing-security/access-control/}{Polaris RBAC} \\
``RBAC của catalog có lọc dòng, che cột'' & OpenFGA và Cedar trong Lakekeeper không biểu diễn được; row filter và column mask do engine/nền tảng thực thi & \href{https://docs.lakekeeper.io/docs/nightly/authorization/}{Lakekeeper Authorization}, \href{https://docs.databricks.com/aws/en/data-governance/unity-catalog/filters-and-masks/}{Databricks} \\
``Để secret S3 trong table property cũng được'' & Polaris dặn không để secret trong property vì ai đọc property sẽ thấy & \href{https://polaris.apache.org/in-dev/unreleased/managing-security/access-control/}{Polaris RBAC} \\
``\code{DELETE FROM} là đã xoá dữ liệu cá nhân'' & Snapshot, tag, branch cũ còn giữ file tới khi expire và dọn file & \href{https://iceberg.apache.org/docs/latest/branching/}{Branching}, \href{https://iceberg.apache.org/docs/latest/maintenance/}{Maintenance} \\
``Công cụ lineage tự tìm ra lineage'' & OpenLineage cần bên phát như Spark listener hay Airflow provider & \href{https://openlineage.io/getting-started/}{OpenLineage} \\
``Glue chỉ nói được giao thức riêng của AWS'' & Glue Data Catalog có Iceberg REST endpoint chuẩn, xác thực SigV4 & \href{https://docs.aws.amazon.com/glue/latest/dg/connect-glu-iceberg-rest.html}{AWS Glue} \\
\end{xltabular}



\partpage{III}{Xử lý và vận hành pipeline}{Spark, các pattern nạp dữ liệu tăng dần, điều phối bằng Airflow, test và DataOps. Kết thúc bằng mốc capstone 3: pipeline tăng dần có test chặn.}
% ---------------- m07_10 ----------------
% ============================================================
% Module 7, 8, 9, 10
% ============================================================

\modulechapter{7}{Apache Spark và tính toán phân tán}
\moduleinfo{Trung cấp}{30 giờ}{Module 1, 5}

Khi dữ liệu và phép tính vượt sức một máy, bạn chia dữ liệu thành partition và chạy song song trên nhiều executor. Spark là engine làm việc đó: nó biến DataFrame hoặc câu SQL thành kế hoạch thực thi gồm job, stage và task. Trong lakehouse của lộ trình này, Spark là engine ghi chính, đưa dữ liệu từ raw qua staging tới marts trên bảng Iceberg. Học xong, bạn đọc được \code{EXPLAIN} và Spark UI để biết vì sao một truy vấn chậm, và viết được \code{MERGE INTO} Iceberg có kế hoạch hợp lý.

\section{Bạn sẽ học được}
\begin{itemize}
  \item Giải thích vai trò của driver, executor và cluster manager, và cách Spark Connect tách client khỏi driver.
  \item Phân biệt transformation với action, phép biến đổi narrow với wide, và chỉ ra ranh giới stage nằm ở shuffle.
  \item Đọc \code{EXPLAIN} ở các chế độ khác nhau để thấy Exchange, kiểu join và bộ lọc được đẩy xuống.
  \item Chọn chiến lược join và cấu hình partition dựa trên số liệu, kể cả khi AQE đang bật.
  \item Chẩn đoán data skew trên Spark UI và chọn cách xử lý phù hợp.
  \item Viết \code{MERGE INTO} và overwrite partition trên bảng Iceberg từ Spark một cách an toàn.
  \item Nêu những thay đổi hành vi chính khi chuyển từ Spark 3.5 lên 4.x.
\end{itemize}

\section{Nội dung chính}

\subsection{Kiến trúc và mô hình thực thi}
\begin{enumerate}
  \item \textbf{Vì sao phân tán.} Một máy có giới hạn về CPU, bộ nhớ và đĩa; mở rộng ngang nghĩa là chia dữ liệu thành partition và xử lý song song. Cái giá phải trả là dữ liệu phải di chuyển qua mạng khi một phép tính cần các dòng cùng khoá nằm chung một chỗ.
  \item \textbf{Driver, executor, cluster manager.} Driver chạy hàm \code{main} và lập kế hoạch; executor chạy task và giữ dữ liệu cache; cluster manager (standalone, YARN, Kubernetes) cấp tài nguyên. Mỗi driver có một web UI, mặc định ở cổng 4040 (\href{https://spark.apache.org/docs/latest/cluster-overview.html}{Cluster Mode Overview}, \href{https://spark.apache.org/docs/latest/web-ui.html}{Web UI}).
  \item \textbf{SparkSession, DataFrame và Spark SQL.} DataFrame và SQL đi qua cùng một bộ tối ưu, nên hai cách viết cho cùng kế hoạch (\href{https://spark.apache.org/docs/latest/sql-programming-guide.html}{Spark SQL Guide}).
  \item \textbf{Transformation, action và đánh giá lười.} Transformation chỉ ghi lại kế hoạch; chỉ khi gọi action như \code{count()}, \code{collect()} hay ghi bảng thì Spark mới chạy job (\href{https://spark.apache.org/docs/latest/rdd-programming-guide.html}{RDD Programming Guide}). Phép narrow như \code{filter} xử lý từng partition riêng; phép wide như \code{groupBy} hay join cần shuffle.
  \item \textbf{Job, stage, task.} Một action tạo một hoặc nhiều job; mỗi job chia thành stage tại các ranh giới shuffle; mỗi stage có một task cho mỗi partition. Bạn thấy cấu trúc này trong tab Jobs và Stages của Spark UI.
  \item \textbf{Catalyst và kế hoạch.} Truy vấn đi qua logical plan, logical plan đã tối ưu, rồi physical plan. \code{EXPLAIN} mặc định chỉ in physical plan; \code{EXTENDED} in cả các logical plan, \code{FORMATTED} in phần khung và chi tiết từng node, \code{COST} in thống kê nếu có, \code{CODEGEN} in mã sinh ra (\href{https://spark.apache.org/docs/latest/sql-ref-syntax-qry-explain.html}{EXPLAIN}). Trong physical plan, node \code{Exchange} chính là một shuffle.
\end{enumerate}

\subsection{Hiệu năng: shuffle, join, AQE, skew}
\begin{enumerate}[resume]
  \item \textbf{Partition sau shuffle.} \code{spark.sql.\allowbreak shuffle.partitions} mặc định là 200 cho join và aggregation (\href{https://spark.apache.org/docs/latest/sql-performance-tuning.html}{Performance Tuning}). \code{repartition(n)} luôn shuffle; \code{coalesce(n)} gộp partition mà không shuffle toàn bộ, chỉ dùng để giảm số partition.
  \item \textbf{Chiến lược join.} Spark có broadcast hash join, sort-merge join, shuffled hash join và nested loop join. Bảng nhỏ hơn \code{spark.sql.\allowbreak autoBroadcast\allowbreak JoinThreshold} (mặc định 10485760 byte, tức 10 MB) được broadcast; đặt \code{-1} để tắt. Hint join gồm \code{BROADCAST}, \code{MERGE}, \code{SHUFFLE\_HASH}, \code{SHUFFLE\_REPLICATE\_NL} (\href{https://spark.apache.org/docs/latest/sql-ref-syntax-qry-select-hints.html}{Hints}).
  \item \textbf{Adaptive Query Execution (AQE).} AQE tối ưu lại kế hoạch lúc chạy dựa trên thống kê thật sau shuffle, và bật mặc định từ Spark 3.2.0. Ba tính năng chính: gộp partition sau shuffle, đổi sort-merge join sang broadcast hash join hoặc shuffled hash join, và tách task bị skew trong sort-merge join (\href{https://spark.apache.org/docs/latest/sql-performance-tuning.html}{Performance Tuning}).
  \item \textbf{Data skew.} Khi vài khoá chiếm phần lớn dữ liệu, vài task chạy lâu hơn hẳn phần còn lại; tab Stages cho thấy thời gian task lớn nhất so với trung vị. AQE skew join chỉ xử lý skew trong sort-merge join; aggregation lệch thì cần cách khác như salting (thêm hậu tố ngẫu nhiên vào khoá, tính hai bước).
  \item \textbf{Cache và persist.} \code{cache()} cũng lười: dữ liệu chỉ được lưu khi một action chạy. Chọn storage level phù hợp và gọi \code{unpersist()} khi xong (\href{https://spark.apache.org/docs/latest/rdd-programming-guide.html\#rdd-persistence}{RDD Persistence}).
  \item \textbf{Spark UI.} Các tab Jobs, Stages, Storage, Environment, Executors và SQL cho bạn thấy shuffle read/write, spill, GC và thời gian từng task (\href{https://spark.apache.org/docs/latest/web-ui.html}{Web UI}). Nếu cổng 4040 bận, Spark thử 4041, 4042 và ghi cổng đã chọn vào log.
\end{enumerate}

\subsection{Spark Connect, PySpark và phiên bản}
\begin{enumerate}[resume]
  \item \textbf{Spark Connect.} Client đổi thao tác DataFrame thành unresolved logical plan, mã hoá bằng protocol buffers và gửi tới server qua gRPC; client mặc định nối cổng 15002. Client PySpark qua Connect không dùng được \code{SparkContext} và RDD (\href{https://spark.apache.org/docs/latest/spark-connect-overview.html}{Spark Connect Overview}). Server khởi động bằng \code{sbin/\allowbreak start-connect-server.sh}; với Spark 3.5, đọc trang này ở bộ docs 3.5.9 vì cách khởi động khác.
  \item \textbf{Classic và Connect trong Spark 4.x.} Spark 4.0 thêm client Python nhẹ \code{pyspark-client}, cấu hình \code{spark.api.mode} để bật tắt Connect cho ứng dụng, và một tarball phụ bật Connect sẵn (\href{https://spark.apache.org/releases/spark-release-4-0-0.html}{Spark 4.0.0 release notes}). Nếu bạn không dùng các cơ chế Connect, session vẫn chạy như trước, tức chế độ Classic.
  \item \textbf{PySpark và Python UDF.} Với code DataFrame, kế hoạch chạy trong JVM; chi phí riêng của Python nằm chủ yếu ở UDF, nơi dữ liệu phải chuyển sang Python worker. Spark 4.2.0 bật tối ưu Arrow cho Python UDF theo mặc định (\href{https://spark.apache.org/releases/spark-release-4-2-0.html}{Spark 4.2.0 release notes}).
  \item \textbf{Chính sách phiên bản.} Bản feature được bảo trì 6 tháng, bản LTS (bản cuối của mỗi dòng major) 18 tháng. Nhánh 3.5 có LTS mở rộng tới 11/2027 và trong giai đoạn này chỉ nhận vá bảo mật; LTS của dòng 4.x theo kế hoạch là 4.5.0 (\href{https://spark.apache.org/versioning-policy.html}{Versioning policy}).
  \item \textbf{Spark 4.0 đổi gì.} ANSI SQL mode bật mặc định, JDK 17 là mặc định (bỏ JDK 8 và 11), Scala 2.13 là mặc định (bỏ 2.12) (\href{https://spark.apache.org/releases/spark-release-4-0-0.html}{Spark 4.0.0 release notes}). Với ANSI, phép \code{CAST} không hợp lệ báo lỗi thay vì trả NULL; xem \href{https://spark.apache.org/docs/latest/sql-migration-guide.html}{SQL Migration Guide} trước khi nâng cấp.
\end{enumerate}

\subsection{Ghi Iceberg từ Spark và streaming}
\begin{enumerate}[resume]
  \item \textbf{MERGE INTO.} Mỗi dòng của bảng đích chỉ được một dòng nguồn cập nhật; nếu nhiều dòng nguồn khớp cùng một dòng đích, câu lệnh báo lỗi, nên hãy khử trùng nguồn trước. Iceberg khuyên dùng \code{MERGE INTO} thay cho \code{INSERT OVERWRITE} vì nó chỉ viết lại file bị ảnh hưởng (\href{https://iceberg.apache.org/docs/latest/spark-writes/}{Iceberg Spark Writes}). \code{MERGE} và \code{CALL} cần bật Iceberg SQL extensions trên phía Spark chạy truy vấn.
  \item \textbf{Overwrite partition.} Chế độ overwrite mặc định của Spark là static; Iceberg khuyên đặt \code{spark.sql.\allowbreak sources.\allowbreak partitionOverwriteMode=dynamic} để chỉ thay các partition có dòng trong kết quả. Với DataFrame API, \code{writeTo(t).\allowbreak overwritePartitions()} tương đương overwrite động.
  \item \textbf{Phân phối khi ghi.} \code{write.\allowbreak distribution-mode} có các giá trị \code{none}, \code{hash}, \code{range}; \code{hash} là mặc định từ Iceberg 1.2.0 và thêm một shuffle theo giá trị partition trước khi ghi, đổi lại ít file nhỏ hơn (\href{https://iceberg.apache.org/docs/latest/spark-writes/}{Iceberg Spark Writes}).
  \item \textbf{Structured Streaming, chỉ giới thiệu.} Spark có API streaming dùng chung DataFrame; module streaming sau sẽ học kỹ. Xem mục Structured Streaming trong \href{https://spark.apache.org/docs/latest/}{Spark docs}.
\end{enumerate}

\section{Bắt đầu}
\begin{itemize}
  \item \href{https://spark.apache.org/docs/latest/spark-connect-overview.html}{Spark Connect Overview}: chạy server Spark Connect và nối từ PySpark bằng \code{SparkSession.\allowbreak builder.\allowbreak remote(\allowbreak "sc://localhost:15002")}.
  \item \href{https://iceberg.apache.org/spark-quickstart/}{Iceberg Spark Quickstart}: tạo catalog và bảng Iceberg đầu tiên từ Spark.
  \item \href{https://spark.apache.org/docs/latest/sql-performance-tuning.html}{SQL Performance Tuning}: đọc mục AQE và join hint trước khi làm lab 2.
\end{itemize}

\begin{ghichu}
Docs \code{latest} của Spark mô tả bản mới nhất (4.3.0 khi kiểm ngày 02/11/2026). Nếu lab chạy 3.5.x, đọc bộ docs có số phiên bản, ví dụ \code{spark.apache.org/\allowbreak docs/3.5.9/}. Chọn gói \code{iceberg-spark-runtime} đúng phiên bản Spark và Scala theo trang Releases của Iceberg.
\end{ghichu}

\section{Tài liệu chính thức}
\begin{xltabular}{\textwidth}{@{}P{4.6cm} L P{1.8cm}@{}}
\toprule
\thd{tài liệu} & \thd{nội dung cần đọc} & \thd{thời lượng} \\
\midrule
\endhead
\bottomrule
\endfoot
\href{https://spark.apache.org/docs/latest/cluster-overview.html}{Spark: Cluster Mode Overview} & driver, executor, cluster manager, thuật ngữ job, stage, task & 1 giờ \\
\href{https://spark.apache.org/docs/latest/sql-programming-guide.html}{Spark SQL, DataFrames and Datasets Guide} & Getting Started, Data Sources (Parquet, JDBC) & 2 giờ \\
\href{https://spark.apache.org/docs/latest/rdd-programming-guide.html}{RDD Programming Guide} & transformation và action, shuffle, persistence; chỉ đọc khái niệm & 1,5 giờ \\
\href{https://spark.apache.org/docs/latest/sql-performance-tuning.html}{SQL Performance Tuning} & caching, partition, join hint, toàn bộ mục AQE & 2 giờ \\
\href{https://spark.apache.org/docs/latest/sql-ref-syntax-qry-explain.html}{EXPLAIN} & các chế độ \code{EXTENDED}, \code{FORMATTED}, \code{COST}, \code{CODEGEN} & 0,5 giờ \\
\href{https://spark.apache.org/docs/latest/web-ui.html}{Web UI} & Jobs, Stages, Storage, Executors, SQL & 1 giờ \\
\href{https://spark.apache.org/docs/latest/spark-connect-overview.html}{Spark Connect Overview} & kiến trúc, cách khởi động server, API không hỗ trợ & 1 giờ \\
\href{https://spark.apache.org/versioning-policy.html}{Spark Versioning Policy} & thời gian bảo trì, LTS mở rộng của 3.5 & 0,3 giờ \\
\href{https://spark.apache.org/releases/spark-release-4-0-0.html}{Spark 4.0.0 release notes}, \href{https://spark.apache.org/docs/latest/sql-migration-guide.html}{SQL Migration Guide} & ANSI mode, JDK, Scala, Spark Connect; mục nâng cấp từ 3.5 & 1,5 giờ \\
\href{https://iceberg.apache.org/docs/latest/spark-writes/}{Iceberg: Spark Writes} & MERGE INTO, INSERT OVERWRITE, Writing Distribution Modes, Controlling File Sizes & 1,5 giờ \\
\href{https://iceberg.apache.org/docs/latest/spark-procedures/}{Iceberg: Spark Procedures} & \code{rewrite\_data\_files}, \code{expire\_snapshots} & 0,5 giờ \\
\end{xltabular}

\section{Khoá học miễn phí}
\begin{xltabular}{\textwidth}{@{}P{4.2cm} L P{2.6cm}@{}}
\toprule
\thd{khoá} & \thd{nội dung} & \thd{miễn phí} \\
\midrule
\endhead
\bottomrule
\endfoot
\href{https://pdos.csail.mit.edu/6.824/}{MIT 6.5840 Distributed Systems (Spring 2026)} & Lab 1 viết MapReduce bằng Go; lịch bài giảng ở trang Schedule. Làm Lab 1 để hiểu map, shuffle, reduce trước khi học Spark sâu & trang khoá và đề lab công khai \\
\end{xltabular}

\section{Lab thực hành}

\subsection{Lab 7.1: Spark Connect và đánh giá lười trên Olist}
\begin{enumerate}
  \item Thêm vào Docker Compose của stack tham chiếu một dịch vụ Spark chạy Spark Connect server, mở cổng 15002 và 4040, cấu hình catalog Iceberg REST và storage RustFS.
  \item Từ máy của bạn, cài \code{pyspark-client} (Spark 4.x) hoặc \code{pyspark[connect]} đúng phiên bản server, rồi nối bằng \code{SparkSession.builder.\allowbreak remote(...)}.
  \item Đọc \code{olist\_orders\_dataset.csv} và \code{olist\_order\_items\_dataset.csv}. Viết năm transformation (lọc, chọn cột, đổi kiểu thời gian) mà không gọi action; mở tab Jobs và xác nhận chưa có job nào.
  \item Gọi \code{count()}; xem job, stage, task xuất hiện. Sau đó chạy \code{groupBy(\allowbreak "order\_status").\allowbreak count()} và đếm số stage.
  \item Thử gọi \code{spark.sparkContext} từ client Connect và ghi lại lỗi nhận được.
\end{enumerate}
\hoanthanh{bạn có ảnh chụp tab Jobs trước và sau action, và giải thích được vì sao \code{groupBy} tạo thêm một stage.}

\subsection{Lab 7.2: EXPLAIN và chiến lược join}
\begin{enumerate}
  \item Join \code{order\_items} với \code{products} theo \code{product\_id}, rồi aggregate doanh thu theo \code{product\_category\_name}.
  \item Chạy \code{EXPLAIN FORMATTED}. Tìm node \code{Exchange} và kiểu join; ghi lại kích thước ước lượng của \code{products}.
  \item Đặt \code{spark.sql.\allowbreak autoBroadcast\allowbreak JoinThreshold=-1}, chạy lại và so kiểu join. Thử các hint \code{BROADCAST}, \code{MERGE}, \code{SHUFFLE\_HASH}.
  \item Tắt rồi bật \code{spark.sql.\allowbreak adaptive.enabled}; so physical plan cuối cùng trong tab SQL (dòng \code{AdaptiveSparkPlan} cho biết kế hoạch cuối).
\end{enumerate}
\hoanthanh{bạn có bảng so sánh bốn cấu hình (mặc định, tắt broadcast, từng hint, tắt AQE) với kiểu join, số Exchange và thời gian chạy.}

\subsection{Lab 7.3: skew và số partition}
\begin{enumerate}
  \item Đếm số đơn theo \code{customer\_state} (join \code{orders} với \code{customers}) và ghi lại bang chiếm nhiều dòng nhất.
  \item Tạo dữ liệu skew có kiểm soát: nhân bản \code{order\_items} sao cho 90\% dòng có cùng một \code{seller\_id}, rồi join với \code{sellers} và aggregate.
  \item Trong tab Stages, so thời gian task lớn nhất với trung vị. Bật tắt \code{spark.sql.\allowbreak adaptive.\allowbreak skewJoin.enabled} và so.
  \item Với aggregation lệch, thử salting: thêm cột \code{salt} ngẫu nhiên từ 0 tới 15, aggregate theo khoá và salt, rồi aggregate lần hai theo khoá.
  \item So \code{spark.sql.\allowbreak shuffle.partitions} 200 với 8 khi AQE tắt và khi AQE bật.
\end{enumerate}
\hoanthanh{bạn chỉ ra được trên Spark UI task bị skew trước và sau khi sửa, và giải thích vì sao AQE xử lý được trường hợp join nhưng không xử lý được trường hợp aggregation.}

\subsection{Lab 7.4: ghi Iceberg và MERGE INTO}
\begin{enumerate}
  \item Tạo bảng Iceberg \code{lake.staging.orders} phân vùng theo \code{months(\allowbreak order\_purchase\_\allowbreak timestamp)}, nạp toàn bộ \code{orders}.
  \item Tạo một DataFrame cập nhật: 1.000 đơn đổi \code{order\_status} và 100 đơn mới. Cố ý để vài \code{order\_id} trùng trong nguồn, chạy \code{MERGE INTO} và đọc lỗi; sau đó khử trùng nguồn bằng \code{row\_number()} và chạy lại.
  \item Chạy \code{EXPLAIN} cho câu \code{MERGE}; chỉ ra kiểu join và Exchange trước bước ghi.
  \item Đếm data file qua bảng metadata \code{lake.staging.orders.files} và snapshot qua \code{.snapshots} (\href{https://iceberg.apache.org/docs/latest/spark-queries/}{Iceberg Spark Queries}). Đặt \code{write.\allowbreak distribution-mode} là \code{none} rồi \code{hash}, ghi lại và so số file.
\end{enumerate}
\hoanthanh{câu \code{MERGE} chạy lại lần hai không đổi số dòng, và bạn giải thích được số file khác nhau giữa hai chế độ phân phối.}

\section{Kế hoạch theo tuần}
\begin{xltabular}{\textwidth}{@{}P{1.4cm} L P{1.4cm}@{}}
\toprule
\thd{tuần} & \thd{việc} & \thd{giờ} \\
\midrule
\endhead
\bottomrule
\endfoot
1 & Cluster Mode Overview, Spark SQL Guide, RDD Programming Guide phần khái niệm, Spark Connect Overview; MIT 6.5840 Lab 1 nếu muốn; Lab 7.1 & 10 \\
2 & Performance Tuning, EXPLAIN, Hints, Web UI; Lab 7.2 và 7.3 & 10 \\
3 & Iceberg Spark Writes, Spark 4.0 release notes, Migration Guide; Lab 7.4; thử thách kỹ năng & 10 \\
thêm & Iceberg Spark Procedures (\code{rewrite\_data\_files}); thử cùng lab trên cả Spark 3.5.x và 4.x & 4 \\
\end{xltabular}

\begin{thuthach}
Bảng mart \code{lake.marts.revenue\_by\_category\_month} được tính từ \code{orders}, \code{order\_items}, \code{products} và \code{product\_category\_name\_translation}, rồi ghi đè theo tháng vào Iceberg. Tự viết truy vấn, đo thời gian lần đầu, sau đó làm nó nhanh hơn ít nhất 30\% mà kết quả không đổi. Mỗi thay đổi phải có bằng chứng từ \code{EXPLAIN} hoặc Spark UI, và lần chạy lại chỉ được thay các partition tháng có dữ liệu.
\hoanthanh{bạn có báo cáo một trang gồm thời gian trước và sau, kế hoạch trước và sau, lý do của từng thay đổi, và kiểm tra rằng tổng doanh thu không đổi.}
\end{thuthach}

\section{Câu hỏi tự kiểm tra}
\begin{enumerate}
  \item Driver làm gì, executor làm gì? Vì sao \code{collect()} trên bảng lớn có thể làm hỏng job?
  \item Cái gì tạo ra một stage mới? Chỉ ra dấu hiệu của nó trong \code{EXPLAIN}.
  \item Một stage có 200 task, 199 task xong trong 5 giây, một task mất 4 phút. Nguyên nhân có thể là gì, xem ở đâu, và kể hai cách sửa.
  \item Khi nào Spark chọn broadcast hash join, và khi nào broadcast gây hại?
  \item AQE quyết định được gì lúc chạy mà bộ tối ưu không quyết được lúc lập kế hoạch?
  \item \code{repartition(n)} khác \code{coalesce(n)} thế nào?
  \item Khi nào \code{.cache()} là ý tồi?
  \item Vì sao client PySpark qua Spark Connect không gọi được \code{spark.sparkContext}?
  \item \code{MERGE INTO} Iceberg báo lỗi khi nào, và bạn chuẩn bị dữ liệu nguồn ra sao để tránh?
  \item Overwrite static và dynamic khác nhau thế nào khi bạn chỉ muốn ghi lại một tháng?
  \item Kể ba thay đổi trong Spark 4.0 có thể làm hỏng một job viết cho 3.5.
\end{enumerate}

\section{Điểm dễ nhầm}
\begin{xltabular}{\textwidth}{@{}P{4.2cm} L P{3.4cm}@{}}
\toprule
\thd{điểm dễ nhầm} & \thd{thực tế} & \thd{docs chính thức} \\
\midrule
\endhead
\bottomrule
\endfoot
Nhiều người nghĩ Spark chạy trong bộ nhớ nên không chạm đĩa & Shuffle ghi dữ liệu ra file và cache có thể tràn xuống đĩa; Web UI hiển thị shuffle write và spill & \href{https://spark.apache.org/docs/latest/web-ui.html}{Web UI} \\
Càng nhiều partition càng nhanh & Quá nhiều task nhỏ làm tăng chi phí lập lịch và sinh nhiều file nhỏ; AQE gộp partition sau shuffle vì lý do này & \href{https://spark.apache.org/docs/latest/sql-performance-tuning.html}{Performance Tuning} \\
\code{cache()} lưu dữ liệu ngay & \code{cache()} lười; dữ liệu chỉ được lưu khi có action & \href{https://spark.apache.org/docs/latest/rdd-programming-guide.html\#rdd-persistence}{RDD Persistence} \\
AQE sửa mọi loại skew & AQE skew join chỉ xử lý skew trong sort-merge join & \href{https://spark.apache.org/docs/latest/sql-performance-tuning.html}{Performance Tuning} \\
Spark Connect là một engine khác & Cùng engine; client gửi unresolved logical plan qua gRPC và server lập kế hoạch như bình thường & \href{https://spark.apache.org/docs/latest/spark-connect-overview.html}{Spark Connect Overview} \\
Từ Spark 4.0, Connect là chế độ mặc định & Spark 4.0 thêm \code{pyspark-client}, \code{spark.api.mode} và một tarball phụ bật Connect sẵn; không dùng cơ chế Connect thì session vẫn chạy Classic & \href{https://spark.apache.org/releases/spark-release-4-0-0.html}{Spark 4.0.0 release notes} \\
\code{CAST} sai kiểu luôn trả NULL & Từ Spark 4.0, ANSI mode bật mặc định nên \code{CAST} không hợp lệ báo lỗi; muốn NULL thì dùng \code{try\_cast} một cách có chủ đích & \href{https://spark.apache.org/docs/latest/sql-migration-guide.html}{SQL Migration Guide} \\
\code{INSERT OVERWRITE} chỉ thay partition có dữ liệu mới & Chế độ mặc định của Spark là static; cần \code{partitionOverwriteMode=dynamic} & \href{https://iceberg.apache.org/docs/latest/spark-writes/}{Iceberg Spark Writes} \\
\end{xltabular}


% ============================================================

\modulechapter{8}{Ingestion và các pattern pipeline}
\moduleinfo{Trung cấp}{20 giờ}{Module 5, 7}

Ingestion là bước đưa dữ liệu từ hệ thống nguồn vào tầng raw của lakehouse. Bài toán khó nhất không phải đọc dữ liệu một lần, mà đọc lặp lại mỗi ngày sao cho không sót, không trùng và chạy lại được khi lỗi. Module này đặt tên cho các pattern chuẩn: nạp toàn bộ và nạp tăng dần, cursor và high-water mark, idempotency, dữ liệu tới muộn, backfill, CDC và schema drift. Học xong, bạn xây được một pipeline tăng dần từ PostgreSQL vào Iceberg mà chạy lại bao nhiêu lần cũng cho cùng kết quả.

\section{Bạn sẽ học được}
\begin{itemize}
  \item Chọn giữa nạp toàn bộ và nạp tăng dần, và gọi đúng tên các sync mode thường gặp.
  \item Thiết kế cursor hoặc high-water mark, và chỉ ra những trường hợp nó bỏ sót thay đổi.
  \item Viết bước ghi idempotent bằng \code{MERGE} hoặc overwrite partition, và lưu tiến độ đúng thời điểm.
  \item Xử lý dữ liệu tới muộn bằng cửa sổ nhìn lại, và lập kế hoạch backfill.
  \item Giải thích CDC dựa trên log khác truy vấn theo \code{updated\_at} ở đâu, và đọc được một sự kiện thay đổi của Debezium.
  \item Phân loại thay đổi schema và chọn phản ứng phù hợp ở từng tầng.
\end{itemize}

\section{Nội dung chính}
\begin{enumerate}
  \item \textbf{Nguồn và tần suất.} Nguồn có thể là cơ sở dữ liệu OLTP, API, file hay luồng sự kiện; bạn kéo (pull) theo lịch hoặc nhận đẩy (push). Batch và streaming là hai đầu của một dải tần suất; batch chạy mỗi vài phút đã gần với micro-batch.
  \item \textbf{Nạp toàn bộ và nạp tăng dần.} Airbyte đặt tên năm sync mode: Incremental Append, Incremental Append + Deduped, Full Refresh Append, Full Refresh Overwrite, Full Refresh Overwrite + Deduped; các mode Deduped khử trùng theo primary key (\href{https://docs.airbyte.com/platform/using-airbyte/core-concepts/sync-modes}{Airbyte Sync Modes}). Dù dùng công cụ nào, đây là bộ từ vựng chung.
  \item \textbf{Cursor và high-water mark.} Nạp tăng dần lưu giá trị lớn nhất đã nạp của một cột (cursor) rồi lần sau chỉ đọc phần lớn hơn. dlt lưu giá trị này trong state và đưa ra dưới dạng \code{start\_value} ở lần chạy sau; biên dưới mặc định là bao gồm, nên dòng ở đúng giá trị cũ được đọc lại và dlt khử trùng chúng (\href{https://dlthub.com/docs/general-usage/incremental/cursor}{dlt: cursor-based incremental loading}). Đừng nhầm high-water mark của batch với watermark của streaming, thứ ước lượng tiến độ theo event time.
  \item \textbf{Cursor bỏ sót gì.} Cursor trên cột thời điểm tạo (ví dụ \code{order\_purchase\_timestamp}) chỉ bắt dòng mới, không bắt dòng bị sửa. Cursor trên \code{updated\_at} bắt được sửa nếu nguồn luôn cập nhật cột này, nhưng không bao giờ thấy dòng bị xoá cứng. Chốt cận trên của mỗi lần chạy (ví dụ thời điểm bắt đầu run) để hai lần chạy không chồng nhau một cách khó đoán.
  \item \textbf{Lưu tiến độ sau khi dữ liệu đã commit.} Nếu ghi cursor trước rồi mới ghi dữ liệu, một lần lỗi giữa chừng làm mất dữ liệu. Ghi dữ liệu trước, cursor sau: lỗi giữa chừng chỉ làm đọc lại một phần, và đích idempotent biến việc đọc lại thành vô hại. Đây là at-least-once cộng ghi idempotent.
  \item \textbf{Idempotency.} Task idempotent cho cùng kết quả mỗi lần chạy lại. Docs Airflow khuyên coi task như transaction, không dùng \code{INSERT} khi chạy lại mà dùng UPSERT, đọc và ghi đúng một partition xác định, và không dùng \code{datetime.now()} cho phép tính quan trọng (\href{https://airflow.apache.org/docs/apache-airflow/stable/best-practices.html}{Airflow Best Practices}). Trên Iceberg, hai cách chuẩn là \code{MERGE INTO} theo khoá và overwrite partition động (\href{https://iceberg.apache.org/docs/latest/spark-writes/}{Iceberg Spark Writes}).
  \item \textbf{Dữ liệu tới muộn.} Dòng có thể tới sau khi cursor đã vượt qua nó. Cách phổ biến là cửa sổ nhìn lại: mỗi lần đọc lại một khoảng trước cursor rồi \code{MERGE}. dlt có tuỳ chọn \code{lag} cho việc này, tính bằng giây với cursor kiểu datetime, bằng ngày với cursor kiểu date (\href{https://dlthub.com/docs/general-usage/incremental/lag}{dlt: lag}).
  \item \textbf{Backfill và xử lý lại.} Backfill là chạy lại pipeline cho một khoảng thời gian trong quá khứ, ví dụ sau khi sửa logic. Nó chỉ an toàn khi mỗi lần chạy ghi đè đúng phần của nó. Airflow 3 có backfill do scheduler quản lý (\href{https://airflow.apache.org/docs/apache-airflow/stable/core-concepts/backfill.html}{Airflow Backfill}); Module 9 sẽ dùng nó.
  \item \textbf{CDC dựa trên log.} Thay vì truy vấn bảng, CDC đọc log của cơ sở dữ liệu. Connector PostgreSQL của Debezium dùng logical decoding với plugin \code{pgoutput} (có sẵn từ PostgreSQL 10) qua một replication slot; lần đầu nó chụp snapshot, mỗi dòng snapshot là sự kiện \code{op=r}, rồi đọc tiếp từ WAL (\href{https://debezium.io/documentation/reference/stable/connectors/postgresql.html}{Debezium PostgreSQL connector}, \href{https://www.postgresql.org/docs/18/logicaldecoding.html}{PostgreSQL: Logical Decoding}).
  \item \textbf{Cấu trúc sự kiện thay đổi.} Payload có \code{before}, \code{after}, \code{source}, \code{op} và \code{ts\_ms}; \code{op} là \code{c} (tạo), \code{u} (sửa), \code{d} (xoá), \code{r} (đọc snapshot), \code{t} (truncate). Nội dung \code{before} phụ thuộc \code{REPLICA IDENTITY} của bảng: mặc định chỉ có khoá chính, \code{FULL} thì có mọi cột. Sau sự kiện xoá, connector gửi thêm một tombstone (khoá giữ nguyên, value null) để Kafka log compaction dọn được khoá đó (\href{https://debezium.io/documentation/reference/stable/connectors/postgresql.html}{Debezium PostgreSQL connector}).
  \item \textbf{Kafka Connect.} Kafka Connect là công cụ truyền dữ liệu giữa Kafka và hệ thống khác một cách có thể mở rộng và tin cậy; nó chạy ở chế độ standalone hoặc distributed và tự quản lý việc commit offset cho connector (\href{https://kafka.apache.org/44/kafka-connect/overview/}{Kafka Connect Overview}). Debezium chạy như một source connector trên Kafka Connect (\href{https://debezium.io/documentation/reference/stable/architecture.html}{Debezium Architecture}).
  \item \textbf{Schema drift.} Thêm cột là thay đổi tương thích; đổi kiểu, đổi tên hay xoá cột có thể phá vỡ tầng sau. Iceberg hỗ trợ schema evolution an toàn vì cột được theo dõi bằng ID chứ không bằng tên (\href{https://iceberg.apache.org/docs/latest/evolution/}{Iceberg Evolution}). Ở tầng raw thường cho phép evolve; ở staging và marts thì nên cố định và kiểm bằng test hoặc contract (Module 10).
  \item \textbf{ETL và ELT.} ETL biến đổi trước khi nạp vào đích; ELT nạp nguyên bản vào tầng raw rồi biến đổi bên trong lakehouse. Kiến trúc raw, staging, marts của lộ trình này là ELT.
\end{enumerate}

\section{Bắt đầu}
\begin{itemize}
  \item \href{https://docs.airbyte.com/platform/using-airbyte/core-concepts/sync-modes}{Airbyte: Sync Modes}: đọc 15 phút để nắm từ vựng.
  \item \href{https://dlthub.com/docs/general-usage/incremental/cursor}{dlt: cursor-based incremental loading}: chạy ví dụ trong trang với destination DuckDB.
  \item \href{https://debezium.io/documentation/reference/stable/tutorial.html}{Debezium Tutorial}: chạy theo tutorial chính thức để thấy sự kiện create, update, delete.
\end{itemize}

\section{Tài liệu chính thức}
\begin{xltabular}{\textwidth}{@{}P{4.6cm} L P{1.8cm}@{}}
\toprule
\thd{tài liệu} & \thd{nội dung cần đọc} & \thd{thời lượng} \\
\midrule
\endhead
\bottomrule
\endfoot
\href{https://docs.airbyte.com/platform/using-airbyte/core-concepts/sync-modes}{Airbyte: Sync Modes} & năm sync mode và cách khử trùng theo primary key & 0,5 giờ \\
\href{https://dlthub.com/docs/general-usage/incremental/cursor}{dlt: cursor-based incremental loading}, \href{https://dlthub.com/docs/general-usage/incremental/lag}{lag} & \code{dlt.sources.incremental}, state, khử trùng ở biên, \code{row\_order}, \code{lag} & 2 giờ \\
\href{https://airflow.apache.org/docs/apache-airflow/stable/best-practices.html}{Airflow: Best Practices} & mục Creating a task: task như transaction, UPSERT, partition xác định & 0,5 giờ \\
\href{https://iceberg.apache.org/docs/latest/spark-writes/}{Iceberg: Spark Writes} & MERGE INTO, INSERT OVERWRITE động & 1 giờ \\
\href{https://spark.apache.org/docs/latest/sql-data-sources-jdbc.html}{Spark: JDBC To Other Databases} & các option \code{url}, \code{query}, \code{dbtable}, driver & 0,5 giờ \\
\href{https://debezium.io/documentation/reference/stable/architecture.html}{Debezium: Architecture} & Debezium trên Kafka Connect, Debezium Server & 0,5 giờ \\
\href{https://debezium.io/documentation/reference/stable/connectors/postgresql.html}{Debezium: PostgreSQL connector} & How the connector works, snapshot, data change events & 2 giờ \\
\href{https://kafka.apache.org/44/kafka-connect/overview/}{Kafka Connect} & Overview và User Guide: standalone, distributed, offset & 1 giờ \\
\href{https://iceberg.apache.org/docs/latest/evolution/}{Iceberg: Evolution} & schema evolution, partition evolution & 0,5 giờ \\
\href{https://airflow.apache.org/docs/apache-airflow/stable/core-concepts/backfill.html}{Airflow: Backfill} & backfill do scheduler quản lý & 0,5 giờ \\
\end{xltabular}

\section{Khoá học miễn phí}
\begin{xltabular}{\textwidth}{@{}P{4.2cm} L P{2.6cm}@{}}
\toprule
\thd{khoá} & \thd{nội dung} & \thd{miễn phí} \\
\midrule
\endhead
\bottomrule
\endfoot
\href{https://github.com/DataTalksClub/data-engineering-zoomcamp}{DE Zoomcamp (DataTalksClub)} & khoá data engineering mở, có workshop nạp dữ liệu với dlt & toàn bộ tài liệu trên GitHub \\
\href{https://developer.confluent.io/courses/}{Confluent Developer: courses} & các khoá về Kafka, Kafka Connect và Schema Registry; học sau khi làm lab Debezium & khoá tự học trên trang \\
\end{xltabular}

\section{Lab thực hành}

\subsection{Lab 8.1: dựng nguồn và nạp toàn bộ}
\begin{enumerate}
  \item Nạp \code{olist\_orders\_dataset.csv} vào PostgreSQL thành bảng \code{src.orders}, thêm cột \code{updated\_at timestamptz} và đặt bằng \code{order\_purchase\_timestamp} cho mọi dòng hiện có.
  \item Để giả lập nguồn đang sống, chỉ giữ lại các đơn tới hết tháng 06/2017; giữ phần còn lại trong một bảng chờ \code{src.orders\_future}.
  \item Từ Spark, đọc \code{src.orders} qua JDBC (driver PostgreSQL phải có trên classpath của Spark server) và ghi nguyên bản vào \code{lake.raw.orders} bằng overwrite.
  \item Chạy lại bước 3 hai lần; so số dòng giữa nguồn và đích.
\end{enumerate}
\hoanthanh{số dòng của \code{lake.raw.orders} bằng \code{src.orders} sau mỗi lần chạy, và bạn giải thích được vì sao nạp toàn bộ bằng overwrite là idempotent nhưng tốn kém khi bảng lớn.}

\subsection{Lab 8.2: nạp tăng dần với high-water mark và MERGE}
Tạo bảng trạng thái \code{lake.ops.load\_state(source STRING, hwm TIMESTAMP)}. Mỗi lần chạy làm đúng thứ tự sau:
\begin{enumerate}
  \item Đọc \code{hwm} của \code{src.orders}; chốt cận trên \code{hi} bằng thời điểm bắt đầu run.
  \item Đọc từ PostgreSQL các dòng có \code{order\_purchase\_timestamp > hwm} và \code{<= hi}.
  \item Khử trùng theo \code{order\_id}, rồi \code{MERGE} vào \code{lake.staging.orders} như dưới.
  \item Chỉ sau khi \code{MERGE} commit xong mới ghi \code{hwm} mới vào \code{lake.ops.load\_state}.
\end{enumerate}
\begin{codeblock}
MERGE INTO lake.staging.orders t
USING batch_orders s
ON t.order_id = s.order_id
WHEN MATCHED AND s.updated_at > t.updated_at THEN UPDATE SET *
WHEN NOT MATCHED THEN INSERT *
\end{codeblock}
Chuyển một tháng từ \code{src.orders\_future} sang \code{src.orders}, chạy pipeline; lặp lại ba tháng. Sau đó chạy lại cùng một lần chạy hai lần, và giả lập lỗi bằng cách dừng job giữa bước 3 và bước 4.
\hoanthanh{sau mọi lần chạy lại và lần lỗi giả lập, số dòng và tổng \code{order\_id} khác nhau của \code{lake.staging.orders} khớp với \code{src.orders}; bạn giải thích được vì sao việc \code{MERGE} và ghi \code{hwm} nằm ở hai commit khác nhau vẫn an toàn.}

\subsection{Lab 8.3: dữ liệu tới muộn và dòng bị sửa}
\begin{enumerate}
  \item Chèn vào \code{src.orders} một đơn có \code{order\_purchase\_timestamp} cũ hơn \code{hwm} hiện tại (giả lập dòng tới muộn). Chạy pipeline và xác nhận nó bị bỏ sót.
  \item Đổi \code{order\_status} của 50 đơn cũ từ \code{shipped} sang \code{delivered} và đặt \code{updated\_at = now()}. Chạy pipeline và xác nhận cursor theo thời điểm mua không thấy thay đổi.
  \item Sửa pipeline: đổi cursor sang \code{updated\_at} và thêm cửa sổ nhìn lại 3 ngày (đọc từ \code{hwm - 3 days}). Chạy lại và kiểm cả hai trường hợp đã được bắt.
  \item Xoá cứng một đơn ở nguồn; ghi lại rằng không cursor nào bắt được việc này, và nêu cách bạn sẽ xử lý (đối soát toàn bộ định kỳ, xoá mềm ở nguồn, hoặc CDC).
\end{enumerate}
\hoanthanh{bạn có bảng ba cột: thay đổi ở nguồn, cursor theo \code{order\_purchase\_timestamp} có bắt được không, cursor theo \code{updated\_at} cộng nhìn lại có bắt được không.}

\subsection{Lab 8.4 (tuỳ chọn): CDC với Debezium}
\begin{enumerate}
  \item Chạy \href{https://debezium.io/documentation/reference/stable/tutorial.html}{Debezium Tutorial} theo đúng image và phiên bản trong trang.
  \item Tạo, sửa, xoá một dòng; đọc các sự kiện và chỉ ra \code{before}, \code{after}, \code{op} và tombstone.
  \item So với Lab 8.3: liệt kê thay đổi nào CDC bắt được mà cursor không bắt được.
\end{enumerate}
\hoanthanh{bạn chỉ ra được sự kiện \code{op=d} và tombstone đi sau nó cho cùng một khoá.}

\section{Kế hoạch theo tuần}
\begin{xltabular}{\textwidth}{@{}P{1.4cm} L P{1.4cm}@{}}
\toprule
\thd{tuần} & \thd{việc} & \thd{giờ} \\
\midrule
\endhead
\bottomrule
\endfoot
1 & Airbyte Sync Modes, dlt cursor và lag, Airflow Best Practices mục Creating a task, Iceberg Spark Writes, Spark JDBC; Lab 8.1 và 8.2 & 10 \\
2 & Debezium Architecture và PostgreSQL connector, Kafka Connect, Iceberg Evolution, Airflow Backfill; Lab 8.3; thử thách kỹ năng & 10 \\
thêm & Lab 8.4 với Debezium; khoá Kafka Connect trên Confluent Developer & 5 \\
\end{xltabular}

\begin{thuthach}
Xây pipeline tăng dần cho \code{order\_items} và \code{order\_payments} của Olist vào tầng staging. Hai bảng này không có cột thời gian riêng; bạn tự chọn cursor và khoá (gợi ý: khoá tổng hợp), tự quyết cách xử lý dữ liệu tới muộn và cách backfill lại một tháng bất kỳ. Pipeline phải chạy lại được nhiều lần mà không trùng dòng.
\hoanthanh{chạy toàn bộ pipeline ba lần liên tiếp và một lần backfill tháng 03/2018 cho cùng số dòng và cùng tổng \code{payment\_value} như khi tính trực tiếp trên CSV.}
\end{thuthach}

\section{Câu hỏi tự kiểm tra}
\begin{enumerate}
  \item Bạn ghi dữ liệu vào đích rồi job chết trước khi lưu cursor. Lần sau chuyện gì xảy ra, và vì sao chấp nhận được?
  \item Nếu đảo thứ tự, lưu cursor trước rồi mới ghi dữ liệu, rủi ro là gì?
  \item Kể ba loại thay đổi ở nguồn mà cursor trên \code{updated\_at} bỏ sót hoặc dễ bỏ sót.
  \item \code{INSERT INTO staging SELECT * FROM raw\_new} có idempotent không? Viết lại cho idempotent.
  \item Khi nào chọn \code{MERGE}, khi nào chọn overwrite partition động?
  \item Cửa sổ nhìn lại giải quyết vấn đề gì, và cái giá là gì?
  \item High-water mark của batch khác watermark của streaming thế nào?
  \item CDC dựa trên log cho bạn gì mà \code{SELECT ... WHERE updated\_at > ?} không cho?
  \item Vì sao \code{REPLICA IDENTITY} ảnh hưởng tới trường \code{before} trong sự kiện Debezium?
  \item Nguồn thêm một cột và đổi tên một cột khác. Tầng raw, staging và marts nên phản ứng thế nào?
\end{enumerate}

\section{Điểm dễ nhầm}
\begin{xltabular}{\textwidth}{@{}P{4.2cm} L P{3.4cm}@{}}
\toprule
\thd{điểm dễ nhầm} & \thd{thực tế} & \thd{docs chính thức} \\
\midrule
\endhead
\bottomrule
\endfoot
Nhiều người nghĩ idempotent nghĩa là nguồn không có bản trùng & Idempotent nghĩa là chạy lại cho cùng kết quả; nguồn vẫn có thể gửi lại dữ liệu, đích phải chịu được & \href{https://airflow.apache.org/docs/apache-airflow/stable/best-practices.html}{Airflow Best Practices} \\
Nạp tăng dần theo \code{updated\_at} bắt mọi thay đổi & Dòng bị xoá cứng không bao giờ xuất hiện trong truy vấn; CDC dựa trên log thấy được sự kiện xoá & \href{https://debezium.io/documentation/reference/stable/connectors/postgresql.html}{Debezium PostgreSQL} \\
Retry làm pipeline an toàn & Retry chỉ an toàn khi task idempotent; nếu không, mỗi lần retry có thể thêm bản trùng & \href{https://airflow.apache.org/docs/apache-airflow/stable/best-practices.html}{Airflow Best Practices} \\
\code{MERGE INTO} tự khử trùng dữ liệu nguồn & Nhiều dòng nguồn khớp một dòng đích thì câu lệnh báo lỗi; phải khử trùng nguồn trước & \href{https://iceberg.apache.org/docs/latest/spark-writes/}{Iceberg Spark Writes} \\
Khử trùng ở biên cursor luôn bật & Trong dlt, khử trùng ở biên tắt khi dùng \code{lag} hoặc khi đặt \code{primary\_key} rỗng & \href{https://dlthub.com/docs/general-usage/incremental/cursor}{dlt cursor} \\
Backfill chỉ là chạy lại DAG & Backfill chỉ đúng khi mỗi lần chạy ghi đè đúng partition của nó và không phụ thuộc thời điểm hiện tại & \href{https://airflow.apache.org/docs/apache-airflow/stable/best-practices.html}{Airflow Best Practices} \\
Debezium tự giữ dữ liệu cũ của dòng bị sửa & Trường \code{before} đầy đủ chỉ khi bảng đặt \code{REPLICA IDENTITY FULL} & \href{https://debezium.io/documentation/reference/stable/connectors/postgresql.html}{Debezium PostgreSQL} \\
\end{xltabular}


% ============================================================

\modulechapter{9}{Orchestration với Apache Airflow}
\moduleinfo{Trung cấp}{20 giờ}{Module 8}

Pipeline ở Module 8 cần chạy theo lịch, theo đúng thứ tự, có retry, có giới hạn đồng thời, và báo khi hỏng. Đó là việc của orchestrator. Apache Airflow lên lịch và theo dõi các task; phần tính toán nặng vẫn chạy ở Spark. Module này dạy Airflow 3, bản có nhiều thay đổi phá vỡ so với tutorial Airflow 2 trên mạng. Học xong, bạn có ba DAG nối raw, staging và marts của Olist bằng asset, có pool giới hạn số job ghi Iceberg cùng lúc.

\section{Bạn sẽ học được}
\begin{itemize}
  \item Mô tả các thành phần của Airflow 3 và việc mỗi thành phần làm.
  \item Viết DAG bằng TaskFlow với import từ \code{airflow.sdk}.
  \item Giải thích \code{logical\_date} và data interval dưới \code{CronTrigger\allowbreak Timetable} mặc định, và chọn timetable phù hợp.
  \item Dùng retry, pool và \code{max\_active\_runs} để chạy an toàn.
  \item Nối DAG bằng asset thay vì đoán giờ chạy.
  \item Viết test cho DAG và phát hiện code kiểu Airflow 2 đã hỏng.
\end{itemize}

\section{Nội dung chính}
\begin{enumerate}
  \item \textbf{Orchestrator làm gì.} Airflow lên lịch, chạy theo phụ thuộc, retry và ghi log. Docs của Astronomer khuyên chuyển job xử lý dữ liệu lớn sang framework chuyên dụng như Spark, thay vì chạy trong worker của Airflow (\href{https://www.astronomer.io/docs/airflow/dag-best-practices}{Astronomer: DAG best practices}).
  \item \textbf{Kiến trúc Airflow 3.} Scheduler, \code{dag-processor}, \code{api-server}, triggerer, executor với worker, và metadata DB. Từ Airflow 3, \code{airflow api-server} thay webserver, \code{airflow dag-processor} phải chạy riêng kể cả ở local, và code của task không được dùng session metadata DB trực tiếp (\href{https://airflow.apache.org/docs/apache-airflow/stable/installation/upgrading_to_airflow3.html}{Upgrading to Airflow 3}).
  \item \textbf{DAG, task, operator, TaskFlow.} DAG là tập task và phụ thuộc. TaskFlow viết task bằng hàm Python với \code{@dag} và \code{@task}. Import mới nằm ở \code{airflow.sdk}; đường import cũ bị deprecate và sẽ bị bỏ (\href{https://airflow.apache.org/docs/apache-airflow/stable/core-concepts/dags.html}{Dags}, \href{https://airflow.apache.org/docs/apache-airflow/stable/tutorial/taskflow.html}{TaskFlow tutorial}).
  \item \textbf{Lịch chạy với chuỗi cron.} Trong Airflow 3, \code{create\_cron\_\allowbreak data\_intervals} mặc định là \code{False}, nên chuỗi cron trần dùng \code{CronTrigger\allowbreak Timetable}: run chạy đúng giờ cron, \code{logical\_date} là thời điểm kích hoạt, \code{data\_interval\_start} và \code{data\_interval\_end} đều bằng thời điểm đó. \code{CronData\allowbreak IntervalTimetable} (mặc định của Airflow 2) tạo khoảng giữa hai mốc cron và chạy ở cuối khoảng. Ví dụ \code{0 0 * * *} bật lúc 15 giờ ngày 31/01: timetable trigger tạo run lúc nửa đêm 01/02; timetable data interval tạo run ngay cho khoảng bắt đầu nửa đêm 31/01 (\href{https://airflow.apache.org/docs/apache-airflow/stable/authoring-and-scheduling/timetable.html}{Timetables}).
  \item \textbf{Ghim timetable.} Đổi cấu hình \code{create\_cron\_\allowbreak data\_intervals} làm đổi mọi DAG dùng chuỗi cron trần; nếu task cần data interval, khai báo tường minh một instance timetable. \code{CronTrigger\allowbreak Timetable} có tham số \code{interval} để có data interval kết thúc tại giờ kích hoạt (\href{https://airflow.apache.org/docs/apache-airflow/stable/authoring-and-scheduling/timetable.html}{Timetables}).
  \item \textbf{Catchup và run trigger tay.} \code{catchup\_by\_default} giờ là \code{False}. Với run trigger tay, data interval có thể không khớp \code{logical\_date} bạn nhập; cần ngày người dùng chọn thì đọc thẳng \code{logical\_date} (\href{https://airflow.apache.org/docs/apache-airflow/stable/installation/upgrading_to_airflow3.html}{Upgrading to Airflow 3}).
  \item \textbf{Context và template.} Các khoá \code{execution\_date}, \code{prev\_ds}, \code{next\_ds}, \code{yesterday\_ds}, \code{tomorrow\_ds} đã bị bỏ; \code{xcom\_pull()} không truyền \code{task\_ids} giờ chỉ kéo từ task hiện tại (\href{https://airflow.apache.org/docs/apache-airflow/stable/installation/upgrading_to_airflow3.html}{Upgrading to Airflow 3}).
  \item \textbf{Retry và idempotency.} Airflow có thể retry task khi lỗi, nên task phải cho cùng kết quả mỗi lần chạy lại (\href{https://airflow.apache.org/docs/apache-airflow/stable/best-practices.html}{Best Practices}). Astronomer gợi ý \code{retries=2} cho phần lớn lỗi thường gặp của hệ phân tán, và mỗi task chỉ làm một việc chạy lại được độc lập (\href{https://www.astronomer.io/docs/airflow/dag-best-practices}{Astronomer}).
  \item \textbf{Đồng thời: pool.} Pool giới hạn số task chạy song song trên một tập task tuỳ ý. Task không gán pool rơi vào \code{default\_pool} có 128 slot; mỗi task chiếm 1 slot, đổi bằng \code{pool\_slots}; pool có tuỳ chọn tính cả task deferred (\href{https://airflow.apache.org/docs/apache-airflow/stable/administration-and-deployment/pools.html}{Pools}). Kết hợp với \code{max\_active\_runs} của DAG để hai run không chồng nhau.
  \item \textbf{XCom.} XCom dành cho metadata nhỏ giữa các task, không phải để truyền DataFrame; trạng thái bền như high-water mark nên nằm ở bảng của bạn (\href{https://airflow.apache.org/docs/apache-airflow/stable/core-concepts/xcoms.html}{XComs}).
  \item \textbf{Sensor và deferrable operator.} Sensor chờ một điều kiện. Operator deferrable nhường slot worker trong lúc chờ và giao việc chờ cho triggerer (\href{https://airflow.apache.org/docs/apache-airflow/stable/authoring-and-scheduling/deferring.html}{Deferrable Operators \& Triggers}).
  \item \textbf{Asset.} Asset là nhóm dữ liệu logic định danh bằng URI. Task khai báo \code{outlets} để phát sự kiện cho asset; DAG khác đặt \code{schedule} theo asset để chạy khi asset được cập nhật. Airflow không kiểm nội dung hay vị trí dữ liệu, URI chỉ là chuỗi. Decorator \code{@asset} tạo asset, DAG và task cùng lúc (\href{https://airflow.apache.org/docs/apache-airflow/stable/authoring-and-scheduling/assets.html}{Assets}).
  \item \textbf{Deadline Alerts thay SLA.} SLA đã bị bỏ trong Airflow 3. Deadline Alerts có từ 3.1, còn ở mức experimental, cần ba tham số \code{reference}, \code{interval}, \code{callback}; ví dụ reference là \code{DAGRUN\_QUEUED\_AT}, \code{DAGRUN\_LOGICAL\_DATE}, \code{FIXED\_DATETIME} (\href{https://airflow.apache.org/docs/apache-airflow/stable/howto/deadline-alerts.html}{Deadline Alerts}).
  \item \textbf{REST API.} \code{/api/v1} được thay bằng \code{/api/v2} (\href{https://airflow.apache.org/docs/apache-airflow/stable/installation/upgrading_to_airflow3.html}{Upgrading to Airflow 3}).
  \item \textbf{Test DAG.} Docs chia thành Dag Loader Test, unit test, self-check và môi trường staging; tránh code nặng ở mức đầu file vì nó chạy mỗi lần DAG được parse (\href{https://airflow.apache.org/docs/apache-airflow/stable/best-practices.html}{Best Practices}). Ruff 0.13.1 trở lên với rule \code{AIR301}, \code{AIR302} bắt các thay đổi phá vỡ của Airflow 3.
  \item \textbf{Triển khai.} Docker Compose chính thức chỉ để học và khám phá, cần ít nhất 4 GB bộ nhớ cho Docker, tốt nhất 8 GB; production dùng \href{https://airflow.apache.org/docs/helm-chart/stable/index.html}{Helm Chart} (\href{https://airflow.apache.org/docs/apache-airflow/stable/howto/docker-compose/index.html}{Running Airflow in Docker}).
\end{enumerate}

\section{Bắt đầu}
\begin{itemize}
  \item \href{https://airflow.apache.org/docs/apache-airflow/stable/howto/docker-compose/index.html}{Running Airflow in Docker}: dựng Airflow 3 bằng Docker Compose chính thức.
  \item \href{https://airflow.apache.org/docs/apache-airflow/stable/tutorial/taskflow.html}{TaskFlow tutorial}: viết DAG đầu tiên.
  \item \href{https://airflow.apache.org/docs/apache-airflow/stable/installation/upgrading_to_airflow3.html}{Upgrading to Airflow 3}: đọc trước khi mở bất kỳ tutorial Airflow 2 nào.
\end{itemize}

\section{Tài liệu chính thức}
\begin{xltabular}{\textwidth}{@{}P{4.6cm} L P{1.8cm}@{}}
\toprule
\thd{tài liệu} & \thd{nội dung cần đọc} & \thd{thời lượng} \\
\midrule
\endhead
\bottomrule
\endfoot
\href{https://airflow.apache.org/docs/apache-airflow/stable/core-concepts/index.html}{Airflow: Core Concepts} & Dags, Dag Runs, Tasks, Operators, XComs & 3 giờ \\
\href{https://airflow.apache.org/docs/apache-airflow/stable/installation/upgrading_to_airflow3.html}{Upgrading to Airflow 3} & toàn trang & 1 giờ \\
\href{https://airflow.apache.org/docs/apache-airflow/stable/authoring-and-scheduling/timetable.html}{Timetables}, \href{https://airflow.apache.org/docs/apache-airflow/stable/authoring-and-scheduling/cron.html}{Cron} & so sánh hai timetable cron, \code{interval}, \code{MultipleCron\allowbreak TriggerTimetable} & 1 giờ \\
\href{https://airflow.apache.org/docs/apache-airflow/stable/authoring-and-scheduling/assets.html}{Assets} & outlets, lịch theo asset, \code{@asset} & 1 giờ \\
\href{https://airflow.apache.org/docs/apache-airflow/stable/administration-and-deployment/pools.html}{Pools} & \code{default\_pool}, \code{pool}, \code{pool\_slots} & 0,5 giờ \\
\href{https://airflow.apache.org/docs/apache-airflow/stable/authoring-and-scheduling/deferring.html}{Deferrable Operators \& Triggers} & triggerer, vì sao tiết kiệm slot worker & 0,5 giờ \\
\href{https://airflow.apache.org/docs/apache-airflow/stable/howto/deadline-alerts.html}{Deadline Alerts} & reference, interval, callback & 0,5 giờ \\
\href{https://airflow.apache.org/docs/apache-airflow/stable/best-practices.html}{Best Practices} & Creating a task, Top level Python Code, Testing a Dag & 1,5 giờ \\
\href{https://airflow.apache.org/docs/apache-airflow/stable/core-concepts/backfill.html}{Backfill} & tạo backfill, chạy thử & 0,5 giờ \\
\href{https://www.astronomer.io/docs/airflow/dag-best-practices}{Astronomer: DAG best practices} & Review idempotency, Set retries, Keep tasks atomic, Avoid top-level code & 1 giờ \\
\end{xltabular}

\section{Khoá học miễn phí}
\begin{xltabular}{\textwidth}{@{}P{4.2cm} L P{2.6cm}@{}}
\toprule
\thd{khoá} & \thd{nội dung} & \thd{miễn phí} \\
\midrule
\endhead
\bottomrule
\endfoot
\href{https://academy.astronomer.io}{Astronomer Academy} & các module ngắn như Airflow: Taskflow API, Airflow: Dynamic Task Mapping, Airflow: Dynamic DAGs; lộ trình Airflow 101 (Airflow 3) & các module có nhãn Free; kiểm giá của lộ trình khi đăng ký \\
\href{https://academy.astronomer.io}{Chứng chỉ Astronomer cho Airflow 3} & kỳ thi chứng chỉ Airflow; ôn bằng Airflow 101 & không miễn phí: 150 USD mỗi lần thi theo trang Academy ngày 02/11/2026 \\
\end{xltabular}

\section{Lab thực hành}

\subsection{Lab 9.1: Airflow local và ngữ nghĩa lịch chạy}
\begin{enumerate}
  \item Dựng Airflow theo Docker Compose chính thức, cấp ít nhất 4 GB (tốt nhất 8 GB) bộ nhớ cho Docker; nối nó vào cùng network với Spark của stack tham chiếu.
  \item Viết DAG \code{schedule\_probe} với \code{schedule="*/10 * * * *"}; task in ra \code{logical\_date}, \code{data\_interval\_start}, \code{data\_interval\_end}.
  \item Viết bản thứ hai với \code{schedule=\allowbreak CronData\allowbreak IntervalTimetable(\allowbreak "*/10 * * * *", \allowbreak timezone="UTC")}, rồi so log của hai DAG. Trigger tay một run và so thêm lần nữa.
  \item Chạy \code{ruff check dags/ -{}-select AIR301,AIR302} và sửa mọi cảnh báo.
\end{enumerate}
\hoanthanh{bạn có bảng so sánh ba giá trị thời gian của hai DAG ở run theo lịch và run trigger tay, và giải thích được khác biệt.}

\subsection{Lab 9.2: ba DAG Olist nối bằng asset}
Viết ba DAG: \code{olist\_raw} nạp CSV và bảng PostgreSQL vào tầng raw; \code{olist\_staging} chạy pipeline tăng dần của Module 8; \code{olist\_marts} tính mart doanh thu. DAG đầu chạy theo cron, hai DAG sau chạy theo asset. Task gọi Spark qua Spark Connect, Airflow chỉ điều phối.
\begin{codeblock}
import pendulum
from airflow.sdk import Asset, dag, task

RAW_ORDERS = Asset("s3://lakehouse/raw/olist_orders")
STG_ORDERS = Asset("s3://lakehouse/staging/olist_orders")

@dag(schedule="0 2 * * *", start_date=pendulum.datetime(2026, 1, 1, tz="UTC"),
     catchup=False, max_active_runs=1, tags=["olist"])
def olist_raw():
    @task(outlets=[RAW_ORDERS], pool="lakehouse_writers", retries=2)
    def load_raw_orders():
        ...  # mở SparkSession.builder.remote(...) và ghi lake.raw.orders
    load_raw_orders()

@dag(schedule=[RAW_ORDERS], start_date=pendulum.datetime(2026, 1, 1, tz="UTC"),
     catchup=False, max_active_runs=1, tags=["olist"])
def olist_staging():
    @task(outlets=[STG_ORDERS], pool="lakehouse_writers", retries=2)
    def merge_orders():
        ...  # pipeline tăng dần + MERGE của Lab 8.2
    merge_orders()

olist_raw()
olist_staging()
\end{codeblock}
\hoanthanh{một run của \code{olist\_raw} tự kích hoạt \code{olist\_staging} rồi \code{olist\_marts}, thấy được trong giao diện asset, và cả ba DAG xanh.}

\subsection{Lab 9.3: pool, retry và chạy lại an toàn}
\begin{enumerate}
  \item Tạo pool \code{lakehouse\_writers} 1 slot trong giao diện (Admin, Pools) và gán mọi task ghi Iceberg vào pool này.
  \item Trigger tay \code{olist\_raw} và một job bảo trì Iceberg (ví dụ task gọi \code{rewrite\_data\_files}) cùng lúc; xác nhận trên giao diện rằng chỉ một task ghi chạy tại một thời điểm.
  \item Cho task \code{merge\_orders} lỗi ngẫu nhiên sau khi \code{MERGE} commit nhưng trước khi ghi \code{hwm}; để retry chạy và kiểm không có dòng trùng.
  \item Tạo backfill cho \code{olist\_raw} trên ba ngày bằng cơ chế backfill của Airflow 3, và xác nhận số dòng không nhân đôi.
\end{enumerate}
\hoanthanh{log cho thấy các task ghi chạy nối tiếp, và sau retry cùng backfill số dòng staging vẫn khớp nguồn.}

\subsection{Lab 9.4: test DAG}
\begin{enumerate}
  \item Viết test pytest kiểm mỗi DAG: import không lỗi, có \code{tags}, mọi task có \code{retries} và task ghi Iceberg nằm trong pool \code{lakehouse\_writers}.
  \item Chạy \code{dag.test()} cho \code{olist\_marts} ở local.
  \item Thêm một DAG cố ý dùng \code{execution\_date} và import cũ; xác nhận Ruff với \code{AIR301}, \code{AIR302} bắt được.
\end{enumerate}
\hoanthanh{lệnh pytest và Ruff chạy xanh trên ba DAG thật và đỏ trên DAG cố ý sai.}

\section{Kế hoạch theo tuần}
\begin{xltabular}{\textwidth}{@{}P{1.4cm} L P{1.4cm}@{}}
\toprule
\thd{tuần} & \thd{việc} & \thd{giờ} \\
\midrule
\endhead
\bottomrule
\endfoot
1 & Upgrading to Airflow 3, Core Concepts, TaskFlow tutorial, Timetables, Docker Compose; Lab 9.1 và 9.2 & 10 \\
2 & Assets, Pools, Deferring, Deadline Alerts, Backfill, Best Practices, Astronomer best practices; Lab 9.3 và 9.4; thử thách kỹ năng & 10 \\
thêm & các module Free của Astronomer Academy; lộ trình Airflow 101 nếu ôn chứng chỉ & 6 \\
\end{xltabular}

\begin{thuthach}
Thêm vào hệ ba DAG một nhánh cho \code{order\_reviews}: nạp vào raw, \code{MERGE} vào staging, rồi cập nhật một mart điểm đánh giá trung bình theo bang. Mart chỉ được chạy khi cả staging orders và staging reviews đều đã cập nhật. Thêm một Deadline Alert báo khi DAG marts chưa xong sau 2 giờ kể từ khi run được xếp hàng. Không được dùng sensor chờ theo giờ.
\hoanthanh{một run trọn vẹn từ raw tới mart reviews chạy tự động theo asset, cảnh báo deadline kích hoạt được khi bạn cố ý làm chậm một task, và test của Lab 9.4 vẫn xanh.}
\end{thuthach}

\section{Câu hỏi tự kiểm tra}
\begin{enumerate}
  \item Kể các thành phần của Airflow 3 và việc của từng thành phần.
  \item Với \code{schedule="0 0 * * *"} trong Airflow 3, \code{logical\_date} và \code{data\_interval\_start} là gì? Lấy lại hành vi Airflow 2 bằng cách nào?
  \item Vì sao nên khai báo timetable tường minh thay vì dựa vào \code{create\_cron\_\allowbreak data\_intervals}?
  \item Vì sao task phải idempotent trước khi bật retry?
  \item Hai cấu hình nào bảo đảm hai job ghi cùng một bảng Iceberg không chạy chồng nhau?
  \item Cái gì nên để trong XCom, cái gì không?
  \item Triggerer làm gì, và vì sao operator deferrable tiết kiệm slot worker?
  \item Asset có di chuyển hay kiểm tra dữ liệu không? Nó cho bạn gì?
  \item Vì sao code nặng ở mức đầu file DAG gây hại?
  \item Thứ gì thay SLA trong Airflow 3, và cần những tham số nào?
\end{enumerate}

\section{Điểm dễ nhầm}
\begin{xltabular}{\textwidth}{@{}P{4.2cm} L P{3.4cm}@{}}
\toprule
\thd{điểm dễ nhầm} & \thd{thực tế} & \thd{docs chính thức} \\
\midrule
\endhead
\bottomrule
\endfoot
Nhiều người nghĩ Airflow xử lý dữ liệu & Airflow điều phối; job xử lý lớn nên chạy ở framework chuyên dụng & \href{https://www.astronomer.io/docs/airflow/dag-best-practices}{Astronomer} \\
\code{logical\_date} luôn là đầu một data interval như Airflow 2 & Với chuỗi cron trần, Airflow 3 dùng \code{CronTrigger\allowbreak Timetable}: \code{logical\_date} là thời điểm kích hoạt và data interval có độ dài bằng 0 & \href{https://airflow.apache.org/docs/apache-airflow/stable/authoring-and-scheduling/timetable.html}{Timetables} \\
Catchup bật mặc định & \code{catchup\_by\_default} là \code{False} trong Airflow 3 & \href{https://airflow.apache.org/docs/apache-airflow/stable/installation/upgrading_to_airflow3.html}{Upgrading to Airflow 3} \\
Tutorial Airflow 2 vẫn chạy & Khoá context như \code{execution\_date}, SLA, \code{/api/v1} và nhiều đường import đã bị bỏ hoặc deprecate & \href{https://airflow.apache.org/docs/apache-airflow/stable/installation/upgrading_to_airflow3.html}{Upgrading to Airflow 3} \\
XCom truyền được DataFrame & XCom dành cho metadata nhỏ & \href{https://airflow.apache.org/docs/apache-airflow/stable/core-concepts/xcoms.html}{XComs} \\
Asset kiểm dữ liệu đã thật sự có & Airflow không giả định gì về nội dung hay vị trí dữ liệu; asset là tín hiệu lên lịch & \href{https://airflow.apache.org/docs/apache-airflow/stable/authoring-and-scheduling/assets.html}{Assets} \\
Docker Compose dùng được cho production & Docs ghi nó chỉ để học và khám phá; production dùng Helm Chart & \href{https://airflow.apache.org/docs/apache-airflow/stable/howto/docker-compose/index.html}{Running Airflow in Docker} \\
\end{xltabular}


% ============================================================

\modulechapter{10}{Chất lượng dữ liệu, test và DataOps}
\moduleinfo{Trung cấp}{20 giờ}{Module 8, 9}

Một pipeline chạy xanh vẫn có thể ghi dữ liệu sai. Module này thêm các lớp phòng vệ: unit test cho logic biến đổi, data test cho dữ liệu thật, contract cho hình dạng bảng, và write-audit-publish (WAP) để dữ liệu chưa kiểm không bao giờ tới người đọc. Bên cạnh đó là các thói quen DataOps: Git, pre-commit và CI. Học xong, bạn có cổng chất lượng chặn marts khi staging sai, và một CI tự chạy test mỗi lần đẩy code.

\section{Bạn sẽ học được}
\begin{itemize}
  \item Phân biệt unit test, data test, contract và test đầu cuối, và đặt mỗi loại đúng chỗ.
  \item Viết các kiểm tra unique, not null, accepted values và relationships cho bảng core.
  \item Làm WAP trên một branch Iceberg và publish bằng \code{fast\_forward}.
  \item Viết unit test PySpark bằng pytest và chạy chúng trong GitHub Actions.
  \item Chọn công cụ chất lượng dữ liệu, có xét license và tình trạng dự án năm 2026.
  \item Mô tả một data contract theo Open Data Contract Standard.
\end{itemize}

\section{Nội dung chính}

\subsection{Các loại kiểm tra}
\begin{enumerate}
  \item \textbf{Bốn loại kiểm tra.} Unit test kiểm logic trên đầu vào tĩnh, trước khi build. Data test kiểm nội dung dữ liệu thật sau khi build. Contract kiểm hình dạng đầu ra (tên cột, kiểu, ràng buộc) ngay lúc build. Test đầu cuối chạy cả pipeline trên môi trường giống thật (\href{https://docs.getdbt.com/docs/build/unit-tests}{dbt unit tests}, \href{https://docs.getdbt.com/docs/mesh/govern/model-contracts}{dbt model contracts}).
  \item \textbf{Các chiều chất lượng.} GX chia use case thành Distribution, Freshness, Integrity, Missingness, Schema, Uniqueness, Volume và Unstructured; dùng làm từ vựng khi thiết kế kiểm tra (\href{https://docs.greatexpectations.io/docs/reference/learn/data_quality_use_cases/dq_use_cases_lp/}{GX: data quality use cases}).
  \item \textbf{Data test kiểu dbt.} dbt có bốn generic test dựng sẵn: \code{unique}, \code{not\_null}, \code{accepted\_values}, \code{relationships}; khai báo dưới khoá \code{data\_tests}. Ngoài ra có singular test (một câu SQL trả về dòng lỗi), mức \code{severity}, và \code{store\_failures} để lưu dòng lỗi (\href{https://docs.getdbt.com/docs/build/data-tests}{dbt data tests}). Dù không dùng dbt, đây là mô hình tốt: mỗi test là một truy vấn trả về dòng vi phạm, rỗng nghĩa là qua.
  \item \textbf{Unit test của dbt.} Có từ dbt 1.8; dbt Labs khuyên chỉ chạy unit test ở môi trường dev và CI, vì đầu vào là tĩnh, không cần chạy lại trên production (\href{https://docs.getdbt.com/docs/build/unit-tests}{dbt unit tests}).
  \item \textbf{Model contract của dbt.} Đặt \code{contract: enforced: true}, rồi khai báo \code{name} và \code{data\_type} cho mọi cột. dbt chạy preflight check so tên và kiểu cột; không khớp thì model không build. Ràng buộc phụ thuộc nền tảng: phần lớn chỉ thực thi \code{not\_null}, còn \code{primary\_key}, \code{unique}, \code{foreign\_key} chỉ được khai báo. Contract hỗ trợ \code{table}, \code{view} và \code{incremental} (với \code{on\_schema\_change} là \code{append\_new\_columns} hoặc \code{fail}), không hỗ trợ \code{ephemeral} hay Python model (\href{https://docs.getdbt.com/docs/mesh/govern/model-contracts}{dbt model contracts}).
\end{enumerate}

\subsection{Write-audit-publish trên Iceberg}
\begin{enumerate}[resume]
  \item \textbf{WAP.} Ghi dữ liệu mới vào một branch, kiểm trên branch đó, rồi mới publish lên \code{main}. Trước khi publish, người đọc \code{main} không thấy dữ liệu của branch (\href{https://iceberg.apache.org/docs/latest/branching/}{Iceberg Branching}).
  \item \textbf{Cách làm trong Spark.} Bật thuộc tính bảng \code{write.wap.enabled}, tạo branch bằng \code{ALTER TABLE ... CREATE BRANCH}, đặt \code{spark.wap.branch} rồi ghi vào tên bảng gốc; branch phải có sẵn, và không được vừa đặt \code{spark.wap.branch} vừa dùng định danh \code{branch\_<tên>} trong cùng lệnh ghi (\href{https://iceberg.apache.org/docs/latest/spark-writes/}{Iceberg Spark Writes}, \href{https://iceberg.apache.org/docs/latest/spark-ddl/}{Iceberg Spark DDL}).
  \item \textbf{Publish.} \code{CALL <catalog>.system.\allowbreak fast\_forward('<bảng>', 'main', '<branch>')} đưa \code{main} tới snapshot mới nhất của branch; \code{cherrypick\_snapshot} là cách khác để publish một snapshot cụ thể (\href{https://iceberg.apache.org/docs/latest/spark-procedures/}{Iceberg Spark Procedures}).
\end{enumerate}

\subsection{Công cụ chất lượng dữ liệu năm 2026}
\begin{enumerate}[resume]
  \item \textbf{GX Core.} Thư viện Python mã nguồn mở, license Apache-2.0. Ngày 06/05/2026, Great Expectations thông báo FICO mua GX Cloud, còn Fivetran trở thành bên chăm sóc cộng đồng mã nguồn mở và dự án GX Core (\href{https://www.greatexpectations.io/blog/an-update-from-great-expectations}{GX blog}, \href{https://www.fivetran.com/press/fivetran-to-become-steward-of-the-great-expectations-open-source-community-and-gx-core-project}{thông cáo của Fivetran}). Mã nguồn và bản phát hành nằm ở \href{https://github.com/fivetran/great_expectations/releases}{\code{fivetran/\allowbreak great\_expectations}}.
  \item \textbf{Soda Core.} Soda đổi license Soda Core từ Apache 2.0 sang Elastic License 2.0 (bài ngày 27/01/2026). Theo Soda, bạn vẫn dùng tự do trong tổ chức, kể cả production; điều bị cấm là cung cấp Soda Core cho bên thứ ba dưới dạng dịch vụ hosted hoặc managed (\href{https://soda.io/blog/soda-core-license-update-moving-to-elastic-license}{Soda: license update}). Kiểm license của đúng phiên bản bạn cài.
  \item \textbf{Bắt đầu rẻ nhất.} Với lab này, kiểm tra bằng Spark SQL ngay trong WAP là đủ; thêm thư viện khi số kiểm tra tăng hoặc cần báo cáo đẹp hơn.
\end{enumerate}

\subsection{DataOps: Git, pre-commit, CI}
\begin{enumerate}[resume]
  \item \textbf{Git.} Commit nhỏ, branch cho mỗi thay đổi, merge qua pull request. Pro Git là sách chính thức trên git-scm.com, đọc miễn phí (\href{https://git-scm.com/book/en/v2}{Pro Git}).
  \item \textbf{Lint và pre-commit.} pre-commit chạy hook (như Ruff) trước mỗi commit để lỗi bị chặn sớm (\href{https://pre-commit.com/}{pre-commit}, \href{https://docs.astral.sh/ruff/}{Ruff}).
  \item \textbf{Unit test PySpark.} pytest dùng fixture để tạo \code{SparkSession} dùng chung (\href{https://docs.pytest.org/en/stable/how-to/fixtures.html}{pytest fixtures}); PySpark có \code{pyspark.testing.\allowbreak assertDataFrameEqual} để so DataFrame (\href{https://spark.apache.org/docs/latest/api/python/getting_started/testing_pyspark.html}{Testing PySpark}).
  \item \textbf{CI với GitHub Actions.} Workflow chạy trên mỗi push và pull request: cài Python, cache dependency, chạy Ruff và pytest (\href{https://docs.github.com/en/actions/tutorials/build-and-test-code/python}{GitHub Actions: Python}). Test cần nhiều dịch vụ Docker (Spark Connect, catalog, RustFS) thường chạy ở máy hoặc workflow chạy tay.
  \item \textbf{Data contract.} Open Data Contract Standard (ODCS) là chuẩn trung lập về nhà cung cấp để mô tả data contract bằng YAML: schema, quan hệ, quy tắc chất lượng, SLA, kênh hỗ trợ. ODCS thuộc dự án Bitol, ở mức incubation trong LF AI \& Data; v3.1.0 ra tháng 12/2025 (\href{https://bitol.io/bitol-announces-odcs-v3-1-0-stronger-smarter-and-stricter/}{Bitol: ODCS v3.1.0}, \href{https://bitol-io.github.io/open-data-contract-standard/v3.1.0/}{ODCS docs}).
  \item \textbf{Quan sát pipeline.} Airflow phát metric qua StatsD hoặc OpenTelemetry, ví dụ thời gian DAG run và số task lỗi; đó là điểm bắt đầu cho dashboard và cảnh báo (\href{https://airflow.apache.org/docs/apache-airflow/stable/administration-and-deployment/logging-monitoring/metrics.html}{Airflow Metrics}).
\end{enumerate}

\section{Bắt đầu}
\begin{itemize}
  \item \href{https://docs.getdbt.com/docs/build/data-tests}{dbt: Add data tests}: đọc để có mô hình bốn generic test.
  \item \href{https://spark.apache.org/docs/latest/api/python/getting_started/testing_pyspark.html}{Testing PySpark}: chạy ví dụ dùng pytest.
  \item \href{https://docs.github.com/en/actions/tutorials/build-and-test-code/python}{GitHub Actions: Building and testing Python}: tạo workflow đầu tiên.
\end{itemize}

\section{Tài liệu chính thức}
\begin{xltabular}{\textwidth}{@{}P{4.6cm} L P{1.8cm}@{}}
\toprule
\thd{tài liệu} & \thd{nội dung cần đọc} & \thd{thời lượng} \\
\midrule
\endhead
\bottomrule
\endfoot
\href{https://docs.getdbt.com/docs/build/data-tests}{dbt: data tests}, \href{https://docs.getdbt.com/reference/resource-properties/data-tests}{data\_tests property} & generic và singular test, \code{severity}, \code{store\_failures} & 1,5 giờ \\
\href{https://docs.getdbt.com/docs/build/unit-tests}{dbt: unit tests} & khi nào chạy, đầu vào tĩnh & 1 giờ \\
\href{https://docs.getdbt.com/docs/mesh/govern/model-contracts}{dbt: model contracts} & preflight check, ràng buộc theo nền tảng, contract và data test khác nhau & 1 giờ \\
\href{https://docs.greatexpectations.io/docs/reference/learn/data_quality_use_cases/dq_use_cases_lp/}{GX: data quality use cases} & tám use case & 1 giờ \\
\href{https://docs.greatexpectations.io/docs/core/introduction/try_gx/}{GX Core: Try GX} & quickstart với pandas & 1 giờ \\
\href{https://iceberg.apache.org/docs/latest/branching/}{Iceberg: Branching and Tagging} & mục Audit Branch & 0,5 giờ \\
\href{https://iceberg.apache.org/docs/latest/spark-procedures/}{Iceberg: Spark Procedures} & \code{fast\_forward}, \code{cherrypick\_snapshot} & 0,5 giờ \\
\href{https://git-scm.com/book/en/v2}{Pro Git} & chương 1 tới 3 & 5 giờ \\
\href{https://pre-commit.com/}{pre-commit} & cài đặt, file cấu hình, chạy trên mọi file & 0,5 giờ \\
\href{https://docs.pytest.org/en/stable/how-to/fixtures.html}{pytest: fixtures} & scope, yield fixture, \code{conftest.py} & 1 giờ \\
\href{https://spark.apache.org/docs/latest/api/python/getting_started/testing_pyspark.html}{Testing PySpark} & \code{assertDataFrameEqual}, dùng với pytest & 1 giờ \\
\href{https://docs.github.com/en/actions/tutorials/build-and-test-code/python}{GitHub Actions: Python} & chọn phiên bản Python, cache dependency, pytest, Ruff & 1 giờ \\
\href{https://bitol-io.github.io/open-data-contract-standard/v3.1.0/}{ODCS v3.1.0} & cấu trúc một contract, ví dụ YAML & 1 giờ \\
\end{xltabular}

\section{Khoá học miễn phí}
\begin{xltabular}{\textwidth}{@{}P{4.2cm} L P{2.6cm}@{}}
\toprule
\thd{khoá} & \thd{nội dung} & \thd{miễn phí} \\
\midrule
\endhead
\bottomrule
\endfoot
\href{https://missing.csail.mit.edu/2026/version-control/}{MIT Missing Semester 2026: Version Control and Git} & mô hình dữ liệu của Git, branch, remote; có bài tập & toàn bộ bài giảng và bài tập \\
\href{https://missing.csail.mit.edu/2026/code-quality/}{MIT Missing Semester 2026: Code Quality} & testing, pre-commit hook, continuous integration & toàn bộ \\
\end{xltabular}

\section{Lab thực hành}

\subsection{Lab 10.1: cổng chất lượng chặn marts}
\begin{enumerate}
  \item Viết một module Python chứa các kiểm tra dạng SQL, mỗi kiểm tra trả về dòng vi phạm:
  \begin{itemize}
    \item \code{unique} và \code{not\_null} cho \code{order\_id} của \code{staging.orders}, \code{customer\_id} của \code{staging.\allowbreak customers}, \code{product\_id} của \code{staging.\allowbreak products}, \code{seller\_id} của \code{staging.sellers}; cặp \code{(order\_id, \allowbreak order\_item\_id)} của \code{staging.order\_items}.
    \item \code{accepted\_values} cho \code{order\_status}: liệt kê giá trị khác nhau trong dữ liệu trước, rồi cố định danh sách.
    \item \code{relationships}: \code{order\_items.order\_id} thuộc \code{orders}, \code{order\_items.product\_id} thuộc \code{products}, \code{order\_items.seller\_id} thuộc \code{sellers}, \code{orders.customer\_id} thuộc \code{customers}.
  \end{itemize}
  \item Chạy các kiểm tra trên dữ liệu hiện có và ghi lại kết quả; nếu một khoá bạn tưởng là duy nhất lại có trùng (thử \code{review\_id} của reviews), quyết định đó là lỗi dữ liệu hay giả định sai.
  \item Thêm task \code{quality\_gate} vào DAG \code{olist\_staging} của Module 9, chạy sau \code{MERGE} và trước khi phát asset; có dòng vi phạm thì task fail.
  \item Cố ý chèn một dòng \code{order\_items} có \code{product\_id} không tồn tại; chạy DAG.
\end{enumerate}
\hoanthanh{task \code{quality\_gate} fail, asset staging không được phát, và \code{olist\_marts} không chạy; xoá dòng xấu thì toàn bộ chạy xanh.}

\subsection{Lab 10.2: write-audit-publish với branch Iceberg}
Làm WAP cho \code{lake.staging.orders} trong Spark SQL:
\begin{codeblock}
ALTER TABLE lake.staging.orders SET TBLPROPERTIES ('write.wap.enabled'='true');
ALTER TABLE lake.staging.orders CREATE BRANCH wap;

SET spark.wap.branch = wap;
MERGE INTO lake.staging.orders t USING batch_orders s
ON t.order_id = s.order_id
WHEN MATCHED THEN UPDATE SET *
WHEN NOT MATCHED THEN INSERT *;
RESET spark.wap.branch;

-- kiểm trên branch, so với main
SELECT count(*) FROM lake.staging.orders.branch_wap;
SELECT count(*) FROM lake.staging.orders;

-- chỉ khi mọi kiểm tra qua
CALL lake.system.fast_forward('staging.orders', 'main', 'wap');
\end{codeblock}
\begin{enumerate}
  \item Chạy các kiểm tra của Lab 10.1 trên \code{branch\_wap} thay vì trên \code{main}.
  \item Lần một: dữ liệu sạch, kiểm tra qua, \code{fast\_forward}; xác nhận \code{main} có dữ liệu mới.
  \item Lần hai: chèn dòng trùng \code{order\_id} vào lô; kiểm tra fail, không publish; xác nhận người đọc \code{main} không thấy dòng xấu. Bỏ branch hỏng bằng \code{ALTER TABLE ... DROP BRANCH} và tạo lại cho lần sau.
  \item Đưa ba bước ghi, kiểm, publish vào DAG \code{olist\_staging} thay cho cổng của Lab 10.1.
\end{enumerate}
\hoanthanh{ở lần hai, số dòng của \code{main} không đổi và bảng snapshot cho thấy snapshot mới chỉ nằm trên branch.}

\subsection{Lab 10.3: Git, pre-commit, unit test và CI}
\begin{enumerate}
  \item Đưa code DAG và job Spark vào Git, thêm \code{.gitignore} bỏ dữ liệu, file \code{.env} và volume Docker; đẩy lên một kho GitHub của bạn.
  \item Thêm pre-commit với hook Ruff.
  \item Tách hàm biến đổi (ví dụ khử trùng theo \code{order\_id} giữ bản \code{updated\_at} mới nhất) thành hàm thuần nhận và trả DataFrame; viết unit test với fixture \code{SparkSession} local và \code{assertDataFrameEqual}, có ca biên hai dòng cùng \code{updated\_at}.
  \item Thêm workflow GitHub Actions như dưới; kiểm phiên bản mới nhất của từng action trong docs GitHub.
\end{enumerate}
\begin{codeblock}
name: ci
on: [push, pull_request]
jobs:
  test:
    runs-on: ubuntu-latest
    steps:
      - uses: actions/checkout@v5
      - uses: actions/setup-java@v4
        with:
          distribution: temurin
          java-version: "17"
      - uses: actions/setup-python@v6
        with:
          python-version: "3.12"
          cache: pip
      - run: pip install -r requirements-dev.txt
      - run: ruff check .
      - run: pytest -q
\end{codeblock}
\hoanthanh{một pull request có test fail bị đánh đỏ, sửa xong thì xanh, và \code{ruff check} chạy cả rule \code{AIR301}, \code{AIR302} cho thư mục DAG.}

\subsection{Lab 10.4 (tuỳ chọn): viết lại kiểm tra bằng GX Core}
\begin{enumerate}
  \item Chạy \href{https://docs.greatexpectations.io/docs/core/introduction/try_gx/}{Try GX} với pandas.
  \item Xuất một mẫu \code{staging.orders} ra pandas và viết lại các kiểm tra \code{not\_null}, \code{unique}, \code{accepted\_values} thành expectation.
  \item So số dòng code, thông báo lỗi và thời gian chạy với cách SQL của Lab 10.1.
\end{enumerate}
\hoanthanh{cùng một dòng xấu bị cả hai cách bắt, và bạn viết được ba câu về khi nào nên dùng thư viện.}

\section{Kế hoạch theo tuần}
\begin{xltabular}{\textwidth}{@{}P{1.4cm} L P{1.4cm}@{}}
\toprule
\thd{tuần} & \thd{việc} & \thd{giờ} \\
\midrule
\endhead
\bottomrule
\endfoot
1 & dbt data tests, unit tests, model contracts, GX use cases, Iceberg Branching và Spark Procedures; Lab 10.1 và 10.2 & 10 \\
2 & Missing Semester (Git, Code Quality), Pro Git chương 1 tới 3, pre-commit, pytest, Testing PySpark, GitHub Actions, ODCS; Lab 10.3; thử thách kỹ năng & 10 \\
thêm & Lab 10.4 với GX Core; Airflow Metrics & 4 \\
\end{xltabular}

\begin{thuthach}
Viết một data contract theo ODCS v3.1.0 cho bảng mart doanh thu theo danh mục và tháng: schema, khoá, quy tắc chất lượng (gồm đối soát tổng doanh thu mart với tổng \code{price} của \code{staging.order\_items} trong cùng tháng) và cam kết độ tươi. Sau đó tự hiện thực các quy tắc đó thành kiểm tra trong WAP của DAG marts, và thêm chúng vào CI ở dạng unit test trên dữ liệu mẫu nhỏ.
\hoanthanh{file contract YAML hợp lệ theo ODCS, một lần cố ý làm lệch tổng doanh thu bị WAP chặn trước khi publish, và CI xanh trên nhánh chính.}
\end{thuthach}

\section{Câu hỏi tự kiểm tra}
\begin{enumerate}
  \item Unit test, data test, contract và test đầu cuối khác nhau thế nào? Cho mỗi loại một ví dụ trên Olist.
  \item Vì sao dbt khuyên chỉ chạy unit test ở dev và CI, còn data test chạy cả ở production?
  \item Kể bốn generic test của dbt và một kiểm tra Olist cho mỗi loại.
  \item Contract của dbt bảo đảm gì, và vì sao khai báo \code{unique} trong contract thường không đủ?
  \item Trong WAP trên Iceberg, người đọc \code{main} thấy gì khi branch đang giữ dữ liệu mới?
  \item Nếu kiểm tra fail, bạn làm gì với branch, và vì sao không cần rollback \code{main}?
  \item Ghép mỗi kiểm tra của Lab 10.1 với một chiều chất lượng trong danh sách use case của GX.
  \item Kiểm tra nào chạy được trong GitHub Actions, kiểm tra nào chỉ chạy ở môi trường có đủ dịch vụ, và vì sao?
  \item License Elastic License 2.0 của Soda Core cho phép và cấm điều gì?
  \item Năm 2026, ai chăm sóc GX Core, và GX Cloud thuộc về ai?
  \item Một data contract theo ODCS gồm những phần chính nào?
\end{enumerate}

\section{Điểm dễ nhầm}
\begin{xltabular}{\textwidth}{@{}P{4.2cm} L P{3.4cm}@{}}
\toprule
\thd{điểm dễ nhầm} & \thd{thực tế} & \thd{docs chính thức} \\
\midrule
\endhead
\bottomrule
\endfoot
Nhiều người nghĩ data test và unit test là một & Unit test kiểm logic trên đầu vào tĩnh trước khi build; data test kiểm dữ liệu thật sau khi build & \href{https://docs.getdbt.com/docs/build/unit-tests}{dbt unit tests} \\
Khai báo ràng buộc \code{unique} trong contract là đủ & Trên phần lớn nền tảng chỉ \code{not\_null} được thực thi; \code{unique}, \code{primary\_key}, \code{foreign\_key} chỉ được khai báo, nên vẫn cần data test & \href{https://docs.getdbt.com/docs/mesh/govern/model-contracts}{dbt model contracts} \\
WAP cần một hệ thống versioning riêng & Branch Iceberg cộng \code{fast\_forward} làm được WAP ngay trong một bảng & \href{https://iceberg.apache.org/docs/latest/branching/}{Iceberg Branching} \\
Đặt \code{spark.wap.branch} sẽ tự tạo branch & Branch phải tồn tại trước khi ghi & \href{https://iceberg.apache.org/docs/latest/spark-writes/}{Iceberg Spark Writes} \\
Soda Core vẫn là Apache 2.0 & Soda đã chuyển Soda Core sang Elastic License 2.0; kiểm license của phiên bản bạn dùng & \href{https://soda.io/blog/soda-core-license-update-moving-to-elastic-license}{Soda} \\
GX Core bị bán cùng GX Cloud & FICO mua GX Cloud; GX Core vẫn mã nguồn mở, do Fivetran chăm sóc & \href{https://www.fivetran.com/press/fivetran-to-become-steward-of-the-great-expectations-open-source-community-and-gx-core-project}{Fivetran} \\
Pipeline xanh nghĩa là dữ liệu đúng & Trạng thái task chỉ nói code chạy xong; cần data test chặn trước khi publish & \href{https://docs.getdbt.com/docs/build/data-tests}{dbt data tests} \\
\end{xltabular}



\partpage{IV}{Phục vụ, streaming và cloud}{Trino và BI, xử lý stream với Kafka và Flink, đưa lakehouse lên cloud với chi phí kiểm soát được. Kết thúc bằng mốc capstone 4: dashboard và bản trên AWS.}
% ---------------- m11_13 ----------------
% ============================================================
% Module 11, 12, 13 và chương Nghề nghiệp
% ============================================================

% ==================================================================
\modulechapter{11}{Phục vụ dữ liệu: Trino và BI}
\moduleinfo{Trung cấp}{15 giờ}{Module 3, 5, 6}

Sau khi các bảng marts (gold) đã nằm trong Iceberg, người phân tích cần hỏi dữ liệu bằng SQL và xem dashboard mà không chạm được vào dữ liệu gốc. Module này đặt Trino làm query engine chỉ đọc trên cùng Iceberg REST catalog mà Spark đang ghi, rồi nối Metabase lên trên. Bạn sẽ chứng minh điều kiện lọc được đẩy xuống bảng Iceberg, khoá quyền ghi ở tầng Trino, và biết khi nào nên tính sẵn (bảng tổng hợp, materialized view) hay định nghĩa metric bằng code (semantic layer). Học xong, bạn dựng được một lớp phục vụ dữ liệu nhỏ nhưng đúng nguyên tắc cho bộ dữ liệu Olist.

\section{Bạn sẽ học được}
\begin{itemize}
  \item Giải thích vai trò của coordinator, worker, connector, catalog và split trong Trino.
  \item Cấu hình connector Iceberg của Trino với REST catalog và object storage tương thích S3.
  \item Đọc \code{EXPLAIN} và \code{EXPLAIN ANALYZE} để chứng minh pushdown và việc chỉ đọc file cần thiết.
  \item Thiết lập phân quyền chỉ đọc và phân quyền theo file cho người dùng phân tích.
  \item Kết nối Metabase bản mã nguồn mở với Trino và dựng dashboard trên marts Olist.
  \item Chọn cách phục vụ phù hợp: bảng tổng hợp sẵn, materialized view của Trino, DuckDB ở local, hoặc semantic layer với MetricFlow.
\end{itemize}

\section{Nội dung chính}
\begin{enumerate}
  \item \textbf{Trino là query engine, không phải cơ sở dữ liệu.} Trino chạy SQL phân tán trên dữ liệu nằm ở nơi khác; mỗi nguồn được nối qua một connector và xuất hiện dưới dạng một catalog (\href{https://trino.io/docs/current/overview/concepts.html}{Trino concepts}).
  \item \textbf{Coordinator, worker, split.} Coordinator nhận câu truy vấn, lập kế hoạch và điều phối; worker thực thi task và xử lý dữ liệu. Split là một phần dữ liệu mà một task xử lý; câu truy vấn được chia thành stage, task và split (\href{https://trino.io/docs/current/overview/concepts.html}{Trino concepts}).
  \item \textbf{Connector Iceberg và loại catalog.} Thuộc tính \code{iceberg.\allowbreak catalog.\allowbreak type} nhận \code{hive\_\allowbreak metastore} (mặc định), \code{glue}, \code{jdbc}, \code{rest}, \code{nessie}, \code{snowflake}. Với REST catalog, đặt \code{iceberg.\allowbreak catalog.\allowbreak type=rest} và \code{iceberg.\allowbreak rest-catalog.\allowbreak uri}; có thể thêm \code{iceberg.\allowbreak rest-catalog.\allowbreak warehouse}, \code{iceberg.\allowbreak rest-catalog.\allowbreak security} (\code{NONE} mặc định, \code{OAUTH2}, \code{SIGV4}, \code{GOOGLE}) và \code{iceberg.\allowbreak rest-catalog.\allowbreak vended-credentials-enabled} (\href{https://trino.io/docs/current/connector/iceberg.html}{Iceberg connector}, \href{https://trino.io/docs/current/object-storage/metastores.html}{Metastores}).
  \item \textbf{Đọc object storage.} Bật S3 gốc của Trino bằng \code{fs.\allowbreak s3.\allowbreak enabled=true}, rồi đặt \code{s3.endpoint}, \code{s3.region}, \code{s3.\allowbreak path-style-access} và khoá truy cập; với RustFS ở local cần path-style (\href{https://trino.io/docs/current/object-storage/file-system-s3.html}{S3 file system support}).
  \item \textbf{Pushdown.} Connector có thể nhận một phần công việc: điều kiện lọc, cột cần đọc, aggregation, limit. Mức hỗ trợ khác nhau theo connector (\href{https://trino.io/docs/current/optimizer/pushdown.html}{Pushdown}). Iceberg lưu đường dẫn data file trong metadata, nên Trino chỉ hỏi file system về những file phải đọc (\href{https://trino.io/docs/current/connector/iceberg.html}{Iceberg connector}).
  \item \textbf{Đọc kế hoạch truy vấn.} \code{EXPLAIN (TYPE DISTRIBUTED)} cho thấy fragment và exchange; \code{EXPLAIN ANALYZE} chạy câu thật và báo số dòng, thời gian ở từng toán tử (\href{https://trino.io/docs/current/sql/explain.html}{EXPLAIN}). Bảng metadata như \code{"orders\$snapshots"}, \code{"orders\$files"}, \code{"orders\$partitions"} giúp bạn đối chiếu (\href{https://trino.io/docs/current/connector/iceberg.html}{Iceberg connector}).
  \item \textbf{Phân quyền hệ thống.} Đặt thuộc tính \code{access-control.\allowbreak name} trong \code{etc/\allowbreak access-control.\allowbreak properties} trên coordinator. Giá trị dựng sẵn: \code{default}, \code{allow-all}, \code{read-only}, \code{file}, \code{opa}, \code{ranger}. \code{read-only} cho phép \code{SELECT}, \code{SHOW} và đặt session property, cấm mọi thao tác ghi dữ liệu hay metadata như \code{CREATE}, \code{INSERT}, \code{DELETE} (\href{https://trino.io/docs/current/security/built-in-system-access-control.html}{System access control}). \code{file} đọc luật theo user, catalog, schema, bảng từ file JSON (\href{https://trino.io/docs/current/security/file-system-access-control.html}{File-based access control}).
  \item \textbf{Đồng thời cho BI.} Resource group giới hạn số truy vấn chạy cùng lúc và xếp hàng phần còn lại, để dashboard không chiếm hết cluster (\href{https://trino.io/docs/current/admin/resource-groups.html}{Resource groups}).
  \item \textbf{Materialized view trên Iceberg.} Connector Iceberg lưu kết quả vào một storage table và làm mới bằng \code{REFRESH MATERIALIZED VIEW}. Incremental refresh bật mặc định (\code{iceberg.\allowbreak incremental-refresh-enabled}); nó quét ít dữ liệu nguồn hơn nhưng dùng append nên có thể sinh thêm data file; đặt \code{false} để luôn refresh toàn bộ (\href{https://trino.io/docs/current/connector/iceberg.html}{Iceberg connector}).
  \item \textbf{Metabase.} Driver Starburst của Metabase ``also works for connections to a Trino database''; các trường cần điền là host, port, catalog, schema (tuỳ chọn), username, role (tuỳ chọn), SSL (\href{https://www.metabase.com/docs/latest/databases/connections/starburst}{Metabase: Starburst}). Trên docs v0.65, chính sách cache Duration và Schedule, cache theo dashboard, cache theo câu hỏi và tự làm mới cache chỉ có ở gói Pro và Enterprise (\href{https://www.metabase.com/docs/latest/configuring-metabase/caching}{Metabase caching}). Apache Superset là phương án BI mã nguồn mở khác (\href{https://superset.apache.org/docs/intro/}{Superset docs}).
  \item \textbf{Semantic layer.} MetricFlow sinh SQL từ metric định nghĩa bằng YAML, là nền của dbt Semantic Layer, cấp phép Apache 2.0. Khái niệm: semantic model, entity (khoá join), dimension, metric (simple, ratio, derived, cumulative, conversion) (\href{https://docs.getdbt.com/docs/build/about-metricflow}{About MetricFlow}). Tự host thì cài \code{metricflow} bằng pip và chạy lệnh \code{mf}; truy vấn metric động từ công cụ downstream qua dbt Semantic Layer cần tài khoản dbt Starter, Enterprise hoặc Enterprise+ (\href{https://docs.getdbt.com/docs/build/metricflow-commands}{MetricFlow commands}). Cú pháp YAML đổi từ dbt v1.12, nên đọc docs đúng phiên bản.
  \item \textbf{DuckDB làm lớp phục vụ local.} DuckDB chạy nhúng trong tiến trình, đọc được bảng Iceberg qua extension \code{iceberg}, hợp để thử câu hỏi, test và trích dữ liệu nhỏ; Trino hợp hơn khi nhiều người cùng truy vấn qua mạng (\href{https://duckdb.org/docs/2.0/core_extensions/iceberg/overview.html}{DuckDB Iceberg extension}).
\end{enumerate}

\section{Bắt đầu}
\begin{enumerate}
  \item Chạy Trino bằng container theo \href{https://trino.io/docs/current/installation/containers.html}{Trino in a Docker container}, ghim image \code{trinodb/\allowbreak trino:484}.
  \item Thêm một catalog Iceberg dùng REST theo ví dụ tối thiểu ở \href{https://trino.io/docs/current/object-storage/metastores.html}{Metastores} (\code{connector.\allowbreak name=iceberg}, \code{iceberg.\allowbreak catalog.\allowbreak type=rest}, \code{iceberg.\allowbreak rest-catalog.\allowbreak uri}).
  \item Chạy Metabase bản mã nguồn mở theo \href{https://www.metabase.com/start/oss/}{Metabase OSS} và nối vào Trino bằng driver Starburst.
\end{enumerate}

\section{Tài liệu chính thức}
\begin{xltabular}{\textwidth}{@{}P{4.2cm} L P{1.8cm}@{}}
\toprule
\thd{tài liệu} & \thd{nội dung cần đọc} & \thd{thời lượng} \\
\midrule
\endhead
\bottomrule
\endfoot
\href{https://trino.io/docs/current/overview/concepts.html}{Trino 484: Concepts} & server, coordinator, worker, catalog, connector, stage, task, split & 1 giờ \\
\href{https://trino.io/docs/current/connector/iceberg.html}{Trino: Iceberg connector} & cấu hình chung, metadata table, materialized views & 1,5 giờ \\
\href{https://trino.io/docs/current/object-storage/metastores.html}{Trino: Metastores} & mục REST catalog: URI, warehouse, security, vended credentials & 0,5 giờ \\
\href{https://trino.io/docs/current/object-storage/file-system-s3.html}{Trino: S3 file system support} & \code{fs.\allowbreak s3.\allowbreak enabled}, endpoint, path-style & 0,5 giờ \\
\href{https://trino.io/docs/current/optimizer/pushdown.html}{Trino: Pushdown}, \href{https://trino.io/docs/current/sql/explain.html}{EXPLAIN} & các loại pushdown; đọc kế hoạch phân tán & 1 giờ \\
\href{https://trino.io/docs/current/security/built-in-system-access-control.html}{Trino: System access control}, \href{https://trino.io/docs/current/security/file-system-access-control.html}{File-based} & \code{read-only}, luật catalog, schema, bảng & 1 giờ \\
\href{https://trino.io/docs/current/admin/resource-groups.html}{Trino: Resource groups} & giới hạn đồng thời, hàng đợi & 0,5 giờ \\
\href{https://www.metabase.com/docs/latest/databases/connections/starburst}{Metabase: Starburst/Trino}, \href{https://www.metabase.com/docs/latest/configuring-metabase/caching}{Caching} & kết nối; tính năng cache theo gói & 1 giờ \\
\href{https://docs.getdbt.com/docs/build/about-metricflow}{dbt: About MetricFlow}, \href{https://docs.getdbt.com/docs/build/metricflow-commands}{MetricFlow commands} & semantic model, metric, lệnh \code{mf} & 1 giờ \\
\href{https://duckdb.org/docs/2.0/core_extensions/iceberg/overview.html}{DuckDB: Iceberg extension} & đọc bảng Iceberg, nối REST catalog & 0,5 giờ \\
\end{xltabular}

\section{Khoá học miễn phí}
\begin{itemize}
  \item \href{https://trino.io/trino-the-definitive-guide.html}{Trino: The Definitive Guide, bản 2} (O'Reilly, 10/2022, 319 trang): trang chính thức của Trino có link ``Free digital copy from Starburst''. Sách viết trước Trino 484, nên kiểm lại tên thuộc tính với docs hiện hành.
  \item \href{https://15445.courses.cs.cmu.edu/fall2025/schedule.html}{CMU 15-445 Fall 2025}: bài giảng và slide miễn phí; xem nhóm bài về thực thi và tối ưu truy vấn nếu muốn hiểu sâu join và cost-based optimizer.
\end{itemize}

\section{Lab thực hành}

\subsection{Lab 1: Trino trên Iceberg REST catalog}
\begin{enumerate}
  \item Thêm service \code{trinodb/\allowbreak trino:484} vào Docker Compose của stack tham chiếu, cùng mạng với RustFS và REST catalog.
  \item Tạo file \code{etc/catalog/lake.properties} (thay giá trị theo catalog của bạn; với Polaris hay Lakekeeper, lấy URI, warehouse và thông tin OAuth2 từ quickstart của dự án):
\end{enumerate}
\begin{codeblock}
connector.name=iceberg
iceberg.catalog.type=rest
iceberg.rest-catalog.uri=http://iceberg-rest:8181
iceberg.rest-catalog.warehouse=s3://warehouse/
# iceberg.rest-catalog.security=OAUTH2
# iceberg.rest-catalog.oauth2.credential=<client_id>:<client_secret>
fs.s3.enabled=true
s3.endpoint=http://rustfs:9000
s3.region=us-east-1
s3.path-style-access=true
s3.aws-access-key=<access-key>
s3.aws-secret-key=<secret-key>
\end{codeblock}
\begin{enumerate}[resume]
  \item Khởi động lại Trino, mở CLI trong container và chạy \code{SHOW SCHEMAS FROM lake}, \code{SHOW TABLES FROM lake.marts}.
  \item Truy vấn một mart Olist do Spark ghi (ví dụ doanh thu theo tháng và bang của khách), và \code{SELECT * FROM lake.marts."<bảng>\$snapshots"}.
\end{enumerate}
\hoanthanh{Trino trả cùng số dòng cho một mart như Spark hoặc PyIceberg, và snapshot mới nhất trong \code{\$snapshots} trùng snapshot mà Spark vừa commit.}

\subsection{Lab 2: Chứng minh pushdown}
\begin{enumerate}
  \item Chọn bảng sự kiện đơn hàng có partition theo ngày (hoặc tháng) của \code{order\_\allowbreak purchase\_\allowbreak timestamp}.
  \item Chạy \code{EXPLAIN (TYPE DISTRIBUTED)} cho một câu lọc theo khoảng ngày và một câu join đơn hàng với khách hàng. Tìm toán tử table scan, điều kiện lọc gắn với nó, và các exchange.
  \item Chạy \code{EXPLAIN ANALYZE} cho câu lọc theo ngày và câu lọc theo một cột không phải partition; so số dòng và dữ liệu đầu vào của table scan.
  \item Đếm file của các partition liên quan qua \code{"<bảng>\$files"} để giải thích vì sao lượng đọc khác nhau.
\end{enumerate}
\hoanthanh{Bạn có một bảng ngắn: câu truy vấn, số dòng đầu vào của scan, số file liên quan, và một câu giải thích vì sao câu lọc theo partition đọc ít hơn.}

\subsection{Lab 3: Người dùng phân tích chỉ đọc}
\begin{enumerate}
  \item Đặt \code{access-control.name=read-only} trong \code{etc/\allowbreak access-control.\allowbreak properties} của coordinator; thử \code{CREATE TABLE} và \code{INSERT}, ghi lại thông báo lỗi.
  \item Đổi sang \code{access-control.name=file} với \code{security.config-file} trỏ tới file luật. Ví dụ: \code{admin} toàn quyền, \code{analyst} chỉ \code{SELECT} schema \code{marts}:
\end{enumerate}
\begin{codeblock}
{
  "catalogs": [
    {"user": "admin", "catalog": ".*", "allow": "all"},
    {"user": "analyst", "catalog": "lake", "allow": "read-only"}
  ],
  "schemas": [
    {"user": "admin", "schema": ".*", "owner": true},
    {"user": "analyst", "catalog": "lake", "schema": "marts", "owner": false}
  ],
  "tables": [
    {"user": "admin", "privileges": ["SELECT", "INSERT", "DELETE", "UPDATE", "OWNERSHIP"]},
    {"user": "analyst", "catalog": "lake", "schema": "marts", "privileges": ["SELECT"]}
  ]
}
\end{codeblock}
\begin{enumerate}[resume]
  \item Đăng nhập CLI với \code{--user analyst}; kiểm \code{SHOW SCHEMAS FROM lake} và thử đọc một bảng ở \code{raw}.
  \item Thêm một resource group cho user \code{analyst} với số truy vấn đồng thời thấp; chạy 5 truy vấn cùng lúc và quan sát hàng đợi trên Web UI.
\end{enumerate}
\begin{luuy}
Trong lab không có xác thực, nên client tự khai tên user. Ở môi trường thật, bật authentication cho Trino trước khi dựa vào luật theo user. Phân quyền ở Metabase không thay được phân quyền ở Trino, vì người dùng có thể nối thẳng vào Trino.
\end{luuy}
\hoanthanh{User \code{analyst} đọc được \code{marts}, không thấy hoặc không đọc được \code{raw}, và mọi lệnh ghi đều bị từ chối.}

\subsection{Lab 4: Dashboard Metabase trên marts Olist}
\begin{enumerate}
  \item Chạy Metabase OSS bằng Docker; thêm database loại Starburst trỏ vào Trino, catalog \code{lake}, schema \code{marts}, username \code{analyst}.
  \item Dựng 4 câu hỏi: doanh thu theo tháng, 10 danh mục sản phẩm bán chạy, thời gian giao hàng trung bình theo bang, phân bố điểm review.
  \item Ghép thành một dashboard có bộ lọc theo khoảng ngày.
  \item Mở trang cài đặt cache của Metabase, ghi lại tuỳ chọn nào có ở bản mã nguồn mở và tuỳ chọn nào cần Pro hay Enterprise.
\end{enumerate}
\hoanthanh{Dashboard trả lời 4 câu hỏi mà người xem không phải viết join, và kết nối dùng user Trino chỉ đọc.}

\section{Kế hoạch theo tuần}
\begin{xltabular}{\textwidth}{@{}P{1.6cm} L P{1.8cm}@{}}
\toprule
\thd{tuần} & \thd{việc} & \thd{giờ} \\
\midrule
\endhead
\bottomrule
\endfoot
1 & Đọc Concepts, Iceberg connector, Metastores, S3 file system, Pushdown, EXPLAIN; làm Lab 1 và Lab 2 & 8 \\
2 & Đọc access control, resource groups, Metabase; làm Lab 3, Lab 4 và thử thách kỹ năng & 7 \\
thêm & Materialized view với incremental refresh; MetricFlow trên một project dbt Core; DuckDB đọc cùng bảng; sách Trino TDG; CMU 15-445 & 8 trở lên \\
\end{xltabular}

\begin{thuthach}
Bộ phận tài chính cần một dashboard ``doanh thu và phí vận chuyển theo tháng, theo phương thức thanh toán'' cho 2017--2018. Hãy tự thiết kế: một user Trino mới chỉ đọc đúng các bảng cần thiết; một bảng tổng hợp sẵn hoặc một materialized view phục vụ dashboard; bằng chứng từ \code{EXPLAIN ANALYZE} rằng dashboard đọc ít dữ liệu hơn khi truy vấn bảng chi tiết; và một dashboard Metabase dùng user đó.
\hoanthanh{Dashboard chạy bằng user mới, user này không đọc được bảng ngoài danh sách, và bạn trình bày được số liệu đọc của hai phương án (bảng chi tiết và bảng tính sẵn).}
\end{thuthach}

\section{Câu hỏi tự kiểm tra}
\begin{enumerate}
  \item Coordinator làm gì mà worker không làm? Split là gì?
  \item Vì sao Trino truy vấn được bảng Iceberg mà không cần nạp dữ liệu vào nó?
  \item Những thuộc tính nào tối thiểu để Trino đọc một Iceberg REST catalog trên RustFS?
  \item Trong \code{EXPLAIN ANALYZE}, bạn nhìn vào đâu để biết điều kiện lọc đã tới table scan?
  \item Vì sao lọc theo cột partition đọc ít file hơn lọc theo cột thường, dù kết quả cùng số dòng?
  \item \code{read-only} cho phép và cấm những thao tác nào? Nó được cấu hình ở node nào?
  \item Vì sao phân quyền trong Metabase không đủ để bảo vệ dữ liệu?
  \item Mười người cùng mở một dashboard lúc 9 giờ sáng: kể ba cách giảm tải, cách nào dùng được trên Metabase bản mã nguồn mở?
  \item Incremental refresh của materialized view trên Iceberg đánh đổi điều gì?
  \item Semantic layer giải quyết vấn đề gì mà một bảng marts không giải quyết?
  \item Khi nào bạn chọn DuckDB thay vì Trino để trả lời một câu hỏi phân tích?
\end{enumerate}

\section{Điểm dễ nhầm}
\begin{xltabular}{\textwidth}{@{}P{4.0cm} L P{3.2cm}@{}}
\toprule
\thd{điểm dễ nhầm} & \thd{thực tế} & \thd{docs chính thức} \\
\midrule
\endhead
\bottomrule
\endfoot
Nhiều người nghĩ Trino lưu dữ liệu như một cơ sở dữ liệu & Trino là query engine; dữ liệu nằm ở nguồn, được đọc qua connector & \href{https://trino.io/docs/current/overview/concepts.html}{Trino concepts} \\
Join giữa hai catalog khác nhau là ``miễn phí'' & Mức pushdown tuỳ connector; phần không đẩy xuống được thì dữ liệu phải chuyển qua Trino & \href{https://trino.io/docs/current/optimizer/pushdown.html}{Pushdown} \\
Đặt chỉ đọc trong Metabase là đủ & Người dùng có thể nối thẳng vào Trino; khoá ghi bằng \code{read-only} hoặc luật theo file ở Trino & \href{https://trino.io/docs/current/security/built-in-system-access-control.html}{System access control} \\
Metabase bản mã nguồn mở có đủ cache như bản trả phí & Duration, Schedule, cache theo dashboard, theo câu hỏi và tự làm mới cache chỉ có ở Pro và Enterprise & \href{https://www.metabase.com/docs/latest/configuring-metabase/caching}{Metabase caching} \\
Cần driver Trino riêng cho Metabase & Driver Starburst dùng được cho Trino & \href{https://www.metabase.com/docs/latest/databases/connections/starburst}{Metabase: Starburst} \\
Incremental refresh luôn tốt hơn & Nó quét ít dữ liệu nguồn hơn nhưng có thể tạo thêm data file, cần bảo trì bảng & \href{https://trino.io/docs/current/connector/iceberg.html}{Iceberg connector} \\
Semantic layer của dbt hoàn toàn miễn phí & MetricFlow là Apache 2.0 và chạy local bằng \code{mf}; truy vấn động từ công cụ downstream qua dbt Semantic Layer cần gói dbt trả phí & \href{https://docs.getdbt.com/docs/build/metricflow-commands}{MetricFlow commands} \\
\end{xltabular}


% ==================================================================
\modulechapter{12}{Streaming: Kafka và Flink}
\moduleinfo{Nâng cao}{30 giờ}{Module 5, 7, 8}

Pipeline batch xử lý dữ liệu theo đợt; streaming xử lý từng sự kiện khi nó tới. Module này dạy hai công cụ phổ biến nhất: Apache Kafka làm log sự kiện bền vững, và Apache Flink làm bộ xử lý luồng có trạng thái. Bạn sẽ đẩy sự kiện đơn hàng Olist vào Kafka, tính tổng theo cửa sổ thời gian bằng Flink SQL dựa trên event time và watermark, ghi kết quả vào Iceberg, rồi khởi động lại job để kiểm exactly-once. Học xong, bạn giải thích đúng Kafka bảo đảm gì, Flink bảo đảm gì, và chi phí của việc ghi lakehouse từ luồng.

\section{Bạn sẽ học được}
\begin{itemize}
  \item Giải thích topic, partition, offset, consumer group và phạm vi bảo đảm thứ tự của Kafka.
  \item Phân biệt at-most-once, at-least-once, exactly-once, và biết idempotent producer, transaction, \code{read\_committed} làm gì.
  \item Phân biệt event time với processing time và giải thích watermark của Flink.
  \item Viết truy vấn Flink SQL dùng window TVF trên luồng Kafka.
  \item Giải thích checkpoint của Flink và cách sink Iceberg dựa vào nó để đạt exactly-once.
  \item Ước lượng số snapshot và file mà một luồng ghi Iceberg sinh ra, và chọn việc bảo trì bảng.
\end{itemize}

\section{Nội dung chính}

\subsection{Kafka}
\begin{enumerate}
  \item \textbf{Topic, partition, offset.} Topic được chia thành các partition có thứ tự hoàn toàn (``totally ordered''); Kafka không bảo đảm thứ tự giữa các partition. Message có cùng key đi vào cùng partition nên giữ thứ tự theo key. Vị trí của consumer trong mỗi partition chỉ là một số nguyên (offset), và consumer có thể chủ động tua lại offset cũ để đọc lại (\href{https://kafka.apache.org/44/design/design/}{Kafka 4.4 Design}).
  \item \textbf{Message không bị xoá khi đã đọc.} Kafka giữ message theo cấu hình retention; mục ``Constant Time Suffices'' của trang Design giải thích vì sao giữ lâu không làm chậm hệ thống. Log compaction giữ ít nhất giá trị cuối cùng của mỗi key, hợp với topic ``trạng thái hiện tại'' (\href{https://kafka.apache.org/44/design/design/}{Kafka 4.4 Design}).
  \item \textbf{Consumer group.} Mỗi partition, tại một thời điểm, được đúng một consumer trong mỗi group đọc. Giao thức consumer group mới (KIP-848) GA từ Kafka 4.0, bật sẵn phía server; client chọn bằng \code{group.protocol=consumer} (\href{https://kafka.apache.org/blog/2025/03/18/apache-kafka-4.0.0-release-announcement/}{Kafka 4.0 announcement}).
  \item \textbf{Share group.} Share group là một loại group khác, chạy song song với consumer group: nhiều consumer cùng đọc và một partition có thể giao cho nhiều consumer, record được acknowledge riêng lẻ (\href{https://kafka.apache.org/44/design/design/}{Kafka 4.4 Design}, mục ``The Share Consumer''). Tính năng ở dạng early access trong 4.0 và được công bố production-ready từ 4.2 (\href{https://kafka.apache.org/blog/2026/02/17/apache-kafka-4.2.0-release-announcement/}{Kafka 4.2 announcement}).
  \item \textbf{Replication.} Leader theo dõi tập replica đồng bộ (ISR); với \code{min.insync.replicas}, partition chỉ nhận ghi khi ISR đủ lớn. Từ 0.11.0.0, mặc định Kafka chờ một replica nhất quán thay vì bầu leader ``unclean'' (\href{https://kafka.apache.org/44/design/design/}{Kafka 4.4 Design}).
  \item \textbf{Chỉ còn KRaft.} Kafka 4.0 là bản lớn đầu tiên chạy hoàn toàn không cần ZooKeeper (\href{https://kafka.apache.org/blog/2025/03/18/apache-kafka-4.0.0-release-announcement/}{Kafka 4.0 announcement}). Tài liệu và video cũ có bước cài ZooKeeper không còn áp dụng. Bản mới nhất là 4.4.0 (\href{https://kafka.apache.org/blog/2026/10/22/apache-kafka-4.4.0-release-announcement/}{Kafka 4.4.0 announcement}).
  \item \textbf{Ngữ nghĩa giao nhận.} At-most-once: có thể mất, không lặp. At-least-once: không mất, có thể lặp. Exactly-once: mỗi message được xử lý đúng một lần. Idempotent producer (từ 0.11.0.0) bảo đảm gửi lại không tạo bản trùng trong log; transaction ghi nhiều message theo kiểu tất cả hoặc không; consumer \code{isolation.\allowbreak level=\allowbreak read\_committed} chỉ trả message của transaction đã commit. Exactly-once sang hệ đích khác ``generally requires cooperation with such systems'' (\href{https://kafka.apache.org/44/design/design/}{Kafka 4.4 Design}, mục ``Message Delivery Semantics'', ``Using Transactions'').
  \item \textbf{Hệ tương thích Kafka API (tuỳ chọn).} Redpanda là broker nói Kafka API nhưng là implementation riêng; kiểm tính năng cần dùng trên \href{https://docs.redpanda.com/}{Redpanda docs} trước khi thay thế.
\end{enumerate}

\subsection{Flink}
\begin{enumerate}[resume]
  \item \textbf{Xử lý luồng có trạng thái.} Flink giữ state cho từng key (ví dụ tổng doanh thu của cửa sổ đang mở) và chụp state định kỳ bằng checkpoint để phục hồi khi lỗi (\href{https://nightlies.apache.org/flink/flink-docs-stable/docs/learn-flink/fault_tolerance/}{Learn Flink: Fault Tolerance}).
  \item \textbf{Event time và processing time.} Processing time là giờ máy khi sự kiện được xử lý; event time là thời điểm sự kiện xảy ra, nằm trong chính bản ghi. Phân tích dựa trên event time cho cùng kết quả khi phát lại dữ liệu cũ (\href{https://nightlies.apache.org/flink/flink-docs-stable/docs/concepts/time/}{Timely Stream Processing}).
  \item \textbf{Watermark.} Watermark với thời điểm \code{t} khẳng định luồng đã tới event time \code{t}, tức không còn chờ sự kiện cũ hơn \code{t}. Sự kiện tới sau watermark là sự kiện trễ; watermark là heuristic nên vẫn có thể có dữ liệu trễ. Với nhiều luồng vào song song, watermark của toán tử là giá trị nhỏ nhất của các luồng vào (\href{https://nightlies.apache.org/flink/flink-docs-stable/docs/concepts/time/}{Timely Stream Processing}). Trong Flink SQL, khai báo bằng \code{WATERMARK FOR ts AS ts - INTERVAL '5' SECOND} trong \code{CREATE TABLE}.
  \item \textbf{Window TVF.} Flink SQL có các window table-valued function \code{TUMBLE}, \code{HOP}, \code{CUMULATE}, \code{SESSION}; kết quả có cột \code{window\_start}, \code{window\_end}, rồi \code{GROUP BY window\_start, window\_end} (\href{https://nightlies.apache.org/flink/flink-docs-stable/docs/sql/reference/queries/window-agg/}{Window Aggregation}).
  \item \textbf{Checkpoint và exactly-once đầu cuối.} Khi lỗi, Flink khôi phục state và vị trí đọc nguồn từ checkpoint gần nhất. Để kết quả ở hệ đích không bị trùng, sink phải hợp tác: commit theo checkpoint (giao dịch) hoặc ghi idempotent. Sink Iceberg của Flink commit data file khi checkpoint hoàn tất, nên khoảng checkpoint quyết định độ tươi dữ liệu và số snapshot (\href{https://iceberg.apache.org/docs/latest/flink/}{Iceberg: Flink}, \href{https://iceberg.apache.org/docs/latest/flink-writes/}{Flink writes}).
  \item \textbf{Phiên bản.} Flink 2.3.0 là bản stable (25/06/2026), nhưng Flink Kafka connector mới nhất 5.0.0 ghi hỗ trợ Flink 2.1.x và 2.2.x (\href{https://flink.apache.org/downloads/}{Flink downloads}). Trước khi dựng lab, chọn một bản Flink mà cả Kafka connector và Iceberg Flink runtime cùng hỗ trợ (\href{https://iceberg.apache.org/multi-engine-support/}{Iceberg multi-engine support}).
\end{enumerate}

\subsection{Ghi lakehouse từ luồng}
\begin{enumerate}[resume]
  \item \textbf{Hai đường ghi Iceberg.} Sink Flink (ghi từ job xử lý) và Iceberg Kafka Connect sink (chép topic vào bảng, có control topic để điều phối commit tập trung) (\href{https://iceberg.apache.org/docs/latest/kafka-connect/}{Iceberg Kafka Connect}).
  \item \textbf{Khoảng commit và số file.} Kafka Connect sink mặc định commit mỗi 5 phút (\code{iceberg.\allowbreak control.\allowbreak commit.\allowbreak interval-ms} = 300000), tức 288 commit mỗi ngày cho mỗi bảng (\href{https://iceberg.apache.org/docs/latest/kafka-connect/}{Iceberg Kafka Connect}). Mỗi commit thêm một snapshot và file mới, nên bảng ghi từ luồng cần compaction và expire snapshot định kỳ (\href{https://iceberg.apache.org/docs/latest/maintenance/}{Iceberg Maintenance}).
  \item \textbf{Khi nào cần streaming.} Nếu nguồn chỉ có dữ liệu theo đợt (ví dụ API kéo 15 phút một lần), streaming thường chỉ thêm độ phức tạp. Streaming đáng giá khi nguồn tự đẩy sự kiện và người dùng cần kết quả trong vài giây tới vài phút.
\end{enumerate}

\section{Bắt đầu}
\begin{enumerate}
  \item Đọc phần Introduction và Quickstart trong \href{https://kafka.apache.org/documentation/}{Kafka docs}, chạy một broker KRaft một node bằng image Docker chính thức \code{apache/kafka}, ghim bản \code{4.4.0}.
  \item Đọc \href{https://nightlies.apache.org/flink/flink-docs-stable/docs/learn-flink/overview/}{Learn Flink: Overview}.
  \item Làm \href{https://iceberg.apache.org/flink-quickstart/}{Flink and Iceberg Quickstart} của Iceberg để có Flink SQL ghi được bảng Iceberg.
\end{enumerate}

\section{Tài liệu chính thức}
\begin{xltabular}{\textwidth}{@{}P{4.2cm} L P{1.8cm}@{}}
\toprule
\thd{tài liệu} & \thd{nội dung cần đọc} & \thd{thời lượng} \\
\midrule
\endhead
\bottomrule
\endfoot
\href{https://kafka.apache.org/documentation/}{Kafka docs} & Introduction, Quickstart, Kafka Connect & 2 giờ \\
\href{https://kafka.apache.org/44/design/design/}{Kafka 4.4: Design} & The Consumer, Message Delivery Semantics, Using Transactions, The Share Consumer, Replication, Log Compaction & 3 giờ \\
\href{https://kafka.apache.org/blog/2025/03/18/apache-kafka-4.0.0-release-announcement/}{Kafka 4.0}, \href{https://kafka.apache.org/blog/2026/10/22/apache-kafka-4.4.0-release-announcement/}{4.4.0 announcement} & bỏ ZooKeeper, KIP-848, thay đổi mới nhất & 0,5 giờ \\
\href{https://nightlies.apache.org/flink/flink-docs-stable/docs/concepts/time/}{Flink: Timely Stream Processing} & event time, watermark, lateness & 1,5 giờ \\
\href{https://nightlies.apache.org/flink/flink-docs-stable/docs/learn-flink/fault_tolerance/}{Flink: Fault Tolerance} & state backend, checkpoint, savepoint, exactly-once & 1,5 giờ \\
\href{https://nightlies.apache.org/flink/flink-docs-stable/docs/sql/reference/queries/window-agg/}{Flink SQL: Window Aggregation} & window TVF, \code{GROUP BY window\_start, window\_end} & 1 giờ \\
\href{https://iceberg.apache.org/docs/latest/flink/}{Iceberg: Flink}, \href{https://iceberg.apache.org/docs/latest/flink-writes/}{Flink writes} & tạo catalog, ghi từ Flink SQL, upsert, bảo trì sau commit & 1,5 giờ \\
\href{https://iceberg.apache.org/docs/latest/kafka-connect/}{Iceberg Kafka Connect} & features, control topic, cấu hình commit & 1 giờ \\
\href{https://iceberg.apache.org/docs/latest/maintenance/}{Iceberg Maintenance} & expire snapshots, rewrite data files & 0,5 giờ \\
\href{https://flink.apache.org/downloads/}{Flink downloads} & bản Flink và connector tương thích & 0,25 giờ \\
\end{xltabular}

\section{Khoá học miễn phí}
\begin{itemize}
  \item \href{https://developer.confluent.io/courses/}{Confluent Developer courses}: các khoá nhập môn miễn phí như Apache Kafka 101 và Apache Flink 101. Một số bài dùng dịch vụ Confluent Cloud; phần khái niệm áp dụng cho Kafka và Flink mã nguồn mở.
  \item \href{https://github.com/DataTalksClub/data-engineering-zoomcamp}{DE Zoomcamp}: module streaming (module 7) miễn phí, có bài thực hành Kafka.
  \item \href{https://github.com/apache/flink-training}{apache/flink-training}: bài tập chính thức của dự án Flink, viết bằng Java, dùng nhánh khớp bản Flink bạn chạy.
\end{itemize}

\section{Lab thực hành}

\subsection{Lab 1: Sự kiện đơn hàng Olist vào Kafka}
\begin{enumerate}
  \item Chạy broker \code{apache/\allowbreak kafka:4.4.0} một node (trong image, các script nằm ở \code{/opt/\allowbreak kafka/\allowbreak bin}). Tạo topic \code{olist.orders} với 3 partition bằng \code{kafka-topics.sh --create}.
  \item Viết producer Python: ghép \code{orders} với \code{order\_payments} (DuckDB là đủ), sắp theo \code{order\_\allowbreak purchase\_\allowbreak timestamp}, gửi mỗi đơn thành một message JSON gồm \code{order\_id}, \code{customer\_id}, \code{order\_status}, \code{payment\_value}, \code{event\_time}; key là \code{customer\_id}.
  \item Chạy hai consumer cùng group; dùng \code{kafka-consumer-\allowbreak groups.sh --describe} xem partition giao cho ai và lag.
  \item Dừng consumer, chạy \code{kafka-consumer-\allowbreak groups.sh --reset-offsets --to-earliest --execute} rồi đọc lại.
\end{enumerate}
\hoanthanh{Bạn chỉ ra được mọi message của một \code{customer\_id} nằm cùng một partition theo đúng thứ tự gửi, và đọc lại được toàn bộ topic sau khi reset offset.}

\subsection{Lab 2: Lặp, mất và transaction}
\begin{enumerate}
  \item Viết vòng consume, transform, produce đọc \code{olist.orders} và ghi \code{olist.orders\_enriched}; bản đầu dùng auto-commit offset.
  \item Giết tiến trình giữa chừng nhiều lần; đếm bản trùng và bản thiếu ở topic đích so với nguồn.
  \item Bản hai: producer có \code{transactional.id}, gửi offset của consumer trong cùng transaction; consumer kiểm tra đọc với \code{isolation.\allowbreak level=\allowbreak read\_committed}.
\end{enumerate}
\hoanthanh{Bảng so sánh hai bản: số message trùng, số message thiếu, và một câu giải thích vì sao bản transaction không trùng khi đọc bằng \code{read\_committed}.}

\subsection{Lab 3: Cửa sổ theo event time bằng Flink SQL}
\begin{enumerate}
  \item Chọn bản Flink có Kafka connector và Iceberg runtime tương thích (xem mục Phiên bản); chạy JobManager, TaskManager và SQL Client bằng Docker.
  \item Khai báo bảng nguồn có watermark và chạy tổng theo ngày (dữ liệu Olist trải 2016--2018, nên dùng cửa sổ 1 ngày trên event time gốc):
\end{enumerate}
\begin{codeblock}
CREATE TABLE orders_events (
  order_id      STRING,
  customer_id   STRING,
  order_status  STRING,
  payment_value DECIMAL(12, 2),
  event_time    TIMESTAMP(3),
  WATERMARK FOR event_time AS event_time - INTERVAL '5' SECOND
) WITH (
  'connector' = 'kafka',
  'topic' = 'olist.orders',
  'properties.bootstrap.servers' = 'kafka:9092',
  'properties.group.id' = 'flink-orders',
  'scan.startup.mode' = 'earliest-offset',
  'format' = 'json'
);

SELECT window_start, window_end, COUNT(*) AS orders, SUM(payment_value) AS revenue
FROM TABLE(TUMBLE(TABLE orders_events, DESCRIPTOR(event_time), INTERVAL '1' DAY))
GROUP BY window_start, window_end;
\end{codeblock}
\begin{enumerate}[resume]
  \item Gửi thêm vài sự kiện có \code{event\_time} lùi 2 giây và lùi 2 ngày so với sự kiện mới nhất; ghi lại sự kiện nào được tính vào cửa sổ, sự kiện nào bị bỏ.
\end{enumerate}
\hoanthanh{Tổng số đơn và doanh thu của một ngày bất kỳ đã đóng cửa sổ khớp với kết quả DuckDB tính trực tiếp từ CSV, và bạn giải thích được số phận của từng sự kiện trễ.}

\subsection{Lab 4: Ghi Iceberg và kiểm exactly-once khi khởi động lại}
\begin{enumerate}
  \item Tạo catalog Iceberg trong Flink SQL trỏ vào REST catalog của stack:
\end{enumerate}
\begin{codeblock}
SET 'execution.checkpointing.interval' = '30s';

CREATE CATALOG lake WITH (
  'type' = 'iceberg',
  'catalog-type' = 'rest',
  'uri' = 'http://iceberg-rest:8181',
  'warehouse' = 's3://warehouse/',
  'io-impl' = 'org.apache.iceberg.aws.s3.S3FileIO',
  's3.endpoint' = 'http://rustfs:9000',
  's3.path-style-access' = 'true'
);

CREATE DATABASE IF NOT EXISTS lake.streaming;
CREATE TABLE IF NOT EXISTS lake.streaming.daily_revenue (
  window_start TIMESTAMP(3), window_end TIMESTAMP(3),
  orders BIGINT, revenue DECIMAL(18, 2));

INSERT INTO lake.streaming.daily_revenue
SELECT window_start, window_end, COUNT(*), SUM(payment_value)
FROM TABLE(TUMBLE(TABLE default_catalog.default_database.orders_events,
                  DESCRIPTOR(event_time), INTERVAL '1' DAY))
GROUP BY window_start, window_end;
\end{codeblock}
\begin{enumerate}[resume]
  \item Cho producer gửi chậm (ví dụ 50 đơn mỗi giây). Sau vài checkpoint, \code{docker kill} container TaskManager rồi chạy lại nó; theo dõi trên Web UI việc job khôi phục từ checkpoint.
  \item Khi producer gửi xong và watermark đã qua ngày cuối, đọc bảng bằng Trino hoặc DuckDB: kiểm không có \code{window\_start} nào xuất hiện hai lần, và tổng \code{orders} bằng số đơn trong CSV (trừ các ngày chưa đóng cửa sổ).
  \item Đếm snapshot qua bảng metadata \code{\$snapshots}; tính số snapshot mỗi ngày nếu checkpoint mỗi 30 giây, và chạy expire snapshot theo Iceberg Maintenance.
\end{enumerate}
\hoanthanh{Sau ít nhất một lần giết TaskManager, bảng Iceberg không có cửa sổ trùng và tổng số đơn khớp nguồn; bạn nêu được số snapshot đã sinh và lệnh bảo trì đã chạy.}

\section{Kế hoạch theo tuần}
\begin{xltabular}{\textwidth}{@{}P{1.6cm} L P{1.8cm}@{}}
\toprule
\thd{tuần} & \thd{việc} & \thd{giờ} \\
\midrule
\endhead
\bottomrule
\endfoot
1 & Kafka docs Introduction, Quickstart; Kafka 4.4 Design các mục The Consumer, Replication, Log Compaction; Lab 1 & 10 \\
2 & Design: Message Delivery Semantics, Using Transactions, The Share Consumer; Flink Timely Stream Processing, Window Aggregation; Lab 2, Lab 3 & 10 \\
3 & Flink Fault Tolerance; Iceberg Flink, Flink writes, Kafka Connect, Maintenance; Lab 4 và thử thách kỹ năng & 10 \\
thêm & Iceberg Kafka Connect sink chép \code{olist.orders} vào bảng raw; Redpanda thay broker; flink-training; Confluent và Zoomcamp & 10 trở lên \\
\end{xltabular}

\begin{thuthach}
Đội vận hành muốn biết gần như ngay lập tức khi tỉ lệ đơn bị huỷ của một bang vượt 5\% trong cửa sổ trượt 1 ngày, bước 1 giờ. Tự thiết kế topic và key, truy vấn Flink SQL (gợi ý: \code{HOP}), độ trễ watermark, cách xử lý sự kiện trễ, bảng Iceberg nhận kết quả, khoảng checkpoint, và lịch bảo trì bảng. Chạy trên dữ liệu Olist phát lại.
\hoanthanh{Job chạy liên tục, chịu được một lần khởi động lại mà không ghi trùng, và bạn trình bày được các bang vượt ngưỡng cùng lý do chọn từng tham số (độ trễ watermark, khoảng checkpoint, lịch compaction).}
\end{thuthach}

\section{Câu hỏi tự kiểm tra}
\begin{enumerate}
  \item Vì sao Kafka bảo đảm thứ tự trong một partition mà không bảo đảm trong cả topic? Key giúp gì?
  \item Consumer group lưu cái gì để biết đã đọc tới đâu? Điều gì xảy ra với message đã xử lý nhưng chưa commit offset khi tiến trình chết?
  \item Share group khác consumer group ở điểm nào? Khi nào bạn chọn nó?
  \item Retention và compaction: cái nào hợp với topic chứa trạng thái mới nhất của mỗi khách hàng?
  \item Idempotent producer bảo đảm gì, và transaction thêm gì?
  \item Vì sao exactly-once bên trong Kafka không tự áp sang PostgreSQL hay Iceberg?
  \item Cho một ví dụ trên dữ liệu Olist mà event time và processing time lệch nhau nhiều ngày.
  \item Watermark tại \code{2017-11-24 00:00} khẳng định điều gì? Vì sao sự kiện cũ hơn vẫn có thể tới?
  \item Checkpoint của Flink chứa gì, và sink Iceberg dùng nó thế nào để tránh ghi trùng?
  \item Với khoảng commit mặc định của Iceberg Kafka Connect sink, mỗi bảng có bao nhiêu snapshot mỗi ngày? Cần bảo trì gì?
  \item Vì sao phải kiểm bảng tương thích phiên bản trước khi chọn bản Flink cho lab?
\end{enumerate}

\section{Điểm dễ nhầm}
\begin{xltabular}{\textwidth}{@{}P{4.0cm} L P{3.2cm}@{}}
\toprule
\thd{điểm dễ nhầm} & \thd{thực tế} & \thd{docs chính thức} \\
\midrule
\endhead
\bottomrule
\endfoot
Nhiều người nghĩ Kafka xoá message khi đã đọc & Message ở lại tới khi retention hay compaction bỏ; consumer chỉ dời offset và có thể tua lại & \href{https://kafka.apache.org/44/design/design/}{Kafka Design} \\
Topic có thứ tự toàn cục & Thứ tự chỉ đúng trong một partition & \href{https://kafka.apache.org/44/design/design/}{Kafka Design} \\
Kafka cần ZooKeeper & Từ 4.0, Kafka chạy hoàn toàn bằng KRaft & \href{https://kafka.apache.org/blog/2025/03/18/apache-kafka-4.0.0-release-announcement/}{Kafka 4.0} \\
Exactly-once đầu cuối tự có khi bật transaction & Exactly-once sang hệ khác cần hệ đích hợp tác: commit theo checkpoint hoặc ghi idempotent & \href{https://kafka.apache.org/44/design/design/}{Kafka Design} \\
Watermark bảo đảm không còn dữ liệu trễ & Watermark là ước lượng tiến độ event time; vẫn cần chính sách cho sự kiện trễ & \href{https://nightlies.apache.org/flink/flink-docs-stable/docs/concepts/time/}{Flink time} \\
Window theo processing time là đủ cho báo cáo & Kết quả khi đó phụ thuộc lúc dữ liệu tới; phát lại cho kết quả khác & \href{https://nightlies.apache.org/flink/flink-docs-stable/docs/concepts/time/}{Flink time} \\
Ghi Iceberg từ luồng mỗi vài giây không sao & Mỗi commit thêm snapshot và file; cần compaction và expire snapshot & \href{https://iceberg.apache.org/docs/latest/maintenance/}{Iceberg Maintenance} \\
Bản Flink mới nhất luôn dùng được với mọi connector & Kafka connector 5.0.0 ghi hỗ trợ Flink 2.1.x và 2.2.x; kiểm trước khi nâng & \href{https://flink.apache.org/downloads/}{Flink downloads} \\
\end{xltabular}


% ==================================================================
\modulechapter{13}{Cloud và vận hành production}
\moduleinfo{Nâng cao}{30 giờ}{Module 5, 6, 9}

Đến đây toàn bộ lakehouse chạy trên laptop. Module này đưa một lát của nó lên AWS theo cách an toàn chi phí: cảnh báo chi phí có trước mọi tài nguyên, mọi tài nguyên dựng bằng Terraform và xoá được bằng \code{terraform destroy}. Bạn sẽ đặt dữ liệu Olist lên S3, khai báo bảng trong Glue Data Catalog, truy vấn bằng Athena, rồi thử bảng Iceberg qua endpoint Iceberg REST của Glue hoặc S3 Tables. Cuối module có bảng ánh xạ sang Google Cloud và thông tin hai chứng chỉ Data Engineer phổ biến.

\section{Bạn sẽ học được}
\begin{itemize}
  \item Giải thích gói Free và gói Paid của AWS Free Tier hiện hành, và dựng cảnh báo chi phí trước khi tạo tài nguyên.
  \item Viết IAM policy theo quyền tối thiểu cho một job đọc một prefix và ghi một prefix khác.
  \item Chọn storage class và lifecycle rule cho S3, tránh bẫy với object nhỏ.
  \item Truy vấn dữ liệu trên S3 bằng Athena qua Glue Data Catalog và giảm byte quét.
  \item Nối engine Iceberg vào endpoint Iceberg REST của Glue hoặc S3 Tables, biết giới hạn của từng endpoint.
  \item Dựng và xoá toàn bộ tài nguyên bằng Terraform.
  \item Ánh xạ các dịch vụ AWS sang Google Cloud và đọc exam guide của DEA-C01 và Professional Data Engineer.
\end{itemize}

\section{Nội dung chính}

\subsection{Tài khoản và chi phí}
\begin{enumerate}
  \item \textbf{AWS Free Tier hiện hành.} Khi đăng ký, bạn nhận 100 USD credit và có thể nhận thêm tới 100 USD qua các hoạt động; bạn chọn gói Free hoặc Paid. Gói Free không tính tiền, chỉ dùng một số dịch vụ, và kết thúc sau 6 tháng hoặc khi hết credit, tuỳ điều nào tới trước; khi đó tài khoản tự đóng, AWS giữ dữ liệu 90 ngày, nâng lên gói Paid trong thời gian này thì giữ được tài khoản và phần credit còn lại (\href{https://docs.aws.amazon.com/awsaccountbilling/latest/aboutv2/free-tier-plans.html}{Choosing a Free Tier plan}). Theo \href{https://aws.amazon.com/free/terms}{AWS Free Tier Terms}, credit hết hạn sau 12 tháng kể từ ngày mở tài khoản; nên credit chỉ dùng tới tháng thứ 12 nếu bạn đã chuyển sang gói Paid.
  \item \textbf{Dịch vụ trong gói Free.} Danh sách cho kiểu đăng ký mới có S3, Glue, Athena, Lake Formation, VPC, IAM, Budgets, Cost Explorer; Amazon EMR chỉ có ở gói Paid (\href{https://docs.aws.amazon.com/accounts/latest/reference/supported-services-sign-up-new.md}{Supported services}). Kiểm lại danh sách trước khi dùng một dịch vụ mới, kể cả EMR Serverless và S3 Tables.
  \item \textbf{Budgets.} AWS Budgets theo dõi chi tiêu và gửi cảnh báo; cảnh báo không tự dừng chi tiêu, trừ khi bạn cấu hình budget action (\href{https://docs.aws.amazon.com/cost-management/latest/userguide/budgets-managing-costs.html}{AWS Budgets}, \href{https://aws.amazon.com/aws-cost-management/aws-budgets/pricing/}{Budgets pricing}). Bật thêm \href{https://docs.aws.amazon.com/cost-management/latest/userguide/getting-started-ad.html}{Cost Anomaly Detection}.
  \item \textbf{IAM.} Bảo vệ root bằng MFA và không tạo access key cho root; dùng danh tính có MFA cho người, và role với credential tạm thời cho workload; cấp quyền tối thiểu (\href{https://docs.aws.amazon.com/IAM/latest/UserGuide/best-practices.html}{IAM best practices}). Role có trust policy (ai được assume) tách khỏi permissions policy (được làm gì) (\href{https://docs.aws.amazon.com/IAM/latest/UserGuide/id_roles.html}{IAM roles}).
\end{enumerate}

\subsection{Lưu trữ và truy vấn}
\begin{enumerate}[resume]
  \item \textbf{S3 storage class.} Standard-IA tính object dưới 128 KB như 128 KB và tính tối thiểu 30 ngày lưu, nên không hợp với file nhỏ hay sống ngắn (\href{https://docs.aws.amazon.com/AmazonS3/latest/userguide/storage-class-intro.html}{S3 storage classes}).
  \item \textbf{Lifecycle.} Từ 09/2024, lifecycle mặc định không chuyển object dưới 128 KB sang storage class khác; muốn chuyển phải thêm bộ lọc kích thước (\code{ObjectSizeGreaterThan}, \code{ObjectSizeLessThan}), và mỗi lần chuyển là một transition request có tính phí (\href{https://docs.aws.amazon.com/AmazonS3/latest/userguide/lifecycle-transition-general-considerations.html}{Lifecycle transitions}). Bật \href{https://docs.aws.amazon.com/AmazonS3/latest/userguide/Versioning.html}{versioning} và dùng lifecycle để xoá version cũ.
  \item \textbf{Glue Data Catalog.} Catalog chứa database và bảng mà Athena, Glue, EMR dùng chung; có thể khai báo bảng bằng tay, bằng Terraform, hoặc bằng crawler (\href{https://docs.aws.amazon.com/glue/latest/dg/catalog-and-crawler.html}{Data Catalog và crawler}).
  \item \textbf{Endpoint Iceberg REST của Glue.} Dạng \code{https://glue.<region>.\allowbreak amazonaws.com/\allowbreak iceberg}; \code{warehouse} là account ID; engine ký SigV4 với signing name \code{glue}. Chỉ hỗ trợ namespace một cấp; mặc định tạo bảng spec v2; lưu được trên S3 thường và S3 Tables; với client OAuth, Glue áp quyền Lake Formation ở mỗi request (\href{https://docs.aws.amazon.com/glue/latest/dg/connect-glu-iceberg-rest.html}{Glue Iceberg REST endpoint}).
  \item \textbf{S3 Tables.} Table bucket lưu bảng Iceberg với bảo trì do AWS quản lý. Endpoint REST riêng có dạng \code{https://s3tables.<region>.\allowbreak amazonaws.com/\allowbreak iceberg}, \code{warehouse} là ARN của table bucket, ký SigV4 với signing name \code{s3tables}. Giới hạn: chỉ namespace một cấp, không có API view, không hỗ trợ \code{stage-create} nên \code{CREATE TABLE AS SELECT} không chạy, \code{DROP TABLE} phải có \code{PURGE}, không hỗ trợ \code{metadata.json} trên 50 MB; compaction và quản lý snapshot làm qua API bảo trì của S3 Tables. AWS khuyên dùng endpoint của Glue nếu cần quản trị tập trung và phân quyền chi tiết (\href{https://docs.aws.amazon.com/AmazonS3/latest/userguide/s3-tables-integrating-open-source.html}{S3 Tables REST endpoint}).
  \item \textbf{Athena.} Athena đọc bảng trong Glue Data Catalog, kể cả bảng Iceberg (\href{https://docs.aws.amazon.com/athena/latest/ug/querying-iceberg.html}{Querying Iceberg tables}). Giá tính theo byte quét, làm tròn lên MB, tối thiểu 10 MB mỗi câu; ví dụ trên trang giá dùng 5 USD mỗi TB; câu lỗi và câu DDL không tính (\href{https://aws.amazon.com/athena/pricing/}{Athena pricing}). Parquet, partition và chọn đúng cột làm giảm byte quét.
  \item \textbf{Tính toán Spark có quản lý.} Glue 6.0 (21/08/2026) chạy Spark 4.1, Python 3.13, hỗ trợ đầy đủ Iceberg v3, giảm giá 30\% (\href{https://aws.amazon.com/about-aws/whats-new/2026/08/aws-glue-6-0-price-reduction-iceberg-v3/}{What's New: Glue 6.0}). Glue job cấu hình theo \code{MaxCapacity} cần từ 2 tới 100 DPU, mặc định 10; timeout mặc định 480 phút với Glue 5.0 trở lên, nên luôn đặt timeout ngắn cho job thử (\href{https://docs.aws.amazon.com/glue/latest/dg/aws-glue-api-jobs-job.html}{Glue Jobs API}). EMR Serverless chạy Spark mà không quản lý cluster (\href{https://docs.aws.amazon.com/emr/latest/EMR-Serverless-UserGuide/getting-started.html}{EMR Serverless}).
  \item \textbf{Mạng.} NAT gateway tính theo giờ khi còn tồn tại, giờ lẻ tính tròn, cộng phí mỗi GB xử lý; mỗi địa chỉ public IPv4 cũng tính theo giờ (\href{https://aws.amazon.com/vpc/pricing/}{VPC pricing}). Trong lab, tránh tạo NAT gateway.
\end{enumerate}

\subsection{Hạ tầng dưới dạng code và vận hành}
\begin{enumerate}[resume]
  \item \textbf{Terraform.} Viết tài nguyên bằng HCL, \code{terraform plan} xem thay đổi, \code{apply} tạo, \code{destroy} xoá những gì có trong state; thứ tạo bằng tay ngoài state vẫn còn và vẫn tính tiền (\href{https://developer.hashicorp.com/terraform/tutorials/aws-get-started}{Get Started - AWS}). State chứa thông tin nhạy cảm; với nhóm, lưu state ở backend từ xa như \href{https://developer.hashicorp.com/terraform/language/backend/s3}{S3 backend}. Từ 10/08/2023, HashiCorp chuyển các bản phát hành mới của sản phẩm sang Business Source License 1.1; người dùng cuối vẫn được dùng, sửa, phân phối cho mục đích thương mại và phi thương mại, trừ trường hợp cung cấp sản phẩm cạnh tranh với HashiCorp (\href{https://www.hashicorp.com/en/blog/hashicorp-adopts-business-source-license}{HashiCorp adopts BSL}, \href{https://www.hashicorp.com/en/license-faq}{Licensing FAQ}).
  \item \textbf{Nguyên tắc kiến trúc.} \href{https://docs.aws.amazon.com/wellarchitected/latest/analytics-lens/analytics-lens.html}{Well-Architected Data Analytics Lens} gom nguyên tắc theo các trụ cột (vận hành, bảo mật, tin cậy, hiệu năng, chi phí, bền vững) cho hệ thống phân tích; dùng nó làm danh sách rà soát thiết kế. Bản này có trước S3 Tables.
  \item \textbf{Sao lưu và khôi phục.} Snapshot Iceberg nằm cùng storage và bị expire, nên không thay được bản sao lưu; kết hợp versioning, sao chép sang nơi khác và đặt mục tiêu RPO, RTO (\href{https://docs.aws.amazon.com/whitepapers/latest/disaster-recovery-workloads-on-aws/disaster-recovery-workloads-on-aws.html}{Disaster Recovery of Workloads on AWS}).
\end{enumerate}

\subsection{Ánh xạ sang Google Cloud}
\begin{xltabular}{\textwidth}{@{}P{3.6cm} P{4.4cm} L@{}}
\toprule
\thd{AWS} & \thd{Google Cloud} & \thd{ghi chú} \\
\midrule
\endhead
\bottomrule
\endfoot
S3 & Cloud Storage & object storage \\
Glue Data Catalog, Iceberg REST & Lakehouse runtime catalog (trước gọi là BigLake metastore) & từ 20/04/2026, BigLake đổi tên thành Lakehouse; API, thư viện và lệnh \code{gcloud biglake} giữ tên cũ (\href{https://docs.cloud.google.com/lakehouse/docs/release-notes}{Lakehouse release notes}) \\
Athena & BigQuery & BigQuery sandbox không cần thẻ: 10 GiB lưu trữ, 1 TiB truy vấn mỗi tháng, bảng hết hạn sau 60 ngày (\href{https://cloud.google.com/bigquery/docs/sandbox}{BigQuery sandbox}) \\
EMR, EMR Serverless & Managed Service for Apache Spark (trước là Dataproc) & \href{https://docs.cloud.google.com/managed-spark/docs}{Managed Service for Apache Spark} \\
MWAA & Cloud Composer & Airflow có quản lý \\
\end{xltabular}

\subsection{Chứng chỉ}
Module này phủ phần dịch vụ cloud của hai chứng chỉ AWS Certified Data Engineer -- Associate và Google Cloud Professional Data Engineer. Chi tiết đề thi và cách các module ánh xạ vào từng mục của exam guide nằm ở chương ``Ánh xạ tới chứng chỉ'' trong Phần V. Với người mới, hãy dùng exam guide làm đề cương học trước, thi sau khi đã có kinh nghiệm thực tế.

\section{Bắt đầu}
\begin{enumerate}
  \item Đọc \href{https://docs.aws.amazon.com/awsaccountbilling/latest/aboutv2/free-tier-plans.html}{Choosing a Free Tier plan} trước khi mở tài khoản.
  \item Làm \href{https://developer.hashicorp.com/terraform/tutorials/docker-get-started}{Terraform: Get Started - Docker} trên laptop (không tốn tiền cloud), rồi \href{https://developer.hashicorp.com/terraform/tutorials/aws-get-started}{Get Started - AWS}.
  \item Làm \href{https://docs.aws.amazon.com/athena/latest/ug/getting-started.html}{Athena: Getting started} sau khi đã có budget.
\end{enumerate}

\section{Tài liệu chính thức}
\begin{xltabular}{\textwidth}{@{}P{4.2cm} L P{1.8cm}@{}}
\toprule
\thd{tài liệu} & \thd{nội dung cần đọc} & \thd{thời lượng} \\
\midrule
\endhead
\bottomrule
\endfoot
\href{https://docs.aws.amazon.com/awsaccountbilling/latest/aboutv2/free-tier-plans.html}{AWS Free Tier plans}, \href{https://aws.amazon.com/free/terms}{Terms} & gói Free, gói Paid, thời hạn credit & 0,5 giờ \\
\href{https://docs.aws.amazon.com/IAM/latest/UserGuide/best-practices.html}{IAM best practices} & MFA, role, credential tạm, quyền tối thiểu & 1,5 giờ \\
\href{https://docs.aws.amazon.com/cost-management/latest/userguide/budgets-managing-costs.html}{AWS Budgets} & budget, ngưỡng cảnh báo, budget action & 0,5 giờ \\
\href{https://docs.aws.amazon.com/AmazonS3/latest/userguide/storage-class-intro.html}{S3 storage classes}, \href{https://docs.aws.amazon.com/AmazonS3/latest/userguide/lifecycle-transition-general-considerations.html}{Lifecycle transitions} & mức tính tối thiểu, mặc định 128 KB, bộ lọc kích thước & 1,5 giờ \\
\href{https://docs.aws.amazon.com/glue/latest/dg/catalog-and-crawler.html}{Glue Data Catalog}, \href{https://docs.aws.amazon.com/glue/latest/dg/connect-glu-iceberg-rest.html}{Glue Iceberg REST} & database, bảng, crawler; endpoint REST & 1,5 giờ \\
\href{https://docs.aws.amazon.com/athena/latest/ug/querying-iceberg.html}{Athena: Querying Iceberg}, \href{https://aws.amazon.com/athena/pricing/}{Pricing} & tạo và truy vấn bảng Iceberg; cách tính tiền & 1,5 giờ \\
\href{https://docs.aws.amazon.com/AmazonS3/latest/userguide/s3-tables-integrating-open-source.html}{S3 Tables REST endpoint} & cấu hình và giới hạn & 1 giờ \\
\href{https://docs.aws.amazon.com/emr/latest/EMR-Serverless-UserGuide/getting-started.html}{EMR Serverless: Getting started} & application, job run & 0,5 giờ \\
\href{https://docs.aws.amazon.com/prescriptive-guidance/latest/apache-iceberg-on-aws/introduction.html}{Using Apache Iceberg on AWS} & Athena, Glue, EMR, best practices & 2 giờ \\
\href{https://developer.hashicorp.com/terraform/tutorials/aws-get-started}{Terraform: Get Started - AWS} & chuỗi tutorial cơ bản, state & 2 giờ \\
\href{https://docs.aws.amazon.com/wellarchitected/latest/analytics-lens/analytics-lens.html}{Data Analytics Lens} & nguyên tắc theo trụ cột, kịch bản batch & 1 giờ \\
\href{https://docs.aws.amazon.com/aws-certification/latest/data-engineer-associate-01/data-engineer-associate-01.html}{DEA-C01 exam guide}, \href{https://cloud.google.com/learn/certification/data-engineer}{Google PDE} & domain và kỹ năng & 1 giờ \\
\end{xltabular}

\section{Khoá học miễn phí}
\begin{itemize}
  \item \href{https://developer.hashicorp.com/terraform/tutorials/aws-get-started}{HashiCorp tutorials: Get Started - AWS}: miễn phí; chỉ phần tài nguyên AWS tạo ra mới có thể tính tiền.
  \item \href{https://github.com/DataTalksClub/data-engineering-zoomcamp}{DE Zoomcamp}: module 1 ``Containerization and Infrastructure as Code'' dùng Docker và Terraform trên Google Cloud; miễn phí.
\end{itemize}

\section{Lab thực hành}

Các lab đi theo thang an toàn chi phí: từ không tốn tiền tới có thể tốn tiền. Mỗi lab trên AWS kết thúc bằng \code{terraform destroy}, và hôm sau bạn mở Billing console kiểm lại.

\subsection{Lab 1: Terraform trên laptop}
Làm chuỗi Get Started - Docker: tạo một container bằng Terraform, đổi cấu hình, \code{terraform destroy}.
\hoanthanh{\code{terraform state list} trả rỗng sau \code{destroy}, và \code{docker ps} không còn container của lab.}

\subsection{Lab 2: An toàn tài khoản trước tiên}
\begin{enumerate}
  \item Mở tài khoản, chọn gói Free nếu chỉ học; bật MFA cho root, không tạo access key cho root.
  \item Tạo một danh tính quản trị có MFA cho việc hằng ngày; cấu hình AWS CLI bằng credential tạm thời.
  \item Tạo budget chi phí hằng tháng nhỏ (ví dụ 5 USD) với email cảnh báo ở 50\%, 80\%, 100\%; bật Cost Anomaly Detection.
\end{enumerate}
\hoanthanh{Bạn nhận được email xác nhận đăng ký cảnh báo, và \code{aws sts get-caller-\allowbreak identity} trả danh tính không phải root.}

\subsection{Lab 3: S3, Glue và Athena trên Olist bằng Terraform}
\begin{enumerate}
  \item Viết Terraform cho: một bucket có versioning, chặn public access, lifecycle xoá version cũ sau 7 ngày; một Glue database \code{olist}; một Athena workgroup có giới hạn byte quét mỗi câu và vị trí lưu kết quả.
\end{enumerate}
\begin{codeblock}
resource "aws_s3_bucket" "lake" {
  bucket        = "olist-lake-<suffix>"
  force_destroy = true   # cho phép destroy cả khi còn object (chỉ dùng cho lab)
}

resource "aws_glue_catalog_database" "olist" {
  name = "olist"
}

resource "aws_athena_workgroup" "lab" {
  name          = "olist-lab"
  force_destroy = true
  configuration {
    bytes_scanned_cutoff_per_query = 1073741824   # 1 GB
    result_configuration {
      output_location = "s3://${aws_s3_bucket.lake.bucket}/athena-results/"
    }
  }
}
\end{codeblock}
\begin{enumerate}[resume]
  \item Tải lên \code{orders} dạng CSV và dạng Parquet (Parquet tạo bằng DuckDB, có thêm bản partition theo năm và tháng mua).
  \item Khai báo ba bảng trong Glue (bằng DDL của Athena hoặc Terraform) và chạy cùng một câu ``số đơn theo trạng thái trong 2018'' trên cả ba.
  \item Ghi lại ``Data scanned'' của mỗi câu.
\end{enumerate}
\hoanthanh{Bảng ba dòng (CSV, Parquet, Parquet có partition) với byte quét và chi phí ước tính theo 5 USD mỗi TB, tối thiểu 10 MB mỗi câu.}

\subsection{Lab 4: Bảng Iceberg trên AWS, rồi dọn sạch}
\begin{enumerate}
  \item Tạo bảng Iceberg \code{olist.\allowbreak orders\_iceberg} bằng Athena (\code{TBLPROPERTIES ('table\_type'=\allowbreak 'ICEBERG')}) và nạp từ bảng Parquet; chạy một \code{UPDATE} hoặc \code{MERGE} rồi xem lịch sử snapshot.
  \item Từ Spark hoặc PyIceberg trên laptop, nối vào endpoint Iceberg REST của Glue (SigV4, signing name \code{glue}, \code{warehouse} là account ID) và đọc cùng bảng.
  \item Tuỳ chọn: tạo table bucket S3 Tables, ghi một bảng Olist qua endpoint REST của S3 Tables, đối chiếu với danh sách giới hạn; xoá table bucket khi xong.
  \item \code{terraform destroy}; xoá tay mọi thứ ngoài state; hôm sau kiểm Billing và Cost Explorer.
\end{enumerate}
\hoanthanh{Cùng một bảng Iceberg được đọc từ Athena và từ engine trên laptop; sau \code{destroy}, không còn bucket, database hay workgroup của lab, và Billing không có dịch vụ lạ.}

\section{Kế hoạch theo tuần}
\begin{xltabular}{\textwidth}{@{}P{1.6cm} L P{1.8cm}@{}}
\toprule
\thd{tuần} & \thd{việc} & \thd{giờ} \\
\midrule
\endhead
\bottomrule
\endfoot
1 & Free Tier plans, IAM best practices, Budgets; Terraform Docker và AWS; Lab 1, Lab 2 & 10 \\
2 & S3 storage classes, lifecycle; Glue Data Catalog; Athena Iceberg và pricing; Lab 3 & 10 \\
3 & Glue Iceberg REST, S3 Tables, Iceberg on AWS, Analytics Lens; Lab 4; thử thách kỹ năng; \code{terraform destroy} và kiểm Billing & 10 \\
thêm & EMR Serverless trên cùng bảng; Lake Formation; DR whitepaper; đọc kỹ hai exam guide; Zoomcamp module 1 trên Google Cloud & 10 trở lên \\
\end{xltabular}

\begin{thuthach}
Dựng bằng một lệnh \code{terraform apply} một ``lát cloud'' cho Olist: bucket với prefix \code{raw/} và \code{curated/}, một IAM role cho job (chỉ đọc \code{raw/}, đọc và ghi \code{curated/}, không có quyền gì khác), Glue database, Athena workgroup có giới hạn byte quét, và budget có cảnh báo. Chứng minh role không ghi được vào \code{raw/}. Sau đó xoá toàn bộ bằng một lệnh \code{terraform destroy}.
\hoanthanh{Chạy \code{apply} rồi \code{destroy} hai lần liên tiếp trên tài khoản sạch đều thành công; lệnh ghi vào \code{raw/} bằng role bị từ chối; IAM Access Analyzer policy validation không báo lỗi hay cảnh báo bảo mật.}
\end{thuthach}

\section{Câu hỏi tự kiểm tra}
\begin{enumerate}
  \item Gói Free của AWS kết thúc khi nào, và điều gì xảy ra với dữ liệu khi nó kết thúc?
  \item Vì sao một budget alert không ngăn được hoá đơn lớn?
  \item Trust policy và permissions policy của một IAM role khác nhau thế nào?
  \item Vì sao job chạy trên AWS nên dùng role thay vì access key lâu dài?
  \item Vì sao lifecycle rule chuyển file Parquet 50 KB sang Standard-IA mặc định không làm gì, và ép chuyển thì có lợi không?
  \item Athena tính tiền một câu truy vấn thế nào? Kể ba cách giảm byte quét trên bảng orders.
  \item Endpoint Iceberg REST của Glue và của S3 Tables khác nhau ở điểm nào (warehouse, signing name, giới hạn)?
  \item Vì sao \code{CREATE TABLE AS SELECT} không chạy qua endpoint REST của S3 Tables?
  \item Cấu hình nào ngăn một Glue job thử chạy lâu và tốn nhiều DPU?
  \item \code{terraform destroy} xoá gì và không xoá gì?
  \item Vì sao snapshot Iceberg không phải bản sao lưu?
  \item Dịch vụ Google Cloud nào tương ứng với Glue Data Catalog, Athena và EMR?
\end{enumerate}

\section{Điểm dễ nhầm}
\begin{xltabular}{\textwidth}{@{}P{4.0cm} L P{3.2cm}@{}}
\toprule
\thd{điểm dễ nhầm} & \thd{thực tế} & \thd{docs chính thức} \\
\midrule
\endhead
\bottomrule
\endfoot
Nhiều người nghĩ Free Tier là ``12 tháng miễn phí'' & Tài khoản mới chọn gói Free (tối đa 6 tháng, dùng credit) hoặc gói Paid; credit hết hạn sau 12 tháng & \href{https://docs.aws.amazon.com/awsaccountbilling/latest/aboutv2/free-tier-plans.html}{Free Tier plans}, \href{https://aws.amazon.com/free/terms}{Terms} \\
Budget alert dừng việc tiêu tiền & Nó chỉ báo, trừ khi cấu hình budget action & \href{https://aws.amazon.com/aws-cost-management/aws-budgets/pricing/}{Budgets pricing} \\
Infrequent Access luôn rẻ hơn & Không với object nhỏ hay sống ngắn: tính tối thiểu 128 KB và 30 ngày; lifecycle mặc định không chuyển object dưới 128 KB & \href{https://docs.aws.amazon.com/AmazonS3/latest/userguide/storage-class-intro.html}{S3 storage classes}, \href{https://docs.aws.amazon.com/AmazonS3/latest/userguide/lifecycle-transition-general-considerations.html}{Lifecycle} \\
Athena miễn phí khi câu trả ít dòng & Tính theo byte quét, tối thiểu 10 MB; chỉ câu lỗi và DDL không tính & \href{https://aws.amazon.com/athena/pricing/}{Athena pricing} \\
Endpoint REST của S3 Tables chạy mọi câu lệnh Iceberg & Không có \code{stage-create}, view API, namespace nhiều cấp; \code{DROP TABLE} cần \code{PURGE} & \href{https://docs.aws.amazon.com/AmazonS3/latest/userguide/s3-tables-integrating-open-source.html}{S3 Tables REST} \\
Xoá EC2 là hết chi phí mạng & NAT gateway tính theo giờ tới khi bị xoá; public IPv4 cũng tính theo giờ & \href{https://aws.amazon.com/vpc/pricing/}{VPC pricing} \\
\code{terraform destroy} xoá mọi thứ & Chỉ xoá tài nguyên có trong state & \href{https://developer.hashicorp.com/terraform/tutorials/aws-get-started}{Terraform AWS} \\
DEA-C01 có 65 câu tính điểm & 50 câu tính điểm, 15 câu không tính điểm và không được đánh dấu & \href{https://docs.aws.amazon.com/aws-certification/latest/data-engineer-associate-01/data-engineer-associate-01.html}{DEA-C01 exam guide} \\
\end{xltabular}


% ==================================================================


\partpage{V}{Capstone, chứng chỉ và nghề nghiệp}{Gom mọi module thành một dự án portfolio, ánh xạ sang chứng chỉ của AWS và Google Cloud, chuẩn bị đi làm.}
% ---------------- capstone ----------------
% ====================== CAPSTONE ======================
\plainchapter{Dự án capstone: lakehouse cho Olist}

Dự án capstone gom mọi module thành một hệ thống chạy thật, có thể đưa vào portfolio. Bạn xây nó dần qua bốn mốc, mỗi mốc ứng với một phần của lộ trình, giống cách learning path chính thức kết thúc từng chặng bằng một skill badge.

\section{Mục tiêu sản phẩm}
Một lakehouse cho dữ liệu Olist mà người khác clone về là chạy được bằng một lệnh, gồm:
\begin{itemize}
  \item tầng \code{raw}, \code{staging}, \code{core}, \code{marts} là các bảng Iceberg trên object storage, quản lý bằng một REST catalog;
  \item pipeline tăng dần, chạy lại cho cùng một kết quả, điều phối bằng Airflow;
  \item test dữ liệu chặn tầng \code{marts} khi \code{core} sai;
  \item dashboard trên Trino và Metabase trả lời được các câu hỏi nghiệp vụ;
  \item tài liệu đủ để người khác vận hành: README, sơ đồ kiến trúc, runbook.
\end{itemize}

\section{Bốn mốc}

\subsection{Mốc 1: mô hình dữ liệu (sau Phần I)}
\begin{itemize}
  \item Viết 5 câu hỏi nghiệp vụ, ví dụ doanh thu theo tháng và bang, thời gian giao hàng so với dự kiến, điểm đánh giá theo danh mục sản phẩm.
  \item Thiết kế star schema trả lời các câu đó: ghi rõ grain của từng bảng fact, khoá của từng bảng chiều, cách xử lý \code{customer\_unique\_id}.
  \item Viết truy vấn SQL kiểm grain trên CSV gốc bằng DuckDB.
\end{itemize}
\hoanthanh{mỗi bảng có một câu mô tả grain, và truy vấn kiểm khoá trả về 0 dòng trùng.}

\subsection{Mốc 2: raw và staging trên Iceberg (sau Phần II)}
\begin{itemize}
  \item Dựng stack tham chiếu bằng Docker Compose; tạo namespace \code{raw}, \code{staging}, \code{core}, \code{marts}.
  \item Nạp CSV vào \code{raw} dạng chỉ thêm, có cột thời điểm nạp và tên file nguồn.
  \item Viết bảng \code{staging} đã ép kiểu, chọn partition và sort order có lý do.
  \item Tạo hai principal trên catalog: một người ghi pipeline, một người chỉ đọc \code{marts}.
\end{itemize}
\hoanthanh{bạn chỉ ra được \code{metadata.json}, manifest list và manifest của một bảng trên object storage, và người dùng chỉ đọc bị từ chối khi ghi.}

\subsection{Mốc 3: pipeline tăng dần có test (sau Phần III)}
\begin{itemize}
  \item Chia CSV \code{orders} theo ngày mua để giả lập dữ liệu đến dần; nạp tăng dần bằng watermark và \code{MERGE} theo khoá.
  \item Viết DAG Airflow cho chuỗi raw, staging, core, marts; giới hạn số task ghi Iceberg cùng lúc bằng pool.
  \item Thêm test unique, not null, relationships trên \code{core}; test lỗi thì \code{marts} không được build.
  \item Bật CI chạy test mỗi lần push.
\end{itemize}
\hoanthanh{chạy lại cùng một ngày hai lần không làm đổi số dòng, và một dòng trùng cố ý chèn vào làm pipeline dừng trước \code{marts}.}

\subsection{Mốc 4: phục vụ và cloud (sau Phần IV)}
\begin{itemize}
  \item Nối Trino vào REST catalog, tạo dashboard Metabase trả lời 5 câu hỏi ở mốc 1.
  \item Tuỳ chọn: thêm nhánh streaming, đẩy sự kiện đơn hàng qua Kafka, Flink ghi vào Iceberg.
  \item Dựng bản trên AWS bằng Terraform: S3, Glue Data Catalog, Athena; đặt AWS Budgets trước khi tạo bất cứ thứ gì; xoá sạch bằng \code{terraform destroy}.
\end{itemize}
\hoanthanh{dashboard chạy trên dữ liệu của chính pipeline, và hoá đơn AWS của mốc này nằm trong ngân sách bạn đặt.}

\section{Tiêu chí tự đánh giá}
\begin{xltabular}{\textwidth}{@{}P{3.4cm} L L@{}}
\toprule
\thd{tiêu chí} & \thd{đạt} & \thd{tốt} \\
\midrule
\endhead
\bottomrule
\endfoot
Chạy được & \code{docker compose up} cộng một lệnh khởi tạo là chạy hết pipeline & có \code{make} target cho từng tầng, có dữ liệu mẫu nhỏ để chạy nhanh \\
Chạy lại an toàn & chạy lại không sinh trùng & backfill được một khoảng ngày bất kỳ \\
Chất lượng & có test trên khoá và quan hệ & test chặn tầng sau, có write-audit-publish cho một bảng \\
Mô hình dữ liệu & grain và khoá được ghi rõ & có SCD loại 2 cho ít nhất một bảng chiều \\
Vận hành & có lịch compaction và expire snapshot & có runbook xử lý lỗi thường gặp \\
Tài liệu & README, sơ đồ kiến trúc & ghi lại các quyết định thiết kế và lý do \\
Chi phí cloud & có budget và teardown & ghi lại chi phí thật của mỗi lần chạy \\
\end{xltabular}

% ====================== CHỨNG CHỈ ======================
\plainchapter{Ánh xạ tới chứng chỉ}

Lộ trình dạy nguyên lý bằng công cụ mã nguồn mở, nên không thay được phần ôn dịch vụ cụ thể của từng cloud. Bảng dưới cho biết module nào phủ phần nào của hai chứng chỉ data engineer phổ biến, để bạn biết cần bổ sung gì khi ôn.

\section{AWS Certified Data Engineer -- Associate (DEA-C01)}
Theo \href{https://docs.aws.amazon.com/aws-certification/latest/data-engineer-associate-01/data-engineer-associate-01.html}{exam guide chính thức}: 50 câu tính điểm và 15 câu không tính điểm, điểm đạt 720 trên thang 100 tới 1.000, chấm theo tổng chứ không yêu cầu đạt từng phần. Ứng viên mục tiêu có kinh nghiệm data engineering tương đương 2 tới 3 năm và ít nhất 1 tới 2 năm thực hành trên AWS. Theo \href{https://aws.amazon.com/certification/certified-data-engineer-associate/}{trang chứng chỉ}: thi 130 phút, 65 câu, phí 150 USD, có tiếng Anh, Nhật, Hàn và Trung giản thể.

\begin{xltabular}{\textwidth}{@{}P{5.2cm} r P{2.4cm} L@{}}
\toprule
\thd{domain} & \thd{tỷ trọng} & \thd{module} & \thd{cần bổ sung khi ôn} \\
\midrule
\endhead
\bottomrule
\endfoot
1. Data Ingestion and Transformation & 34\% & 7, 8, 9, 12 & Glue ETL, Kinesis, MSK, Step Functions, MWAA \\
2. Data Store Management & 26\% & 3, 4, 5, 6 & lựa chọn giữa S3, Redshift, DynamoDB; Lake Formation \\
3. Data Operations and Support & 22\% & 9, 10, 11, 13 & CloudWatch, Athena, giám sát pipeline \\
4. Data Security and Governance & 18\% & 6, 13 & IAM, KMS, Lake Formation, Macie, CloudTrail \\
\end{xltabular}

\section{Google Cloud Professional Data Engineer}
Theo \href{https://cloud.google.com/learn/certification/data-engineer}{trang chứng chỉ}: thi 2 giờ, 40 tới 50 câu, phí 200 USD, bằng tiếng Anh hoặc tiếng Nhật, hiệu lực 2 năm, không có điều kiện tiên quyết; Google khuyến nghị từ 3 năm kinh nghiệm, trong đó ít nhất 1 năm thiết kế và quản lý giải pháp trên Google Cloud. Trang này báo đề thi sắp được cập nhật theo tên sản phẩm mới, nên hãy đối chiếu \href{https://services.google.com/fh/files/misc/professional_data_engineer_exam_guide_english.pdf}{exam guide} trước khi ôn. Tên sản phẩm trong đề có thể còn là tên cũ. Google có \href{https://www.skills.google/paths/16}{learning path chính thức} gồm khoá học, lab và skill badge cho chứng chỉ này.

\begin{xltabular}{\textwidth}{@{}P{5.2cm} P{2.4cm} L@{}}
\toprule
\thd{mục của exam guide} & \thd{module} & \thd{cần bổ sung khi ôn} \\
\midrule
\endhead
\bottomrule
\endfoot
1. Design data processing systems & 2, 3, 5, 13 & lựa chọn dịch vụ, bảo mật và tuân thủ trên Google Cloud \\
2. Ingest and process the data & 7, 8, 9, 12 & Dataflow, Pub/Sub, Managed Service for Apache Spark (trước là Dataproc), Cloud Composer \\
3. Store the data & 4, 5, 6 & BigQuery, Cloud Storage, Bigtable, Spanner, Lakehouse (trước là BigLake) \\
4. Prepare and use data for analysis & 3, 11 & BigQuery cho phân tích, chia sẻ dữ liệu, BigQuery ML \\
5. Maintain and automate data workloads & 9, 10, 13 & giám sát, tối ưu chi phí và độ tin cậy \\
\end{xltabular}

% ---------------- career ----------------
\plainchapter{Nghề nghiệp: portfolio, phỏng vấn và cộng đồng}

Chương này không phải một module có giờ cố định; nó chạy song song suốt lộ trình, khoảng 1 giờ mỗi tuần, cộng 30--40 giờ ôn tập trước khi ứng tuyển. Mục tiêu là biến những gì bạn dựng trong các module thành bằng chứng mà nhà tuyển dụng kiểm được: một project chạy lại được, vài bài viết kỹ thuật, một đóng góp mã nguồn mở và khả năng giải thích đánh đổi khi phỏng vấn.

\section{Portfolio: một project lakehouse chạy lại được}
\begin{enumerate}
  \item \textbf{Một project, nhiều lát.} Gom lakehouse Olist qua các module vào một kho GitHub công khai: raw → staging/core → marts trên Iceberg, orchestration bằng Airflow, test chất lượng, lớp phục vụ Trino và Metabase, tuỳ chọn một luồng streaming và một lát cloud.
  \item \textbf{README là sản phẩm.} Nêu bài toán kinh doanh, sơ đồ kiến trúc, mô hình dữ liệu, ``cách chạy'' bằng ít lệnh trên máy sạch, phần chi phí và đánh đổi, và cách kiểm kết quả.
  \item \textbf{Dùng một rubric công khai.} \href{https://github.com/DataTalksClub/data-engineering-zoomcamp/blob/main/projects/README.md}{Rubric project của DE Zoomcamp} có 7 tiêu chí, mỗi tiêu chí tối đa 4 điểm: mô tả bài toán, cloud (đạt điểm tối đa khi dùng cloud cùng công cụ IaC), ingestion có orchestration, bảng warehouse được tối ưu, biến đổi dữ liệu, dashboard (tối đa khi có hai ô), khả năng chạy lại. Test, \code{make} và CI là gợi ý ``going the extra mile'', không được tính điểm.
  \item \textbf{Viết 2--3 bài kỹ thuật ngắn.} Một quyết định thiết kế (ví dụ chọn partition cho bảng orders), một lỗi hay sự cố bạn tìm ra và cách sửa, một số đo hiệu năng hay chi phí (ví dụ byte quét của Athena ở Module 13).
\end{enumerate}

\section{Thị trường việc làm ở Việt Nam: đọc số liệu đúng cách}
\href{https://itviec.com/bao-cao/luong-it-va-thi-truong-tuyen-dung-it-vietnam}{Báo cáo Lương \& Thị trường Tuyển dụng IT 2025--2026 của ITviec} khảo sát 1.839 người (hai đợt trực tuyến: 27/05--31/07 và 21/10--11/11/2025). Bảng lương ghi Data Engineer có trung vị 41,3 triệu đồng (đơn vị của bảng là triệu VND); theo số năm kinh nghiệm, chỉ nhóm 3--4 năm có số (56,9 triệu đồng), các nhóm dưới 1 năm, 1--2 năm, 5--8 năm và trên 8 năm để trống. Bản đầy đủ cần điền form để nhận qua email.

Vì vậy:
\begin{itemize}
  \item Trung vị 41,3 triệu gộp mọi cấp, không phải mức cho người mới.
  \item Báo cáo không có số cho nhóm dưới 3 năm, nên đừng suy ra mức junior từ báo cáo này.
  \item Cách thực tế hơn là đọc vài tin tuyển dụng thật và ghi lại kỹ năng lặp lại (xem bài tập ở dưới).
\end{itemize}

\section{Ôn phỏng vấn bằng tài liệu chính thức và miễn phí}
\begin{xltabular}{\textwidth}{@{}P{3.0cm} P{5.2cm} L@{}}
\toprule
\thd{mảng} & \thd{tài liệu} & \thd{cách dùng} \\
\midrule
\endhead
\bottomrule
\endfoot
SQL & \href{https://www.postgresql.org/docs/current/tutorial-window.html}{PostgreSQL: Window Functions}, \href{https://www.postgresql.org/docs/current/queries-with.html}{WITH Queries} & tự viết lại các câu trên Olist: top-N theo nhóm, đơn thứ hai của mỗi khách, chuỗi ngày liên tiếp (gaps and islands), doanh thu luỹ kế \\
Mô hình dữ liệu & \href{https://www.kimballgroup.com/data-warehouse-business-intelligence-resources/kimball-techniques/dimensional-modeling-techniques/}{Kimball Dimensional Modeling Techniques} & fact, dimension, grain, SCD loại 1 và 2 trên Olist \\
Spark & \href{https://spark.apache.org/docs/latest/sql-performance-tuning.html}{Spark SQL Performance Tuning} & shuffle, broadcast join, AQE, đọc Spark UI \\
Thiết kế hệ thống dữ liệu & \href{https://docs.aws.amazon.com/wellarchitected/latest/analytics-lens/scenarios.html}{Data Analytics Lens: Scenarios} & khung trình bày thiết kế batch và streaming: yêu cầu, luồng dữ liệu, đánh đổi \\
Đề cương kỹ năng & \href{https://docs.aws.amazon.com/aws-certification/latest/data-engineer-associate-01/data-engineer-associate-01.html}{DEA-C01 exam guide}, \href{https://cloud.google.com/learn/certification/data-engineer}{Google PDE} & dùng danh sách domain làm checklist ôn tập, kể cả khi chưa định thi \\
Streaming & \href{https://kafka.apache.org/44/design/design/}{Kafka 4.4 Design}, \href{https://nightlies.apache.org/flink/flink-docs-stable/docs/concepts/time/}{Flink time} & giải thích thứ tự, offset, exactly-once, watermark bằng ví dụ của chính bạn \\
\end{xltabular}

Với câu hỏi hành vi, chuẩn bị 3 câu chuyện theo cấu trúc tình huống, nhiệm vụ, hành động, kết quả, lấy từ chính project: một lỗi pipeline, một đánh đổi thiết kế, một lần tối ưu chi phí hay hiệu năng. Một sự cố thật trong pipeline của bạn thuyết phục hơn ví dụ trong sách.

\section{Đóng góp mã nguồn mở: bắt đầu từ PyIceberg}
PyIceberg hợp với stack tham chiếu vì bộ test dùng đúng các công cụ bạn đã chạy (\href{https://py.iceberg.apache.org/contributing/}{PyIceberg contributing}):
\begin{enumerate}
  \item \code{make install} cài uv nếu cần, tạo \code{.venv} và cài dependency; chọn bản Python bằng \code{PYTHON=3.12 make install}.
  \item \code{make test} chạy pytest có đo coverage; test S3 và ADLS bị bỏ qua mặc định: \code{make test-s3} cần RustFS đang chạy, \code{make test-adls} cần azurite.
  \item \code{make test-integration} khởi động lại container rồi chạy bộ test tích hợp; test Spark chạy Spark trong container để tạo bảng cho PyIceberg đọc.
  \item \code{make lint} chạy prek để format và lint; \code{prek install} cài pre-commit hook.
  \item Tìm việc trong các issue có nhãn \href{https://github.com/apache/iceberg-python/issues?q=is\%3Aissue+is\%3Aopen+label\%3A\%22good+first+issue\%22}{good first issue} và bình luận trước khi làm. Docs và test cũng là đóng góp có giá trị.
\end{enumerate}
\begin{luuy}
Không bao giờ chạy test tích hợp của PyIceberg vào catalog thật: bộ test xoá dữ liệu trên toàn catalog, docs ghi rõ ``Do not run against production catalogs'' (\href{https://py.iceberg.apache.org/contributing/}{PyIceberg contributing}).
\end{luuy}
Thảo luận thiết kế và bỏ phiếu của Apache Iceberg diễn ra trên mailing list; đăng ký dev list bằng thư tới \code{dev-subscribe@iceberg.apache.org} và xem các kênh khác ở \href{https://iceberg.apache.org/community/}{Iceberg community}.

\section{Cộng đồng, khoá học có lịch và sự kiện}
\begin{itemize}
  \item \href{https://github.com/DataTalksClub/data-engineering-zoomcamp}{DE Zoomcamp}: tự học bất cứ lúc nào; lớp có hạn nộp tiếp theo bắt đầu 01/2027. Tham gia lớp mới có chấm bài, chấm chéo project và chứng chỉ; tự học thì không.
  \item \href{https://meetup.com/ho-chi-minh-city-data-meetup-group}{HCMC Data Meetup}: nhóm meetup dữ liệu ở TP.HCM, từng tổ chức buổi bằng tiếng Việt.
  \item \href{https://kubernetes.io/vi/docs/tutorials/kubernetes-basics/}{Kubernetes Basics} có bản tiếng Việt chính thức, hữu ích khi bạn bắt đầu học vận hành container.
  \item Sự kiện của các dự án trong lộ trình: \href{https://www.icebergsummit.org/}{Iceberg Summit} (bản 2026 diễn ra 8--9/04/2026 ở San Francisco), \href{https://current.confluent.io/}{Current} (San Francisco, 4--5/11/2026), \href{https://www.databricks.com/dataaisummit}{Data + AI Summit} (21--24/06/2027). Phần lớn talk được đăng lại miễn phí sau sự kiện; xem talk về công cụ bạn đang dùng thay vì cố theo hết.
\end{itemize}

\section{Bài tập}
\begin{xltabular}{\textwidth}{@{}P{0.6cm} L P{1.8cm}@{}}
\toprule
\thd{\#} & \thd{bài} & \thd{giờ} \\
\midrule
\endhead
\bottomrule
\endfoot
1 & Tự chấm project theo 7 tiêu chí của rubric Zoomcamp; với mỗi tiêu chí dưới 4 điểm, ghi việc cần làm & 1 \\
2 & 30--50 câu SQL trên Olist bằng PostgreSQL hoặc DuckDB, có window function, CTE đệ quy và gaps and islands & 15--20 \\
3 & Mô hình hoá Olist thành star schema có một dimension SCD loại 2; viết một trang giải thích grain & 4 \\
4 & Thiết kế thử trong 45 phút: ``pipeline đơn hàng hằng ngày từ API vào lakehouse tới dashboard''; viết lại kèm đánh đổi theo Data Analytics Lens & 2 \\
5 & Ba câu chuyện hành vi lấy từ project & 3 \\
6 & Đọc 5 tin tuyển dụng Data Engineer thật: ba kỹ năng nào xuất hiện ở cả năm tin, project của bạn chứng minh được kỹ năng nào? & 1 \\
7 & Chạy \code{make install}, \code{make test} và \code{make test-s3} của PyIceberg với RustFS; chọn một good first issue & 3--5 \\
\end{xltabular}

\section{Danh sách tự kiểm tra}
\begin{enumerate}
  \item Một người lạ chạy được project từ README trên máy sạch trong khoảng 30 phút chưa?
  \item Với mỗi tiêu chí rubric bạn tự chấm 4 điểm, bằng chứng nằm ở file hay ảnh chụp nào?
  \item Bạn giải thích được vì sao trung vị lương trong báo cáo ITviec không phải mức lương cho người mới?
  \item Bạn viết được câu SQL trả đơn thứ hai của mỗi khách mà không dùng \code{LIMIT} hay \code{OFFSET} theo nhóm chưa?
  \item Bạn giải thích được khi nào chọn SCD loại 2 thay loại 1 cho dimension khách hàng?
  \item Bạn chỉ ra được trên Spark UI chỗ nào có shuffle và vì sao?
  \item Bạn làm ingestion Olist idempotent bằng cách nào, và chứng minh ra sao?
  \item Vì sao không bao giờ trỏ test tích hợp của PyIceberg vào catalog thật?
  \item Quyết định thiết kế của Iceberg được thảo luận và bỏ phiếu ở đâu?
  \item Bạn đã có ít nhất một bài viết kỹ thuật và một pull request hay issue công khai chưa?
\end{enumerate}


% ---------------- appendix ----------------
% ====================== PHỤ LỤC ======================
\plainchapter{Phụ lục A: Thuật ngữ dễ nhầm}

Nhiều từ trong data engineering có hai, ba nghĩa tuỳ ngữ cảnh. Khi đọc docs, hãy xác định bạn đang ở ngữ cảnh nào.

\begin{xltabular}{\textwidth}{@{}P{2.2cm} L L@{}}
\toprule
\thd{từ} & \thd{nghĩa trong lakehouse của lộ trình} & \thd{nghĩa khác hay gặp} \\
\midrule
\endhead
\bottomrule
\endfoot
warehouse & thư mục gốc trên object storage mà catalog Iceberg dùng để đặt bảng (thuộc tính \code{warehouse} khi cấu hình catalog) & data warehouse: hệ thống lưu trữ và truy vấn phân tích \\
watermark & mốc cao nhất đã nạp trong tải tăng dần (high-water mark) & mốc event time trong streaming, cho biết engine coi dữ liệu trước mốc đó là đã đủ (\href{https://nightlies.apache.org/flink/flink-docs-stable/docs/concepts/time/}{Flink: Timely Stream Processing}) \\
catalog & Iceberg catalog: giữ con trỏ tới metadata hiện hành của từng bảng (\href{https://iceberg.apache.org/spec/}{Iceberg spec}) & data catalog để tìm kiếm và mô tả dữ liệu cho người dùng \\
schema & namespace chứa bảng trong Spark SQL và Trino & danh sách cột và kiểu của một bảng \\
partition & partition của bảng Iceberg, có thể ẩn qua transform (\href{https://iceberg.apache.org/docs/latest/partitioning/}{Iceberg partitioning}) & phần dữ liệu trong bộ nhớ Spark khi tính toán; partition của topic Kafka \\
snapshot & một phiên bản của bảng Iceberg tại một lần commit & periodic snapshot fact table của Kimball; snapshot của dbt dùng cho SCD loại 2 (\href{https://docs.getdbt.com/docs/build/snapshots}{dbt snapshots}) \\
staging & tầng dữ liệu đã ép kiểu, nằm giữa raw và core & môi trường staging trước production; lớp model staging của dbt \\
commit & một lần đổi con trỏ metadata của bảng Iceberg & git commit; commit offset của Kafka consumer \\
lineage & đường đi của dữ liệu từ nguồn tới dashboard (\href{https://openlineage.io/docs/}{OpenLineage}) & lineage của Spark dùng để tính lại dữ liệu bị mất \\
contract & data contract giữa bên sinh và bên dùng dữ liệu (\href{https://bitol.io/}{Open Data Contract Standard}) & model contract của dbt: kiểm tên và kiểu cột khi build (\href{https://docs.getdbt.com/docs/mesh/govern/model-contracts}{dbt model contracts}) \\
\end{xltabular}

\plainchapter{Phụ lục B: Sách đọc thêm}

Sách không bắt buộc: mọi module đều học được chỉ bằng docs chính thức. Sách hữu ích khi bạn muốn một mạch trình bày dài, có ví dụ và lý giải thiết kế. Thông tin ấn bản được kiểm ngày 02/11/2026.

\begin{xltabular}{\textwidth}{@{}L P{4.3cm} P{1.8cm}@{}}
\toprule
\thd{sách} & \thd{ấn bản} & \thd{module} \\
\midrule
\endhead
\bottomrule
\endfoot
\href{https://martin.kleppmann.com/2026/03/24/designing-data-intensive-applications-2e.html}{Designing Data-Intensive Applications}, Martin Kleppmann và Chris Riccomini & bản 2, O'Reilly, 2026, 14 chương & 4, 5, 7, 8, 12 \\
\href{https://www.oreilly.com/library/view/fundamentals-of-data/9781098108298/}{Fundamentals of Data Engineering}, Joe Reis và Matt Housley & O'Reilly, 2022 & 2, 8 \\
\href{https://oreilly.com/library/view/the-data-warehouse/9781118530801}{The Data Warehouse Toolkit}, Ralph Kimball và Margy Ross & bản 3, Wiley, 2013 & 3 \\
\href{https://db-book.com/}{Database System Concepts} & bản 7, McGraw-Hill, 2019 & 1 \\
\href{https://oreilly.com/library/view/sql-antipatterns-volume/9798888650011}{SQL Antipatterns, Volume 1}, Bill Karwin & bản 2, Pragmatic Bookshelf, 2022 & 1 \\
Learning Spark & bản 2, O'Reilly, 2020 & 7 \\
\href{https://oreilly.com/library/view/high-performance-spark/9781098145842}{High Performance Spark} & bản 2, O'Reilly, 2026 & 7 \\
\href{https://oreilly.com/library/view/apache-iceberg-the/9781098148614}{Apache Iceberg: The Definitive Guide} & O'Reilly, 2024 & 5, 6 \\
\href{https://manning.com/books/data-pipelines-with-apache-airflow-second-edition}{Data Pipelines with Apache Airflow} & bản 2, Manning, 2026, viết cho Airflow 3 & 9 \\
\href{https://trino.io/trino-the-definitive-guide.html}{Trino: The Definitive Guide} & bản 2, O'Reilly, 2022; tải miễn phí qua Starburst & 11 \\
\href{https://oreilly.com/library/view/~/9781492043072}{Kafka: The Definitive Guide} & bản 2, O'Reilly, 2021 & 12 \\
\end{xltabular}

\plainchapter{Về tài liệu này}

\begin{itemize}
  \item Phiên bản 2.0 nháp, cập nhật ngày 06/11/2026. Phần I được viết lại thành tutorial thực hành ngày 06/11/2026, còn Phần II tới IV vẫn giữ nội dung bản 1.0, nên tên module trong mục lục có thể khác bản đồ lộ trình cho tới khi được viết lại. Chương capstone cũng vẫn ghi bốn mốc của bản 1.0. Giờ trong bảng module là số của bản 2.0. Tiêu đề giờ, khung và mục \textit{Lướt được} của các chương M4 đến M13 vẫn theo bản 1.0.
  \item Phiên bản phần mềm, thông tin chứng chỉ và link của bản 1.0 được kiểm với nguồn chính thức ngày 02/11/2026.
  \item Kiến thức kỹ thuật được đối chiếu với docs, spec và release notes chính thức của từng dự án và nhà cung cấp. Bài viết của kỹ sư có tên tuổi và sách được trích khi chúng giúp nối với kinh nghiệm SWE, và luôn được ghi rõ là ``bài viết'' hay ``sách''.
  \item Thời lượng là ước tính cho phần lõi; tốc độ thật phụ thuộc nền tảng của bạn.
  \item Khi một link hỏng hoặc một hành vi thay đổi ở phiên bản mới, nguồn đúng luôn là docs của phiên bản bạn đang chạy.
\end{itemize}


\end{document}
